% આ ભાગનો પૂર્ણ સંપાદનીય પાઠ છે. શૈલી, ફોન્ટ, ચિત્રો અને પ્રકરણસંદર્ભ માટે સંપૂર્ણ સ્રોત ZIP તથા reader/full-reader.tex જુઓ. આ એકલી સ્વતંત્ર કમ્પાઇલ થતી મુખ્ય ફાઇલ નથી.

% BEGIN OLP-0274
\olpart{inc}{અપૂર્ણતા}
\phantomsection\label{guunit:OLP-0274}


\begin{editorial}
 આ ભાગમાં અપૂર્ણતાનાં પ્રમેયો આવે છે. તે પ્રથમ-ક્રમના
 તર્કશાસ્ત્રના ભાગો (ખાસ કરીને સાબિતી-તંત્ર) અને સંગણનીયતા
 ભાગમાંના પુનરાવર્તી વિધેયો અંગેની સામગ્રી પર આધાર રાખે છે.
 આ લખાણ જેરેમી એવિગાડની નોંધો પર આધારિત છે; રિચર્ડ ઝાકે
 તેમાં સુધારા કર્યા છે.
\end{editorial}

\begingroup
\makeatletter\def\ol@sectioncs{section}\makeatother

% BEGIN OLP-0275
\olchapter{inc}{int}{અપૂર્ણતાનો પરિચય}
\olchapterid{inc}{int}
\phantomsection\label{guunit:OLP-0275}


\begingroup
\makeatletter\def\ol@sectioncs{section}\makeatother

% BEGIN OLP-0276
\olfileid{inc}{int}{bgr}

\olsection{ઐતિહાસિક પૃષ્ઠભૂમિ}
\phantomsection\label{guunit:OLP-0276}


આ વિભાગમાં અપૂર્ણતાનાં પ્રમેયોને યોગ્ય સંદર્ભમાં સમજવામાં
મદદરૂપ થતી ઐતિહાસિક ઘટનાઓની ટૂંકી ચર્ચા કરીશું.
ખાસ કરીને, ગાણિતિક તર્કશાસ્ત્રના ઇતિહાસનો ખૂબ જ
સંક્ષિપ્ત આલેખ આપીશું અને પછી ગણિતના પાયા વિશેના
ઇતિહાસની થોડી વાત કરીશું.

\begin{digress}
  ``ગાણિતિક તર્કશાસ્ત્ર'' શબ્દસમૂહના બે અર્થ થઈ શકે.
  ``ગાણિતિક'' શબ્દ વિષયને વર્ણવે છે એમ લઈએ, તો
  ``ગણિતનો તર્ક''નો અર્થ થાય: ગણિતીય દલીલોના સિદ્ધાંતો.
  એ શબ્દ પદ્ધતિને વર્ણવે છે એમ લઈએ, તો ``તર્કનું ગણિત''નો
  અર્થ થાય: દલીલોના સિદ્ધાંતોનો ગાણિતિક અભ્યાસ.
  આગળનું વર્ણન ઘણી વાર એકસાથે આ બંને અર્થમાં ગાણિતિક
  તર્કશાસ્ત્ર વિશે છે.
\end{digress}

તર્કશાસ્ત્રનો અભ્યાસ મૂળે એરિસ્ટોટલથી શરૂ થયો; તે
લગભગ 384--322 \textsc{bce} દરમિયાન જીવ્યા.
તેમનાં \emph{Categories}, \emph{Prior
  analytics} અને \emph{Posterior analytics}માં
વૈજ્ઞાનિક દલીલોના સિદ્ધાંતોનો વ્યવસ્થિત અભ્યાસ છે;
તેમાં નિગમનત્રયનો પણ વિગતવાર અને વ્યવસ્થિત અભ્યાસ છે.

મધ્યયુગના પંડિતીય દર્શનમાં એરિસ્ટોટલનું તર્કશાસ્ત્ર
પ્રભુત્વ ધરાવતું હતું. અઢારમી સદી જેટલા મોડા સમયમાં પણ
કાન્ટ માનતા હતા કે એરિસ્ટોટલનું તર્કશાસ્ત્ર સંપૂર્ણ છે
અને તેમાં સુધારાની જરૂર નથી. પરંતુ નિગમનત્રયનો સિદ્ધાંત
ગણિતીય દલીલના માત્ર ઉપરછલ્લા પાસાંઓનું જ નિરૂપણ કરી
શકે છે. તેની એક સદી પહેલાં, ન્યૂટનના સમકાલીન લાઇબનિટ્ઝે
તાર્કિક દલીલ માટે સંપૂર્ણ ``કલન''ની કલ્પના કરી હતી અને
એવું કલન રચવા તરફ કેટલાક પ્રાથમિક પગલાં પણ લીધાં;
તેમણે મૂળે વિધાનાત્મક તર્કશાસ્ત્રનું એક સ્વરૂપ વર્ણવ્યું.

ઓગણીસમી સદી તર્કશાસ્ત્ર માટે નિર્ણાયક બની. 1854માં
જ્યોર્જ બૂલે \emph{The Laws of Thought} લખ્યું;
તેમાં વિધાનાત્મક તર્કશાસ્ત્રનો વિસ્તૃત બીજગણિતીય અભ્યાસ
છે, જે આજની રજૂઆતોથી બહુ દૂર નથી. 1879માં ગોટલોબ
ફ્રેગેએ \emph{Begriffsschrift} (સંકેતલેખન) પ્રકાશિત
કર્યું. તે વિધાનાત્મક તર્કશાસ્ત્રમાં પરિમાણકો અને સંબંધો
ઉમેરીને પ્રથમ-ક્રમનું તર્કશાસ્ત્ર પણ સમાવે છે. વાસ્તવમાં,
ફ્રેગેની તાર્કિક પદ્ધતિઓમાં ઉચ્ચ-ક્રમનું તર્કશાસ્ત્ર અને
તેથી વધુ પણ હતું. પોતાના \emph{Basic Laws of
  Arithmetic}માં તેમણે બતાવવાનો પ્રયત્ન કર્યો કે
અંકગણિતનું સમગ્ર જ્ઞાન કેવળ તાર્કિક પૂર્વધારણાઓમાંથી
તેમના \emph{Begriffsschrift}માં વ્યૂત્પન્ન થઈ શકે.
દુર્ભાગ્યે, 1902માં રસેલે બતાવ્યું તેમ આ પૂર્વધારણાઓ
અસુસંગત નીવડી. તે અસુસંગત સ્વયંસિદ્ધ વાક્ય બાજુએ
રાખીએ તો ફ્રેગેએ લગભગ એકલા હાથે આધુનિક તર્કશાસ્ત્ર
શોધ્યું—આ નોંધપાત્ર સિદ્ધિ છે. બૂલ પછી બીજગણિતીય
અભિગમ ધરાવતા પીયર્સ અને શ્રોડર સહિતના વિચારકોએ પણ
પરિમાણકયુક્ત તર્કશાસ્ત્ર સ્વતંત્ર રીતે વિકસાવ્યું.

હવે ગણિતના પાયા અંગેના વિકાસ તરફ વળીએ. ગણિતમાં
તર્કશાસ્ત્રની મહત્ત્વની ભૂમિકા હોવાથી ઉપર વર્ણવેલા
વિકાસો સાથે તેનો ઘણો સંબંધ છે. દાખલા તરીકે, સમગ્ર
ગણિત માત્ર પોતાની તાર્કિક ગોઠવણી પર આધારિત છે એવું
બતાવવાના સ્પષ્ટ હેતુથી ફ્રેગેએ તેમનું તર્કશાસ્ત્ર
વિકસાવ્યું. ખાસ કરીને, કાન્ટના મત મુજબનાં અનુભવથી
સ્વતંત્ર \emph{સંશ્લેષક} સત્યોને બદલે ગણિત અનુભવથી
સ્વતંત્ર \emph{વિશ્લેષક} સત્યોનું બનેલું છે એમ તેઓ બતાવવા ઇચ્છતા હતા.

ઘણા લોકો શુદ્ધ ગણિતનો જન્મ ગ્રીક લોકો પાસે થયો એમ
માને છે. લગભગ 300 B.C.માં લખાયેલું યુક્લિડનું
\emph{Elements} ચોકસાઈ અને કડક દલીલ પર ભાર મૂકતું
ગ્રીક ગણિતનું પરિપક્વ ઉદાહરણ છે. તેના વ્યાખ્યાઓ અને
સાબિતીઓ આજના માધ્યમિક શાળાના ભૂમિતિનાં પાઠ્યપુસ્તકોમાં
લગભગ મૂળ સ્વરૂપે જ છે (જ્યાં સુધી ભૂમિતિ હજી શીખવાય છે).
ગણિતમાં જ નહીં, દર્શનની શાખાઓમાં પણ આ ગાણિતિક દલીલની
રીતને કડક દલીલનું આદર્શ માનવામાં આવ્યું છે. સ્પિનોઝાએ તો
નૈતિક અને ધાર્મિક દલીલો પણ યુક્લિડની શૈલીમાં રજૂ કરી હતી—
જે જોવામાં વિચિત્ર લાગે છે!

સત્તરમી સદીમાં ન્યૂટન અને લાઇબનિટ્ઝે કલન શોધ્યું.
(સદીઓ સુધી પ્રાથમિકતાનો ઉગ્ર વિવાદ રહ્યો, પરંતુ આજે
મોટાભાગના વિદ્વાનો માને છે કે બંને વિકાસ મુખ્યત્વે
સ્વતંત્ર હતા.) કલનમાં, દાખલા તરીકે, અનંત સૂક્ષ્મ
રાશિઓના અનંત સરવાળા વિશે દલીલો આવે છે. આ લક્ષણોને
કારણે બિશપ બર્કલીએ તેની ટીકા કરી: તેમના મતે ઈશ્વરમાં
માનવું તેમના સમયના ગણિત કરતાં ઓછું તર્કસંગત નહોતું.
અઢારમી સદીમાં કલનની પદ્ધતિઓનો વ્યાપક ઉપયોગ થયો;
લિયોનાર્ડ ઓયલરે અનંત સરવાળાઓવાળી ગણતરીઓથી નોંધપાત્ર
પરિણામો મેળવ્યાં.

ઓગણીસમી સદીમાં ગણિતશાસ્ત્રીઓએ કલનને વધુ મજબૂત પાયો
આપીને બર્કલીની ટીકાનો જવાબ આપવાનો પ્રયત્ન કર્યો.
કોશી, વાયરસ્ટ્રાસ, બોલ્ઝાનો વગેરેના પ્રયાસોથી
``એપ્સિલોન અને ડેલ્ટા'' વડે લક્ષ, સાતત્ય,
અવકલન અને સંકલનની આજની વ્યાખ્યાઓ મળી; એટલે કે
તેમાં અનંતસૂક્ષ્મ રાશિઓનો કોઈ ઉલ્લેખ નથી. સદીના
પછીના ભાગમાં ગણિતશાસ્ત્રીઓ આગળ વધીને વાસ્તવિક
સંખ્યાઓ સહિત સમગ્ર કલનને પ્રાકૃતિક સંખ્યાઓની મદદથી
સમજાવવા પ્રયત્નશીલ બન્યા. (ક્રોનેકર: ``પૂર્ણ સંખ્યાઓ
ઈશ્વરે બનાવી; બાકી બધું માણસનું કામ છે.'') 1872માં
ડેડકિન્ડે ``Continuity and the irrational numbers''
લખ્યું. તેમાં તેમણે સંમેય સંખ્યાઓના ગણ તરીકે વાસ્તવિક
સંખ્યાઓ કેવી રીતે ``રચવી'' તે બતાવ્યું (અને સંમેય
સંખ્યાઓને, આપણે જાણીએ છીએ તેમ, પ્રાકૃતિક સંખ્યાઓની
જોડીઓ તરીકે જોઈ શકાય). 1888માં તેમણે ``Was sind und
was sollen die Zahlen'' (લગભગ, ``પ્રાકૃતિક સંખ્યાઓ
શું છે અને તે કેવી હોવી જોઈએ?'') લખ્યું; તેમાં
પ્રાકૃતિક સંખ્યાઓને કેવળ ``તાર્કિક'' શબ્દોમાં સમજાવવાનો
હેતુ હતો. 1887માં ક્રોનેકરે ``\"Uber den Zahlbegriff''
(``સંખ્યાની સંકલ્પના વિશે'') લખ્યું; તેમાં તેમણે તમામ
ગણિતીય પદાર્થોને પૂર્ણ સંખ્યાઓ વડે નિરૂપવાની વાત કરી.
1889માં જ્યુસેપ્પે પિયાનોને પ્રાકૃતિક સંખ્યાઓ માટે
ઔપચારિક, સાંકેતિક સ્વયંસિદ્ધ વાક્યો આપ્યાં.

ઓગણીસમી સદીના અંતે અનંત અંગે વધુ સાહસિક અભિગમ પણ
દેખાયો. તે પહેલાં, અનંત પદાર્થો અને સંરચનાઓને
(જેમ કે પ્રાકૃતિક સંખ્યાઓનો ગણ) સાવચેતીથી જોવામાં
આવતાં: ``અનંત સંખ્યામાં'' એટલે ``જેટલાં જોઈએ તેટલાં'',
અને ``લક્ષ તરફ જાય'' એટલે ``જેટલા જોઈએ તેટલા નજીક
જાય'' એવો અર્થ લેવાતો. પરંતુ જ્યોર્જ કૅન્ટૉરે બતાવ્યું
કે અનંતને ખરેખર અનંત તરીકે લઈ શકાય. કૅન્ટૉર,
ડેડકિન્ડ અને અન્યના કાર્યથી આજે વ્યાપક સ્વીકૃત
ગણસિદ્ધાંત આધારિત ગણિતની સામાન્ય સમજ વિકસી.

આથી આપણે વીસમી સદીના તર્કશાસ્ત્ર અને પાયાના વિકાસો
સુધી પહોંચીએ છીએ. 1902માં રસેલે ફ્રેગેની તાર્કિક
પદ્ધતિમાં વિરોધાભાસ શોધ્યો. 1904માં ઝર્મેલોએ
કહેવાતા ``પસંદગીના સ્વયંસિદ્ધ વાક્ય'' વડે કૅન્ટૉરનો
સુવ્યવસ્થિત-ક્રમનો સિદ્ધાંત સાબિત કર્યો; આ સ્વયંસિદ્ધ
વાક્યની માન્યતા અંગે ઘણો વિવાદ થયો. 1910થી 1913
વચ્ચે રસેલ અને વ્હાઇટહેડના \emph{Principia Mathematica}ના
ત્રણ ખંડ પ્રસિદ્ધ થયા; તેમાં ગણિતને તર્કના પાયા પર
ઊભું કરવાના ફ્રેગેના કાર્યક્રમને આગળ વધાર્યો. પરંતુ
રસેલ અને વ્હાઇટહેડે બે એવા સિદ્ધાંતો સ્વીકારવા પડ્યા
જેને કેવળ તાર્કિક ગણવાનું મુશ્કેલ લાગતું હતું:
અનંતનું સ્વયંસિદ્ધ વાક્ય અને ``ઘટાડનીયતા''નું
સ્વયંસિદ્ધ વાક્ય. 1900ના દાયકામાં પોઇનકારે ગણિતમાં
``પોતે વ્યાખ્યાયિત થનારને સમાવતી સમષ્ટિ પર આધારિત
વ્યાખ્યાઓ''ના ઉપયોગની ટીકા કરી;
1910ના દાયકામાં બ્રાઉવરે સમગ્ર ગણિતને ``સ્ફુરણાવાદી''
પાયા પર નવેસરથી સ્થાપવાનો પ્રસ્તાવ મૂક્યો; તેમાં
વર્જિત મધ્યનો નિયમ ($!A \lor \lnot !A$) વપરાતો નથી.

ખરેખર વિચિત્ર સમય!{} સમગ્ર ગણિતને તર્કમાં ઘટાડવાનો
કાર્યક્રમ હવે ``તર્કવાદ'' કહેવાય છે અને ઉપર જણાવેલી
મુશ્કેલીઓને કારણે તેને સામાન્ય રીતે નિષ્ફળ માનવામાં
આવે છે. માનસિક સ્ફુરણાત્મક રચનાઓ વડે ગણિત વિકસાવવાનો
કાર્યક્રમ ``સ્ફુરણાવાદ'' કહેવાય છે; રોજિંદા ગણિત પર
તે અતિ કડક મર્યાદાઓ મૂકે છે એવું મનાય છે. સદીના
વળાંકે ડેવિડ હિલ્બર્ટ—અત્યંત પ્રભાવશાળી ગણિતશાસ્ત્રીઓમાંના
એક—કૅન્ટૉર અને ડેડકિન્ડે રજૂ કરેલી નવી અમૂર્ત
પદ્ધતિઓના પ્રબળ સમર્થક હતા: ``કૅન્ટૉરે આપણા માટે
રચેલા સ્વર્ગમાંથી આપણને કોઈ હાંકી કાઢશે નહીં.''
તેમ છતાં તેઓ આ નવી પદ્ધતિઓની પાયા અંગેની ટીકાઓને
પણ ગંભીરતાથી લેતા (વિચિત્ર રીતે, હવે આ પદ્ધતિઓ
``શાસ્ત્રીય'' કહેવાય છે). તેમણે બંને હેતુ સાધવાનો
નીચેનો રસ્તો સૂચવ્યો:
\begin{enumerate}
\item શાસ્ત્રીય પદ્ધતિઓને ઔપચારિક સ્વયંસિદ્ધ વાક્યો અને
  નિયમો વડે નિરૂપવી; ગણિતીય પ્રશ્નોને સ્વયંસિદ્ધ-તંત્રમાં
  સૂત્રો તરીકે નિરૂપવા.
\item સલામત, ``સાન્ત-પદ્ધતિઓ'' વડે આ ઔપચારિક
  નિષ્પત્તિ-તંત્રો સુસંગત છે તે સાબિત કરવું.
\end{enumerate}

પહેલો હેતુ સિદ્ધ કરવા હિલ્બર્ટનું કાર્ય ઘણું આગળ
વધ્યું. 1899માં તેમના પ્રસિદ્ધ ગ્રંથ \emph{Foundations
of geometry}માં તેમણે ભૂમિતિ માટે આવું કર્યું હતું.
પછીનાં વર્ષોમાં તેમણે અને તેમના વિદ્યાર્થીઓ તથા
સહયોગીઓએ ગણિતની બીજી શાખાઓમાં પણ એવું કરવાનો
પ્રયત્ન કર્યો. હિલ્બર્ટે પોતે અંકગણિત અને વિશ્લેષણ
માટે સ્વયંસિદ્ધ-તંત્રો આપ્યાં. ઝર્મેલોએ ગણસિદ્ધાંતનું
સ્વયંસિદ્ધીકરણ આપ્યું; ફ્રેન્કેલ, સ્કોલેમ,
ફોન નોયમાન અને અન્યોએ તેને વિસ્તાર્યું. 1920ના
દાયકાના મધ્ય સુધી ``સમગ્ર'' ગણિતના સ્વયંસિદ્ધીકરણનો
દાવો કરતા બે અભિગમ હતા: રસેલ અને વ્હાઇટહેડનું
\emph{Principia mathematica} અને જે ઝર્મેલો--ફ્રેન્કેલ
ગણસિદ્ધાંત કહેવાયું તે.

1921માં હિલ્બર્ટે આ તંત્રોની સુસંગતતા સાબિત કરવાનો
હેતુ ધરાવતો સંશોધન-કાર્યક્રમ શરૂ કર્યો. તેમાં તેમના
કેટલાક વિદ્યાર્થીઓ, ખાસ કરીને બર્નેય્સ, એકરમાન
અને પછી જેન્ટ્ઝેને મદદ કરી. મુખ્ય વિચાર આ હતો:
ગણિતમાં અસુસંગતતાની નિષ્પત્તિ શક્ય છે કે નહીં
તે પ્રશ્નને ચિહ્નોની શક્ય શ્રેણીઓ અંગેના સંયોજનાત્મક
પ્રશ્નમાં ફેરવવો; એટલે કે અંકગણિત, વિશ્લેષણ કે
ગણસિદ્ધાંતના સ્વયંસિદ્ધ-તંત્રમાંથી, દાખલા તરીકે,
$!A \land \lnot !A$ની સાચી નિષ્પત્તિ હોવાના
માપદંડને સંતોષતી વાક્યોની શક્ય શ્રેણીઓનો પ્રશ્ન.
આવી ચિહ્ન-શ્રેણી શક્ય નથી તેની સાબિતી પોતે ગાણિતિક
સાબિતી હોવાથી આ સ્વયંસિદ્ધ-તંત્રોમાં ઔપચારિક બનાવી
શકાય. બીજા શબ્દોમાં, દાખલા તરીકે અંકગણિત સુસંગત
છે એવું કહેતો કોઈ વાક્ય~$\OCon$ હશે. વધુમાં,
ખાસ કરીને જો તેની સાબિતી માટે માત્ર ખૂબ મર્યાદિત,
``સાન્ત'' સાધનો જોઈએ, તો તે વાક્ય સંબંધિત તંત્રોમાં
સાબિત થઈ શકવું જોઈએ.

બીજો હેતુ—વિકસાવેલા સ્વયંસિદ્ધ-તંત્રો દરેક ગણિતીય
પ્રશ્ન નક્કી કરે—બે રીતે ચોક્કસ કરી શકાય. એક રીતે
આવું કહી શકાય: ગણિતના સ્વયંસિદ્ધ-તંત્રની ભાષાનાં
દરેક વાક્ય~$!A$ માટે સ્વયંસિદ્ધ વાક્યોમાંથી
કાં તો~$!A$ અથવા~$\lnot !A$ સાબિત થાય.
આ સાચું હોય તો એવું કોઈ વાક્ય નહીં હોય જેને
આ સ્વયંસિદ્ધ વાક્યોના આધારે ન સાબિત કરી શકાય
કે ન નકારી શકાય; તંત્ર કોઈ પ્રશ્ન અનિર્ણીત નહીં
રાખે. આ ગુણધર્મવાળું સ્વયંસિદ્ધ-તંત્ર
\emph{પૂર્ણ} કહેવાય છે. અલબત્ત, કોઈ ચોક્કસ
વાક્ય માટે બેમાંથી કઈ શક્યતા છે તે નક્કી કરવું
હજી મુશ્કેલ હોઈ શકે. પરંતુ સિદ્ધાંતમાં એવું
કરવાની પદ્ધતિ હોવી જોઈએ. વાસ્તવમાં, હિલ્બર્ટે
વિચારેલાં સ્વયંસિદ્ધ અને નિષ્પત્તિ તંત્રો માટે
પૂર્ણતા એવી પદ્ધતિનું અસ્તિત્વ ફલિત કરે—જોકે
હિલ્બર્ટ આ વાત જાણતા નહોતા. પ્રશ્નનો બીજો
અર્થ વધુ મજબૂત શરત છે: આપેલું વાક્ય~$!A$
સ્વયંસિદ્ધ વાક્યોમાંથી વ્યૂત્પન્ન થાય છે કે નહીં
તે નક્કી કરતી યાંત્રિક, સંગણકીય પદ્ધતિ હોય.

1931માં ગડેલે બે ``અપૂર્ણતાનાં પ્રમેયો'' સાબિત
કર્યાં, જે બતાવે છે કે આ કાર્યક્રમ સફળ થઈ
શકે નહીં. ગણિતના પૂરતા શક્તિશાળી,
સુસંગત અને સ્વયંસિદ્ધીકરણીય તંત્રો પૂર્ણ
હોઈ શકતાં નથી. વધુ શરતો હેઠળ આવું તંત્ર
પોતાની સુસંગતતા વ્યક્ત કરતું વાક્ય સાબિત
કરી શકતું નથી.
\sourcecorrection{OLINC-001}{સ્થિર અંગ્રેજી મૂળ અહીં
``ગણિત માટેનું કોઈ તંત્ર પૂર્ણ નથી'' અને સુસંગતતાનું
વાક્ય ``ન સાબિત થાય કે ન નકારી શકાય'' એમ અતિ વ્યાપક
રીતે કહે છે. આ જ ગ્રંથના અપૂર્ણતા-પરિણામોના અવલોકનમાં
પ્રથમ પ્રમેય માટે સુસંગત, સ્વયંસિદ્ધીકરણીય અને
સંગણનીય વિધેયો/નિર્ણેય સંબંધોને નિરૂપતું તંત્ર
ધાર્યું છે; બીજા પ્રમેયમાં પોતાની સુસંગતતા સાબિત ન
થવાની વાત વધુ કડક શરતો હેઠળ છે. મૂળનું નકારી ન
શકાય તેવું સામાન્ય દાવું અહીં રાખ્યું નથી.}

આ હિલ્બર્ટના મૂળ કાર્યક્રમ માટે ભારે આઘાત હતો.
તેમ છતાં, ગણિતમાં ઘણી વાર થાય છે તેમ, તેનાથી
સંશોધન માટે નવી અને રસપ્રદ દિશાઓ પણ ખુલી.
ગણિતનું એક જ સર્વસમાવેશક ઔપચારિક તંત્ર ન હોય,
તો વધુ મર્યાદિત તંત્રો વિકસાવી તેમાં શું સાબિત
કરી શકાય તે તપાસવું યોગ્ય છે. આ તંત્રોની
સુસંગતતા સ્થાપવા માટે ઓછા મર્યાદિત સાબિતી-માર્ગો
વિકસાવવાનું અને તેમની સુસંગતતા સાબિત કરવી કેટલી
મુશ્કેલ છે તે માપવાનું પણ યોગ્ય છે. ગડેલે બતાવ્યું
કે (લગભગ) દરેક ઔપચારિક તંત્રમાં તે નક્કી ન
કરી શકે એવા પ્રશ્નો છે; તેથી કોઈ ચોક્કસ તંત્ર
નક્કી ન કરી શકે એવા ``રસપ્રદ'' પ્રશ્નો શોધવા
અને તેમને નક્કી કરવા માટે તંત્ર કેટલું શક્તિશાળી
હોવું જોઈએ તે સમજવું યોગ્ય છે. આજ સુધી
તર્કશાસ્ત્રીઓ સાબિતીઓના સિદ્ધાંત નામની નવી
ગણિતીય શાખામાં આ પ્રશ્નોનો અભ્યાસ કરે છે.
% END OLP-0276
\endgroup


\begingroup
\makeatletter\def\ol@sectioncs{section}\makeatother

% BEGIN OLP-0277
\olfileid{inc}{int}{def}

\olsection{વ્યાખ્યાઓ}
\phantomsection\label{guunit:OLP-0277}


ગણિતને ઔપચારિક બનાવવાનો અને આવું ઔપચારિકીકરણ
સુસંગત તથા પૂર્ણ છે એમ બતાવવાનો હિલ્બર્ટનો
કાર્યક્રમ હાથ ધરવા પહેલાં ભાષા, તાર્કિક માળખું
અને સ્વયંસિદ્ધ વાક્યોનું તંત્ર પસંદ કરવું પડે.
અહીં આપણે ધારી લઈએ કે ગણિતને પ્રથમ-ક્રમની
ભાષામાં ઔપચારિક બનાવી શકાય: એટલે, અચળો,
ફલનો અને વિધેયોનો કોઈ ગણ એવો
હોય કે પ્રથમ-ક્રમના તર્કશાસ્ત્રના સંયોજકો અને
પરિમાણકો સાથે મળીને તે ગણિતના દાવા વ્યક્ત
કરી શકે. મોટાભાગના લોકો માને છે કે એવી ભાષા
છે: ગણસિદ્ધાંતની ભાષા, જેમાં~$\in$ એકમાત્ર
અતાર્કિક ચિહ્ન છે. આટલી સાદી ભાષા આટલું બધું
વ્યક્ત કરી શકે છે એ પહેલી નજરે અવિશ્વસનીય
લાગે છે; ગણિતનો લગભગ આખો ભાગ આટલી
મર્યાદિત શબ્દાવલિમાં વ્યક્ત થઈ શકે છે તે
સ્થાપવામાં ઘણું કામ લાગ્યું. હમણાં માટે વાત
સરળ રાખવા અંકગણિત સુધી જ ચર્ચા મર્યાદિત
કરીએ, એટલે માત્ર પ્રાકૃતિક સંખ્યાઓ~$\Nat$
સાથે સંબંધિત ગણિત. અંકગણિતના તથ્યો વ્યક્ત
કરવાની સ્વાભાવિક ભાષા~$\Lang L_A$ છે. $\Lang L_A$માં
એક દ્વિસ્થાનિક વિધેય~$<$, એક
અચળ~$\Obj 0$, એક એકસ્થાનિક
ફલન~$\prime$ અને બે દ્વિસ્થાનિક
ફલનો~$+$ તથા~$\times$ છે.

\begin{defn}
વાક્યોનો ગણ~$\Gamma$ \emph{સિદ્ધાંત} છે, જો
તે તાર્કિક ફલિતતા હેઠળ બંધ હોય; એટલે કે જો
$\Gamma = \Setabs{!A}{\Gamma \Entails
!A}$ હોય.
\end{defn}

સિદ્ધાંતો નિર્દિષ્ટ કરવાની બે સરળ રીત છે. એક
રીત કોઈ સંરચનામાં સત્ય હોય એવાં
વાક્યોનો ગણ લેવાની. દાખલા તરીકે,
$\Lang L_A$ માટે એવું સંરચના લો જેમાં
પ્રદેશ~$\Nat$ હોય અને બધાં અતાર્કિક
ચિહ્નોના અર્થઘટન અપેક્ષા મુજબ હોય.

\begin{defn}\ollabel{def:standard-model}
\emph{અંકગણિતનું પ્રમાણભૂત નિદર્શ} નીચે મુજબ
વ્યાખ્યાયિત સંરચના~$\Struct N$ છે:
\begin{enumerate}
\item $\Domain N = \Nat$
\item $\Assign{\Obj 0}{N} = 0$
\item $\Assign{\Obj \prime}{N}(n) = n + 1$, દરેક $n \in \Nat$ માટે
\item $\Assign{\Obj +}{N}(n, m) = n + m$, દરેક $n, m \in \Nat$ માટે
\item $\Assign{\Obj \times}{N}(n, m) = n\cdot m$, દરેક $n, m \in \Nat$ માટે
\item $\Assign{\Obj <}{N} = \Setabs{\tuple{n, m}}{n \in \Nat, m \in
  \Nat, n < m}$
\end{enumerate}
\end{defn}

`$\times$' અને `$\cdot$'નો ભેદ નોંધો: $\times$
અંકગણિતની ભાષાનું ચિહ્ન છે. ગુણાકાર યાદ આવે
તે માટે એ પસંદ કર્યું છે, પણ $\times$ પોતે
ગુણાકારની ક્રિયા નથી; તે દ્વિસ્થાનિક વિધેય-ચિહ્ન
છે (સત્તાવાર રીતે~$\Obj f^2_1$). તેનાથી
વિપરીત~$\cdot$ સામાન્ય ગુણાકારનું વિધેય
\emph{છે}. $n \cdot m$ દેખાય ત્યારે તેનો
અર્થ $n$ અને~$m$ સંખ્યાઓનો ગુણાકાર છે;
$x \times y$ દેખાય ત્યારે અંકગણિતની ભાષાના
પદ વિશે વાત છે. પ્રમાણભૂત નિદર્શમાં વિધેય-ચિહ્ન
$\times$નું અર્થઘટન પ્રાકૃતિક સંખ્યાઓ પરના
વિધેય~$\cdot$ તરીકે થાય છે. સરવાળા માટે
આપણે~$+$ને અંકગણિતની ભાષાના વિધેય-ચિહ્ન
અને પ્રાકૃતિક સંખ્યાઓ પરના સરવાળાના વિધેય,
બંને માટે વાપરીએ છીએ. અહીં અર્થ પ્રસંગથી
નક્કી કરવો પડે.

\begin{defn}
\emph{સાચું અંકગણિત}નો સિદ્ધાંત અંકગણિતના
પ્રમાણભૂત નિદર્શમાં સંતોષાતાં વાક્યોનો ગણ છે;
એટલે કે,
\[
\Th{TA} = \Setabs{!A}{\Sat{N}{!A}}.
\]
\end{defn}

$\Th{TA}$ સિદ્ધાંત છે: $\Th{TA} \Entails !A$
હોય ત્યારે~$!A$, $\Th{TA}$ને સંતોષતા દરેક
સંરચનામાં સંતોષાય છે. $\Sat{N}{\Th{TA}}$
હોવાથી~$\Sat{N}{!A}$ મળે, તેથી
$!A \in \Th{TA}$.

સિદ્ધાંત~$\Gamma$ નિર્દિષ્ટ કરવાની બીજી રીત
તેને વાક્યોના કોઈ ગણ~$\Gamma_0$માંથી ફલિત
થતાં બધાં વાક્યોનો ગણ માનવાની છે.
એ સ્થિતિમાં~$\Gamma$, તાર્કિક ફલિતતા હેઠળ
$\Gamma_0$નું ``બંધ'' છે. આ રીતે સિદ્ધાંત
નિર્દિષ્ટ કરવો ત્યારે જ રસપ્રદ છે જ્યારે
$\Gamma_0$ સ્પષ્ટ રીતે આપવામાં આવ્યો હોય;
દાખલા તરીકે, $\Gamma_0$ના ઘટકોની યાદી હોય.
ઓછામાં ઓછું~$\Gamma_0$ નિર્ણેય તો હોવો
જોઈએ: કોઈ વાક્ય એ~$\Gamma_0$નો
ઘટક છે કે નહીં તેની સંગણકીય તપાસ શક્ય હોવી
જોઈએ. $\Gamma_0$નાં વાક્યોને આપણે
$\Gamma$નાં \emph{સ્વયંસિદ્ધ વાક્યો} અને
$\Gamma$ને~$\Gamma_0$ દ્વારા
\emph{સ્વયંસિદ્ધીકૃત} કહીએ છીએ.

\begin{defn}
સિદ્ધાંત~$\Gamma$ એ~$\Gamma_0$ વડે
\emph{સ્વયંસિદ્ધીકૃત} છે ત્યારે અને માત્ર ત્યારે,
જ્યારે
\[
\Gamma = \Setabs{!A}{\Gamma_0 \Entails !A}
\]
હોય.
\end{defn}

\begin{defn}
નીચેના વાક્યો વડે સ્વયંસિદ્ધીકૃત સિદ્ધાંત
$\Th{Q}$ ``રોબિન્સનનું~$\Th{Q}$'' કહેવાય છે;
તે અંકગણિતનો ખૂબ સરળ સિદ્ધાંત છે.
\begin{align*}
& \lforall[x][\lforall[y][(\eq[x'][y'] \lif \eq[x][y])]] \tag{$!Q_1$}\\
& \lforall[x][\eq/[\Obj 0][x']] \tag{$!Q_2$}\\
& \lforall[x][(\eq[x][\Obj 0] \lor \lexists[y][\eq[x][y']])] \tag{$!Q_3$}\\
& \lforall[x][\eq[(x + \Obj 0)][x]] \tag{$!Q_4$}\\
& \lforall[x][\lforall[y][\eq[(x + y')][(x + y)']]] \tag{$!Q_5$}\\
& \lforall[x][\eq[(x \times \Obj 0)][\Obj 0]] \tag{$!Q_6$}\\
& \lforall[x][\lforall[y][\eq[(x \times y')][((x \times y) + x)]]] \tag{$!Q_7$}\\
& \lforall[x][\lforall[y][(x < y \liff \lexists[z][\eq[(z' + x)][y]])]] \tag{$!Q_8$}
\end{align*}
વાક્યોનો ગણ $\{!Q_1, \dots, !Q_8\}$ એ
$\Th{Q}$નાં સ્વયંસિદ્ધ વાક્યો છે; એટલે આ
વાક્યોમાંથી ફલિત થતાં બધાં વાક્યો
$\Th{Q}$માં છે:
\[
\Th{Q} = \Setabs{!A}{\{!Q_1, \dots, !Q_8\} \Entails !A}.
\]
\end{defn}

\begin{defn}
માનો કે $!A(x)$ એ~$\Lang L_A$નું એવું
સૂત્ર છે જેના મુક્ત ચલ~$x$ અને
$y_1$, \dots, $y_n$ છે. તો નીચેના સ્વરૂપનું
કોઈ પણ વાક્ય
\[
\lforall[y_1][\dots\lforall[y_n][((!A(\Obj 0) \land \lforall[x][(!A(x)
\lif !A(x'))]) \lif \lforall[x][!A(x)])]]
\]
\emph{આગમન યોજનાનું} એક ઉદાહરણ છે.

\emph{પિયાનો અંકગણિત}~$\Th{PA}$ એ
$\Th{Q}$નાં સ્વયંસિદ્ધ વાક્યો અને આગમન
યોજનાનાં બધાં ઉદાહરણોથી સ્વયંસિદ્ધીકૃત
સિદ્ધાંત છે.
\end{defn}

\begin{explain}
આગમન યોજનાનું દરેક ઉદાહરણ~$\Struct{N}$માં
સત્ય છે. સૂત્ર~$!A$માં માત્ર એક મુક્ત
ચલ~$x$ હોય ત્યારે આ સૌથી સહેલું
દેખાય છે. ત્યારે~$!A(x)$, $\Struct{N}$માં
$\Nat$નો ઉપગણ~$X_{!A}$ વ્યાખ્યાયિત કરે છે.
$X_{!A}$ એ એવાં બધાં~$n \in \Nat$નો ગણ છે
જેમના માટે $s(x)=n$ હોય ત્યારે
$\Sat{N}{!A(x)}[s]$ છે. આગમન યોજનાનું
અનુરૂપ ઉદાહરણ આ છે:
\[
((!A(\Obj 0) \land \lforall[x][(!A(x) \lif !A(x'))]) \lif 
  \lforall[x][!A(x)]).
\]
તેનું પૂર્વાંગ~$\Struct{N}$માં સત્ય હોય,
તો $0 \in X_{!A}$ અને $n \in X_{!A}$ હોય
ત્યારે $n+1$ પણ તેમાં છે. $0 \in X_{!A}$
હોવાથી~$1 \in X_{!A}$ મળે; $1 \in X_{!A}$માંથી
$2 \in X_{!A}$ મળે; અને આમ આગળ.
તેથી દરેક~$n \in \Nat$ માટે
$n \in X_{!A}$. આનો અર્થ
$\lforall[x][!A(x)]$ એ~$\Struct{N}$માં
સંતોષાય છે.
\end{explain}

$\Th{Q}$ અને~$\Th{PA}$ બંને
સ્વયંસિદ્ધીકૃત સિદ્ધાંતો છે. મહત્ત્વનો પ્રશ્ન
એ છે કે તે કેટલા શક્તિશાળી છે? દાખલા તરીકે,
$\Th{PA}$, $\Lang L_A$માં વ્યક્ત થઈ શકે
એવાં~$\Nat$ વિશેનાં બધાં સત્યો સાબિત કરી
શકે છે? ચોક્કસ રીતે કહીએ, તો
$\Th{PA}$નાં સ્વયંસિદ્ધ વાક્યો
$\Lang L_A$માં પૂછાયેલા બધા પ્રશ્નો
નક્કી કરે છે?

બીજી રીતે પૂછીએ: શું $\Th{PA} = \Th{TA}$?
$\Th{TA}$ સ્પષ્ટપણે~$\Nat$ વિશેનાં બધાં
સત્યો સાબિત કરે છે (એટલે તેમને સમાવે છે) અને
$\Lang L_A$માં વ્યક્ત થઈ શકે એવા બધા
પ્રશ્નો નક્કી કરે છે. ખરેખર, $!A$ એ
$\Lang L_A$નું વાક્ય હોય તો કાં તો
$\Sat{N}{!A}$ અથવા~$\Sat{N}{\lnot !A}$;
એટલે કાં તો $\Th{TA} \Entails !A$ અથવા
$\Th{TA} \Entails \lnot !A$. આવા સિદ્ધાંતને
\emph{પૂર્ણ} કહીએ.

\begin{defn}
સિદ્ધાંત~$\Gamma$ \emph{પૂર્ણ} છે
ત્યારે અને માત્ર ત્યારે, જ્યારે તેની ભાષાના
દરેક વાક્ય~$!A$ માટે કાં તો
$\Gamma \Entails !A$ અથવા
$\Gamma \Entails \lnot !A$ હોય.
\end{defn}

\begin{explain}
પૂર્ણતાના પ્રમેય મુજબ
$\Gamma \Entails !A$ ત્યારે અને માત્ર ત્યારે
જ્યારે $\Gamma \Proves !A$. તેથી~$\Gamma$
પૂર્ણ છે ત્યારે અને માત્ર ત્યારે, જ્યારે
તેની ભાષાના દરેક વાક્ય~$!A$ માટે
કાં તો $\Gamma \Proves !A$ અથવા
$\Gamma \Proves \lnot !A$ હોય.
\end{explain}

બીજો પ્રશ્ન પણ ઊભો થાય છે: કોઈ
વાક્ય એ~$\Th{TA}$, $\Th{PA}$
અથવા માત્ર~$\Th{Q}$માં પણ છે કે નહીં તે
તપાસવાની સંગણકીય પ્રક્રિયા છે? કોઈ ગણ
(દાખલા તરીકે, વાક્યોનો ગણ)
ક્યારે નિર્ણેય છે તેની વ્યાખ્યા
આપીને પ્રશ્ન વધુ ચોક્કસ કરી શકીએ.

\begin{defn}
ગણ~$X$ \emph{નિર્ણેય} છે ત્યારે
અને માત્ર ત્યારે, જ્યારે એવી સંગણકીય પ્રક્રિયા
હોય જે ઇનપુટ~$x$ માટે $x \in X$ હોય તો
$1$ અને અન્યથા~$0$ આપે.
\end{defn}

એટલે પ્રશ્ન આ છે: $\Th{TA}$
($\Th{PA}$, $\Th{Q}$) નિર્ણેય છે?

આ બધા પ્રશ્નોનો જવાબ ``ના'' છે. આમાંથી
કોઈ સિદ્ધાંત નિર્ણેય નથી. જોકે આ બાબત માત્ર
આ ખાસ સિદ્ધાંતો પૂરતી નથી. ખરેખર,
ચોક્કસ શરતો સંતોષતા \emph{દરેક} સિદ્ધાંતને
એવાં જ પરિણામો લાગુ પડે છે. $\Th{Q}$
અને~$\Th{PA}$ને લાગુ પડતી એવી એક શરત
એ છે કે તેઓ સ્વયંસિદ્ધ વાક્યોના
નિર્ણેય ગણથી સ્વયંસિદ્ધીકૃત છે.

\begin{defn}
સિદ્ધાંત \emph{સ્વયંસિદ્ધીકરણીય} છે, જો
તે સ્વયંસિદ્ધ વાક્યોના નિર્ણેય ગણથી
સ્વયંસિદ્ધીકૃત હોય.
\end{defn}

\begin{ex}
વાક્યોના સાન્ત ગણથી સ્વયંસિદ્ધીકૃત
દરેક સિદ્ધાંત સ્વયંસિદ્ધીકરણીય છે, કારણ કે
દરેક સાન્ત ગણ નિર્ણેય છે. તેથી,
દાખલા તરીકે, $\Th{Q}$ સ્વયંસિદ્ધીકરણીય છે.

$\Th{PA}$ જેવા યોજના વડે સ્વયંસિદ્ધીકૃત
સિદ્ધાંતો પણ સ્વયંસિદ્ધીકરણીય છે. $!B$ એ
$\Th{PA}$ના સ્વયંસિદ્ધ વાક્યોમાંનું છે કે
નહીં તે તપાસવા—એટલે કે $!B$ એ
$\Th{PA}$નું સ્વયંસિદ્ધ વાક્ય હોય તો
$\Char{X}(!B)=1$ અને અન્યથા~$=0$
આપતું વિધેય~$\Char{X}$ ગણવા—આ કરી શકાય:
પહેલાં તપાસો કે $!B$ એ~$\Th{Q}$ના
સ્વયંસિદ્ધ વાક્યોમાંનું એક છે કે નહીં.
હોય તો જવાબ ``હા'' અને
$\Char{X}(!B)=1$. ન હોય તો તે આગમન
યોજનાનું ઉદાહરણ છે કે નહીં તપાસો. આ
તપાસ પદ્ધતિસર થઈ શકે છે. અહીં તે આ રીતે
થાય છે એમ જોવું કદાચ સહેલું છે: આગમન
યોજનાનું કોઈ પણ ઉદાહરણ થોડા સર્વપરિમાણકોથી
શરૂ થાય છે અને પછી શરતી ઉપ-સૂત્ર
આવે છે. તે શરતનું ઉત્તરાંગ
$\lforall[x][!A(x, y_1, \dots, y_n)]$
છે, જ્યાં $x$ અને $y_1$, \dots, $y_n$
એ બધાં~$!A$ના મુક્ત ચલ છે; $!B$ના
પ્રારંભિક પરિમાણકો $y_1$, \dots,~$y_n$
ને બાંધે છે. આ~$!A$ કાઢીને તેના મુક્ત
ચલો પ્રારંભના સર્વપરિમાણકો અને
$\lforall[x]$થી બંધાયેલા ચલો સાથે સરખા
છે એ તપાસ્યા પછી શરતનું પૂર્વાંગ
નીચેના સૂત્ર સાથે મેળ ખાય છે કે નહીં
તે તપાસીએ:
\[
!A(\Obj 0, y_1, \dots, y_n) \land \lforall[x][(!A(x, y_1, \dots, y_n)
\lif !A(x', y_1, \dots, y_n))]
\]
મેળ ખાય તો~$!B$ આગમન યોજનાનું
ઉદાહરણ છે; ન ખાય તો~$!B$ ઉદાહરણ નથી.
\end{ex}

આ પ્રશ્નનો—અને કયા સિદ્ધાંતો પૂર્ણ અથવા
નિર્ણેય છે એના વધુ સામાન્ય પ્રશ્નનો—જવાબ
શોધવામાં નીચેની વ્યાખ્યા પણ ઉપયોગી થશે.
યાદ કરો કે ગણ~$X$ ગણનીય છે ત્યારે
અને માત્ર ત્યારે, જ્યારે તે ખાલી હોય અથવા
કોઈ વ્યાપ્ત વિધેય
$f \colon \Nat \to X$ હોય. આવું વિધેય
$X$ની પરિગણના કહેવાય છે.

\begin{defn}
ગણ~$X$ \emph{સંગણકીય રીતે પરિગણનીય}
(ટૂંકમાં~સંગણકીય રીતે પરિગણનીય) કહેવાય છે ત્યારે અને માત્ર
ત્યારે, જ્યારે તે ખાલી હોય અથવા તેની
સંગણનીય પરિગણના હોય.
\end{defn}

સ્વયંસિદ્ધીકરણીયતા ઉપરાંત અપૂર્ણતાનાં
પ્રમેયો લાગુ પડે એવા સિદ્ધાંતો માટે બીજી શરત
એ છે કે તેઓ સંગણનીય વિધેયો અને
નિર્ણેય સંબંધો વિશેનાં મૂળભૂત તથ્યો
સાબિત કરવા જેટલા શક્તિશાળી હોય. ``મૂળભૂત
તથ્યો'' એટલે વિધેયના દરેક નિવેશ માટે તેનું
મૂલ્ય વ્યક્ત કરતાં વાક્યો. અને
``પૂરતા શક્તિશાળી'' એટલે આ વાક્યો એ
સિદ્ધાંતોના પ્રમેયોમાં હોય. દાખલા તરીકે
સરવાળો એક સામાન્ય સંગણનીય વિધેય છે.
$2$ અને~$3$ નિવેશ માટે~$+$નું મૂલ્ય~$5$
છે, એટલે $2+3=5$. અંકગણિતની ભાષામાં
$2$ અને~$3$ નિવેશ માટે~$+$નું મૂલ્ય~$5$
છે તે વ્યક્ત કરતું વાક્ય છે:
$(\num{2} + \num{3}) = \num{5}$.
દાખલા તરીકે, $\Th{Q}$ આ વાક્ય સાબિત
કરે છે. વધુ સામાન્ય રીતે, દરેક સંગણનીય
વિધેય~$f(x_1, x_2)$ માટે~$\Lang{L_A}$માં
એવું સૂત્ર~$!A_f(x_1, x_2, y)$
હોવું જોઈએ કે જ્યારે $f(n_1, n_2)=m$
હોય ત્યારે $\Th{Q} \Proves
!A_f(\num{n_1}, \num{n_2}, \num{m})$.
આ રીતે $\Th{Q}$ સાબિત કરે છે કે
$n_1$, $n_2$ નિવેશ માટે~$f$નું મૂલ્ય~$m$
છે. વાસ્તવમાં, આપણે થોડું વધુ માગીએ છીએ:
$n_1$, $n_2$ નિવેશ માટે આ મૂલ્ય સિવાય
બીજી કોઈ સંખ્યા~$f$નું મૂલ્ય નથી, તે પણ
સાબિત થાય.
નિર્ણેય સંબંધો માટે પણ આવી જ શરત
છે. આગળની બે વ્યાખ્યાઓમાં આ ચોક્કસ કરીએ.

\begin{defn}
સૂત્ર~$!A(x_1, \dots, x_k, y)$ એ
વિધેય~$f\colon \Nat^k \to \Nat$ને
$\Gamma$માં \emph{નિરૂપણ કરે} છે
ત્યારે અને માત્ર ત્યારે, જ્યારે
$f(n_1, \dots, n_k)=m$ હોય ત્યારે
\begin{enumerate}
\item $\Gamma \Proves !A(\num{n_1}, \dots, \num{n_k}, \num{m})$, અને
\item $\Gamma \Proves \lforall[y](!A(\num{n_1}, \dots, \num{n_k},
y) \lif y = \num{m})$.
\end{enumerate}
\end{defn}

\begin{defn}
સૂત્ર~$!A(x_1, \dots, x_k)$ સંબંધ
$R \subseteq \Nat^k$ને એ જ સિદ્ધાંતમાં
\emph{નિરૂપણ કરે}
છે ત્યારે અને માત્ર ત્યારે,
\begin{enumerate}
\item $R(n_1, \dots, n_k)$ હોય ત્યારે
  $\Gamma \Proves !A(\num{n_1},
\dots, \num{n_k})$, અને
\item $R(n_1, \dots, n_k)$ ન હોય ત્યારે
  $\Gamma \Proves \lnot
!A(\num{n_1}, \dots, \num{n_k})$.
\end{enumerate}
\sourcecorrection{OLINC-002}{સ્થિર અંગ્રેજી મૂળમાં આ
સંબંધના નિરૂપણની વ્યાખ્યામાં સિદ્ધાંત સ્પષ્ટ કહ્યો
નથી, છતાં બંને શરતોમાં~$\Gamma$ વપરાય છે.
અહીં ``એ જ સિદ્ધાંતમાં''થી અગાઉની વ્યાખ્યાનો
સિદ્ધાંત ઉદ્દિષ્ટ કર્યો છે; સૂત્રો બદલ્યાં નથી.}
\end{defn}

સિદ્ધાંત અપૂર્ણતાનાં પ્રમેયો લાગુ પડે એટલો
``પૂરતો શક્તિશાળી'' છે, જો તે બધાં સંગણનીય
વિધેયો અને બધાં નિર્ણેય સંબંધોને
નિરૂપણ કરે. $\Th{Q}$ અને તેના
વિસ્તારો આ શરત સંતોષે છે, પરંતુ આ સ્થાપવામાં
આપણને સમય લાગશે: $\Th{Q}$ કેવાં તથ્યો
સાબિત કરી શકે તે વિશે આ અ-સરળ પરિણામ
છે, અને $\Th{Q}$નાં થોડાં જ સ્વયંસિદ્ધ વાક્યોમાંથી
બધું સાબિત કરવું પડે છે. તેમ છતાં
$\Th{Q}$ ખૂબ નબળો સિદ્ધાંત છે. એટલે
$\Th{Q}$ બધાં સંગણનીય વિધેયોનું
નિરૂપણ કરે છે તે સાબિત કરવું મુશ્કેલ
હોવા છતાં, મોટાભાગના રસપ્રદ સિદ્ધાંતો
$\Th{Q}$ કરતાં વધુ શક્તિશાળી છે, એટલે
$\Th{Q}$ કરતાં વધુ સાબિત કરે છે.
$\Th{Q}$ કંઈક સાબિત કરે તો દરેક વધુ
શક્તિશાળી સિદ્ધાંત પણ કરે. $\Th{Q}$
બધાં સંગણનીય વિધેયોનું નિરૂપણ કરે છે,
તેથી દરેક વધુ શક્તિશાળી સિદ્ધાંત પણ
કરે છે. આમ, ઘણા રસપ્રદ સિદ્ધાંતો
અપૂર્ણતાનાં પ્રમેયોની આ શરત સંતોષે છે.
એટલે આપણું મુશ્કેલ કામ ઉપયોગી નીવડશે:
તે આ પ્રમેયો ઘણા સિદ્ધાંતોને લાગુ પડે
છે તે બતાવે છે. ``સમગ્ર ગણિત''ને
ઔપચારિક બનાવવાનો હેતુ ધરાવતો કોઈ
સિદ્ધાંત $\Th{Q}$ જે સાબિત કરે છે
તે બધું તો સાબિત કરી શકે જ જોઈએ;
ઓછામાં ઓછું પ્રાથમિક સંગણનાઓના
પરિણામોને તે સમાવી શકે. તેથી
``સમગ્ર ગણિત''નો ઉમેદવાર સિદ્ધાંત
અપૂર્ણતાનાં પ્રમેયો લાગુ પડે એવો હશે.
% END OLP-0277
\endgroup


\begingroup
\makeatletter\def\ol@sectioncs{section}\makeatother

% BEGIN OLP-0278
\olfileid{inc}{int}{ovr}

\olsection{અપૂર્ણતાનાં પરિણામોનું અવલોકન}
\phantomsection\label{guunit:OLP-0278}


હિલ્બર્ટને આશા હતી કે ગણિતને એવા
સ્વયંસિદ્ધીકરણીય સિદ્ધાંતમાં ઔપચારિક બનાવી
શકાય જેને પૂર્ણ અને નિર્ણેય
છે એમ સાબિત કરી શકાય. વધુમાં, તે આ
સિદ્ધાંતની સુસંગતતા ખૂબ નબળી, ``સાન્ત''
પદ્ધતિઓથી સાબિત કરવા ઇચ્છતા હતા;
તેના દ્વારા શાસ્ત્રીય ગણિતને સ્ફુરણાવાદના
પડકારોથી બચાવી શકાય. ગડેલનાં અપૂર્ણતાનાં
પ્રમેયોએ બતાવ્યું કે આ હેતુઓ હાંસલ થઈ
શકતા નથી.

ગડેલના પ્રથમ અપૂર્ણતા-પ્રમેયે બતાવ્યું કે
રસેલ અને વ્હાઇટહેડના \emph{Principia
Mathematica}નું એક સ્વરૂપ પૂર્ણ
નથી. પરંતુ સાબિતી ઘણી સામાન્ય હતી અને
વિવિધ પ્રકારના સિદ્ધાંતોને લાગુ પડે છે. આનો
અર્થ કે \emph{Principia Mathematica}
ગણિતને સંપૂર્ણ રીતે સમાવી શક્યું નહીં એટલું
જ નહીં, પરંતુ \emph{કોઈ પણ} સ્વીકાર્ય
સિદ્ધાંત તેમ કરી શકતો નથી. અપૂર્ણતાનાં
પ્રમેયો લાગુ પડવા માટે સિદ્ધાંતનાં કયાં
લક્ષણો પૂરતાં છે તે અલગ તારવવામાં, અને
ગડેલની સાબિતીને માત્ર આ લક્ષણો પર આધારિત
કરી સામાન્ય બનાવવામાં સમય લાગ્યો.
હવે આપણે અંકગણિતની ભાષા~$\Lang L_A$માં
રચાયેલા સિદ્ધાંતો માટે પ્રથમ અપૂર્ણતા-પ્રમેયનું
ખૂબ સામાન્ય સ્વરૂપ કહી શકીએ છીએ.

\begin{thm}
જો $\Gamma$ એ~$\Lang L_A$માં સુસંગત
અને સ્વયંસિદ્ધીકરણીય સિદ્ધાંત હોય, જે બધાં
સંગણનીય વિધેયો અને નિર્ણેય સંબંધોને
નિરૂપણ કરતો હોય, તો $\Gamma$
પૂર્ણ નથી.
\end{thm}

$\Gamma$ પૂર્ણ નથી એનો અર્થ એ
કે ઓછામાં ઓછું એક વાક્ય~$!A$ એવું
છે કે $\Gamma \Proves/ !A$ અને
$\Gamma \Proves/ \lnot !A$. આવું
વાક્ય ($\Gamma)$થી \emph{સ્વતંત્ર}
કહેવાય છે. ખરેખર, સ્વતંત્ર વાક્યો અસ્તિત્વ
ધરાવે છે તે આપણે તદ્દન ઝડપથી સાબિત કરી
શકીએ. પરંતુ ગડેલની આ પ્રમેયની સાબિતીની
શક્તિ એ છે કે તે આવા સ્વતંત્ર વાક્યનું
એક \emph{ચોક્કસ ઉદાહરણ} આપે છે. આ
રસપ્રદ રચનાથી~$\Gamma$ માટેનું
વાક્ય~$!G_\Gamma$ મળે છે, જેને
\emph{ગડેલ-વાક્ય} કહે છે. તે સાબિત ન થઈ
શકે, કારણ કે~$\Gamma$માં $!G_\Gamma$
એવું કહેતા વિધાનને સમકક્ષ છે કે
$!G_\Gamma$ એ~$\Gamma$માં સાબિત
થઈ શકતું નથી. આ રચના \emph{રચનાત્મક}
રીતે થાય છે: $\Gamma$નું સ્વયંસિદ્ધીકરણ
અને નિષ્પત્તિ તંત્રનું વર્ણન આપવામાં
આવે તો સાબિતી $!G_\Gamma$ ખરેખર લખવાની
પદ્ધતિ આપે છે.

ગડેલની સાબિતીની રચનામાં~$\Lang L_A$ના
પદો અને સૂત્રોના ગુણધર્મો તથા તેમના
પરની ક્રિયાઓને $\Lang L_A$માં જ વ્યક્ત કરવાની
રીત જોઈએ. તેમાં ``$!A$ એ વાક્ય
છે'', ``$\delta$ એ~$!A$ની
નિષ્પત્તિ છે'' જેવા ગુણધર્મો અને
$\Subst{!A}{t}{x}$ જેવી ક્રિયાઓ આવે છે.
આ રીતમાં (ક) ચિહ્નો અને તેમના અનુક્રમો—જેમાંથી
પદો, સૂત્રો, નિષ્પત્તિઓ વગેરે
બને છે—તેમને પ્રાકૃતિક સંખ્યાઓમાં ``સંકેતબદ્ધ''
કરી આ ગુણધર્મો અને સંબંધો વ્યક્ત કરવા પડે;
કારણ કે~$\Lang L_A$ પ્રાકૃતિક સંખ્યાઓ
વિશે જ વાત કરી શકે છે. (ખ) સંકેતબદ્ધીકરણ
એવું હોવું જોઈએ કે~$\Gamma$ સંબંધિત તથ્યો
સાબિત કરે. તેથી બતાવવું પડે કે આ ગુણધર્મો
પ્રાકૃતિક સંખ્યાઓના નિર્ણેય ગુણધર્મો
વડે અને ક્રિયાઓ પ્રાકૃતિક સંખ્યાઓ પરના
સંગણનીય વિધેયો વડે સંકેતબદ્ધ થાય છે.
આને ``વાક્યરચનાનું અંકગણિતીકરણ'' કહે છે.

વાક્યરચનાનું અંકગણિતીકરણ કેવી રીતે થાય છે
તે તપાસતાં પહેલાં, $\Gamma$ ``પૂરતો
શક્તિશાળી'' છે—એટલે બધાં સંગણનીય વિધેયો
અને નિર્ણેય સંબંધોને નિરૂપણ
કરે છે—એ શરત જોઈએ. તે માટે
``સંગણનીય''ની ચોક્કસ વ્યાખ્યા આપવી પડશે.
આ અનેક રીતે કરી શકાય; દાખલા તરીકે,
ટ્યુરિંગ મશીનોના નિદર્શથી અથવા કોઈ
સામાન્ય ઉપયોગની પ્રોગ્રામિંગ ભાષામાંના
પ્રોગ્રામો વડે ગણાતાં વિધેયો તરીકે. પરંતુ
આ વિધેયો અને સંબંધોને ભાષા~$\Lang L_A$માં
રચાયેલા સિદ્ધાંતમાં નિરૂપણ કરવાનો
આપણો હેતુ હોવાથી, માત્ર સંખ્યાત્મક વિધેયોની
સંગણનીયતાની સરળ વ્યાખ્યા પસંદ કરવી સારી.
તે \emph{પુનરાવર્તી વિધેય}નો ખ્યાલ છે.
તેથી પહેલાં આપણે પુનરાવર્તી વિધેયોની ચર્ચા
કરીશું. ત્યાર પછી બતાવીશું કે~$\Th{Q}$
પહેલેથી જ બધાં પુનરાવર્તી વિધેયો અને
સંબંધોને નિરૂપણ કરે છે. આમ
$\Th{Q}$ અને~$\Th{PA}$ જેવા ચોક્કસ
સિદ્ધાંતોને અપૂર્ણતા-પ્રમેય લાગુ કરી શકીશું,
કારણ કે તેઓ ``પૂરતા શક્તિશાળી'' સિદ્ધાંતોનાં
ઉદાહરણો છે તે સ્થાપિત થઈ જશે.

વાક્યરચનાના અંકગણિતીકરણનું અંતિમ પરિણામ
સૂત્ર $\OProv[\Gamma](x)$ છે. તે
સૂત્રોને સંખ્યાઓથી સંકેતબદ્ધ કરીને
$\Gamma$નાં સ્વયંસિદ્ધ વાક્યોમાંથી સાબિત
થવાની શક્યતા વ્યક્ત કરે છે. ચોક્કસ રીતે,
$!A$નો સંકેત નંબર~$n$ હોય અને
$\Gamma \Proves !A$ હોય, તો
$\Gamma \Proves \Prov[\Gamma](\num{n})$.
$\Gamma$ માટેનું આ ``સાબિતીક્ષમતા-વિધેય''
આપણને ચોક્કસ અર્થમાં~$\Gamma$ની સુસંગતતા
$\Lang L_A$ના વાક્ય તરીકે વ્યક્ત
કરવા દે છે. $\Gamma$ માટેનું
``સુસંગતતા-વિધાન'' વાક્ય
$\lnot \Prov[\Gamma](\num{n})$ ગણીએ,
જ્યાં~$n$ કોઈ વિરોધાભાસનો સંકેત નંબર,
દાખલા તરીકે~$\lfalse$નો, છે. બીજો
અપૂર્ણતા-પ્રમેય કહે છે કે યોગ્ય વધુ શરતો
સંતોષતા સુસંગત સ્વયંસિદ્ધીકરણીય સિદ્ધાંતો પોતાના
સુસંગતતા-વિધાનને પણ સાબિત કરતા નથી.
આ પ્રમેય લાગુ કરવા માટે માત્ર બધાં
સંગણનીય વિધેયો અને નિર્ણેય
સંબંધોનું નિરૂપણ કરવાની શરત કરતાં
થોડી વધુ કડક શરતો જોઈએ; $\Th{PA}$
તે સંતોષે છે તે આપણે બતાવીશું.
\sourcecorrection{OLINC-003}{સ્થિર અંગ્રેજી મૂળના
પહેલા વાક્યમાં બધા સુસંગત સ્વયંસિદ્ધીકરણીય
સિદ્ધાંતો માટે પોતાની સુસંગતતા સાબિત ન થવાનો
દાવો આવે છે. એ જ ફકરો તરત વધુ કડક શરતો
જરૂરી હોવાનું કહે છે. તેથી અહીં દાવામાં
એ શરતો સ્પષ્ટ કરી છે.}
% END OLP-0278
\endgroup


\begingroup
\makeatletter\def\ol@sectioncs{section}\makeatother

% BEGIN OLP-0279
\olfileid{inc}{int}{dec}

\olsection{અનિર્ણેયતા અને અપૂર્ણતા}
\phantomsection\label{guunit:OLP-0279}


ગડેલની અપૂર્ણતાનાં પ્રમેયોની સાબિતી માટે
વાક્યરચનાનું અંકગણિતીકરણ જરૂરી છે. પરંતુ
તેના વિના પણ સિદ્ધાંત બધા નિર્ણેય
સંબંધોને નિરૂપણ કરે છે એવી ધારણાથી
કેટલાંક સરસ પરિણામો મેળવી શકાય છે. તેની
સાબિતી અટકવાની સમસ્યાની અનિર્ણેયતાની
સાબિતી જેવી વિકર્ણ દલીલ છે.

\begin{thm}
જો $\Gamma$ એવો સુસંગત સિદ્ધાંત હોય જે
દરેક નિર્ણેય સંબંધને નિરૂપણ
કરે છે, તો $\Gamma$ નિર્ણેય નથી.
\end{thm}

\begin{proof}
ધારો કે~$\Gamma$ નિર્ણેય છે. હવે
બતાવીએ કે તે દરેક નિર્ણેય સંબંધને
નિરૂપણ કરતો હોય તો~$\Gamma$ અસંગત થવો જ
પડે.

નિર્ણેય ગુણધર્મો (એકસ્થાનિક સંબંધો)
એક મુક્ત ચલવાળાં સૂત્રો વડે નિરૂપાય
છે. આવા બધાં સૂત્રોની સંગણનીય પરિગણના
$!A_0(x)$, $!A_1(x)$, \dots, લો. હવે
નીચેનો ગણ~$D \subseteq \Nat$ વિચારો:
\[
D = \Setabs{n}{\Gamma \Proves \lnot !A_n(\num{n})}
\]
$D$ નિર્ણેય છે, કારણ કે $n \in D$
છે કે નહીં તે તપાસવા પહેલાં~$!A_n(x)$
અને તેમાંથી~$\lnot !A_n(\num{n})$ ગણી
શકાય. $!A_n(x)$માં $x$ની દરેક મુક્ત
ઘટનાની જગ્યાએ પદ~$\num{n}$ મૂકવું અને
$!A(\num{n})$ પહેલાં~$\lnot$ મૂકવું સ્પષ્ટ
રીતે યાંત્રિક કામ છે. ધારણા મુજબ~$\Gamma$
નિર્ણેય હોવાથી
$\lnot !A(\num{n}) \in \Gamma$ છે કે નહીં
તપાસી શકાય. હોય તો~$n \in D$ અને
ન હોય તો~$n \notin D$. આમ~$D$
પણ નિર્ણેય છે.
\sourcecorrection{OLINC-004}{સ્થિર અંગ્રેજી મૂળમાં
આ ફકરાના પછીના બે સૂત્રોમાં~$!A$ લખ્યું છે;
પરંતુ અહીં પરિગણનામાંથી પસંદ થયેલું
$!A_n$ ઉદ્દિષ્ટ છે. મૂળનાં સૂત્રો જાળવી
આ સૂચન બાજુમાં આપ્યું છે.}

$\Gamma$ બધાં નિર્ણેય ગુણધર્મોને
નિરૂપણ કરતો હોવાથી~$D$ને પણ
નિરૂપણ કરે છે. $\Gamma$માં~$D$ને
નિરૂપણ કરનારાં સૂત્રો પણ
$!A_0(x)$, $!A_1(x)$, \dots,ની યાદીમાં
જ છે. તેથી એવો નંબર~$d$ લો કે
$!A_d(x)$, $\Gamma$માં~$D$ને
નિરૂપણ કરે. જો $d \notin D$ હોય,
તો $!A_d(x)$ એ~$D$ને નિરૂપણ
કરતું હોવાથી $\Gamma \Proves
\lnot !A_d(\num{d})$. પરંતુ તેનો અર્થ
એ કે~$d$ એ~$D$ની વ્યાખ્યાયિત શરત
સંતોષે છે, એટલે~$d \in D$. આ
$d \notin D$થી વિરોધાભાસ છે. તેથી
પરોક્ષ સાબિતીથી~$d \in D$.

$d \in D$ હોવાથી~$D$ની વ્યાખ્યા મુજબ
$\Gamma \Proves \lnot !A_d(\num{d})$.
બીજી તરફ, $!A_d(x)$ એ~$\Gamma$માં
$D$ને નિરૂપણ કરતું હોવાથી
$\Gamma \Proves !A_d(\num{d})$.
આથી~$\Gamma$ અસંગત છે.
\end{proof}

\begin{explain}
પહેલાનો પ્રમેય બતાવે છે કે બધાં
નિર્ણેય સંબંધોને નિરૂપણ
કરતો કોઈ સુસંગત સિદ્ધાંત નિર્ણેય
હોઈ શકે નહીં. આપણે બતાવીશું કે~$\Th{Q}$
ખરેખર બધા નિર્ણેય સંબંધોને
નિરૂપણ કરે છે. તેથી~$\Th{Q}$ને
સમાવતા બધા સિદ્ધાંતો, જેમ કે~$\Th{PA}$
અને~$\Th{TA}$, પણ આવું કરે છે અને તેથી
નિર્ણેય નથી. (આ બધા સિદ્ધાંતો
પ્રમાણભૂત નિદર્શમાં સત્ય હોવાથી સુસંગત છે.)

આ પરિણામથી પ્રથમ અપૂર્ણતા-પ્રમેયનું એક
નબળું સ્વરૂપ પણ મળે છે. જે સિદ્ધાંત
સ્વયંસિદ્ધીકરણીય અને પૂર્ણ બંને
હોય તે નિર્ણેય હોય છે. તેથી
સ્વયંસિદ્ધીકરણીય હોય અને બધા
નિર્ણેય ગુણધર્મોને નિરૂપણ
કરતા સુસંગત સિદ્ધાંતો પૂર્ણ
હોઈ શકતા નથી.
\end{explain}

\begin{thm}
જો $\Gamma$ સ્વયંસિદ્ધીકરણીય અને
પૂર્ણ હોય, તો તે નિર્ણેય છે.
\end{thm}

\begin{proof}
દરેક અસંગત સિદ્ધાંત નિર્ણેય છે,
કારણ કે તે બધાં વાક્યો સમાવે છે.
આથી ``શું $!A \in \Gamma$?''નો જવાબ
હંમેશાં ``હા'' હોય છે; એટલે તેને નક્કી
કરી શકાય છે.

હવે ધારો કે~$\Gamma$ સુસંગત છે અને
સ્વયંસિદ્ધીકરણીય તથા પૂર્ણ પણ છે.
સ્વયંસિદ્ધીકરણીય હોવાથી~$\Gamma$
સંગણકીય રીતે પરિગણનીય છે. ખરેખર,
$\Gamma$નાં સ્વયંસિદ્ધ વાક્યોમાંથી મળતી
બધી સાચી નિષ્પત્તિઓને સંગણનીય
વિધેયથી પરિગણી શકાય છે. સાચી
નિષ્પત્તિમાંથી તે જે વાક્ય
નિષ્પન્ન કરે છે તે ગણી શકાય છે;
એટલે આ બંનેને જોડીને~$\Gamma$ના બધા
પ્રમેયોની પરિગણના કરતું સંગણનીય વિધેય
મળે છે. $\Gamma$ સુસંગત અને પૂર્ણ
હોવાથી વાક્ય એ~$\Gamma$નો પ્રમેય છે ત્યારે અને
માત્ર ત્યારે જ્યારે તેનો નિષેધ~$\lnot !A$ પ્રમેય નથી.
આથી $!A \in \Gamma$ છે કે નહીં તે
આ રીતે નક્કી કરી શકાય: $\Gamma$ના
બધા પ્રમેયોની પરિગણના કરો. યાદીમાં~$!A$
આવે ત્યારે $\Gamma \Proves !A$ જાણીએ.
યાદીમાં~$\lnot !A$ આવે ત્યારે
$\Gamma \Proves/ !A$ જાણીએ. $\Gamma$
પૂર્ણ હોવાથી આ બેમાંથી એક ઘટના
અંતે બને જ છે, એટલે પ્રક્રિયા જવાબ આપે છે.
\end{proof}

\begin{cor}
\ollabel{cor:incompleteness}
જો $\Gamma$ સુસંગત, સ્વયંસિદ્ધીકરણીય
અને દરેક નિર્ણેય ગુણધર્મને
નિરૂપણ કરતો હોય, તો તે
પૂર્ણ નથી.
\end{cor}

\begin{proof}
$\Gamma$ પૂર્ણ હોત તો
પહેલાના પ્રમેય મુજબ તે નિર્ણેય
હોત (કારણ કે તે સ્વયંસિદ્ધીકરણીય
અને સુસંગત છે). પરંતુ~$\Gamma$
દરેક નિર્ણેય ગુણધર્મને
નિરૂપણ કરતો હોવાથી પહેલા
પ્રમેય મુજબ નિર્ણેય નથી.
\end{proof}

\begin{prob}
$\Th{TA} = \Setabs{!A}{\Sat{N}{!A}}$
સ્વયંસિદ્ધીકરણીય નથી તે બતાવો.
$\Th{TA}$ બધા નિર્ણેય ગુણધર્મોનું
નિરૂપણ કરે છે તે ધારી શકો છો.
\end{prob}

દાખલા તરીકે~$\Th{Q}$ બધા
નિર્ણેય ગુણધર્મોને નિરૂપણ
કરે છે તે સ્થાપિત કર્યા પછી, આ પરિણામથી
$\Th{Q}$ અપૂર્ણ છે એમ મળે છે. પરંતુ તેની
સાબિતી સ્વતંત્ર વાક્યનું ઉદાહરણ
આપતી નથી; ફક્ત આવું વાક્ય
અસ્તિત્વ ધરાવે છે તે બતાવે છે. ચોક્કસ
ઉદાહરણ માટે વાક્યરચનાનું અંકગણિતીકરણ
કરી ગડેલની મૂળ સાબિતીનો વિચાર અનુસરવો
પડે. અને પહેલો દાવો—$\Th{Q}$ ખરેખર
બધા નિર્ણેય ગુણધર્મોને
નિરૂપણ કરે છે—હજી બતાવવાનો છે.

દરેક \emph{રસપ્રદ} સિદ્ધાંત અપૂર્ણ કે
અનિર્ણેય હોય એવું નથી. ઘણા સિદ્ધાંતો
રસપ્રદ ગાણિતિક તથ્યો વ્યક્ત કરી શકે
એટલા શક્તિશાળી છે, છતાં ગડેલના પરિણામની
શરતો સંતોષતા નથી. દાખલા તરીકે,
$\Th{Pres} = \Setabs{!A \in \Lang{L_{A^+}}}{\Sat{N}{!A}}$,
એટલે ગુણાકાર-ચિહ્ન~$\times$ વિનાની
અંકગણિતની ભાષાના વાક્યોનો
પ્રમાણભૂત નિદર્શમાં સત્ય હોય એવો ગણ,
પૂર્ણ અને નિર્ણેય બંને છે. આ સિદ્ધાંત
\emph{પ્રેસબર્ગર અંકગણિત} કહેવાય છે અને
માત્ર~$\Obj{0}$, $\prime$ અને~$+$ વડે
વ્યક્ત થઈ શકે એવાં પ્રાકૃતિક સંખ્યાઓ
વિશેનાં બધાં સત્યો સાબિત કરે છે.
% END OLP-0279
\endgroup


\OLEndChapterHook
% END OLP-0275
\endgroup


\begingroup
\makeatletter\def\ol@sectioncs{section}\makeatother

% BEGIN OLP-0280
\olchapter{inc}{art}{વાક્યરચનાનું અંકગણિતીકરણ}
\olchapterid{inc}{art}
\phantomsection\label{guunit:OLP-0280}


\begin{editorial}
  નોંધો કે ચિહ્નિત ટેબ્લો માટેનું અંકગણિતીકરણ હજી ઉપલબ્ધ નથી.
\end{editorial}

\begingroup
\makeatletter\def\ol@sectioncs{section}\makeatother

% BEGIN OLP-0281
\olfileid{inc}{art}{int}
\olsection{પરિચય}
\phantomsection\label{guunit:OLP-0281}


સંગણનીયતા અને તર્કને જોડવા માટે તર્કના
પદાર્થો—ચિહ્નો, પદો, સૂત્રો અને
નિષ્પત્તિઓ—તેમના પરની ક્રિયાઓ તથા
તેમના ગુણધર્મો અને સંબંધો વિશે એવી રીતે
વાત કરવી જોઈએ કે તેમને સંગણકીય રીતે
સંભાળી શકાય. ચિહ્નો, ચિહ્નોના અનુક્રમો
અને તેમાંથી બનેલા બીજા પદાર્થો પરનાં
સંગણનીય વિધેયો અને સંબંધો સીધાં લઈને
આવું કરી શકાય. તાર્કિક વાક્યરચનાના
બધા પદાર્થો સાન્ત છે અને ચિહ્નોના
ગણનીય ગણમાંથી બને છે; તેથી
સંગણનાના કેટલાક નિદર્શોમાં આ શક્ય છે.
પરંતુ બીજા નિદર્શો—જેમ કે પુનરાવર્તી
વિધેયો—સંખ્યાઓ અને તેમના સંબંધો તથા
વિધેયો સુધી મર્યાદિત છે. વધુમાં, અંતે
આપણે અમુક સિદ્ધાંતોની અંદર પણ વાક્યરચના
સંભાળવી છે, ખાસ કરીને અંકગણિતની ભાષામાં
રચાયેલા સિદ્ધાંતોની અંદર. આવી સ્થિતિમાં
વાક્યરચનાને \emph{અંકગણિતીકૃત} કરવી
જરૂરી છે: એટલે વાક્યરચનાના પદાર્થો,
તેમના પરની ક્રિયાઓ અને તેમના સંબંધોને
અનુક્રમે સંખ્યાઓ, અંકગણિતીય વિધેયો અને
અંકગણિતીય સંબંધો તરીકે નિરૂપવા. આ
વિચાર લાઇબનિટ્ઝ સુધી પહોંચે છે: વાક્યરચનાના
પદાર્થોને સંખ્યાઓ આપવી.

ચિહ્નોને તેમના ``સંકેત નંબર'' તરીકે સંખ્યાઓ
આપવી પ્રમાણમાં સરળ છે. કેટલાક ચિહ્નો
થોડો પડકાર ઊભો કરે છે: દાખલા તરીકે,
અનંત ચલો છે અને દરેક
સ્થાનિકતા~$n$ માટે પણ અનંત
ફલનો છે. છતાં ચિહ્નોને પદ્ધતિસર
અલગ સંખ્યાઓ આપી શકાય; જેમ કે
$\Obj v_2$ અને~$\Obj v_3$ના સંકેત
નંબર જુદા હોય. ચિહ્નોના અનુક્રમો—પદો
અને સૂત્રો જેવા—વધુ મોટો પડકાર
છે. પરંતુ સંખ્યાઓના અનુક્રમોને શુદ્ધ
અંકગણિતીય રીતે સંભાળી શકીએ (દાખલા તરીકે,
અવિભાજ્ય સંખ્યાઓની ઘાતો વડે અનુક્રમોને
સંકેતબદ્ધ કરીને), તો એક-એક ચિહ્નના
સંકેતબદ્ધીકરણને ચિહ્નોના અનુક્રમો સુધી,
અને પછી સૂત્રોના અનુક્રમો અથવા
નિષ્પત્તિઓ જેવી બીજી ગોઠવણો સુધી
વિસ્તારી શકીએ. આ વિસ્તૃત સંકેતબદ્ધીકરણને
``ગડેલ-સંખ્યાકરણ'' કહે છે. દરેક પદ,
સૂત્ર અને નિષ્પત્તિને એક
ગડેલ-સંખ્યા અપાય છે.

ચિહ્નોના અનુક્રમોને તેમના સંકેત નંબરોના
અનુક્રમો તરીકે સંકેતબદ્ધ કરીને અને આવા
અનુક્રમોનું સંકેતબદ્ધીકરણ સંગણનીય વિધેયોથી
સંભાળી શકાય એવું પસંદ કરીને, ગડેલ-સંખ્યાઓ
પર પણ સંગણનીય વિધેયો વાપરી શકીએ છીએ.
વ્યવહારમાં જરૂરી બધા વિધેયો આદિમ
પુનરાવર્તી હશે. દાખલા તરીકે, અનુક્રમની
લંબાઈ અને તેના સંકેત નંબર પરથી તેનો
$i$મો ઘટક ગણવાના બંને વિધેયો આદિમ
પુનરાવર્તી છે. અનુક્રમનો સંકેત નંબર,
દાખલા તરીકે, સૂત્ર~$!A$ની
ગડેલ-સંખ્યા હોય, તો સૂત્રની
લંબાઈ અને તે સૂત્રમાંના $i$મા ચિહ્નનો સંકેત
નંબર પણ~$!A$ની ગડેલ-સંખ્યામાંથી
ગણી શકાય છે. યોગ્ય રીતે રચાયેલા પદની
અથવા યોગ્ય નિષ્પત્તિની ગડેલ-સંખ્યા
હોવાનો ગુણધર્મ આદિમ પુનરાવર્તી છે તે
સાબિત કરવું થોડું કઠિન છે. છતાં આવું
શક્ય છે, કારણ કે આપણે જોતા અનુક્રમો—પદો,
સૂત્રો અને નિષ્પત્તિઓ—આગમનાત્મક
રીતે વ્યાખ્યાયિત છે.

ઉદાહરણ તરીકે, પ્રતિસ્થાપનની ક્રિયા લો.
$!A$ સૂત્ર, $x$ ચલ અને $t$ પદ હોય,
તો~$\Subst{!A}{t}{x}$ એ~$!A$માં
$x$ની દરેક મુક્ત ઘટનાને~$t$થી
બદલવાનું પરિણામ છે. હવે ધારો કે~$!A$,
$x$, $t$ને અનુક્રમે ગડેલ-સંખ્યાઓ
$k$, $l$ અને~$m$ આપી છે. એ જ
યોજના~$\Subst{!A}{t}{x}$ને પણ
ગડેલ-સંખ્યા આપે છે, માનો~$n$.
$k$, $l$ અને~$m$માંથી~$n$માં
જતો આ નકશો પ્રતિસ્થાપનની ક્રિયાનો
અંકગણિતીય સમકક્ષ છે. પ્રતિસ્થાપનની
ક્રિયા $!A$, $x$, $t$ને
$\Subst{!A}{t}{x}$માં મોકલે છે;
અંકગણિતીકૃત પ્રતિસ્થાપન વિધેય
ગડેલ-સંખ્યાઓ $k$, $l$, $m$ને
ગડેલ-સંખ્યા~$n$માં મોકલે છે.
આ વિધેય આદિમ પુનરાવર્તી છે તે આપણે
જોઈશું.

વાક્યરચનાનું અંકગણિતીકરણ માત્ર અમૂર્ત
રસની વાત નથી. અલબત્ત, અંકગણિતની ભાષા
જેવી ભાષા પાસે ભાષાઓ વિશે ``વાત કરવાની''
અલગ વ્યવસ્થા ન હોવા છતાં અભિવ્યક્તિઓના
જટિલ ગુણધર્મોને ઔપચારિક બનાવી શકે છે,
એ મૂળે એક નોંધપાત્ર સમજ હતી. પછી એ
ભાષામાં રચાયેલો પિયાનો અંકગણિત જેવો
સિદ્ધાંત પોતાની ભાષા વિશે શું \emph{સાબિત}
કરી શકે—જેમ કે વાક્યો સાબિત
થઈ શકે કે સત્ય છે—એ પૂછવું સહેલું બને.
આથી ગડેલનાં (અસાબિતી વિશેનાં) અને
ટાર્સ્કીનાં (સત્યની અવ્યાખ્યેયતા વિશેનાં)
પ્રસિદ્ધ મર્યાદા-પ્રમેયો તરફ જઈએ છીએ.
પરંતુ વાક્યરચનાને અંકગણિતીકૃત કરવાની
યુક્તિ સંગણનીયતા-સિદ્ધાંતનાં મહત્વનાં
પરિણામો સાબિત કરવા માટે પણ ઉપયોગી છે;
દાખલા તરીકે, સિદ્ધાંતોની સંગણકીય શક્તિ
અથવા સંગણનીયતાનાં જુદાં નિદર્શો વચ્ચેનો
સંબંધ. વાક્યરચનાનું અંકગણિતીકરણ બીજા
પદાર્થો અને ગુણધર્મોને અંકગણિતીકૃત કરવાનો
નમૂનો આપે છે. દાખલા તરીકે, સંરૂપણો
અને સંગણનાઓ—માનો ટ્યુરિંગ મશીનની—પણ
એ રીતે અંકગણિતીકૃત કરી શકાય. તેથી
એક નિદર્શની સંગણનાને, જેમ કે ટ્યુરિંગ
મશીનની, બીજા નિદર્શમાં, જેમ કે પુનરાવર્તી
વિધેયોમાં, અનુકરી શકાય છે.
% END OLP-0281
\endgroup


\begingroup
\makeatletter\def\ol@sectioncs{section}\makeatother

% BEGIN OLP-0282
\olfileid{inc}{art}{cod}
\olsection{ચિહ્નોનું સંકેતબદ્ધીકરણ}
\phantomsection\label{guunit:OLP-0282}


પ્રથમ-ક્રમના તર્કશાસ્ત્રની મૂળભૂત ભાષા~$\Lang L$
નીચેના ચિહ્નો વાપરે છે:
\[
\lfalse \quad \lnot \quad \lor \quad \land \quad \lif \quad \lforall
\quad \lexists \quad \eq \quad ( \quad ) \quad ,
\]
આ સાથે ચલો અને અચળોના
ગણનીય ગણો તથા મનસ્વી સ્થાનિકતાવાળા
ફલનો અને વિધેયોના
ગણનીય ગણો હોય છે. આપણે દરેક
ચિહ્નને \emph{સંકેત નંબર} એવી રીતે આપી
શકીએ કે દરેક ચિહ્નને અનન્ય નંબર મળે અને
બે જુદાં ચિહ્નોને ક્યારેય એ જ નંબર ન મળે.
આ શક્ય છે કારણ કે બધા ચિહ્નોનો ગણ
ગણનીય છે અને તેથી તેની તથા
પ્રાકૃતિક સંખ્યાઓના ગણની વચ્ચે
એક-એક વ્યાપ્ત વિધેય છે. પરંતુ સંકેત નંબરમાંથી
ચિહ્ન તથા તેના વિશેની કેટલીક માહિતી
(દાખલા તરીકે, ફલનની સ્થાનિકતા)
સંગણનીય રીતે પાછી મેળવી શકાય તેની
ખાતરી પણ જોઈએ. એ કરવાની ઘણી રીત છે.
અહીં આદિમ પુનરાવર્તી વિધેયો વાપરતી
એક રીત આપી છે. (યાદ કરો કે
$\tuple{n_0, \dots, n_k}$ એ સંખ્યાઓ
$n_0$, \dots, $n_k$ના અનુક્રમનો સંકેત
નંબર છે.)

\begin{defn}
$s$ એ~$\Lang L$નું ચિહ્ન હોય તો
તેનો \emph{ચિહ્ન-સંકેત}~$\scode s$
આ પ્રમાણે વ્યાખ્યાયિત કરીએ:
\begin{enumerate}
\item $s$ તાર્કિક ચિહ્નોમાંનું હોય તો
  $\scode s$ નીચેના કોષ્ટક મુજબ છે:
\[
\begin{array}{cccccccccc}
  \lfalse & \lnot & \lor & \land & \lif & \lforall \\
\tuple{0, 0} & \tuple{0, 1} & \tuple{0, 2} & \tuple{0, 3} &
\tuple{0, 4} & \tuple{0, 5} \\
\lexists & \eq & ( & ) & ,\\
\tuple{0, 6} & \tuple{0, 7} &
\tuple{0, 8} & \tuple{0, 9} & \tuple{0, 10}
\end{array}
\]
\item $s$ એ~$i$મો ચલ~$\Obj v_i$ હોય, તો $\scode s = \tuple{1, i}$.
\item $s$ એ~$i$મો અચળ~$\Obj c_i$ હોય, તો
  $\scode s = \tuple{2, i}$.
\item $s$ એ~$i$મો $n$-સ્થાનિક ફલન~$\Obj f_i^n$ હોય, તો
  $\scode s = \tuple{3, n, i}$.
\item $s$ એ~$i$મો $n$-સ્થાનિક વિધેય~$\Obj P_i^n$ હોય, તો
  $\scode s = \tuple{4, n, i}$.
\end{enumerate}
\end{defn}

\begin{prop}
નીચેના સંબંધો આદિમ પુનરાવર્તી છે:
\begin{enumerate}
\item $\fn{Fn}(x, n)$ ત્યારે અને માત્ર ત્યારે,
  જ્યારે કોઈ~$i$ માટે~$x$ એ~$\Obj f^n_i$નો
  સંકેત નંબર હોય; એટલે~$x$ એ $n$-સ્થાનિક
  ફલનનો સંકેત નંબર હોય.
\item $\fn{Pred}(x, n)$ ત્યારે અને માત્ર ત્યારે,
  જ્યારે કોઈ~$i$ માટે~$x$ એ~$\Obj P^n_i$નો
  સંકેત નંબર હોય અથવા~$x$ એ~$\eq$નો
  સંકેત નંબર અને~$n = 2$ હોય; એટલે~$x$
  એ $n$-સ્થાનિક વિધેયનો સંકેત
  નંબર હોય.
\end{enumerate}
\end{prop}

\begin{defn}
$s_0, \dots, s_{n-1}$ ચિહ્નોનો અનુક્રમ
હોય, તો તેની \emph{ગડેલ-સંખ્યા}
$\tuple{\scode{s_0}, \dots, \scode{s_{n-1}}}$ છે.
\end{defn}

\begin{explain}
\emph{સંકેત નંબરો} અને \emph{ગડેલ-સંખ્યાઓ}
અલગ વસ્તુઓ છે. દાખલા તરીકે, ચલ~$\Obj v_5$
નો સંકેત નંબર~$\scode{\Obj
  v_5} = \tuple{1, 5} = 2^2\cdot 3^6$ છે.
પરંતુ પદ તરીકે વિચારેલું ચલ~$\Obj v_5$
લંબાઈ~$1$ ધરાવતો ચિહ્નોનો અનુક્રમ પણ
છે. \emph{પદ}~$\Obj v_5$ની
\emph{ગડેલ-સંખ્યા}~$\Gn{\Obj v_5}$ એ
$\tuple{\scode{\Obj v_5}} = 2^{\scode{\Obj
    v_5} + 1} = 2^{2^2\cdot 3^6 + 1}$ છે.
\end{explain}

\begin{ex}
યાદ કરો કે $k_0$, \dots, $k_{n-1}$
સંખ્યાઓનો અનુક્રમ હોય, તો અવિભાજ્ય
સંખ્યાઓની ઘાતો વડે થતા સંકેતબદ્ધીકરણમાં
અનુક્રમ~$\tuple{k_0, \dots, k_{n-1}}$નો
સંકેત નંબર આ છે:
\[
2^{k_0+1}\cdot3^{k_1+1}\cdot \dots \cdot p_{n-1}^{k_{n-1}+1},
\]
અહીં~$p_i$ એ~$i$મી અવિભાજ્ય સંખ્યા છે
(શરૂઆત~$p_0 = 2$થી થાય છે). તેથી,
દાખલા તરીકે, સૂત્ર~$\eq[\Obj v_0][\Obj 0]$
અથવા વધુ સ્પષ્ટ રીતે~${\eq}(\Obj v_0,\Obj c_0)$ની
ગડેલ-સંખ્યા છે:
\[
\tuple{\scode{\eq},\scode{(},\scode{\Obj v_0},\scode{,}, \scode{\Obj
    c_0},\scode{)}}.
\]
અહીં $\scode{\eq}$ એ~$\tuple{0,7} = 2^{0+1}\cdot
3^{7+1}$ અને $\scode{\Obj v_0}$ એ~$\tuple{1,0} = 2^{1+1}\cdot3^{0+1}$,
વગેરે. તેથી~$\Gn{=(\Obj v_0,\Obj c_0)}$ એ
\begin{multline*}
2^{\scode{=} + 1}\cdot 3^{\scode{(}+1}\cdot 5^{\scode{\Obj v_0}+1}
\cdot 7^{\scode{,} + 1} \cdot 11^{\scode{\Obj c_0}+1} \cdot
13^{\scode{)}+1} = \\
2^{2^1\cdot 3^8 + 1}\cdot 3^{2^1\cdot 3^9+1}\cdot 5^{2^2\cdot 3^1+1}
\cdot 7^{2^1\cdot 3^{11} + 1} \cdot 11^{2^3\cdot3^1+1} \cdot
13^{2^1\cdot3^{10}+1} = \\
2^{13\,123}\cdot 3^{39\,367}\cdot 5^{13}\cdot 7^{354\,295}\cdot11^{25}\cdot13^{118\,099}.
\end{multline*}
\end{ex}
% END OLP-0282
\endgroup


\begingroup
\makeatletter\def\ol@sectioncs{section}\makeatother

% BEGIN OLP-0283
\olfileid{inc}{art}{trm}
\olsection{પદોનું સંકેતબદ્ધીકરણ}
\phantomsection\label{guunit:OLP-0283}


\begin{explain}
પદ એ ચિહ્નોના અનુક્રમનો એક ખાસ પ્રકાર
જ છે: પદોની રચનાના નિયમો મુજબ તે
સ્થિરો અને ચલોમાંથી આગમનાત્મક રીતે
બને છે. ચિહ્નોને સંકેતબદ્ધ કરવાની યોજના
અને સંખ્યાઓના અનુક્રમોને સંકેતબદ્ધ કરવાની
રીત વડે ચિહ્નોના અનુક્રમોને સંખ્યાઓ તરીકે
સંકેતબદ્ધ કરી શકાય છે. તેથી પદોને
ગડેલ-સંખ્યાઓ આપવી મુશ્કેલ નથી.
ખરો પડકાર તો એ બતાવવાનો છે કે
યોગ્ય રીતે રચાયેલા પદની ગડેલ-સંખ્યા
હોવાનો ગુણધર્મ સંગણનીય, ખરેખર આદિમ
પુનરાવર્તી છે.
\end{explain}

ચલો અને અચળો સૌથી
સરળ પદો છે. $x$ આવા પદની ગડેલ-સંખ્યા
છે કે નહીં તે તપાસવું સહેલું છે:
કોઈ~$i$ માટે~$x$ એ~$\Gn{\Obj v_i}$
હોય ત્યારે $\fn{Var}(x)$ છે. બીજા
શબ્દોમાં, $x$ એ લંબાઈ~$1$નો અનુક્રમ છે
અને તેનો એકમાત્ર ઘટક~$(x)_0$ એ
કોઈ ચલ~$\Obj v_i$નો સંકેત
નંબર છે; એટલે કોઈ~$i$ માટે~$x$
એ~$\tuple{\tuple{1, i}}$ છે. એ જ રીતે,
કોઈ~$i$ માટે~$x$ એ~$\Gn{\Obj c_i}$
હોય ત્યારે~$\fn{Const}(x)$ છે. આ
બંને સંબંધો આદિમ પુનરાવર્તી છે, કારણ કે
આવો~$i$ હોય તો તે~$< x$ હોવો જ
જોઈએ:
\begin{align*}
  \fn{Var}(x) & \defiff \bexists{i<x}{x = \tuple{\tuple{1, i}}}\\
  \fn{Const}(x) & \defiff \bexists{i<x}{x = \tuple{\tuple{2, i}}}
\end{align*}

\begin{prop}
\ollabel{prop:term-primrec}
$x$ અનુક્રમે પદ અથવા બંધ પદની
ગડેલ-સંખ્યા હોય ત્યારે અને માત્ર ત્યારે
સત્ય થતા સંબંધો~$\fn{Term}(x)$ અને
$\fn{ClTerm}(x)$ આદિમ પુનરાવર્તી છે.
\end{prop}

\begin{proof}
ચિહ્નોનો અનુક્રમ~$s$ પદ છે ત્યારે અને
માત્ર ત્યારે, જ્યારે પદોનો એવો અનુક્રમ
$s_0$, \dots, $s_{k-1} = s$ હોય જે
પદોના રચના-નિયમો મુજબ~$s$ અચળો અને
ચલોમાંથી કેવી રીતે બન્યું તે નોંધે છે.
આવી સંભવિત રચના-શ્રેણી રચના-નિયમો
અનુસરે એ માટે દરેક~$i < k$ માટે
નીચેનામાંથી એક હોવું જોઈએ:
\begin{enumerate}
\item $s_i$ કોઈ ચલ~$\Obj v_j$ છે, અથવા
\item $s_i$ કોઈ અચળ~$\Obj c_j$ છે, અથવા
\item $s_i$ એ સ્થાન~$i$ પહેલાં આવતાં
  $n$ પદો~$t_1$, \dots, $t_n$માંથી
  $n$-સ્થાનિક ફલન~$\Obj f^n_j$
  વાપરીને બનાવેલું છે.
\end{enumerate}
ગડેલ-સંખ્યાઓ પરનો અનુરૂપ સંબંધ આદિમ
પુનરાવર્તી છે તે બતાવવા આ શરતને
આદિમ પુનરાવર્તી રીતે—એટલે આદિમ
પુનરાવર્તી વિધેયો, સંબંધો અને સીમિત
પરિમાણન વડે—વ્યક્ત કરવી પડશે.

માનો કે~$y$ એ અનુક્રમ~$s_0$, \dots,
$s_{k-1}$નો સંકેત નંબર છે; એટલે
$y = \tuple{\Gn{s_0}, \dots, \Gn{s_{k-1}}}$.
દરેક~$i < k$ માટે નીચેની શરતો
સંતોષાય ત્યારે અને માત્ર ત્યારે તે
ગડેલ-સંખ્યા~$x$ ધરાવતા પદનો
રચના-શ્રેણી સંકેતબદ્ધ કરે છે:
\begin{enumerate}
\item $\fn{Var}((y)_i)$, અથવા
\item $\fn{Const}((y)_i)$, અથવા
\item કોઈ~$n$ અને
  નંબર~$z = \tuple{z_1, \dots, z_n}$
  હોય, જ્યાં દરેક~$z_l$ કોઈ~$i' < i$
  માટે~$(y)_{i'}$ જેટલો હોય, અને
\[
(y)_i = \Gn{\Obj f^n_j(} \concat \fn{flatten}(z) \concat \Gn{)},
\]
\end{enumerate}
અને વધુમાં~$(y)_{k-1} = x$ હોય.
(વિધેય~$\fn{flatten}(z)$ અનુક્રમ
$\tuple{\Gn{t_1}, \dots, \Gn{t_n}}$ને
$\Gn{t_1, \dots,
  t_n}$માં ફેરવે છે અને આદિમ પુનરાવર્તી છે.)

કિસ્સા~(3)માં સૂચકાંકો~$j$, $n$, પદો~$t_l$ની
ગડેલ-સંખ્યાઓ~$z_l$ અને
અનુક્રમ~$\tuple{z_1, \dots, z_n}$નો
સંકેત નંબર~$z$—આ બધાં~$y$ કરતાં
નાનાં છે. ઉપરનું~$k$ આપણે~$\len{y}$થી
બદલી શકીએ. તેથી ``$y$ એ ગડેલ-સંખ્યા~$x$
ધરાવતા પદની રચના-શ્રેણીનો સંકેત
નંબર છે'' એ શરતને અનુરૂપ સંબંધ
આદિમ પુનરાવર્તી છે તે સ્પષ્ટ થાય
એવી રીતે વ્યક્ત કરી શકીએ છીએ.

હવે~$y$ માટે આદિમ પુનરાવર્તી મર્યાદા
છે એટલું બતાવવાનું રહે છે. $x$ પદની
ગડેલ-સંખ્યા હોય, તો વધુમાં વધુ~$\len{x}$
પદોની રચના-શ્રેણી મળી શકે છે
(કારણ કે~$s$ની રચના-શ્રેણીનું દરેક પદ
$s$માં કોઈ સ્થાનથી શરૂ થવું જોઈએ,
અને બે ઉપપદો એક જ સ્થાનથી શરૂ
થઈ શકતાં નથી). $s$ના દરેક ઉપપદની
ગડેલ-સંખ્યા નિશ્ચિતપણે~$\le x$ છે.
તેથી હંમેશાં એવી રચના-શ્રેણી છે જેનો
સંકેત નંબર~$\le p_{k-1}^{k(x+1)}$ છે,
જ્યાં~$k=\len{x}$.

$\fn{ClTerm}$ માટે માત્ર
ચલોવાળી શરત કાઢી નાખો.
\end{proof}

\begin{prob}
અનુક્રમ~$\tuple{\Gn{t_1}, \dots, \Gn{t_n}}$ને
$\Gn{t_1, \dots, t_n}$માં ફેરવતું
વિધેય~$\fn{flatten}(z)$ આદિમ
પુનરાવર્તી છે તે બતાવો.
\end{prob}

\begin{prop}\ollabel{prop:num-primrec}
  વિધેય~$\fn{num}(n) = \Gn{\num{n}}$ આદિમ પુનરાવર્તી છે.
\end{prop}

\begin{proof}
  $\fn{num}(n)$ને આદિમ પુનરાવર્તનથી વ્યાખ્યાયિત કરીએ:
  \begin{align*}
    \fn{num}(0) & = \Gn{\Obj 0}\\
    \fn{num}(n+1) & = \Gn{\prime(} \concat \fn{num}(n) \concat \Gn{)}.
  \end{align*}
\end{proof}
% END OLP-0283
\endgroup


\begingroup
\makeatletter\def\ol@sectioncs{section}\makeatother

% BEGIN OLP-0284
\olfileid{inc}{art}{frm}
\olsection{સૂત્રોનું સંકેતબદ્ધીકરણ}
\phantomsection\label{guunit:OLP-0284}


સંબંધ~$\fn{Term}(x)$ને આદિમ
પુનરાવર્તી રીતે વ્યાખ્યાયિત કર્યા પછી,
તેની મદદથી સૂત્રો માટેનો અનુરૂપ
સંબંધ~$\fn{Frm}(x)$ પણ આદિમ
પુનરાવર્તી રીતે વ્યાખ્યાયિત કરી શકીએ.

\begin{prop}
$x$ આણ્વિક સૂત્રની ગડેલ-સંખ્યા
હોય ત્યારે અને માત્ર ત્યારે સત્ય થતો
સંબંધ~$\fn{Atom}(x)$ આદિમ પુનરાવર્તી છે.
\end{prop}

\begin{proof}
નીચેમાંથી એક શરત હોય ત્યારે અને માત્ર
ત્યારે નંબર~$x$ આણ્વિક સૂત્રની
ગડેલ-સંખ્યા છે:
\begin{enumerate}
\item એવાં~$n$, $j < x$ અને~$z < x$
  હોય કે દરેક~$i < n$ માટે
  $\fn{Term}((z)_i)$ અને $x = $
\[
\Gn{\Obj P^n_j(} \concat \fn{flatten}(z) \concat \Gn{)}.
\]
\item એવાં~$z_1, z_2 < x$ હોય કે
  $\fn{Term}(z_1)$, $\fn{Term}(z_2)$
  અને $x = $
\[
\Gn{{\eq}(} \concat z_1 \concat \Gn{,} \concat z_2 \concat{\Gn{)}}.
\]
  \item $x = \Gn{\lfalse}$.
  \item $x = \Gn{\ltrue}$.
\end{enumerate}
\end{proof}

\begin{prop}
\ollabel{prop:frm-primrec}
$x$ કોઈ સૂત્રની ગડેલ-સંખ્યા
હોય ત્યારે અને માત્ર ત્યારે સત્ય થતો
સંબંધ~$\fn{Frm}(x)$ આદિમ પુનરાવર્તી છે.
\end{prop}

\begin{proof}
ચિહ્નોનો અનુક્રમ~$s$ સૂત્ર છે
ત્યારે અને માત્ર ત્યારે, જ્યારે
સૂત્રની એવી રચના-શ્રેણી~$s_0$,
\dots, $s_{k-1} = s$ હોય જે રચના-નિયમો
મુજબ આણ્વિક સૂત્રોમાંથી~$s$
કેવી રીતે બન્યું તે નોંધે છે. બિનજરૂરી
વચગાળાનાં સૂત્રો કાઢીને માત્ર ઉપસૂત્રોની
રચના-શ્રેણી પસંદ કરીએ, તો દરેક~$s_i$નો
સંકેત નંબર અંતિમ~$s$ના સંકેત નંબર~$x$થી
મોટો નથી. સમગ્ર રચના-શ્રેણી
$\tuple{s_0, \dots, s_{k-1}}$ના સંકેત
નંબર માટે પણ અંતિમ નંબરથી ગણાતી
આદિમ પુનરાવર્તી મર્યાદા મળે છે.
\sourcecorrection{OLINC-005}{સ્થિર અંગ્રેજી મૂળમાં
અંતિમ સભ્ય સહિત દરેક~$s_i$નો સંકેત
નંબર અને આખી રચના-શ્રેણીનો સંકેત
નંબર બંને~$x$ કરતાં નાના કહેવાયા છે.
અંતિમ સભ્યનો સંકેત પોતે જ~$x$ છે,
અને આખી શ્રેણીનો સંકેત તેનાથી મોટો
થઈ શકે છે. બિનજરૂરી વચગાળાનાં સૂત્રો
કાઢીને ઉપસૂત્રો પૂરતી રચના-શ્રેણી લેવી
પડે. પહેલાંના પદ-પ્રમાણની જેમ
અહીં જરૂરી વાત~$x$ પરથી ગણાતી આદિમ
પુનરાવર્તી મર્યાદા હોવાની છે.}
\end{proof}

\begin{prob}
પહેલી સાબિતીની રીત અનુસરીને
\olref[inc][art][frm]{prop:frm-primrec}ની
વિગતવાર સાબિતી આપો;
\olref[inc][art][trm]{prop:term-primrec}ની
પહેલી સાબિતી માર્ગદર્શક છે.
\end{prob}

\begin{prop}
\ollabel{prop:freeocc-primrec}
ગડેલ-સંખ્યા~$x$ ધરાવતા સૂત્રનું
$i$મું ચિહ્ન એ ગડેલ-સંખ્યા~$z$
ધરાવતા ચલની મુક્ત ઘટના હોય ત્યારે
અને માત્ર ત્યારે સત્ય થતો સંબંધ
$\fn{FreeOcc}(x, z, i)$ આદિમ
પુનરાવર્તી છે.
\end{prop}

\begin{proof}
અભ્યાસપ્રશ્ન.
\end{proof}

\begin{prob}
\olref[inc][art][frm]{prop:freeocc-primrec}
સાબિત કરો. સૂત્રનો કોઈ
ઉપઅનુક્રમ પણ સૂત્ર હોય, તો
તે મૂળ સૂત્રનું ઉપ-સૂત્ર હોય છે,
એ તથ્ય વાપરી શકો છો.
\end{prob}

\begin{prop}
$x$ કોઈ વાક્યની ગડેલ-સંખ્યા
હોય ત્યારે અને માત્ર ત્યારે સત્ય થતો
ગુણધર્મ~$\fn{Sent}(x)$ આદિમ
પુનરાવર્તી છે.
\end{prop}

\begin{proof}
વાક્ય એ ચલોની મુક્ત
ઘટના વિનાનું સૂત્ર છે. તેથી
$\fn{Sent}(x)$ ત્યારે અને માત્ર ત્યારે,
જ્યારે
\begin{multline*}
\bforall{i<\len{x}}{\bforall{z<x}{}}\\
(\bexists{j<z}{z=\Gn{\Obj v_j}} \lif \lnot\fn{FreeOcc}(x,z,i)).
\end{multline*}
\end{proof}
% END OLP-0284
\endgroup


\begingroup
\makeatletter\def\ol@sectioncs{section}\makeatother

% BEGIN OLP-0285
\olfileid{inc}{art}{sub}
\olsection{પ્રતિસ્થાપન}
\phantomsection\label{guunit:OLP-0285}


યાદ કરો કે પ્રતિસ્થાપન એ
સૂત્ર~$!A$માં ચલ~$u$ની
બધી મુક્ત ઘટનાઓની જગ્યાએ પદ~$t$
મૂકવાની ક્રિયા છે; તેને~$\Subst{!A}{t}{u}$
લખીએ છીએ. ચલો, સૂત્રો
અને પદોની ગડેલ-સંખ્યાઓ પર કરવામાં આવે
ત્યારે આ ક્રિયા આદિમ પુનરાવર્તી છે.

\begin{prop}\ollabel{prop:subst-primrec}
એવું આદિમ પુનરાવર્તી વિધેય
$\fn{Subst}(x, y, z)$ છે કે
\[
\fn{Subst}(\Gn{!A}, \Gn{t}, \Gn{u}) = \Gn{\Subst{!A}{t}{u}}.
\]
\end{prop}

\begin{proof}
આદિમ પુનરાવર્તન વડે વિધેય~$\fn{hSubst}$
આ પ્રમાણે વ્યાખ્યાયિત કરી શકીએ:
\begin{multline*}
\begin{aligned}
\fn{hSubst}(x, y, z, 0) & = \emptyseq \\
\fn{hSubst}(x, y, z, i+1) & =
\end{aligned}\\
\begin{cases}
\fn{hSubst}(x, y, z, i) \concat y & \text{જો $\fn{FreeOcc}(x, z, i)$ હોય} \\
\fn{append}(\fn{hSubst}(x, y, z, i), (x)_{i}) & \text{અન્યથા.}
\end{cases}
\end{multline*}
હવે~$\fn{Subst}(x, y, z)$ને
$\fn{hSubst}(x, y, z, \len{x})$ તરીકે
વ્યાખ્યાયિત કરી શકીએ.
\end{proof}

\begin{prop}
\ollabel{prop:free-for}
ગડેલ-સંખ્યા~$x$ ધરાવતા સૂત્રમાં
ગડેલ-સંખ્યા~$y$ ધરાવતું પદ,
ગડેલ-સંખ્યા~$z$ ધરાવતા ચલની જગ્યાએ
મુક્ત હોય ત્યારે અને માત્ર ત્યારે
સત્ય થતો સંબંધ~$\fn{FreeFor}(x, y, z)$
આદિમ પુનરાવર્તી છે.
\end{prop}

\begin{proof} અભ્યાસપ્રશ્ન. \end{proof}

\begin{prob}
\olref[inc][art][sub]{prop:free-for} સાબિત કરો.
\end{prob}
% END OLP-0285
\endgroup


%
  \begingroup
\makeatletter\def\ol@sectioncs{section}\makeatother

% BEGIN OLP-0286
\olfileid{inc}{art}{plk}
\olsection{$\Log{LK}$માં નિષ્પત્તિઓ}
\phantomsection\label{guunit:OLP-0286}


\begin{explain}
નિષ્પત્તિઓનું અંકગણિતીકરણ કરવા માટે
આપણે નિષ્પત્તિઓને સંખ્યાઓ વડે નિરૂપવાં પડશે.
દરેક નિષ્પત્તિ એ સિક્વન્ટોનું વૃક્ષ છે
અને તેમાં દરેક અનુમાનને એક ચિહ્નન પણ
આપેલું છે. તેથી પુનરાવર્તી નિરૂપણ
સૌથી સ્વાભાવિક છે: નિષ્પત્તિને
એવી ટપલ તરીકે નિરૂપીએ જેના ઘટકો
અંતિમ સિક્વન્ટ, ચિહ્નન અને છેલ્લા અનુમાનના
આધાર-સિક્વન્ટો સુધી પહોંચતી
ઉપ-નિષ્પત્તિઓનાં નિરૂપણ છે.
\end{explain}

\begin{defn}
$\Gamma$ એ વાક્યોની સાન્ત શ્રેણી,
$\Gamma = \tuple{!A_1, \dots, !A_n}$ હોય,
તો $\Gn{\Gamma} = \tuple{\Gn{!A_1},
  \dots, \Gn{!A_n}}$.

$\Gamma \Sequent \Delta$ સિક્વન્ટ હોય,
તો $\Gamma \Sequent \Delta$ની
ગડેલ-સંખ્યા આ છે:
\[
\Gn{\Gamma \Sequent \Delta} = \tuple{\Gn{\Gamma}, \Gn{\Delta}}
\]

$\pi$ એ $\Log{LK}$માં નિષ્પત્તિ હોય,
તો $\Gn{\pi}$ને આ પ્રમાણે વ્યાખ્યાયિત કરીએ:
\begin{enumerate}
\item $\pi$ માત્ર આરંભિક સિક્વન્ટ
  $\Gamma \Sequent \Delta$નું બનેલું હોય,
  તો $\Gn{\pi}$ આ છે:
  \[
    \tuple{0, \Gn{\Gamma \Sequent \Delta}}.
  \]
\item $\pi$નો અંત એક કે બે આધાર-સિક્વન્ટવાળા
  અનુમાનથી થતો હોય, તેનું ફલિત-સિક્વન્ટ
  $\Gamma \Sequent \Delta$ હોય અને
  $\pi_1$ તથા, લાગુ પડે ત્યારે, $\pi_2$
  છેલ્લા અનુમાનના આધાર-સિક્વન્ટોમાં
  પૂરાં થતાં સીધાં ઉપપ્રમાણો હોય,
  તો અનુક્રમે $\Gn{\pi}$ આ છે:
  \begin{align*}
    & \tuple{1, \Gn{\pi_1}, \Gn{\Gamma \Sequent \Delta}, k} \text{ અથવા}\\ 
    & \tuple{2, \Gn{\pi_1}, \Gn{\pi_2}, \Gn{\Gamma \Sequent \Delta}, k}, 
  \end{align*}
  અહીં છેલ્લા અનુમાનમાં વપરાયેલા નિયમ
  પ્રમાણે નીચેના કોષ્ટકમાંથી~$k$ મળે છે:

  \begin{tabular}{lccccccc}
    \text{નિયમ:} & \LeftR{\Weakening} & \RightR{\Weakening} &
    \LeftR{\Contraction} & \RightR{\Contraction} &
    \LeftR{\Exchange} & \RightR{\Exchange} \\
    $k$: & 1 & 2 & 3 & 4 & 5 & 6 \\[2ex]
    \text{નિયમ:} & \LeftR{\lnot} & \RightR{\lnot} &
    \LeftR{\land} & \RightR{\land} & 
    \LeftR{\lor} & \RightR{\lor} \\
    $k$: & 7 & 8 & 9 & 10 & 11 & 12 \\[2ex]
    \text{નિયમ:} & \LeftR{\lif} & \RightR{\lif} &
    \LeftR{\lforall} & \RightR{\lforall} &
    \LeftR{\lexists} & \RightR{\lexists} \\
    $k$: & 13 & 14 & 15 & 16 & 17 & 18 \\[2ex]
    \text{નિયમ:} & \Cut & = \\
    $k$: & 19 & 20
  \end{tabular}
\end{enumerate}
\end{defn}

\begin{ex}
  આ અત્યંત સરળ નિષ્પત્તિ જુઓ:
  \begin{prooftree}
    \Axiom$!A \fCenter !A$
    \RightLabel{\LeftR{\land}}
    \UnaryInf$!A \land !B \fCenter !A$
    \RightLabel{\RightR{\lif}}
    \UnaryInf$\fCenter (!A \land !B) \lif !A$
  \end{prooftree}
  આરંભિક સિક્વન્ટની ગડેલ-સંખ્યા
  $p_0 = \tuple{0,
    \Gn{!A \Sequent !A}}$ થાય.
  $\LeftR{\land}$ના ફલિત-સિક્વન્ટમાં
  પૂરી થતી નિષ્પત્તિની ગડેલ-સંખ્યા
  $p_1 =
    \tuple{1, p_0, \Gn{!A \land !B \Sequent !A}, 9}$
  થાય ($1$ એટલા માટે કે
  $\LeftR{\land}$નો એક આધાર-સિક્વન્ટ છે;
  $!A \land !B \Sequent !A$ના
  ફલિત-સિક્વન્ટની ગડેલ-સંખ્યા લેવાઈ છે;
  અને $9$ એ $\LeftR{\land}$નો સંકેત
  નંબર છે). આખી નિષ્પત્તિની
  ગડેલ-સંખ્યા તેથી
  $\tuple{1, p_1, \Gn{\Sequent (!A \land !B) \lif !A}, 14}$,
  એટલે કે
  \[
  \tuple{1, \tuple{1, \tuple{0, \Gn{!A \Sequent !A}}, \Gn{!A \land !B \Sequent !A}, 9},
    \Gn{\Sequent (!A \land !B) \lif !A}, 14}.
  \]
  \sourcecorrection{OLINC-006}{સ્થિર અંગ્રેજી મૂળમાં
  આ ઉદાહરણની ટપલમાં બે જગ્યાએ
  સિક્વન્ટની અંદર વધારાનું બંધ કૌંસ છે.
  અહીં એ વધારાનાં કૌંસ કાઢ્યાં છે;
  નિયમ-સંકેત~$14$ અને ટપલના બાકીના
  ઘટકો યથાવત્ છે.}
\end{ex}

\begin{explain}
નિષ્પત્તિઓનું નિરૂપણ નક્કી કર્યા પછી,
આવા નિષ્પત્તિઓ પર આદિમ પુનરાવર્તી રીતે ક્રિયા
કરી શકાય છે અને તેમના આવશ્યક ગુણધર્મો
તથા સંબંધો એ જ રીતે વ્યક્ત કરી શકાય છે
તે પણ બતાવવું પડશે. કેટલીક ક્રિયાઓ
સરળ છે. દાખલા તરીકે, નિષ્પત્તિની
ગડેલ-સંખ્યા~$p$ હોય, તો
$\fn{EndSequent}(p) = (p)_{(p)_0+1}$
તેના અંતિમ સિક્વન્ટની ગડેલ-સંખ્યા આપે છે.
આરંભિક સિક્વન્ટ સિવાયના કિસ્સામાં
$\fn{LastRule}(p) = (p)_{(p)_0+2}$
છેલ્લા નિયમનો સંકેત આપે છે.
આરંભિક સિક્વન્ટના કિસ્સામાં આ વિધેયને
ઇચ્છીએ તો શૂન્ય મૂલ્ય આપીને સર્વત્ર
વ્યાખ્યાયિત કરી શકીએ.
\sourcecorrection{OLINC-007}{સ્થિર અંગ્રેજી મૂળ
અહીં અંતિમ સિક્વન્ટના વિધેયને
$\fn{EndSeq}$ કહે છે, પણ આગળ
$\fn{EndSequent}$ વાપરે છે; અહીં
બીજું નામ સતત રાખ્યું છે. આરંભિક
સિક્વન્ટવાળી બે-ઘટક ટપલ માટે
છેલ્લા નિયમનો દર્શાવેલો ઘટક નથી,
એથી એ વિધેયનું દર્શાવેલું સમીકરણ
માત્ર બિનઆરંભિક કિસ્સાને લાગુ પડે છે;
આરંભિક કિસ્સામાં~$0$ મૂકવાથી
પછીનું વિધાન સર્વત્ર અર્થપૂર્ણ બને છે.}
નીચેના વિધાનથી વ્યાખ્યાયિત
$\fn{Sequent}(s)$ ગુણધર્મ જુઓ:
\[
  \len{s} = 2 \land \bforall{i<\len{(s)_0} + \len{(s)_1}}{\fn{Sent}(((s)_0 \concat (s)_1)_i)}
\]
તે $s$ માટે ત્યારે અને માત્ર ત્યારે
સત્ય છે જ્યારે $s$ એ વાક્યોના
સિક્વન્ટની ગડેલ-સંખ્યા હોય. કેટલાક
કિસ્સા ઘણાં કઠિન છે. તેમને કેવી રીતે
સંભાળવા તેનો ઓછામાં ઓછો રૂપરેખાત્મક
ખ્યાલ આપીશું. હેતુ એ બતાવવાનો છે કે
``$\pi$ એ~$\Gamma$માંથી~$!A$ની
નિષ્પત્તિ છે'' એ સંબંધ
$\pi$ અને~$!A$ની ગડેલ-સંખ્યાઓનો
આદિમ પુનરાવર્તી સંબંધ છે.
\end{explain}

\begin{prop}\ollabel{prop:followsby}
  ગડેલ-સંખ્યા~$p$ ધરાવતી
  નિષ્પત્તિ~$\pi$નું છેલ્લું અનુમાન
  યોગ્ય હોય ત્યારે અને માત્ર ત્યારે સત્ય
  થતો ગુણધર્મ~$\fn{Correct}(p)$ આદિમ
  પુનરાવર્તી છે.
\end{prop}

\begin{proof}
  $\Gamma \Sequent \Delta$ આરંભિક સિક્વન્ટ
  છે, જો કાં તો કોઈ વાક્ય~$!A$
  માટે $\Gamma \Sequent \Delta$ એ
  $!A \Sequent !A$ હોય,
  અથવા કોઈ પદ~$t$ માટે
  $\Gamma \Sequent \Delta$ એ
  $\emptyset \Sequent \eq[t][t]$ હોય.
  ગડેલ-સંખ્યાઓની ભાષામાં
  $\fn{InitSeq}(s)$ ત્યારે અને માત્ર ત્યારે
  સત્ય છે જ્યારે
  \begin{align*}
  \bexists{x < s}{} (\fn{Sent}(x) & \land
  s = \tuple{\tuple{x},\tuple{x}})
  \lor {}\\
  \bexists{t<s}{} (\fn{Term}(t) & \land
  s = \tuple{0, \tuple{\Gn{{\eq}(} \concat t \concat \Gn{,} \concat t \concat \Gn{)}}}).
  \end{align*}

  દરેક અનુમાન-નિયમ~$R$ માટે
  $\fn{FollowsBy}_R(p)$ સંબંધ આદિમ
  પુનરાવર્તી છે તે પણ બતાવવું પડે.
  આ $\fn{FollowsBy}_R(p)$ સંબંધ ત્યારે
  અને માત્ર ત્યારે સત્ય છે જ્યારે $p$ એ
  નિષ્પત્તિ~$\pi$ની ગડેલ-સંખ્યા
  હોય અને~$\pi$નું અંતિમ સિક્વન્ટ
  $\pi$ના સીધા ઉપ-નિષ્પત્તિઓમાંથી
  $R$ના યોગ્ય ઉપયોગ વડે ફલિત થતું હોય.

  એક સરળ કિસ્સો \RightR{\land} નિયમનો
  છે. $\pi$નો અંત યોગ્ય
  $\RightR{\land}$ અનુમાનથી થતો હોય,
  તો તે આ પ્રમાણે દેખાય:
  \begin{prooftree}
    \AxiomC{}
    \RightLabel{$\pi_1$}
    \Deduce$\Gamma \fCenter \Delta, !A$
    
    \AxiomC{}
    \RightLabel{$\pi_2$}
    \Deduce$\Gamma \fCenter \Delta, !B$

    \RightLabel{\RightR\land}
    \BinaryInf$\Gamma \fCenter \Delta, !A \land !B$
  \end{prooftree}
  એટલે, નિષ્પત્તિ~$\pi$નું છેલ્લું
  અનુમાન $\RightR{\land}$નો યોગ્ય ઉપયોગ
  ત્યારે અને માત્ર ત્યારે છે જ્યારે
  વાક્યોની શ્રેણીઓ $\Gamma$ અને
  $\Delta$, તેમજ બે વાક્યો
  $!A$ અને~$!B$ એવાં હોય કે
  $\pi_1$નું અંતિમ સિક્વન્ટ
  $\Gamma \Sequent \Delta, !A$,
  $\pi_2$નું અંતિમ સિક્વન્ટ
  $\Gamma \Sequent \Delta, !B$ અને
  $\pi$નું અંતિમ સિક્વન્ટ
  $\Gamma \Sequent \Delta, !A \land !B$ હોય.
  આને ગડેલ-સંખ્યાઓમાં વ્યક્ત કરીએ.
  $s = \Gn{\Gamma \Sequent \Delta}$
  હોય, તો $(s)_0 = \Gn{\Gamma}$
  અને $(s)_1 = \Gn{\Delta}$.
  તેથી $\fn{FollowsBy}_{\RightR{\land}}(p)$
  ત્યારે અને માત્ર ત્યારે સત્ય છે જ્યારે
\begin{align*}
& \bexists{g < p}{\bexists{d < p}{\bexists{a < p}{\bexists{b < p}{\quad}}}} \\
& \qquad \fn{EndSequent}(p) = 
  \tuple{g, d \concat 
    \tuple{\Gn{(} \concat a \concat \Gn{\land} \concat b \concat \Gn{)}}} \land {} \\
& \qquad \fn{EndSequent}((p)_1) = \tuple{g, d \concat \tuple{a}} \land {}\\
& \qquad \fn{EndSequent}((p)_2) = \tuple{g, d \concat \tuple{b}} \land {}\\
& \qquad (p)_0 = 2 \land \fn{LastRule}(p) = 10.
\end{align*}
અનુક્રમે આ પંક્તિઓ કહે છે કે
``ગડેલ-સંખ્યા~$g$ ધરાવતી
શ્રેણી~($\Gamma$), ગડેલ-સંખ્યા~$d$
ધરાવતી શ્રેણી~($\Delta$),
ગડેલ-સંખ્યા~$a$ ધરાવતું
સૂત્ર~($!A$) અને
ગડેલ-સંખ્યા~$b$ ધરાવતું
સૂત્ર~($!B$) અસ્તિત્વમાં છે'';
``$\pi$નું અંતિમ સિક્વન્ટ
$\Gamma \Sequent \Delta, !A \land !B$
છે''; ``$\pi_1$નું અંતિમ સિક્વન્ટ
$\Gamma \Sequent \Delta, !A$ છે'';
``$\pi_2$નું અંતિમ સિક્વન્ટ
$\Gamma \Sequent \Delta, !B$ છે'';
અને ``$\pi$નાં બે સીધાં
ઉપપ્રમાણો છે તથા છેલ્લો અનુમાન-નિયમ
$\RightR\land$ (નંબર~$10$) છે.''

$\pi$નું છેલ્લું અનુમાન
$\RightR\lexists$નો યોગ્ય ઉપયોગ હોય,
તો અને માત્ર તો જ સૂત્ર~$!A$,
ચલ~$x$, પદ~$t$ અને શ્રેણીઓ
$\Gamma$, $\Delta$ એવાં હોય કે
$\pi$નું અંતિમ સિક્વન્ટ
$\Gamma \Sequent \Delta, \lexists[x][!A]$
અને $\pi_1$નું અંતિમ સિક્વન્ટ
$\Gamma \Sequent \Delta,
\Subst{!A}{t}{x}$ હોય. તેથી
ગડેલ-સંખ્યાઓમાં
$\fn{FollowsBy}_{\RightR\lexists}(p)$
ત્યારે અને માત્ર ત્યારે સત્ય છે જ્યારે
\begin{align*}
  & \bexists{g<p}{\bexists{d<p}{\bexists{a<p}{\bexists{x<p}{\bexists{t<p}{\quad}}}}}\\
  & \qquad \fn{EndSequent}(p) = 
  \tuple{
    g,
    d \concat \tuple{\Gn{\lexists} \concat x \concat a}
  } \land {}\\
  & \qquad \fn{EndSequent}((p)_1) = 
  \tuple{
    g,
    d \concat \tuple{\fn{Subst}(a, t, x)}
  } \land {}\\
  & \qquad (p)_0 = 1 \land \fn{LastRule}(p) = 18.
\end{align*}
હવે $\fn{Correct}(p)$ને આ પ્રમાણે
વ્યાખ્યાયિત કરીએ:
\begin{multline*}
  \fn{Sequent}(\fn{EndSequent}(p)) \land {}\\ 
  [(\fn{LastRule}(p) = 1 \land 
  \fn{FollowsBy}_{\LeftR\Weakening}(p)) \lor \dots \lor {}\\
  (\fn{LastRule}(p) = 20 \land \fn{FollowsBy}_{\eq}(p)) \lor {}\\
  (p)_0 = 0 \land \fn{InitSeq}(\fn{EndSequent}(p))]
\end{multline*}
\sourcecorrection{OLINC-008}{સ્થિર અંગ્રેજી મૂળમાં
આરંભિક સિક્વન્ટના ગુણધર્મનું નામ
પહેલાં $\fn{InitSeq}$ છે, પણ આ
છેલ્લી પંક્તિમાં $\fn{InitialSeq}$ છે.
અહીં અગાઉ વ્યાખ્યાયિત નામ જ રાખ્યું છે.
આગળની સમજૂતીમાં અંતિમ સિક્વન્ટને
ભૂલથી~$d$નું કહેવામાં આવ્યું છે;
અહીં યોગ્ય ચલ~$p$ રાખ્યો છે.}
પહેલી પંક્તિથી~$p$નું અંતિમ સિક્વન્ટ
ખરેખર વાક્યોનું સિક્વન્ટ છે
તે નિશ્ચિત થાય છે. છેલ્લી પંક્તિ
$p$ માત્ર આરંભિક સિક્વન્ટ હોય તે
કિસ્સાને આવરી લે છે.
\end{proof}

\begin{prob}
  \olref[inc][art][plk]{prop:followsby}ની
  જેમ નીચેના ગુણધર્મો વ્યાખ્યાયિત કરો:
  \begin{enumerate}
  \item $\fn{FollowsBy}_{\Cut}(p)$,
  \item $\fn{FollowsBy}_{\LeftR\lif}(p)$,
  \item $\fn{FollowsBy}_{\eq}(p)$,
  \item $\fn{FollowsBy}_{\RightR{\lforall}}(p)$.
  \end{enumerate}
  છેલ્લા ગુણધર્મ માટે એ પણ બતાવવું પડશે
  કે ગડેલ-સંખ્યા~$p$ ધરાવતી
  નિષ્પત્તિનું છેલ્લું અનુમાન
  વિશિષ્ટ ચલની શરત સંતોષે છે કે નહીં
  તે આદિમ પુનરાવર્તી રીતે ચકાસી શકાય.
  એટલે $\RightR{\lforall}$નો વિશિષ્ટ
  ચલ~$a$ અંતિમ સિક્વન્ટમાં આવતો નથી.
  \end{prob}
  

\begin{prop}
  \ollabel{prop:deriv}
  યોગ્ય નિષ્પત્તિ~$\pi$ની
  ગડેલ-સંખ્યા~$p$ હોય ત્યારે સત્ય
  થતો સંબંધ~$\fn{Deriv}(p)$ આદિમ
  પુનરાવર્તી છે.
\end{prop}

\begin{proof}
  નિષ્પત્તિ~$\pi$ ત્યારે યોગ્ય
  છે જ્યારે તેનું દરેક અનુમાન કોઈ
  નિયમનો યોગ્ય ઉપયોગ હોય, એટલે કે
  દરેક ઉપ-નિષ્પત્તિનો અંત યોગ્ય
  અનુમાનથી થતો હોય. તેથી
  $\fn{Deriv}(p)$ ત્યારે અને માત્ર ત્યારે,
  જ્યારે
  \[
  \bforall{i<\len{\fn{SubtreeSeq}(p)}}{\fn{Correct}((\fn{SubtreeSeq}(p))_i)}.
  \]
  \sourcecorrection{OLINC-009}{સ્થિર અંગ્રેજી
  મૂળમાં આ અનુચ્છેદના ગદ્યમાં
  $\fn{Deriv}(d)$ લખ્યું છે, જ્યારે
  વિધાન અને સૂત્રનો ચલ~$p$ છે.
  મૂળ સૂત્રમાં $\fn{Correct}$ના ઘટકનું
  બંધ કૌંસ પણ છૂટી ગયું છે; અહીં
  બંને અસંગતતાઓ સુધારી છે.}
\end{proof}


\begin{prop}
$\Gamma$ એ વાક્યોનો આદિમ
પુનરાવર્તી ગણ હોય એમ ધારો. ત્યારે
$\Prf[\Gamma](x, y)$ સંબંધ આદિમ
પુનરાવર્તી છે. તેનો અર્થ એ છે કે
$x$ એ કોઈ સાન્ત
$\Gamma_0 \subseteq \Gamma$ માટે
$\Gamma_0 \Sequent !A$ની
નિષ્પત્તિ~$\pi$નો સંકેત છે અને
$y$ એ~$!A$ની ગડેલ-સંખ્યા છે.
\end{prop}

\begin{proof}
``$y \in \Gamma$'' આદિમ પુનરાવર્તી
વિધાન~$R_\Gamma(y)$ વડે આપવામાં
આવ્યું હોય એમ ધારો. આપણે બતાવવું છે
કે $\Prf[\Gamma](x, y)$ આદિમ
પુનરાવર્તી છે. આ સંબંધ ત્યારે અને
માત્ર ત્યારે સત્ય છે જ્યારે $y$ એ
વાક્ય~$!A$ની ગડેલ-સંખ્યા
હોય અને $x$ એ અંતિમ સિક્વન્ટ
$\Gamma_0 \Sequent !A$ ધરાવતી
$\Log{LK}$-નિષ્પત્તિનો સંકેત હોય.

અગાઉના વિધાન પ્રમાણે, $x$ એ
$\Log{LK}$માં યોગ્ય
નિષ્પત્તિ~$\pi$નો સંકેત હોય
ત્યારે સત્ય થતો
$\fn{Deriv}(x)$ ગુણધર્મ આદિમ
પુનરાવર્તી છે. $x$ એવો સંકેત હોય,
તો $\fn{EndSequent}(x)$ એ~$\pi$ના
અંતિમ સિક્વન્ટનો સંકેત છે. તેથી
$(\fn{EndSequent}(x))_0$ અંતિમ
સિક્વન્ટની ડાબી બાજુનો સંકેત છે
અને $(\fn{EndSequent}(x))_1$
જમણી બાજુનો સંકેત છે. આમ,
``$\pi$ના અંતિમ સિક્વન્ટની જમણી
બાજુ માત્ર~$!A$ છે'' તે
$\len{(\fn{EndSequent}(x))_1} = 1
\land ((\fn{EndSequent}(x))_1)_0 = y$
વડે વ્યક્ત થાય છે.
\sourcecorrection{OLINC-010}{સ્થિર અંગ્રેજી મૂળના
આ ગદ્ય-સૂત્રમાં છેલ્લે $y$ને બદલે
$x$ છે; પરંતુ તરત નીચેની
$\Prf[\Gamma](x,y)$ની વ્યાખ્યામાં
યોગ્ય રીતે $y$ છે. અહીં ગદ્ય-સૂત્રને
એ જ વ્યાખ્યા સાથે મેળમાં કર્યું છે.}
$\pi$ના અંતિમ સિક્વન્ટની ડાબી
બાજુ તો આપોઆપ સાન્ત છે; માત્ર
તેનું દરેક વાક્ય~$\Gamma$માં છે
તે વ્યક્ત કરવું પડે. આમ
$\Prf[\Gamma](x, y)$ને આ પ્રમાણે
વ્યાખ્યાયિત કરી શકીએ:
\begin{align*}
\Prf[\Gamma](x, y) \defiff {}&
\fn{Deriv}(x) \land {} \\
& \bforall{i <
  \len{(\fn{EndSequent}(x))_0}}{R_\Gamma(((\fn{EndSequent}(x))_0)_i)} \land {}\\
& \len{(\fn{EndSequent}(x))_1} = 1 \land ((\fn{EndSequent}(x))_1)_0 = y.
\end{align*}
\end{proof}
% END OLP-0286
\endgroup


%
  \begingroup
\makeatletter\def\ol@sectioncs{section}\makeatother

% BEGIN OLP-0287
\olfileid{inc}{art}{pnd}
\olsection{સ્વાભાવિક નિગમનમાં નિષ્પત્તિઓ}
\phantomsection\label{guunit:OLP-0287}


\begin{explain}
નિષ્પત્તિઓનું અંકગણિતીકરણ કરવા
આપણે નિષ્પત્તિઓને સંખ્યાઓ વડે
નિરૂપવાં પડશે. નિષ્પત્તિઓ એ
સૂત્રોનાં વૃક્ષો છે અને તેમાં
દરેક અનુમાનને એક કે બે ચિહ્નન હોય છે.
તેથી પુનરાવર્તી નિરૂપણ સ્વાભાવિક છે:
નિષ્પત્તિને એવી ટપલ તરીકે નિરૂપીએ
જેના ઘટકો છેલ્લા અનુમાનના આધારવાક્યો
તરફ જતા સીધા ઉપ-નિષ્પત્તિઓની સંખ્યા,
આ ઉપ-નિષ્પત્તિઓનાં નિરૂપણો,
અંતિમ સૂત્ર, છેલ્લા અનુમાનનું
નિવૃત્તિ-ચિહ્નન અને છેલ્લા અનુમાનનો
પ્રકાર દર્શાવતો નંબર છે.
\end{explain}

\begin{defn}
  $\delta$ એ સ્વાભાવિક નિગમનમાં
  નિષ્પત્તિ હોય, તો $\Gn{\delta}$ને
  નીચે પ્રમાણે આગમનાત્મક રીતે વ્યાખ્યાયિત કરીએ:
  \begin{enumerate}
  \item
    $\delta$ માત્ર ધારણા~$!A$નું બનેલું
    હોય, તો $\Gn{\delta}$ એ
    $\tuple{0, \Gn{!A}, n}$ છે.
    ધારણા અનિવૃત્ત હોય તો
    નંબર~$n$ એ~$0$ છે; અન્યથા તે
    ધારણાનું અંકીય ચિહ્નન છે.
  \item
    $\delta$નો અંત શૂન્ય, એક, બે કે ત્રણ
    આધારવાક્યોવાળા અનુમાનથી થતો હોય,
    તો અનુક્રમે $\Gn{\delta}$ આ છે:
    \begin{align*}
      & \tuple{0, \Gn{!A}, n, k}, \\
      & \tuple{1, \Gn{\delta_1}, \Gn{!A}, n, k}, \\
      & \tuple{2, \Gn{\delta_1}, \Gn{\delta_2}, \Gn{!A}, n, k}, \text{ અથવા}\\
      & \tuple{3, \Gn{\delta_1}, \Gn{\delta_2}, \Gn{\delta_3}, \Gn{!A}, n,
        k},
    \end{align*}
    અહીં $\delta_1$, $\delta_2$, $\delta_3$
    એ~$\delta$ના છેલ્લા અનુમાનનાં
    આધારવાક્યોમાં પૂરાં થતાં
    ઉપ-નિષ્પત્તિઓ છે;
    $!A$ એ~$\delta$ના છેલ્લા અનુમાનનું
    ફલિતવાક્ય છે; $n$ એ છેલ્લા અનુમાનનું
    નિવૃત્તિ-ચિહ્નન છે (અનુમાન કોઈ ધારણા
    નિવૃત્ત ન કરતું હોય તો~$0$); અને
    છેલ્લા અનુમાનમાં વપરાયેલા નિયમ પ્રમાણે
    નીચેના કોષ્ટકમાંથી~$k$ મળે છે.
  
    \begin{tabular}{lccccccc}
      \text{નિયમ:} & \Intro{\land} & \Elim{\land} & \Intro{\lor} & \Elim{\lor} \\
      $k$: & 1 & 2 & 3 & 4  \\[1.5ex]
      \text{નિયમ:} &  \Intro{\lif} & \Elim{\lif} & \Intro{\lnot} & \Elim{\lnot} \\
      $k$: & 5 & 6 & 7 & 8 \\[1.5ex]
      \text{નિયમ:}  &  \FalseInt & \FalseCl & \Intro{\lforall} & \Elim{\lforall} \\
      $k$: & 9 & 10 & 11 & 12 \\[1.5ex]
      \text{નિયમ:} & \Intro{\lexists} & \Elim{\lexists} & \Intro{\eq} & \Elim{\eq} \\
      $k$: & 13 & 14 & 15 & 16 
    \end{tabular}
  \end{enumerate}
\end{defn}

\begin{ex}
  આ અત્યંત સરળ નિષ્પત્તિ જુઓ:
  \begin{prooftree}
    \AxiomC{$\Discharge{!A \land !B}{1}$}
    \RightLabel{\Elim{\land}}
    \UnaryInfC{$!A$}
    \DischargeRule{\Intro{\lif}}{1}
    \UnaryInfC{$(!A \land !B) \lif !A$}
  \end{prooftree}
  ધારણાની ગડેલ-સંખ્યા
  $d_0 = \tuple{0,
    \Gn{(!A \land !B)}, 1}$ થાય.
  $\Elim{\land}$ના ફલિતવાક્યમાં
  પૂરી થતી નિષ્પત્તિની ગડેલ-સંખ્યા
  $d_1 = \tuple{1,
     d_0, \Gn{!A}, 0, 2}$ થાય
  ($1$ એટલા માટે કે $\Elim{\land}$ને
  એક આધારવાક્ય છે, ફલિતવાક્ય~$!A$ની
  ગડેલ-સંખ્યા લેવાઈ છે, $0$ એટલા માટે
  કે કોઈ ધારણા નિવૃત્ત થતી નથી અને
  $2$ એ $\Elim{\land}$નો સંકેત નંબર
  છે). આખી નિષ્પત્તિની ગડેલ-સંખ્યા
  તેથી
  $\tuple{1, d_1, \Gn{((!A \land !B) \lif
      !A)}, 1, 5}$, એટલે કે
  \[
  \tuple{1, \tuple{1, \tuple{0, \Gn{(!A \land !B)}, 1}, \Gn{!A}, 0, 2},
    \Gn{((!A \land !B) \lif !A)}, 1, 5}.
  \]
  \sourcecorrection{OLINC-011}{સ્થિર અંગ્રેજી
  મૂળમાં ધારણાના સંકેત~$d_0$માં
  $\Gn{!A \land !B}$ છે, જ્યારે
  એ જ ધારણાની વિસ્તૃત ટપલમાં
  $\Gn{(!A \land !B)}$ છે. વાક્યરચનાની
  એકરૂપ સંકેતશ્રેણી માટે અહીં બંને
  જગ્યાએ બીજા સ્વરૂપને રાખ્યું છે.}
\end{ex}

\begin{explain}
નિષ્પત્તિઓનું નિરૂપણ નક્કી કર્યા પછી
આવા નિષ્પત્તિઓની ગડેલ-સંખ્યાઓ પર
આદિમ પુનરાવર્તી રીતે ક્રિયા કરી શકાય
અને તેમના આવશ્યક ગુણધર્મો તથા સંબંધો
વ્યક્ત કરી શકાય તે પણ બતાવવું પડે.
કેટલીક ક્રિયાઓ સરળ છે: દાખલા તરીકે,
નિષ્પત્તિની ગડેલ-સંખ્યા~$d$ હોય,
તો $\fn{EndFmla}(d) = (d)_{(d)_0+1}$
તેના અંતિમ સૂત્રની ગડેલ-સંખ્યા
આપે છે; $\fn{DischargeLabel}(d) = (d)_{(d)_0 + 2}$
નિવૃત્તિ-ચિહ્નન આપે છે અને બિનધારણા
કિસ્સામાં
$\fn{LastRule}(d) = (d)_{(d)_0 + 3}$
છેલ્લા અનુમાનનો પ્રકાર દર્શાવતો નંબર
આપે છે. માત્ર ધારણા હોય ત્યારે આ
વિધેયને શૂન્ય મૂલ્ય આપીને સર્વત્ર
વ્યાખ્યાયિત કરી શકીએ.
\sourcecorrection{OLINC-012}{સ્થિર અંગ્રેજી
મૂળમાં ત્રણ-ઘટક ધારણા-ટપલ માટે
$\fn{LastRule}(d)$નું દર્શાવેલું
ઘટક અસ્તિત્વમાં નથી. આ સમીકરણ
બિનધારણા કિસ્સે લાગુ પડે છે;
ધારણાના કિસ્સામાં શૂન્ય મૂલ્ય
મૂકવાથી આગળનો ગુણધર્મ સર્વત્ર
અર્થપૂર્ણ થાય છે.}
બીજી કેટલીક ક્રિયાઓ ઘણી કઠિન છે.
તેમને કેવી રીતે સંભાળવી તેનો ઓછામાં
ઓછો રૂપરેખાત્મક ખ્યાલ આપીશું. હેતુ
એ બતાવવાનો છે કે ``$\delta$ એ
$\Gamma$માંથી~$!A$ની
નિષ્પત્તિ છે'' એ સંબંધ
$\delta$ અને~$!A$ની ગડેલ-સંખ્યાઓનો
આદિમ પુનરાવર્તી સંબંધ છે.
\end{explain}

\begin{prop}
  નીચેના સંબંધો આદિમ પુનરાવર્તી છે:
  \begin{enumerate}
  \item $!A$ એ~$\delta$માં
    ચિહ્નન~$n$વાળી ધારણા તરીકે આવે છે.
  \item $\delta$માં ચિહ્નન~$n$વાળી
    બધી ધારણાઓ $!A$ સ્વરૂપની છે
    (એટલે $\delta$માં ચિહ્નન~$n$
    વાપરીને ધારણા~$!A$ને
    નિવૃત્ત કરી શકીએ).
  \end{enumerate}
\end{prop}

\begin{proof}
  સૂત્રોની ગડેલ-સંખ્યાઓ અને
  નિષ્પત્તિઓની ગડેલ-સંખ્યાઓ વચ્ચેના
  અનુરૂપ સંબંધો આદિમ પુનરાવર્તી છે
  તે બતાવવાનું છે.
  \begin{enumerate}
  \item $\fn{Assum}(x, d, n)$ આદિમ
    પુનરાવર્તી છે તે બતાવવું છે. તે
    ત્યારે સત્ય છે જ્યારે $x$ એ
    ગડેલ-સંખ્યા~$d$ ધરાવતી
    નિષ્પત્તિમાં ચિહ્નન~$n$વાળી
    ધારણાની ગડેલ-સંખ્યા હોય.
    ગડેલ-સંખ્યા~$\tuple{0, x,
      n}$ ધરાવતી નિષ્પત્તિ એ~$d$ની
    ઉપ-નિષ્પત્તિ હોય ત્યારે આવું બને.
    નિષ્પત્તિઓનું આપેલું
    સંકેતબદ્ધીકરણ
    \olref[cmp][rec][tre]{sec}માં રજૂ થયેલા
    વૃક્ષોના સંકેતબદ્ધીકરણનો વિશિષ્ટ
    કિસ્સો છે. તેથી આદિમ પુનરાવર્તી
    વિધેય~$\fn{SubtreeSeq}(d)$ આ~$d$ની
    બધી ઉપ-નિષ્પત્તિઓની ગડેલ-સંખ્યાઓની
    શ્રેણી આપે છે (તેની લંબાઈ વધુમાં
    વધુ~$d$ છે). આમ વ્યાખ્યાયિત કરીએ:
    \[
    \fn{Assum}(x, d, n) \defiff \bexists{i<d}{(\fn{SubtreeSeq}(d))_i =
      \tuple{0, x, n}}.
    \]
  \item $\fn{Discharge}(x, d, n)$
    સંબંધ આદિમ પુનરાવર્તી છે તે બતાવવું છે.
    તે ત્યારે સત્ય છે જ્યારે ગડેલ-સંખ્યા~$d$
    ધરાવતી નિષ્પત્તિમાં ચિહ્નન~$n$
    ધરાવતી બધી ધારણાઓ ગડેલ-સંખ્યા~$x$
    ધરાવતા સૂત્ર હોય. આ સંબંધ
    ત્યારે અને માત્ર ત્યારે સત્ય છે જ્યારે
    $\bforall{y<d}{(\fn{Assum}(y, d, n) \lif y
      = x)}$.
  \end{enumerate}
\end{proof}

\begin{prop}
  \ollabel{prop:followsby}
  ગડેલ-સંખ્યા~$d$ ધરાવતી
  નિષ્પત્તિ~$\delta$નું છેલ્લું
  અનુમાન યોગ્ય હોય ત્યારે અને માત્ર
  ત્યારે સત્ય થતો ગુણધર્મ
  $\fn{Correct}(d)$ આદિમ પુનરાવર્તી છે.
\end{prop}

\begin{proof}
  દરેક અનુમાન-નિયમ~$R$ માટે
  $\fn{FollowsBy}_R(d)$ સંબંધ આદિમ
  પુનરાવર્તી છે તે બતાવવાનું છે.
  $\fn{FollowsBy}_R(d)$ ત્યારે અને માત્ર
  ત્યારે સત્ય છે જ્યારે $d$ એ
  નિષ્પત્તિ~$\delta$ની ગડેલ-સંખ્યા
  હોય અને~$\delta$નું અંતિમ સૂત્ર
  $\delta$ના સીધા ઉપ-નિષ્પત્તિઓમાંથી
  $R$ના યોગ્ય ઉપયોગ વડે ફલિત થતું હોય.

  એક સરળ કિસ્સો \Intro{\land} નિયમનો
  છે. $\delta$નો અંત યોગ્ય
  $\Intro{\land}$ અનુમાનથી થતો હોય,
  તો તે આ પ્રમાણે દેખાય:
  \begin{prooftree}
    \AxiomC{}
    \RightLabel{$\delta_1$}
    \DeduceC{$!A$}
    
    \AxiomC{}
    \RightLabel{$\delta_2$}
    \DeduceC{$!B$}

    \RightLabel{\Intro\land}
    \BinaryInfC{$!A \land !B$}
  \end{prooftree}
  ત્યારે $\delta$ની ગડેલ-સંખ્યા~$d$
  એ $\tuple{2, d_1, d_2,
    \Gn{(!A \land !B)}, 0, k}$ છે, જ્યાં
  $\fn{EndFmla}(d_1) = \Gn{!A}$,
  $\fn{EndFmla}(d_2) = \Gn{!B}$,
  $n=0$ અને $k=1$. તેથી
  $\fn{FollowsBy}_{\Intro\land}(d)$ને
  આ પ્રમાણે વ્યાખ્યાયિત કરીએ:
  \begin{multline*}
    (d)_0 = 2 \land \fn{DischargeLabel}(d) = 0 \land \fn{LastRule}(d) = 1 \land {}\\
    \fn{EndFmla}(d) = {}\\
    \Gn{(} \concat \fn{EndFmla}((d)_1) \concat \Gn{\land}
    \concat \fn{EndFmla}((d)_2) \concat \Gn{)}.
  \end{multline*}

  બીજું સરળ ઉદાહરણ $\Intro\eq$ નિયમ છે.
  તેને આધારવાક્યો નથી, એટલે ધારણાની
  જેમ $(d)_0 = 0$. તેને
  નિવૃત્તિ-ચિહ્નન પણ નથી, એટલે
  $n=0$. જોકે $!A$ કોઈ બંધ
  પદ~$t$ માટે $\eq[t][t]$ સ્વરૂપનું
  હોવું જોઈએ. અહીં આદિમ પુનરાવર્તી
  વ્યાખ્યા આ છે:
  \begin{multline*}
    (d)_0 = 0 \land \fn{DischargeLabel}(d) = 0 \land
    \fn{LastRule}(d) = 15 \land {}\\
    \bexists{t<d}{(\fn{ClTerm}(t) \land 
      \fn{EndFmla}(d) = {}\\
      \Gn{{\eq}(} \concat t \concat \Gn{,} \concat t \concat \Gn{)})}.
  \end{multline*}
  \sourcecorrection{OLINC-013}{સ્થિર અંગ્રેજી
  મૂળમાં આ નિયમના સૂત્રમાં
  છેલ્લો નિયમ~$15$ છે તે શરત
  છૂટી છે. વ્યાખ્યા જણાવેલા
  $\Intro\eq$ નિયમને જ ઓળખે
  તે માટે અહીં આ શરત ઉમેરી છે.}

  વધુ જટિલ ઉદાહરણમાં
  $\fn{FollowsBy}_{\Intro{\lif}}(d)$
  ત્યારે અને માત્ર ત્યારે સત્ય છે જ્યારે
  $\delta$નું અંતિમ સૂત્ર
  $(!A \lif !B)$ સ્વરૂપનું હોય,
  $\delta_1$નું અંતિમ સૂત્ર~$!B$
  હોય અને~$\delta$માં ચિહ્નન~$n$
  ધરાવતી દરેક ધારણા $!A$ સ્વરૂપની હોય.
  આને આદિમ પુનરાવર્તી રીતે વ્યક્ત કરીએ:
  \begin{multline*}
    (d)_0 = 1 \land \fn{LastRule}(d) = 5 \land {}\\
    \bexists{a<d}{(\fn{Discharge}(a, (d)_1, \fn{DischargeLabel}(d)) \land {}}\\
      \fn{EndFmla}(d) = (\Gn{(} \concat a \concat \Gn{\lif}
      \concat \fn{EndFmla}((d)_1) \concat \Gn{)}))
  \end{multline*}
  (અહીં $a$ને~$!A$ની ગડેલ-સંખ્યા
  તરીકે વિચારો).
  \sourcecorrection{OLINC-014}{સ્થિર અંગ્રેજી
  મૂળના આ સૂત્રમાં છેલ્લો
  નિયમ~$5$ છે તે શરત છૂટી છે;
  નિયમ ઓળખવા માટે અહીં ઉમેરી છે.}

  હવે \Intro{\lexists}નું ઉદાહરણ લો.
  અહીં~$\delta$નું છેલ્લું અનુમાન
  ત્યારે અને માત્ર ત્યારે યોગ્ય છે
  જ્યારે એવું સૂત્ર~$!A$,
  બંધ પદ~$t$ અને ચલ~$x$
  હોય કે $\Subst{!A}{t}{x}$ એ
  નિષ્પત્તિ~$\delta_1$નું અંતિમ
  સૂત્ર હોય અને
  $\lexists[x][!A]$ છેલ્લા અનુમાનનું
  ફલિતવાક્ય હોય. તેથી
  $\fn{FollowsBy}_{\Intro{\lexists}}(d)$
  ત્યારે અને માત્ર ત્યારે સત્ય છે જ્યારે
  \begin{multline*}
    (d)_0 = 1 \land \fn{DischargeLabel}(d) = 0 \land
    \fn{LastRule}(d) = 13 \land {} \\
    \bexists{a < d}{\bexists{x<d}{\bexists{t<d}{
          (\fn{ClTerm}(t) \land \fn{Var}(x)  \land {}}}}\\
    \fn{Subst}(a,t,x) = \fn{EndFmla}((d)_1) \land
    \fn{EndFmla}(d) = (\Gn{\lexists} \concat x \concat a)).
  \end{multline*}
  \sourcecorrection{OLINC-015}{સ્થિર અંગ્રેજી
  મૂળના આ સૂત્રમાં છેલ્લો
  નિયમ~$13$ છે તે શરત છૂટી છે;
  અહીં એ સ્પષ્ટ કરી છે.}

  હવે $\fn{Correct}(d)$ને આ પ્રમાણે
  વ્યાખ્યાયિત કરીએ:
  \begin{multline*}
    \fn{Sent}(\fn{EndFmla}(d)) \land \bigl[{}\\
    (\fn{LastRule}(d) = 1 \land
    \fn{FollowsBy}_{\Intro\land}(d)) \lor \dots \lor {}\\
    (\fn{LastRule}(d) = 16 \land \fn{FollowsBy}_{\Elim\eq}(d)) \lor {}\\
    \bexists{n<d}{\bexists{x<d}{(d = \tuple{0, x, n})}}\bigr].
  \end{multline*}
  \sourcecorrection{OLINC-016}{સ્થિર અંગ્રેજી
  મૂળના સૂત્રમાં છેલ્લાં વિકલ્પોની
  આસપાસ કૌંસ નથી. સામાન્ય
  સંયોજક-પ્રાથમિકતા પ્રમાણે
  $\fn{Sent}$ની શરત ફક્ત પહેલા
  નિયમને લાગુ પડે તેમ બને છે.
  અહીં આખા વિકલ્પ-સમૂહને કૌંસમાં
  રાખ્યો છે જેથી દરેક કિસ્સામાં
  અંતિમ સૂત્ર વાક્ય હોવું જરૂરી બને.}
  પહેલી પંક્તિથી~$d$નું અંતિમ
  સૂત્ર વાક્ય છે તે નિશ્ચિત
  થાય છે. છેલ્લી પંક્તિ~$d$
  માત્ર ધારણા હોય તે કિસ્સો આવરી લે છે.
\end{proof}

\begin{prob}
  \olref[inc][art][pnd]{prop:followsby}ની
  જેમ નીચેના ગુણધર્મો વ્યાખ્યાયિત કરો:
  \begin{enumerate}
  \item $\fn{FollowsBy}_{\Elim{\lif}}(d)$,
  \item $\fn{FollowsBy}_{\Elim{\eq}}(d)$,
  \item $\fn{FollowsBy}_{\Elim{\lor}}(d)$,
  \item $\fn{FollowsBy}_{\Intro{\lforall}}(d)$.
  \end{enumerate}
  છેલ્લા માટે એ પણ બતાવવું પડશે કે
  ગડેલ-સંખ્યા~$d$ ધરાવતી
  નિષ્પત્તિનું છેલ્લું અનુમાન
  વિશિષ્ટ ચલની શરત સંતોષે છે કે નહીં
  તે આદિમ પુનરાવર્તી રીતે ચકાસી શકાય.
  એટલે $\Intro{\lforall}$ અનુમાનનો
  વિશિષ્ટ ચલ~$a$ ન તો~$d$ના
  અંતિમ સૂત્રમાં આવે, ન તો~$d$ની
  ખુલ્લી ધારણામાં આવે. આ માટે
  \olref[inc][art][pnd]{prop:openassum}નો
  આદિમ પુનરાવર્તી વિધાન
  $\fn{OpenAssum}$ વાપરી શકો.
\end{prob}

\begin{prop}
  \ollabel{prop:deriv}
  યોગ્ય નિષ્પત્તિ~$\delta$ની
  ગડેલ-સંખ્યા~$d$ હોય ત્યારે સત્ય
  થતો સંબંધ~$\fn{Deriv}(d)$ આદિમ
  પુનરાવર્તી છે.
\end{prop}

\begin{proof}
  નિષ્પત્તિ~$\delta$ ત્યારે યોગ્ય
  છે જ્યારે તેનું દરેક અનુમાન કોઈ
  નિયમનો યોગ્ય ઉપયોગ હોય, એટલે કે
  તેની દરેક ઉપ-નિષ્પત્તિનો અંત
  યોગ્ય અનુમાનથી થતો હોય. તેથી
  $\fn{Deriv}(d)$ ત્યારે અને માત્ર
  ત્યારે સત્ય છે જ્યારે
  \[
  \bforall{i<\len{\fn{SubtreeSeq}(d)}}{\fn{Correct}((\fn{SubtreeSeq}(d))_i)}
  \]
\end{proof}

\begin{prop}
  \ollabel{prop:openassum}
  ગડેલ-સંખ્યા~$d$ ધરાવતી
  નિષ્પત્તિ~$\delta$ની
  અનિવૃત્ત ધારણા~$!A$ની
  ગડેલ-સંખ્યા~$z$ હોય ત્યારે સત્ય
  થતો સંબંધ~$\fn{OpenAssum}(z, d)$
  આદિમ પુનરાવર્તી છે.
\end{prop}

\begin{proof}
  ધારણાની કોઈ ઘટના~$\delta$ની એવી
  ઉપ-નિષ્પત્તિમાં ચિહ્નન~$n$
  સાથે આવે જેનો અંત નિવૃત્તિ-ચિહ્નન~$n$
  ધરાવતા નિયમથી થાય, તો તે ઘટના
  નિવૃત્ત કહેવાય. તેથી~$!A$ની
  ઓછામાં ઓછી એક ઘટના~$\delta$માં
  નિવૃત્ત ન હોય તો તે
  $\delta$ની અનિવૃત્ત ધારણા છે.
  સાવચેતી જરૂરી છે: $\delta$માં~$!A$ની
  નિવૃત્ત અને
  અનિવૃત્ત બંને પ્રકારની
  ઘટનાઓ હોઈ શકે.

  $\delta_0$, \dots, $\delta_k$ એવી
  શ્રેણી લો જેમાં $\delta_0 =
  \delta$, $\delta_k$ એ ધારણા
  $\Discharge{!A}{n}$ છે (કોઈ~$n$
  માટે) અને $\delta_{i+1}$ એ
  $\delta_i$ની સીધી ઉપ-નિષ્પત્તિ
  છે. ધારણાનું ચિહ્નન શૂન્ય હોય અથવા એવી શ્રેણી
  હોય કે તેના કોઈ $\delta_i$નો અંત
  નિવૃત્તિ-ચિહ્નન~$n$ ધરાવતા
  અનુમાનથી ન થતો હોય, તો~$!A$ એ
  $\delta$ની અનિવૃત્ત
  ધારણા છે.

  આદિમ પુનરાવર્તી
  વિધેય~$\fn{SubtreeSeq}(d)$ આપણને
  $\delta$ની બધી ઉપ-નિષ્પત્તિઓની
  ગડેલ-સંખ્યાઓની શ્રેણી આપે છે.
  $\delta$ની ઉપ-નિષ્પત્તિઓના
  સંકેતોની કોઈ માર્ગ-શ્રેણીના બધા
  ઘટકો તેમાં આવે છે. ઘટક-સમાવેશ
  આદિમ પુનરાવર્તી સંબંધ છે:
  $\fn{Subseq}(s, s')$ ત્યારે અને
  માત્ર ત્યારે સત્ય છે જ્યારે
  $\bforall{i<\len{s}}{\lexists[j<\len{s'}][(s)_i = (s')_j]}$.
  સીધી ઉપ-નિષ્પત્તિ હોવું પણ
  આદિમ પુનરાવર્તી છે:
  $\fn{Subderiv}(d, d')$ ત્યારે અને
  માત્ર ત્યારે, જ્યારે
  $\bexists{j<(d')_0}{d = (d')_{j+1}}$.
  અંતિમ વૃક્ષમાં માર્ગની લંબાઈ અને
  તેના સભ્યોના સંકેતો મૂળ કોડથી
  મર્યાદિત છે; તેથી અગાઉનું
  $\fn{sequenceBound}$ વિધેય
  વાપરીને $B(d)=\fn{sequenceBound}(d,d+1)+1$
  બધા આવા માર્ગોના સંકેતોની
  આદિમ પુનરાવર્તી મર્યાદા આપે છે.
  આમ $\fn{OpenAssum}(z, d)$ને
  આ પ્રમાણે વ્યાખ્યાયિત કરીએ:
  \begin{multline*}
    \bexists{s<B(d)}{(\fn{Subseq}(s, \fn{SubtreeSeq}(d))
      \land (s)_0 = d \land {}} \\
    \bexists{n<d}{((s)_{\len{s} \tsub 1} = \tuple{0, z, n} \land {}}\\
      \bforall{i<(\len{s} \tsub 1)}{(\fn{Subderiv}((s)_{i+1}, (s)_i) \land {}}\\
      (n=0 \lor \fn{DischargeLabel}((s)_i) \neq n)))).
  \end{multline*}
  \sourcecorrection{OLINC-017}{સ્થિર અંગ્રેજી
  મૂળના માર્ગ-પરીક્ષણમાં ચાર સમસ્યાઓ
  છે. $\fn{Subseq}$નું આપેલું સૂત્ર
  ક્રમયુક્ત ઉપશ્રેણી નહીં, માત્ર
  ઘટક-સમાવેશ તપાસે છે; માર્ગનો ક્રમ
  આગળની સીધી-ઉપપ્રમાણ શરત આપે છે.
  સીધી ઉપપ્રમાણની ટપલનું પ્રથમ સ્થાન
  ગણતરી છે, તેથી સંતાન-સૂચકમાં
  એક ઉમેરવું પડે. આખી
  $\fn{SubtreeSeq}(d)$નો સંકેત
  માર્ગના સંકેતોની મર્યાદા હોવો
  જરૂરી નથી; અહીં અગાઉ વ્યાખ્યાયિત
  $\fn{sequenceBound}$માંથી સ્પષ્ટ
  મર્યાદા લીધી છે. અને ખુલ્લી
  ધારણાનું ચિહ્નન~$0$ હોય ત્યારે
  પૂર્વજનાં શૂન્ય નિવૃત્તિ-ચિહ્નનો
  તેને નિવૃત્ત કરતો નથી; તે કિસ્સો
  અહીં અલગ રાખ્યો છે.}
\end{proof}

\begin{prop}
  \ollabel{prop:prf-prim-rec}
  $\Gamma$ એ વાક્યોનો આદિમ
  પુનરાવર્તી ગણ હોય એમ ધારો.
  ત્યારે $\Prf[\Gamma](x, y)$ આદિમ
  પુનરાવર્તી છે: તેનો અર્થ એ છે કે
  $x$ એ~$\Gamma$ની
  અનિવૃત્ત ધારણાઓમાંથી
  $!A$ની નિષ્પત્તિ~$\delta$નો
  સંકેત છે અને $y$ એ~$!A$ની
  ગડેલ-સંખ્યા છે.
\end{prop}

\begin{proof}
  ``$y \in \Gamma$'' આદિમ પુનરાવર્તી
  વિધાન~$R_\Gamma(y)$ વડે આપવામાં
  આવ્યું હોય એમ ધારો. આપણે બતાવવું છે
  કે $\Prf[\Gamma](x, y)$ આદિમ
  પુનરાવર્તી છે. આ સંબંધ ત્યારે અને
  માત્ર ત્યારે સત્ય છે જ્યારે $y$ એ
  વાક્ય~$!A$ની ગડેલ-સંખ્યા હોય અને
  $x$ એ અંતિમ સૂત્ર~$!A$
  તથા~$\Gamma$માંની બધી
  અનિવૃત્ત ધારણાઓ ધરાવતી
  સ્વાભાવિક નિગમનની
  નિષ્પત્તિનો સંકેત હોય.

  \olref{prop:deriv} પ્રમાણે, $x$ એ
  સ્વાભાવિક નિગમનની યોગ્ય
  નિષ્પત્તિ~$\delta$ની ગડેલ-સંખ્યા
  હોય ત્યારે સત્ય થતો
  $\fn{Deriv}(x)$ ગુણધર્મ આદિમ
  પુનરાવર્તી છે. તેથી
  $\Prf[\Gamma](x, y)$ને આ પ્રમાણે
  વ્યાખ્યાયિત કરી શકીએ:
  \begin{align*}
    \Prf[\Gamma](x, y) \defiff {}
    & \fn{Deriv}(x) \land \fn{EndFmla}(x) = y \land {} \\
    & \bforall{z < x}{(\fn{OpenAssum}(z, x) \lif R_\Gamma(z))}.
  \end{align*}
\end{proof}
% END OLP-0287
\endgroup


  \begingroup
\makeatletter\def\ol@sectioncs{section}\makeatother

% BEGIN OLP-0288
\olfileid{inc}{art}{pax}
\olsection{સ્વયંસિદ્ધિમૂલક નિષ્પત્તિઓ}
\phantomsection\label{guunit:OLP-0288}


\begin{explain}
સ્વયંસિદ્ધિમૂલક નિષ્પત્તિઓનું
અંકગણિતીકરણ કરવા માટે આપણે
નિષ્પત્તિઓને સંખ્યાઓ વડે નિરૂપવાં
પડશે. નિષ્પત્તિઓ તો માત્ર
સૂત્રોની શ્રેણીઓ છે. તેથી
સ્પષ્ટ રીત એ છે કે દરેક
નિષ્પત્તિને તેમાંનાં સૂત્રોના
સંકેતોની શ્રેણીના સંકેત વડે નિરૂપવું.
\end{explain}

\begin{defn}
$\delta$ એ સૂત્રો
$!A_1$, \dots,~$!A_n$થી બનેલી
સ્વયંસિદ્ધિમૂલક નિષ્પત્તિ હોય,
તો $\Gn{\delta}$ આ છે:
\[
\tuple{\Gn{!A_1}, \dots, \Gn{!A_n}}.
\]
\end{defn}

\begin{ex}
  આ અત્યંત સરળ નિષ્પત્તિ જુઓ:
  \begin{derivation}
  1. & $!B \lif (!B \lor !A)$ \\
  2. & $(!B \lif (!B \lor !A)) \lif (!A  \lif (!B \lif (!B \lor !A)))$\\
  3. & $!A  \lif (!B \lif (!B \lor !A))$
  \end{derivation}
  આ નિષ્પત્તિની ગડેલ-સંખ્યા
  નીચે પ્રમાણે થાય:
  \begin{align*}
  \openTuple\, 
    & \Gn{!B \lif (!B \lor !A)}, \\
    &\Gn{(!B \lif (!B \lor !A)) \lif (!A  \lif (!B \lif (!B \lor !A)))},\\
    & \Gn{!A  \lif (!B \lif (!B \lor !A))} \,\closeTuple.
  \end{align*}
\end{ex}

\begin{explain}
નિષ્પત્તિઓનું નિરૂપણ નક્કી કર્યા પછી
આવા નિષ્પત્તિઓ પર આદિમ
પુનરાવર્તી રીતે ક્રિયા કરી શકાય
અને તેમના આવશ્યક ગુણધર્મો તથા
સંબંધો એ જ રીતે વ્યક્ત કરી શકાય
તે પણ બતાવવું પડે. કેટલીક ક્રિયાઓ
સરળ છે: દાખલા તરીકે,
નિષ્પત્તિની ગડેલ-સંખ્યા~$d$
હોય, તો $(d)_{\len{d}-1}$ તેના અંતિમ
સૂત્રની ગડેલ-સંખ્યા આપે છે.
બીજી કેટલીક ક્રિયાઓ ઘણી કઠિન છે.
તેમને કેવી રીતે સંભાળવી તેનો ઓછામાં
ઓછો રૂપરેખાત્મક ખ્યાલ આપીશું.
હેતુ એ બતાવવાનો છે કે
``$\delta$ એ~$\Gamma$માંથી~$!A$ની
નિષ્પત્તિ છે'' એ સંબંધ
$\delta$ અને~$!A$ની ગડેલ-સંખ્યાઓનો
આદિમ પુનરાવર્તી સંબંધ છે.
\end{explain}

\begin{prop}
\ollabel{prop:followsby}
નીચેના સંબંધો આદિમ પુનરાવર્તી છે:
\begin{enumerate}
\item $!A$ સ્વયંસિદ્ધિ છે.
\item $\delta$ની $i$મી પંક્તિ
  મોડસ પોનેન્સ વડે પ્રમાણિત છે.
\item $\delta$ની $i$મી પંક્તિ
  \QR{} વડે પ્રમાણિત છે.
\item $\delta$ યોગ્ય નિષ્પત્તિ છે.
\end{enumerate}
\end{prop}

\begin{proof}
સૂત્રોની ગડેલ-સંખ્યાઓ અને
નિષ્પત્તિઓની ગડેલ-સંખ્યાઓ વચ્ચેના
અનુરૂપ સંબંધો આદિમ પુનરાવર્તી છે
તે બતાવવું છે.
\begin{enumerate}
\item સ્વયંસિદ્ધિ-યોજનાઓની સીમિત યાદી
  આપેલી છે. $!A$ કોઈ યોજના દર્શાવેલા
  સ્વરૂપનું હોય તો સ્વયંસિદ્ધિ છે.
  યોજનાઓની યાદી સીમિત હોવાથી,
  દરેક યોજના માટે $!A$ એ સ્વરૂપનું
  છે કે નહીં તે આદિમ પુનરાવર્તી રીતે
  ચકાસી શકાય તે બતાવવું પૂરતું છે.
  દાખલા તરીકે આ યોજના લો:
  \[
  !B \lif (!C \lif !B).
  \]
  એવાં સૂત્રો $!B$ અને $!C$
  હોય કે `$($' સાથે $!B$, પછી
  `$\lif$', પછી `$($', પછી $!C$,
  પછી `$\lif$', પછી $!B$ અને
  છેલ્લે~`$))$' જોડવાથી $!A$ મળે,
  તો $!A$ આ સ્વયંસિદ્ધિ-યોજનાનું
  દૃષ્ટાંત છે. $!A$ની ગડેલ-સંખ્યા
  $n$ માટે અનુરૂપ ગુણધર્મ ચકાસી
  શકીએ છીએ, કારણ કે શ્રેણીઓનું
  જોડાણ આદિમ પુનરાવર્તી છે. વધુમાં,
  સંબંધ સત્ય હોય ત્યારે $!B$ અને
  $!C$ બંને~$!A$નાં ઉપ-સૂત્રો
  છે; તેથી $!B$ અને $!C$ની ગડેલ-સંખ્યાઓ
  $!A$ની ગડેલ-સંખ્યા કરતાં નાની
  છે. આમ વ્યાખ્યાયિત કરીએ:
  \begin{multline*}
  \fn{IsAx}_{!B \lif (!C \lif !B)}(n) \defiff \bexists{b< n}{
    \bexists{c < n}{(\fn{Sent}(b) \land \fn{Sent}(c) \land {}}}\\ n =
  \Gn{(} \concat b \concat \Gn{\lif} \concat \Gn{(} \concat c \concat
  \Gn{\lif} \concat b \concat \Gn{))}).
  \end{multline*}
  દરેક સ્વયંસિદ્ધિ-યોજના માટે આવી
  વ્યાખ્યા હોય, તો તેમનો વિકલ્પયોગ
  $\fn{IsAx}(n)$ ગુણધર્મ વ્યાખ્યાયિત
  કરે છે: ``$n$ સ્વયંસિદ્ધિની
  ગડેલ-સંખ્યા છે.''
\item $\delta$ની $i$મી પંક્તિ
  મોડસ પોનેન્સથી ત્યારે અને માત્ર
  ત્યારે પ્રમાણિત છે જ્યારે એવી
  પંક્તિઓ $j$ અને $k < i$ હોય કે
  પંક્તિ~$j$ પરનું વાક્ય
  કોઈ સૂત્ર~$!A$ હોય, પંક્તિ~$k$
  પરનું વાક્ય $!A \lif !B$ હોય
  અને પંક્તિ~$i$ પરનું વાક્ય~$!B$ હોય.
  \begin{multline*}
    \fn{MP}(d, i) \defiff \bexists{j < i}{\bexists{k < i}{}}\\
    (d)_k = \Gn{(} \concat (d)_j
      \concat \Gn{\lif} \concat (d)_i \concat \Gn{)}
  \end{multline*}
  સીમિત પરિમાણન, શ્રેણીઓનું જોડાણ
  અને $=$ આદિમ પુનરાવર્તી હોવાથી
  આ સંબંધ પણ આદિમ પુનરાવર્તી છે.
\item $\delta$ની કોઈ પંક્તિ
  \QR{} વડે પ્રમાણિત છે જો તે
  $!B \lif \lforall[x][!A(x)]$
  સ્વરૂપની હોય, કોઈ અગાઉની પંક્તિ
  કોઈ અચળ~$c$ માટે
  $!B \lif !A(c)$ હોય અને $c$
  એ~$!B$માં ન આવતો હોય. આ ત્યારે
  અને માત્ર ત્યારે બને જ્યારે
  \begin{enumerate}
  \item કોઈ વાક્ય~$!B$ હોય,
  \item મુક્ત રીતે માત્ર ચલ~$x$
    ધરાવતું સૂત્ર~$!A(x)$ હોય,
  \item પંક્તિ~$i$ પર
    $!B \lif \lforall[x][!A(x)]$ હોય,
  \item કોઈ સ્થિર પદ~$c$ માટે અગાઉની
    પંક્તિ $j < i$ પર
    $!B \lif \Subst{!A}{c}{x}$ હોય,
  \item અને એ~$!B$માં ન આવતું હોય.
  \end{enumerate}
  આ બધી શરતો આદિમ પુનરાવર્તી રીતે
  ચકાસી શકાય છે: $!B$, $!A(x)$
  અને $x$ની ગડેલ-સંખ્યાઓ પંક્તિ~$i$
  પરના સૂત્રની ગડેલ-સંખ્યા કરતાં
  નાની છે, અને $c$ની સંખ્યા
  પંક્તિ~$j$ પરના સૂત્રની ગડેલ-સંખ્યા
  કરતાં નાની છે:
  \begin{multline*}
    \fn{QR}_1(d, i) \defiff \bexists{j<i}{\bexists{b < (d)_i}{\bexists{x <
        (d)_i}{\bexists{a < (d)_i}{\bexists{c < (d)_j}{(}}}}}
    \\ \fn{Var}(x) \land \fn{Const}(c) \land {} \\
    (d)_i = \Gn{(} \concat
    b \concat \Gn{\lif} \concat \Gn{\lforall} \concat x \concat a
    \concat \Gn{)} \land {}\\
    (d)_j = \Gn{(} \concat b \concat
    \Gn{\lif} \concat \fn{Subst}(a,c,x) \concat \Gn{)} \land {}\\ \fn{Sent}(b)
    \land \fn{Sent}(\fn{Subst}(a,c,x)) \land {} \bforall{k <
      \len{b}}{(b)_k \neq (c)_0})
  \end{multline*}
  \sourcecorrection{OLINC-018}{સ્થિર અંગ્રેજી
  મૂળના આ સૂત્રમાં અગાઉની પંક્તિનો
  સૂચક~$j$ તેની મર્યાદા
  $(d)_j$માં અને પછીના સમીકરણમાં
  વપરાય છે, પરંતુ તેને પરિમાણિત
  કરવામાં આવ્યો નથી. અહીં
  $j<i$ની સીમિત અસ્તિત્વ-શરત
  બહારથી ઉમેરી છે. મૂળ ગદ્યમાં
  સ્થિર પદ~$c$ને એક જગ્યાએ
  ભૂલથી~$a$ લખ્યું છે;
  અહીં~$c$ જ રાખ્યું છે.}
  અહીં $c$ અને $x$ અનુક્રમે
  સ્થિર પદ અને ચલની ગડેલ-સંખ્યાઓ છે;
  માત્ર તેમના ચિહ્ન-સંકેતો નથી.
  $!A(x)$નું એકમાત્ર મુક્ત ચલ~$x$
  છે તે ચકાસવા માટે
  $\Subst{!A(x)}{c}{x}$
  વાક્ય છે કે નહીં તે ચકાસીએ.
  $!B$નું દરેક ચિહ્ન~$c$થી જુદું છે
  એવી શરતથી $c$ એ~$!B$માં નથી
  તેની ખાતરી કરીએ.

  \QR{}નું બીજું સ્વરૂપ અભ્યાસપ્રશ્ન
  તરીકે મૂકીએ છીએ.
\item $d$ યોગ્ય નિષ્પત્તિની
  ગડેલ-સંખ્યા ત્યારે અને માત્ર ત્યારે
  છે જ્યારે તેની દરેક પંક્તિ કાં તો
  સ્વયંસિદ્ધિ હોય અથવા મોડસ પોનેન્સ
  કે~\QR{} વડે પ્રમાણિત હોય. તેથી:
  \[
  \fn{Deriv}(d) \defiff \bforall{i < \len{d}}{(\fn{IsAx}((d)_i) \lor
    \fn{MP}(d,i) \lor \fn{QR}(d, i))}
  \]
\end{enumerate}
\end{proof}

\begin{prob}
\olref[inc][art][pax]{prop:followsby}ની
જેમ નીચેના સંબંધો વ્યાખ્યાયિત કરો:
\begin{enumerate}
\item $\fn{IsAx}_{!A \lif (!B \lif (!A \land !B))}(n)$,
\item $\fn{IsAx}_{\lforall[x][!A(x)] \lif !A(t)}(n)$,
\item $\fn{QR}_{2}(d, i)$
  (\QR{}ના બીજા સ્વરૂપ માટે).
\end{enumerate}
\end{prob}

\begin{prop}
$\Gamma$ એ વાક્યોનો આદિમ
પુનરાવર્તી ગણ હોય એમ ધારો.
ત્યારે $\Prf[\Gamma](x, y)$ આદિમ
પુનરાવર્તી છે: અહીં~$x$ એ
મૂળ નિષ્કર્ષ~$!A$ માટે
$\Gamma$માંથી મળતાં સાન્ત
ધારણાવાક્યોને પૂર્વપક્ષ બનાવીને
રચાયેલા શરતી વાક્યની બંધ
નિષ્પત્તિ~$\delta$નો સંકેત છે
અને $y$ એ મૂળ નિષ્કર્ષ~$!A$ની
ગડેલ-સંખ્યા છે.
\sourcecorrection{OLINC-019}{સ્થિર અંગ્રેજી
મૂળ આ વિધાનમાં~$x$ને
$\Gamma$માંથી~$!A$ના સીધા
નિગમનનો સંકેત કહે છે,
પરંતુ તેની નીચેની વ્યાખ્યામાં~$x$
ખાલી ધારણાઓમાંથી સાન્ત
$\Gamma$-વાક્યોના શરતીકરણના
નિગમનનો સંકેત છે.
આ બે સંકેત-સંબંધો જોડીદીઠ સમાન
નથી. અહીં વિધાનને નીચે ખરેખર
વ્યાખ્યાયિત થયેલા સંબંધ સાથે
મેળમાં કર્યું છે; મૂળ નિષ્કર્ષની
સાધ્યતા સાથેનો સંબંધ નિગમન
પ્રમેય આપે છે.}
\end{prop}

\begin{proof}
``$y \in \Gamma$'' આદિમ પુનરાવર્તી
વિધાન~$R_\Gamma(y)$ વડે આપવામાં
આવ્યું હોય એમ ધારો. આપણે બતાવવું છે
કે $\Prf[\Gamma](x, y)$ સંબંધ આદિમ
પુનરાવર્તી છે. અંગ્રેજી મૂળની
શરૂઆતમાં $\Prf[\Gamma](x, y)$નો
આશય એવો આપ્યો છે કે
$y$ એ વાક્ય~$!A$ની
ગડેલ-સંખ્યા છે અને $x$ એ
$\Gamma$માંથી~$!A$ની
નિષ્પત્તિનો સંકેત છે. નીચેની
રચના જોકે પુરાવા-સંકેત માટે બંધ નિગમનનો
સંકેત લે છે; આ ભેદ ઉપર નોંધ્યો છે.

અગાઉના વિધાન પ્રમાણે,
$\fn{Deriv}(x)$ આદિમ પુનરાવર્તી છે:
તે ત્યારે અને માત્ર ત્યારે સત્ય છે
જ્યારે $x$ યોગ્ય
નિષ્પત્તિ~$\delta$નો સંકેત હોય.
જોકે તે વ્યાખ્યામાં નિષ્પત્તિની પંક્તિઓને
પ્રમાણિત કરવાની વધારાની રીત તરીકે
ધારણાગણ લેવાયો નથી. \QR{} વડે
પંક્તિ પ્રમાણિત છે કે નહીં તેની
આપણી આદિમ પુનરાવર્તી ચકાસણીમાં
સ્થિર પદ~$c$ એ~$\Gamma$માં આવવું
જોઈએ નહીં એવી શરત પણ સામેલ નથી.
વ્યાખ્યામાં સીધો~$\Gamma$ લઈએ તેવો
ફેરફાર શક્ય છે; પરંતુ
$\fn{Deriv}$ અને નિગમન પ્રમેય
વાપરવા સહેલા છે.
$\Gamma \Proves !A$ ત્યારે અને માત્ર
ત્યારે જ્યારે~$\Gamma$માંનાં
વાક્યોની કોઈ સાન્ત યાદી
$!B_1$, \dots, $!B_n \in
\Gamma$ એવી હોય કે
$\{!B_1, \dots, !B_n\} \Proves !A$.
નિગમન પ્રમેય પ્રમાણે આ ત્યારે
સત્ય છે જ્યારે
$\Proves (!B_1 \lif (!B_2 \lif
\cdots (!B_n \lif !A)\cdots))$.
ગડેલ-સંખ્યા~$z$ ધરાવતું
વાક્ય આ સ્વરૂપનું છે કે
નહીં તે આદિમ પુનરાવર્તી રીતે
ચકાસી શકાય છે. તેથી $x$ને
ગડેલ-સંખ્યા~$y$ ધરાવતા
વાક્યની \emph{$\Gamma$માંથી}
નિષ્પત્તિનો સંકેત લેવા બદલે,
આ~$x$ને ઉપરના સ્વરૂપના નેસ્ટેડ
શરતી વાક્યની~$\emptyset$માંથી
નિષ્પત્તિનો સંકેત લઈએ છીએ.

પહેલાં, વાક્યોની શ્રેણી હોય,
તો એ બધાં વાક્યો પૂર્વપક્ષમાં
અને આપેલું વાક્ય ઉત્તરપક્ષમાં
રાખીને આદિમ પુનરાવર્તી રીતે
શરતી વાક્ય બનાવી શકીએ:
\begin{align*}
  \fn{hCond}(s, y, 0) & = y \\
  \fn{hCond}(s, y, n+1) & = \Gn{(} \concat (s)_{n} \concat \Gn{\lif}
  \concat \fn{hCond}(s, y, n) \concat \Gn{)}\\
  \fn{Cond}(s, y) & = \fn{hCond}(s, y, \len{s})\\
  \intertext{તેથી $\Prf[\Gamma](x, y)$ને આ પ્રમાણે વ્યાખ્યાયિત કરી શકીએ:}
  \Prf[\Gamma](x, y) & \defiff \bexists{s < \fn{sequenceBound}(x,x)}{(} \\
  &\qquad (x)_{\len{x}-1} = \fn{Cond}(s,y) \land {} \\
  &\qquad\bforall{i<\len{s}}{R_\Gamma((s)_i)} \land {} \\
  &\qquad\fn{Deriv}(x)).
\end{align*}
\sourcecorrection{OLINC-020}{સ્થિર અંગ્રેજી
મૂળમાં $\fn{hCond}$ના પુનરાવર્તનના
પગલે ત્રણ દલીલોવાળું
$\fn{Cond}(s,y,n)$ લખ્યું છે,
જ્યારે $\fn{Cond}$ની વ્યાખ્યા
નીચે માત્ર બે દલીલો માટે છે.
અહીં પુનરાવર્તનમાં યોગ્ય
$\fn{hCond}(s,y,n)$ રાખ્યું છે.}
\sourcecorrection{OLINC-021}{સ્થિર અંગ્રેજી
મૂળના અંતિમ સૂત્રમાં
$(s)_i \in \Gamma$ અનૌપચારિક
સભ્યત્વ-લેખન છે; ઉપરથી આપેલા
આદિમ પુનરાવર્તી વિધાન
$R_\Gamma((s)_i)$ વડે જ તેની
સીમિત ચકાસણી વ્યક્ત થાય છે.
અહીં એ જ વિધાન વાપર્યું છે.}
$s$ માટેની મર્યાદા એ હકીકત પરથી
આવે છે કે દરેક $(s)_i$ એ
નિષ્પત્તિની છેલ્લી પંક્તિના
કોઈ ઉપ-સૂત્રની ગડેલ-સંખ્યા
છે; એટલે તે $(x)_{\len{x}-1}$થી
નાની છે. $!B \in \Gamma$
પૂર્વપક્ષોની સંખ્યા, એટલે~$s$ની
લંબાઈ, $x$ની છેલ્લી પંક્તિની
લંબાઈ કરતાં નાની છે.
\end{proof}
% END OLP-0288
\endgroup


\OLEndChapterHook
% END OLP-0280
\endgroup


\begingroup
\makeatletter\def\ol@sectioncs{section}\makeatother

% BEGIN OLP-0289
\olchapter{inc}{req}{$\Th{Q}$માં નિરૂપ્યતા}
\olchapterid{inc}{req}
\phantomsection\label{guunit:OLP-0289}


\begingroup
\makeatletter\def\ol@sectioncs{section}\makeatother

% BEGIN OLP-0290
\olfileid{inc}{req}{int}
\olsection{પરિચય}
\phantomsection\label{guunit:OLP-0290}


અપૂર્ણતાનાં પ્રમેયો એવા સિદ્ધાંતોને લાગુ
પડે છે જેમાં સંગણનીય વિધેયો વિશેનાં
પાયાનાં તથ્યો વ્યક્ત અને સાબિત કરી
શકાય. આવા અત્યંત ન્યૂનતમ સિદ્ધાંતનું
વર્ણન કરીશું; તેનું નામ
``$\Th{Q}$'' છે (રાફેલ રોબિન્સનના
નામ પરથી તેને ક્યારેક
``રોબિન્સનનું $Q$'' પણ કહે છે).
વિધેય $\Th{Q}$માં \emph{નિરૂપ્ય}
હોય તેનો અર્થ જણાવીશું અને પછી આ
સાબિત કરીશું:
\begin{quote}
  વિધેય $\Th{Q}$માં નિરૂપ્ય હોય
  ત્યારે અને માત્ર ત્યારે તે સંગણનીય છે.
\end{quote}
એક રીતે આ આપણને સંગણનીયતાનું બીજું
નિદર્શ આપે છે. પરંતુ અટકવાની
સમસ્યાને તેમાં રૂપાંતરિત કરીને
$\Setabs{!A}{\Th{Q} \Proves !A}$
ગણ નિર્ણેય નથી તે બતાવવા પણ તેને
વાપરીશું. અંતે આના કરતાં ઘણાં વધુ
શક્તિશાળી પરિણામો સાબિત કર્યાં હશે.

$\Th{Q}$ની ભાષા અંકગણિતની ભાષા છે.
$\Th{Q}$ નીચેના સ્વયંસિદ્ધ વાક્યોનું
બનેલું છે (તેમની સાથે
તાદાત્મ્યવાળા પ્રથમ-ક્રમના
તર્કશાસ્ત્રનાં બીજા સ્વયંસિદ્ધ વાક્યો
અને નિયમો પણ વાપરવાના છે):
\begin{align*}
& \lforall[x][\lforall[y][(\eq[x'][y'] \lif \eq[x][y])]] \tag{$!Q_1$}\\
& \lforall[x][\eq/[\Obj 0][x']] \tag{$!Q_2$}\\
& \lforall[x][(\eq[x][\Obj 0] \lor \lexists[y][\eq[x][y']])] \tag{$!Q_3$}\\
& \lforall[x][\eq[(x + \Obj 0)][x]] \tag{$!Q_4$}\\
& \lforall[x][\lforall[y][\eq[(x + y')][(x + y)']]] \tag{$!Q_5$}\\
& \lforall[x][\eq[(x \times \Obj 0)][\Obj 0]] \tag{$!Q_6$}\\
& \lforall[x][\lforall[y][\eq[(x \times y')][((x \times y) + x)]]] \tag{$!Q_7$}\\
& \lforall[x][\lforall[y][(x < y \liff \lexists[z][\eq[(z' + x)][y]])]] \tag{$!Q_8$}
\end{align*}
દરેક પ્રાકૃતિક સંખ્યા~$n$ માટે
અંકનિરૂપણ~$\num{n}$ને એવું
પદ~$\Obj{0}^{\prime\prime\ldots\prime}$
વ્યાખ્યાયિત કરીએ જેમાં કુલ~$n$
ઉપરચિહ્નો હોય. આમ $\num{0}$
માત્ર અચળ~$\Obj{0}$ છે,
$\num{1}$ એ $\Obj{0}'$ છે,
$\num{2}$ એ $\Obj{0}''$ છે, વગેરે.

અંકગણિતના સિદ્ધાંત તરીકે
$\Th{Q}$ \emph{અત્યંત} નબળો છે.
દાખલા તરીકે,
$\lforall[x][\eq/[x][x']]$ અથવા
$\lforall[x][\lforall[y][(x + y) = (y +
    x)]]$ જેવાં ઘણાં સરળ વિધાનો પણ
તેમાં સાબિત થઈ શકતાં નથી. જોકે
$\Th{Q}$ આટલો નબળો છે \emph{એટલા
માટે જ} તે રસપ્રદ છે, એવું આપણે
જોઈશું. હકીકતમાં અપૂર્ણતા પ્રમેય
લાગુ પડે એટલો તે માંડ શક્તિશાળી છે.
$\Th{Q}$ રસપ્રદ હોવાનું બીજું કારણ
તેનો સ્વયંસિદ્ધ વાક્યોનો ગણ
\emph{સાન્ત} હોવો છે.

$\Th{Q}$માં આગમનની યોજના ઉમેરીએ,
તો $\Th{Q}$ કરતાં વધુ શક્તિશાળી સિદ્ધાંત
(\emph{પિયાનો અંકગણિત}~$\Th{PA}$)
મળે છે:
\[
(!A(\Obj 0) \land \lforall[x][(!A(x) \lif !A(x'))]) \lif \lforall[x][!A(x)]
\]
અહીં $!A(x)$ કોઈપણ સૂત્ર હોઈ શકે.
$!A(x)$માં $x$ સિવાયના મુક્ત
ચલો હોય, તો એ બધાંને બાંધવા
આગળ સર્વપરિમાણકો ઉમેરીએ છીએ
(જેથી આગમન-યોજનાનું અનુરૂપ દૃષ્ટાંત
વાક્ય બને). દાખલા તરીકે,
$!A(x, y)$માં ચલ~$y$
પણ મુક્ત હોય, તો અનુરૂપ દૃષ્ટાંત આ છે:
\[
\lforall[y][((!A(\Obj 0) \land \lforall[x][(!A(x) \lif !A(x'))]) \lif
  \lforall[x][!A(x)])]
\]
આગમન-યોજનાનાં દૃષ્ટાંતો વાપરીને
$\Th{PA}$નાં સ્વયંસિદ્ધ વાક્યોમાંથી
$\Th{Q}$નાં વાક્યો કરતાં ઘણું વધુ
સાબિત કરી શકાય છે. હકીકતમાં,
પ્રાકૃતિક સંખ્યાઓ વિશેનાં એવાં
``સ્વાભાવિક'' વિધાનો શોધવા ખાસ્સી
મહેનત પડે છે જે પિયાનો અંકગણિતમાં
સાબિત ન થઈ શકે!{}

\begin{defn}
\ollabel{defn:representable-fn}
  પ્રાકૃતિક સંખ્યાઓમાંથી પ્રાકૃતિક
  સંખ્યાઓમાં જતું વિધેય
  $f(x_0,\ldots,x_k)$
  {$\Th{Q}$માં \emph{નિરૂપ્ય}}
  કહેવાય, જો એવું સૂત્ર
  $!A_f(x_0,\dots,x_k,y)$ હોય કે
  જ્યારે પણ $f(n_0,\dots,n_k) = m$
  હોય, ત્યારે $\Th{Q}$ નીચેના
  બંને વિધાનો સાબિત કરે:
\begin{enumerate}
\item\ollabel{defn:rep:a} $!A_f(\num{n_0}, \dots, \num{n_k}, \num{m})$
\item\ollabel{defn:rep:b} $\lforall[y][(!A_f(\num{n_0}, \dots,
\num{n_k}, y) \lif \num{m} = y)]$.
\end{enumerate}
\end{defn}

વ્યાખ્યા બીજી રીતે પણ કહી શકાય.
દાખલા તરીકે, સમકક્ષ રીતે એવી
શરત મૂકી શકાય કે $\Th{Q}$
$\lforall[y][(!A_f(\num{n_0}, \dots,
\num{n_k}, y) \liff \eq[y][\num{m}])]$
સાબિત કરે.

\begin{thm}
\ollabel{thm:representable-iff-comp}
વિધેય $\Th{Q}$માં નિરૂપ્ય હોય ત્યારે
અને માત્ર ત્યારે તે સંગણનીય છે.
\end{thm}

પ્રમેય સાબિત કરવાની બે દિશાઓ છે.
વાક્યરચનાનું અંકગણિતીકરણ હાથવગું
હોય, તો ડાબેથી જમણી દિશા ઘણી
સીધી છે. બીજી દિશા વધુ મહેનત માગે
છે. મૂળ વિચાર આ છે: ``સંગણનીય''નો
ચોક્કસ અર્થ આપવા ``સામાન્ય
પુનરાવર્તી'' પસંદ કરીએ અને દરેક
સામાન્ય પુનરાવર્તી વિધેય
$\Th{Q}$માં નિરૂપ્ય છે તે બતાવીએ.
યાદ કરો કે વિધેય સામાન્ય પુનરાવર્તી
છે જો તે $\Zero$, ઉત્તરગામી
વિધેય~$\Succ$ અને પ્રક્ષેપ
વિધેયો~$\Proj{n}{i}$માંથી સંયોજન,
આદિમ પુનરાવર્તન અને નિયમિત
લઘુત્તમીકરણ વાપરીને વ્યાખ્યાયિત
કરી શકાય. તેથી દરેક સામાન્ય
પુનરાવર્તી વિધેય $\Th{Q}$માં
નિરૂપ્ય છે તે બતાવવાની એક રીત
એ છે કે પાયાનાં વિધેયો નિરૂપ્ય છે
તે સાબિત કરીએ અને નિરૂપ્ય
વિધેયોમાંથી સંયોજન, આદિમ
પુનરાવર્તન તથા નિયમિત
લઘુત્તમીકરણ વડે વ્યાખ્યાયિત
થતાં વિધેયો પણ નિરૂપ્ય છે તે
સાબિત કરીએ. બીજા શબ્દોમાં,
પાયાનાં વિધેયો નિરૂપ્ય છે અને
નિરૂપ્ય વિધેયોનો વર્ગ સંયોજન,
આદિમ પુનરાવર્તન તથા નિયમિત
લઘુત્તમીકરણ હેઠળ ``બંધ'' છે
તે બતાવી શકીએ. આમ દરેક સામાન્ય
પુનરાવર્તી વિધેયની નિરૂપ્યતા મળે.

પરંતુ નિરૂપ્ય વિધેયો આદિમ
પુનરાવર્તન હેઠળ બંધ છે તે બતાવવાનું
પગલું મુશ્કેલ છે. તેને ટાળવા
પહેલાં બતાવીએ કે ખરેખર આદિમ
પુનરાવર્તન વિના પણ કામ ચાલે છે.
એટલે કે દરેક સામાન્ય પુનરાવર્તી
વિધેય પાયાનાં વિધેયોમાંથી માત્ર
સંયોજન અને નિયમિત લઘુત્તમીકરણ
વાપરીને વ્યાખ્યાયિત થઈ શકે છે.
આ માટે બતાવીએ કે આદિમ પુનરાવર્તન
કોઈ ચોક્કસ નિયમિત લઘુત્તમીકરણ
વડે કરી શકાય છે. જોકે આ રીત
કાર્ય કરે તે માટે પાયાનાં વિધેયોમાં
કેટલાક વધુ વિધેયો ઉમેરવા પડે:
સરવાળો, ગુણાકાર અને સમતા
સંબંધનું લાક્ષણિક વિધેય~$\Char{=}$.
પછી આ \emph{બધાં} પાયાનાં વિધેયો
$\Th{Q}$માં નિરૂપ્ય છે અને નિરૂપ્ય
વિધેયો સંયોજન તથા નિયમિત
લઘુત્તમીકરણ હેઠળ બંધ છે તે
બતાવીને પ્રમેય સાબિત કરી શકીએ.
% END OLP-0290
\endgroup


\begingroup
\makeatletter\def\ol@sectioncs{section}\makeatother

% BEGIN OLP-0291
\olfileid{inc}{req}{rpc}
\olsection{$\Th{Q}$માં નિરૂપ્ય વિધેયો સંગણનીય છે}
\phantomsection\label{guunit:OLP-0291}


$\Th{Q}$માં નિરૂપ્ય દરેક વિધેય
સંગણનીય છે તે સાબિત કરીશું.
પહેલાં~$\Th{Q}$માં નિરૂપ્ય વિધેયો
વિશે એક લેમ્મા સ્થાપવો પડશે.

\begin{lem}\ollabel{lem:rep-q}
  $f(x_0,
  \dots, x_k)$ વિધેય~$\Th{Q}$માં
  નિરૂપ્ય હોય, તો એવું
  સૂત્ર~$!A_f(x_0, \dots, x_k, y)$
  છે કે
  \sourcecorrection{OLINC-022}{સ્થિર અંગ્રેજી
  મૂળના આ વાક્યમાં સૂત્રનું નામ
  $!A$ લખાયું છે, પરંતુ અગાઉની
  વ્યાખ્યા અને આ લેમ્માના તરત
  પછીના વિધાનમાં તે $!A_f$ છે.
  અહીં એ જ નામ સતત રાખ્યું છે.}
  \[
  \Th{Q} \Proves !A_f(\num{n_0}, \dots, \num{n_k}, \num{m})
  \quad\text{ત્યારે અને માત્ર ત્યારે}\quad m = f(n_0, \dots, n_k).
  \]
\end{lem}

\begin{proof}
  ``જો'' દિશા
  \olref[int]{defn:representable-fn}\olref[int]{defn:rep:a}
  પરથી મળે છે. ``માત્ર જો'' દિશા
  આ પ્રમાણે મળે: માનો કે
  $\Th{Q} \Proves !A_f(\num{n_0},
  \dots, \num{n_k}, \num{m})$,
  પરંતુ $m \neq f(n_0, \dots, n_k)$.
  $l = f(n_0, \dots, n_k)$ લો.
  \olref[int]{defn:representable-fn}\olref[int]{defn:rep:a}
  પ્રમાણે $\Th{Q}
  \Proves !A_f(\num{n_0}, \dots,
  \num{n_k}, \num{l})$.
  \olref[int]{defn:representable-fn}\olref[int]{defn:rep:b}
  પ્રમાણે
  $\lforall[y][(!A_f(\num{n_0}, \dots,
  \num{n_k}, y) \lif \num{l} =
  y)]$. તર્કના નિયમો અને
  $\Th{Q} \Proves
  !A_f(\num{n_0}, \dots,
  \num{n_k}, \num{m})$ની ધારણા વડે
  $\Th{Q} \Proves
  \eq[\num{l}][\num{m}]$ મળે છે.
  બીજી તરફ,
  \olref[bre]{lem:q-proves-neq} પ્રમાણે
  $\Th{Q} \Proves
  \eq/[\num{l}][\num{m}]$.
  એટલે $\Th{Q}$ અસંગત થાય.
  પરંતુ એવું શક્ય નથી, કારણ કે
  પ્રમાણભૂત નિદર્શ $\Th{Q}$ને સંતોષે
  છે (જુઓ
  \olref[int][def]{def:standard-model}),
  $\Sat{N}{\Th{Q}}$; અને યથાર્થતા
  પ્રમેય પ્રમાણે સંતોષી શકાય એવો
  સિદ્ધાંત હંમેશાં સુસંગત હોય છે
  (\tagrefs{prfAX/{fol:axd:sou:cor:consistency-soundness},
  prfSC/{fol:seq:sou:cor:consistency-soundness},
  prfND/{fol:ntd:sou:cor:consistency-soundness},
  prfTab/{fol:tab:sou:cor:consistency-soundness}}).
\end{proof}

\begin{lem}
$\Th{Q}$માં નિરૂપ્ય દરેક વિધેય
સંગણનીય છે.
\end{lem}

\begin{proof}
આ દાવો સાચો કેમ છે તેનો પહેલાં
સહજ વિચાર જોઈએ. $f$ ગણવા માટે
આ પ્રમાણે કરીએ. અંકગણિતની
ભાષામાં શક્ય બધી
નિષ્પત્તિઓ~$\delta$ની યાદી
બનાવીએ. આ યાંત્રિક રીતે શક્ય છે.
દરેક માટે તપાસીએ કે તે
$!A_f(\num{n_0}, \dots,
\num{n_k}, \num{m})$ સ્વરૂપના
સૂત્રની નિષ્પત્તિ
છે કે નહીં ($\Th{Q}$માં~$f$નું
નિરૂપણ કરતું આ સૂત્ર
\olref{lem:rep-q}માંથી મળે છે).
એવું હોય, તો \olref{lem:rep-q}
પ્રમાણે $m = f(n_0, \dots, n_k)$
અને~$f$નું મૂલ્ય મળી ગયું.
શોધ પૂરી થાય છે, કારણ કે
$\Th{Q} \Proves
!A_f(\num{n_0}, \dots,
\num{n_k}, \num{f(n_0, \dots,
n_k)})$; તેથી આખરે યોગ્ય
$\delta$ મળી જ રહેશે.

આ દલીલ હજી સંપૂર્ણ ચોક્કસ નથી:
આપણી પ્રક્રિયા માત્ર સંખ્યાઓને
બદલે નિષ્પત્તિઓ અને
સૂત્રો પર કામ કરે છે;
વળી ``શક્ય બધી નિષ્પત્તિઓની
યાદી'' યાંત્રિક રીતે બનાવી
શકાય તે પણ હજી સમજાવ્યું નથી.
પરંતુ પદો, સૂત્રો અને
નિષ્પત્તિઓને ગડેલ-સંખ્યાઓ
આપી શકાય છે તે આપણે જોયું છે.
સંગણનાનું ચોક્કસ નિદર્શ—સામાન્ય
પુનરાવર્તી વિધેયો—પણ રજૂ કર્યું છે.
અને $\Prf[\Th{Q}](d,y)$ સંબંધ,
જે ત્યારે અને માત્ર ત્યારે સત્ય છે
જ્યારે $d$ એ~$\Th{Q}$નાં
સ્વયંસિદ્ધ વાક્યોમાંથી ગડેલ-સંખ્યા~$y$
ધરાવતા સૂત્રની
નિષ્પત્તિની ગડેલ-સંખ્યા
હોય, (આદિમ) પુનરાવર્તી છે તે પણ
જોયું છે. આગળ જરૂરી અન્ય આદિમ
પુનરાવર્તી વિધેયો $\fn{num}$
(\olref[art][trm]{prop:num-primrec})
અને $\fn{Subst}$
(\olref[art][sub]{prop:subst-primrec}) છે.
આમાંથી લઘુત્તમીકરણ વડે~$f$
વ્યાખ્યાયિત કરી શકાય; તેથી~$f$
પુનરાવર્તી છે.

પહેલાં આ વ્યાખ્યા આપો:
\begin{multline*}
  A(n_0, \dots, n_k, m) = \\
  \fn{Subst}(\fn{Subst}(\dots\fn{Subst}(\Gn{!A_f}, \fn{num}(n_0), \Gn{x_0}),\\ \dots),
  \fn{num}(n_k),  \Gn{x_k}), \fn{num}(m), \Gn{y})
\end{multline*}
આ જટિલ દેખાય છે, પરંતુ તે માત્ર
$A(n_0, \dots, n_k,
m) = \Gn{!A_f(\num{n_0}, \dots,
\num{n_k}, \num{m})}$ વિધેય છે.

હવે $R(n_0, \dots, n_k, s)$
સંબંધ લો. $(s)_0$ એ
$\Th{Q}$માંથી
$!A_f(\num{n_0}, \dots,
\num{n_k}, \num{(s)_1})$ની
નિષ્પત્તિની ગડેલ-સંખ્યા
હોય ત્યારે આ સંબંધ સત્ય છે:
\[
R(n_0, \dots, n_k, s) \quad\text{ત્યારે અને માત્ર ત્યારે}\quad \Prf[\Th{Q}]((s)_0, A(n_0,
\dots, n_k, (s)_1))
\]
$R(n_0, \dots, n_k, s)$ સત્ય થાય
એવું~$s$ શોધીએ, તો સંખ્યાઓની
જોડ~$(s)_0$ અને~$(s)_1$ મળે છે,
જેમાં $(s)_0$ એ
$A_f(\num{n_0}, \dots,
\num{n_k}, \num{(s)_1})$ની
નિષ્પત્તિની ગડેલ-સંખ્યા છે.
\sourcecorrection{OLINC-023}{સ્થિર અંગ્રેજી મૂળના
આ ગદ્યમાં છેલ્લી દલીલ
$(s)_1$ લખાઈ છે, જ્યારે તેના
પહેલાંનું વિધાન અને સંબંધ બંનેમાં
સૂત્રની અંદર અંકનિરૂપણ
$\num{(s)_1}$ જરૂરી છે.
અહીં ગદ્યને એ જ સૂત્ર સાથે
મેળમાં કર્યું છે.}
આમ~$s$ શોધવું એ અનૌપચારિક
સાબિતીમાં~$d$ અને~$m$ની જોડ
શોધવા જેવું છે. એવું~$s$
``શોધતું'' સંગણનીય વિધેય નિયમિત
લઘુત્તમીકરણ વડે વ્યાખ્યાયિત થાય.
યાદ રાખો કે $R$ નિયમિત છે:
દરેક $n_0$, \dots, $n_k$
માટે $\Th{Q} \Proves
!A_f(\num{n_0}, \dots,
\num{n_k}, \num{f(n_0, \dots,
n_k)})$નું કોઈ
નિષ્પત્તિ~$\delta$ છે;
તેથી $s = \tuple{\Gn{\delta},
f(n_0, \dots, n_k)}$ માટે
$R(n_0, \dots, n_k, s)$ સત્ય છે.
એટલે $f$ને આ પ્રમાણે લખી શકીએ:
\[
f(n_0,\dots,n_{k}) = (\umin{s}{R(n_0, \dots, n_k, s)})_1.
\]
\end{proof}
% END OLP-0291
\endgroup


\begingroup
\makeatletter\def\ol@sectioncs{section}\makeatother

% BEGIN OLP-0292
\olfileid{inc}{req}{bet}
\olsection{બીટા વિધેયનું લેમ્મા}
\phantomsection\label{guunit:OLP-0292}



સરવાળો, ગુણાકાર અને $\Char{=}$
ઉપલબ્ધ હોય ત્યારે આદિમ પુનરાવર્તન
કરી શકાય છે તે બતાવવા આપણે શ્રેણીઓ
સંભાળતાં વિધેયો વિકસાવવાં પડશે.
(ઘાતાંકન પણ ઉપલબ્ધ હોત, તો આ કામ
સરળ બનત.) આદિમ પુનરાવર્તન
ઉપલબ્ધ હતું ત્યારે ``$n$મો
અવિભાજ્ય'' જેવા વિધેયો વ્યાખ્યાયિત
કરીને પ્રમાણમાં સીધું સંકેતીકરણ
પસંદ કરી શકતા હતા. પરંતુ અહીં
આદિમ પુનરાવર્તન ઉપલબ્ધ નથી—હકીકતમાં
લઘુત્તમીકરણ વડે આદિમ પુનરાવર્તન
કરી શકાય છે તે જ બતાવવાનું છે—તેથી
વધુ ચતુર રીત જોઈએ.

\begin{lem}
\ollabel{lem:beta}
એવું વિધેય $\beta(d,i)$ છે કે
દરેક શ્રેણી $a_0$, \dots,~$a_n$
માટે એવી સંખ્યા~$d$ મળે છે કે
દરેક $i \le n$ માટે
$\beta(d,i) = a_i$. વળી, $\beta$
મૂળભૂત વિધેયોમાંથી માત્ર સંયોજન
અને નિયમિત લઘુત્તમીકરણ વડે
વ્યાખ્યાયિત કરી શકાય છે.
\end{lem}

$d$ને શ્રેણી $\tuple{a_0, \dots,
a_n}$નો સંકેત અને $\beta(d,i)$ને
તેનો $i$મો ઘટક પરત આપતું વિધેય
સમજો. (ધ્યાન રહે કે આ
``સંકેતીકરણ'' આપણે અગાઉ જાણેલા
અવિભાજ્ય સંખ્યાઓની ઘાતોવાળા સંકેતીકરણનો
ઉપયોગ \emph{કરતું નથી}!) લેમ્મા
ખૂબ મર્યાદિત દાવો કરે છે: શ્રેણીઓ
જોડી શકાય કે ઘટકો ઉમેરી શકાય
એવું તે કહેતું નથી. સંયોજન અને
નિયમિત લઘુત્તમીકરણ વડે વ્યાખ્યાયિત
વિધેયોનો ઉપયોગ કરીને $a_0$,
\dots,~$a_n$માંથી~$d$ \emph{ગણી}
શકાય એવું પણ તે કહેતું નથી.
તે માત્ર એટલું કહે છે કે એવું
``સંકેત ઉકેલતું'' વિધેય છે જેના માટે
દરેક શ્રેણી ``સંકેતિત'' થાય છે.

$\beta$ સંકેતનો ઉપયોગ ગડેલે કર્યો
હતો. ફરી કહીએ તો, લેમ્માની
સાબિતીમાં અઘરું કામ દેખીતી રીતે
મર્યાદિત સાધનો—માત્ર સંયોજન અને
લઘુત્તમીકરણ—વડે યોગ્ય~$\beta$
વ્યાખ્યાયિત કરવાનું છે. જોકે સરવાળો,
ગુણાકાર અને~$\Char{=}$ વાપરી
શકીએ છીએ. આ લેમ્માને સાબિત
કરવાની અનેક રીતો છે; તેમાંની એક
સૌથી સરળ રીત ગડેલની મૂળ રીત છે.
તેમાં સુનઝીનું પ્રમેય કહેવાતું
સંખ્યાસિદ્ધાંતનું તથ્ય વપરાય છે
(પરંપરાગત નામ ``ચીની શેષફળ પ્રમેય'').

\begin{defn}
બે પ્રાકૃતિક સંખ્યાઓ $a$ અને $b$
તેમનો મહત્તમ સામાન્ય અવયવ~$1$
હોય ત્યારે અને માત્ર ત્યારે
\emph{પરસ્પર અવિભાજ્ય} કહેવાય છે;
બીજા શબ્દોમાં, ૧ સિવાયનો
કોઈ અવયવ બંનેમાં સામાન્ય નથી.
\end{defn}

\begin{defn}
પ્રાકૃતિક સંખ્યાઓ $a$ અને $b$
\emph{મોડ્યુલો~$c$ પ્રમાણે સમશેષ}
છે, $a \equiv b \mod c$, ત્યારે
અને માત્ર ત્યારે $c \mid (a-b)$;
એટલે કે~$c$ વડે ભાગતાં
$a$ અને~$b$નું
શેષફળ સમાન મળે છે.
\end{defn}

સુનઝીનું પ્રમેય આ છે:
\begin{thm}
માનો કે $x_0$, \dots,~$x_n$
કોઈપણ બે પરસ્પર અવિભાજ્ય છે.
$y_0$, \dots,~$y_n$ કોઈપણ
સંખ્યાઓ લો. તો એવી સંખ્યા $z$
છે કે
\begin{align*}
z & \equiv y_0 \mod x_0 \\
z & \equiv y_1 \mod x_1 \\
& \vdots  \\
z & \equiv y_n \mod x_n.
\end{align*}
\end{thm}

સુનઝીના પ્રમેયનો ઉપયોગ આમ કરીશું:
જો $x_0$, \dots,~$x_n$ અનુક્રમે
$y_0$, \dots,~$y_n$ કરતાં મોટાં
હોય, તો $z$ને શ્રેણી
$\tuple{y_0, \dots,y_n}$નો
સંકેત લઈ શકાય. $y_i$ પાછું
મેળવવા $z$ને~$x_i$ વડે ભાગીને
શેષફળ લેવું પૂરતું છે. આ
સંકેતીકરણ માટે $x_0$,
\dots,~$x_n$નાં યોગ્ય મૂલ્યો
શોધવાં પડશે.

બે અવલોકનો તેમાં મદદ કરશે.
$y_0$, \dots,~$y_n$ આપેલાં હોય,
તો આ લો:
\begin{align*}
j &= \max(n, y_0 + 1, \dots, y_n + 1), \\
m &= \lcm(1,\dots,j),
\end{align*}
અને પછી આ લો:
\begin{align*}
x_0 & = 1 + m \\
x_1 & = 1 + 2 \cdot m \\
x_2 & = 1 + 3 \cdot m \\
& \vdots  \\
x_n & = 1 + (n+1) \cdot m
\end{align*}
તો બે બાબતો સાચી છે:
\begin{enumerate}
\item\ollabel{rel-prime} $x_0,\dots,x_n$માંથી કોઈપણ બે પરસ્પર અવિભાજ્ય છે.
\item\ollabel{less} દરેક $i$ માટે $y_i < x_i$.
\end{enumerate}
\olref{rel-prime} કેમ સાચું છે તે
જોવા માનો કે $i\ne k$, $p$ અવિભાજ્ય
છે અને
$p \mid x_i$ તથા $p \mid x_k$.
\sourcecorrection{OLINC-024}{સ્થિર અંગ્રેજી
મૂળની આ દલીલમાં $i$ અને~$k$
જુદા છે એમ સ્પષ્ટ કરવું જરૂરી
છે; નહીં તો $i-k=0$ પરથી
$p\mid m$ નીકળતું નથી.
પરસ્પર અવિભાજ્યતાના દાવા
મુજબ અહીં $i\ne k$ લીધું છે.}
તો $p \mid 1 + (i+1) m$ અને
$p \mid 1 + (k+1) m$; તેથી $p$
તેમના તફાવતનો પણ અવયવ છે:
\[
(1 + (i+1)m) - (1+ (k+1)m) = (i-k) m.
\]
$p$ એ $1 + (i+1)m$નો અવયવ
હોવાથી તે~$m$નો અવયવ હોઈ
શકે નહીં (નહીં તો પહેલું પદ~$p$
વડે ભાગતાં શેષ~$1$ રહે).
હવે $p$ એ $(i-k)m$નો અવયવ
છે અને ઉપર પ્રમાણે બીજા
ગુણકનો નથી, તેથી તે
$i-k$નો અવયવ છે. પરંતુ
$\left|i-k\right|$ વધુમાં વધુ~$n$
છે અને $j \geq n$ પસંદ કર્યું છે;
એટલે $p \mid m$, ફરી વિરોધાભાસ.
આથી $x_i$ અને $x_k$ બંનેનો
અવયવ હોય એવી અવિભાજ્ય સંખ્યા નથી.
\olref{less} સરળ છે:
$y_i < j \leq m < x_i$.

હવે $\beta$-વિધેયનું લેમ્મા સાબિત
કરીએ. યાદ રાખો કે $0$, ઉત્તરગામી,
સરવાળો, ગુણાકાર, $\Char{=}$,
પ્રક્ષેપણો અને આમાંથી સંયોજન તથા
નિયમિત વિધેયો પર લાગુ કરેલા
લઘુત્તમીકરણ વડે વ્યાખ્યાયિત
કોઈપણ વિધેય વાપરી શકીએ છીએ.
જે સંબંધનું લાક્ષણિક વિધેય આ
રીતે વ્યાખ્યાયિત થાય તે સંબંધ પણ
વાપરી શકીએ. અગાઉની જેમ,
આ સંબંધો બૂલિયન સંયોજનો અને
સીમિત પરિમાણન હેઠળ બંધ છે તે
બતાવી શકીએ; દાખલા તરીકે:
\begin{align*}
\fn{not}(x) & \defis \Char{=}(x,0)\\
\bmin{x \leq z}{R(x,y)} & \defis \umin{x}{(R(x,y) \lor x = z)}\\
\bexists{x \leq z}{R(x,y)} & \defiff R(\bmin{x \leq z}{R(x,y)}, y)
\end{align*}
પછી આ બધું પણ આદિમ પુનરાવર્તન
વિના વ્યાખ્યાયિત કરી શકાય છે
તે બતાવી શકીએ:
\begin{enumerate}
\item જોડ-નિરૂપણ વિધેય,
  $J(x,y) = \frac{1}{2}[(x+y)(x+y+1)] + x$;
\item પ્રક્ષેપણ વિધેયો
\begin{align*}
K(z) & = \bmin{x \leq z}{\bexists{y \leq z}{z = J(x,y)}},\\
L(z) & = \bmin{y \leq z}{\bexists{x \leq z}{z = J(x,y)}};
\end{align*}
\item ન્યૂનતા સંબંધ $x < y$;
\item વિભાજ્યતા સંબંધ $x \mid y$;
\item $y$ને~$x$ વડે ભાગતાં
  મળતું શેષફળ આપતું વિધેય
  $\fn{rem}(x,y)$.
\end{enumerate}
હવે આ વ્યાખ્યાયિત કરો:
\begin{align*}
\beta^*(d_0,d_1,i) & = \fn{rem}(1+(i+1) d_1,d_0) \text{અને}\\
\beta(d,i) & = \beta^*(K(d),L(d),i).
\end{align*}
આ જ આપણને જોઈતું વિધેય છે.
ઉપરની જેમ $a_0,\dots,a_n$
આપેલાં હોય, તો આ લો:
\[
j = \max(n,a_0+1,\dots,a_n+1),
\]
અને $d_1 = \lcm(1,\dots,j)$ લો.
ઉપરના \olref{rel-prime} પ્રમાણે
$1+d_1$, $1+2 d_1$, \dots,
$1+(n+1) d_1$ જોડી-જોડી
પરસ્પર અવિભાજ્ય છે; અને~\olref{less}
પ્રમાણે તે અનુક્રમે
$a_0,\dots,a_n$ કરતાં મોટાં છે.
સુનઝીના પ્રમેય પ્રમાણે એવું~$d_0$
છે કે દરેક~$i$ માટે
\[
d_0 \equiv a_i \mod (1+(i+1)d_1)
\]
અને તેથી ($d_1$ એ~$a_i$
કરતાં મોટું હોવાથી)
\[
a_i = \fn{rem}(1+(i+1)d_1,d_0).
\]
હવે $d = J(d_0,d_1)$ લો.
તો દરેક $i \le n$ માટે
\begin{align*}
\beta(d,i) & =  \beta^*(d_0,d_1,i) \\
& =  \fn{rem}(1+(i+1) d_1,d_0) \\
& =  a_i
\end{align*}
મળે છે, અને એ જ જરૂરી હતું.
આથી $\beta$-વિધેયના લેમ્માની
સાબિતી પૂરી થાય છે.

\begin{prob}
  સંબંધો $x < y$ અને $x \mid y$
  તથા વિધેય~$\fn{rem}(x,y)$
  આદિમ પુનરાવર્તન વિના
  વ્યાખ્યાયિત કરી શકાય છે તે
  બતાવો. તમે $0$, ઉત્તરગામી,
  સરવાળો, ગુણાકાર, $\Char{=}$,
  પ્રક્ષેપણો તથા સીમિત
  લઘુત્તમીકરણ અને પરિમાણન
  વાપરી શકો છો.
\end{prob}
% END OLP-0292
\endgroup


\begingroup
\makeatletter\def\ol@sectioncs{section}\makeatother

% BEGIN OLP-0293
\olfileid{inc}{req}{pri}
\olsection{આદિમ પુનરાવર્તનનું અનુસરણ}
\phantomsection\label{guunit:OLP-0293}


હવે બીટા વિધેય વડે આદિમ
પુનરાવર્તનથી થતી વ્યાખ્યાનું
નિયમિત લઘુત્તમીકરણથી ``અનુસરણ''
કરી શકાય છે તે બતાવી શકીએ.
માનો કે આપણી પાસે $f(\vec x)$
અને $g(\vec x, y, z)$ છે.
તો $f$ અને~$g$માંથી આદિમ
પુનરાવર્તન વડે વ્યાખ્યાયિત વિધેય
$h(\vec x,y)$ આ છે:
\sourcecorrection{OLINC-025}{સ્થિર અંગ્રેજી
મૂળના આ ગદ્યમાં $h(x,\vec z)$
લખાયું છે, જ્યારે તરત પછીનાં
બંને સમીકરણો અને બાકીની
સાબિતીમાં $h(\vec x,y)$ છે.
અહીં ચલોના એ જ ક્રમનો ઉપયોગ
કર્યો છે.}
\begin{align*}
h(\vec x, 0) & =  f(\vec x) \\
h(\vec x, y+1) & =  g(\vec x, y, h(\vec x, y)).
\end{align*}
$f$ અને~$g$માંથી માત્ર સંયોજન
અને નિયમિત લઘુત્તમીકરણ વડે~$h$
વ્યાખ્યાયિત થાય છે તે બતાવવું છે.
એ માટે મૂળભૂત વિધેયો અને
તેમનામાંથી સંયોજન તથા નિયમિત
લઘુત્તમીકરણ વડે વ્યાખ્યાયિત
વિધેયો (જેમ કે~$\beta$) વાપરીશું.

\begin{lem}
\ollabel{lem:prim-rec}
$h$ને $f$ અને $g$માંથી આદિમ
પુનરાવર્તન વડે વ્યાખ્યાયિત
કરી શકાય, તો $f$, $g$ અને
વિધેયો $\Zero$, $\Succ$,
$\Proj{n}{i}$, $\Add$, $\Mult$,
$\Char{=}$માંથી સંયોજન અને
નિયમિત લઘુત્તમીકરણ વડે પણ
વ્યાખ્યાયિત કરી શકાય છે.
\end{lem}

\begin{proof}
પહેલાં સહાયક વિધેય $\hat h(\vec x, y)$
વ્યાખ્યાયિત કરો. તે નીચેની
શરતો સંતોષતી શ્રેણીને સંકેતિત
કરતી ન્યૂનતમ સંખ્યા~$d$ આપે છે;
આ~$d$ શ્રેણીનો સંકેત છે:
\begin{enumerate}
\item $(d)_0 = f(\vec x)$, અને
\item દરેક $i < y$ માટે $(d)_{i+1} = g(\vec x, i, (d)_i)$,
\end{enumerate}
અહીં હવે $(d)_i$ એ $\beta(d,i)$નું
ટૂંકું લખાણ છે. બીજા શબ્દોમાં,
$\hat h$ એ
$\tuple{h(\vec x, 0), h(\vec x, 1),
\dots, h(\vec x, y)}$ શ્રેણીનો
સંકેત આપે છે.
\sourcecorrection{OLINC-026}{સ્થિર અંગ્રેજી
મૂળમાં $\hat h$ શ્રેણી પોતે
આપે છે એવું કહેવાયું છે;
પરંતુ અગાઉની વ્યાખ્યા મુજબ
તે શ્રેણીનો સંકેત~$d$ આપે છે.
અહીં એ ભેદ સ્પષ્ટ કર્યો છે.}
$\hat h$ને આ પ્રમાણે લખી શકીએ:
\[
\hat h(\vec x, y) = \umin{d}{(\beta(d,0) = f(\vec x) \land \bforall{i <
  y}{\beta(d,i+1) = g(\vec x, i,\beta(d,i)})}.
\]
ધ્યાન આપો કે અહીં આદિમ
પુનરાવર્તન જરૂરી નથી, માત્ર
લઘુત્તમીકરણ પૂરતું છે. બીટા
વિધેયના લેમ્મા
\olref[bet]{lem:beta}ને કારણે
લઘુત્તમીકરણ હેઠળનું વિધેય
નિયમિત છે.

હવે, આ મળે છે:
\[
h(\vec x, y) = \beta(\hat h(\vec x, y), y),
\]
આથી~$h$ને મૂળભૂત વિધેયોમાંથી
માત્ર સંયોજન અને નિયમિત
લઘુત્તમીકરણ વડે વ્યાખ્યાયિત
કરી શકાય છે.
\end{proof}
% END OLP-0293
\endgroup


\begingroup
\makeatletter\def\ol@sectioncs{section}\makeatother

% BEGIN OLP-0294
\olfileid{inc}{req}{bre}
\olsection{મૂળભૂત વિધેયો~$\Th{Q}$માં નિરૂપ્ય છે}
\phantomsection\label{guunit:OLP-0294}


પહેલાં બધા મૂળભૂત વિધેયો
$\Th{Q}$માં નિરૂપ્ય છે તે બતાવવું
પડશે. છેવટે દરેક $k$-ચલવાળા
મૂળભૂત વિધેય $f(x_0,\dots,x_{k-1})$
માટે તેનું નિરૂપણ કરતું
સૂત્ર
$!A_f(x_0,\dots,x_{k-1},y)$
કેવી રીતે આપવું તે બતાવવાનું છે.

શૂન્ય વિધેય, ઉત્તરગામી વિધેય,
સરવાળો, ગુણાકાર, સમતા માટેનું
લાક્ષણિક વિધેય અને પ્રક્ષેપણોનું
નિરૂપણ કરી શકીશું. દરેક કિસ્સામાં
યોગ્ય નિરૂપક સૂત્ર સીધું છે:
દાખલા તરીકે, શૂન્ય વિધેયને
$y = \Obj 0$ સૂત્ર, ઉત્તરગામીને
$x_0' = y$ સૂત્ર, અને
સરવાળાને $(x_0 + x_1) = y$
સૂત્ર નિરૂપે છે. કામ એ
બતાવવાનું છે કે સંબંધિત
વાક્યો~$\Th{Q}$માં
સાબિત થાય છે. દાખલા તરીકે,
ઉપરનું સૂત્ર સરવાળાને
નિરૂપે છે એમ કહેવા માટે દરેક
પ્રાકૃતિક સંખ્યાની જોડ $m$ અને
$n$ માટે $\Th{Q}$માં આ બંને
સાબિત થાય છે તે બતાવવું પડે:
\begin{align*}
& \eq[\num n + \num m][\num {n+m}] \text{અને}\\
& \lforall[y][(\eq[(\num n + \num m)][y] \lif \eq[y][\num{n+m}])].
\end{align*}

\begin{prop}
\ollabel{prop:rep-zero}
શૂન્ય વિધેય $\Zero(x) = 0$ને
$!A_{\Zero}(x,y) \ident \eq[y][\Obj 0]$
સૂત્ર~$\Th{Q}$માં નિરૂપે છે.
\end{prop}

\begin{prop}
\ollabel{prop:rep-succ}
ઉત્તરગામી વિધેય $\Succ(x) = x+1$ને
$!A_{\Succ}(x,y) \ident \eq[y][x']$
સૂત્ર~$\Th{Q}$માં નિરૂપે છે.
\end{prop}

\begin{prop}
\ollabel{prop:rep-proj}
પ્રક્ષેપણ વિધેય
$\Proj{n}{i}(x_0, \dots, x_{n-1}) = x_i$ને
$\Th{Q}$માં આ સૂત્ર નિરૂપે છે:
\[!A_{\Proj{n}{i}}(x_0, \dots, x_{n-1}, y) \ident \eq[y][x_i].\]
\end{prop}

\begin{prob}
$\eq[y][\Obj 0]$, $\eq[y][x']$
અને $\eq[y][x_i]$ અનુક્રમે
$\Zero$, $\Succ$ અને
$\Proj{n}{i}$ને નિરૂપે છે
તે સાબિત કરો.
\end{prob}

\begin{prop}
\ollabel{prop:rep-id}
સમતા~$=$નું લાક્ષણિક વિધેય
\[
\Char{=}(x_0, x_1) =
\begin{cases}
  1 & \text{જો } x_0 =x_1\\
  0 & \text{અન્યથા}
\end{cases}
\]
ને~$\Th{Q}$માં આ સૂત્ર નિરૂપે છે:
\[
  !A_{\Char{=}}(x_0, x_1, y) \ident (\eq[x_0][x_1] \land \eq[y][\num{1}]) \lor (\eq/[x_0][x_1] \land
\eq[y][\num{0}]).
\]
\end{prop}

સાબિતી માટે આ લેમ્મા જરૂરી છે.

\begin{lem}
\ollabel{lem:q-proves-neq}
પ્રાકૃતિક સંખ્યાઓ $n$ અને $m$
માટે, જો $n \neq m$ હોય, તો
$\Th{Q} \Proves \eq/[\num n][\num m]$.
\end{lem}

\begin{proof}
$n$ પર આગમન વડે સાબિત કરો કે
દરેક $m$ માટે, $n \neq m$ હોય
તો $Q \Proves \eq/[\num n][\num m]$.

પાયાના કિસ્સામાં $n = 0$.
$m$ એ $0$ની બરાબર ન હોય, તો
કોઈ પ્રાકૃતિક સંખ્યા $k$ માટે
$m = k + 1$. આપણી પાસે
$\lforall[x][\eq/[0][x']]$
કહેતું સ્વયંસિદ્ધ વાક્ય છે.
પરિમાણકના સ્વયંસિદ્ધ વાક્યથી
$x$ને $\num k$ વડે બદલીને
$\eq/[0][\num k']$ તારવી
શકીએ. પરંતુ $\num k'$ એ જ
$\num m$ છે.

આગમનના પગથિયામાં $n$ માટે
દાવો સાચો છે એમ માનીને $n+1$
વિચારીએ. $m$ કોઈપણ પ્રાકૃતિક
સંખ્યા લો. બે શક્યતાઓ છે:
$m = 0$, અથવા કોઈ $k$ માટે
$m = k+1$. પહેલો કિસ્સો ઉપર
પ્રમાણે ઉકેલાય છે. બીજા કિસ્સામાં
માનો કે $n+1 \neq k+1$.
તો $n \neq k$. $n$ માટેની
આગમન-ધારણાથી
$\Th{Q} \Proves \eq/[\num n][\num k]$.
આપણી પાસે આ સ્વયંસિદ્ધ વાક્ય છે:
$\lforall[x][\lforall[y][\eq[x'][y'] \lif \eq[x][y]]]$.
પરિમાણકના સ્વયંસિદ્ધ વાક્યથી
$\eq[\num n'][\num k'] \lif \eq[\num n][\num
  k]$ મળે છે. વિધાનાત્મક તર્કથી
$\Th{Q}$માં
$\eq/[\num n][\num k] \lif \eq/[\num n'][\num k']$
તારવી શકીએ. મોડસ પોનેન્સથી
$\eq/[\num n'][\num k']$ મળે,
અને એ જ જોઈએ છે, કારણ કે
$\num k'$ એ $\num m$ છે.
\end{proof}

\begin{explain}
લેમ્માનો દાવો બહુ મોટો નથી:
તેનો સાર એટલો છે કે જુદાં
અંકનિરૂપણો જુદા પદાર્થો દર્શાવે
છે તે $\Th{Q}$ સાબિત કરી શકે છે.
દાખલા તરીકે, $\Th{Q}$ આ સાબિત
કરે છે: $0'' \neq 0'''$.
પરંતુ આ સામાન્ય રીતે સાચું છે
તે બતાવવા થોડી કાળજી જોઈએ.
આગમન વાપરીએ છીએ, પણ એ
$\Th{Q}$ની \emph{બહાર}નું
આગમન છે, એ પણ ધ્યાન રાખો.
\end{explain}

\begin{proof}[\olref{prop:rep-id}ની સાબિતી]
$n = m$ હોય, તો $\num{n}$ અને
$\num{m}$ એક જ પદ છે અને
$\Char{=}(n, m) = 1$.
પરંતુ $\Th{Q} \Proves
(\eq[\num{n}][\num{m}] \land
\eq[\num{1}][\num{1}])$,
તેથી $!A_=(\num{n}, \num{m},
\num{1})$ પણ સાબિત થાય છે.
$n \neq m$ હોય, તો
$\Char=(n, m) = 0$.
\olref{lem:q-proves-neq} પ્રમાણે
$\Th{Q} \Proves \eq/[\num{n}][\num{m}]$;
તેથી
$(\eq/[\num{n}][\num{m}] \land \Obj 0 = \Obj 0)$
પણ સાબિત થાય છે. આમ
$\Th{Q} \Proves !A_=(\num{n},
\num{m}, \num{0})$.

બીજા ભાગમાં પણ બે કિસ્સા છે.
$n = m$ હોય, તો
$\Th{Q} \Proves \lforall[y][(!A_=(\num{n}, \num{m}, y)
  \lif \eq[y][\num{1}])]$
બતાવવાનું છે. અનૌપચારિક રીતે
દલીલ કરીએ: માનો કે
$!A_=(\num{n}, \num{m}, y)$,
એટલે કે
\[
(\eq[\num{n}][\num{n}] \land \eq[y][\num{1}]) \lor
(\eq/[\num{n}][\num{n}] \land \eq[y][\num{0}])
\]
ડાબો વિકલ્પ તર્કથી
$\eq[y][\num{1}]$ આપે છે;
જમણો વિકલ્પ તર્કથી સાબિત થતા
$\eq[\num{n}][\num{n}]$ સાથે
વિરોધાભાસમાં છે.

બીજી તરફ, માનો કે $n \neq m$.
ત્યારે $!A_=(\num{n},
\num{m}, y)$ આ છે:
\[
(\eq[\num{n}][\num{m}] \land \eq[y][\num{1}]) \lor
(\eq/[\num{n}][\num{m}] \land \eq[y][\num{0}])
\]
અહીં ડાબો વિકલ્પ
$\eq/[\num{n}][\num{m}]$ સાથે
વિરોધાભાસમાં છે; આ વિધાન
$\Th{Q}$માં
\olref{lem:q-proves-neq}થી
સાબિત થાય છે. જમણા
વિકલ્પમાંથી $\eq[y][\num{0}]$
નીકળે છે.
\end{proof}

\begin{prop}
\ollabel{prop:rep-add}
સરવાળાનું વિધેય
$\Add(x_0, x_1) = x_0+x_1$ને
$\Th{Q}$માં આ સૂત્ર નિરૂપે છે:
\[
  !A_{\Add}(x_0, x_1, y) \ident \eq[y][(x_0 + x_1)].
\]
\end{prop}

\begin{lem}
\ollabel{lem:q-proves-add}
$\Th{Q} \Proves \eq[(\num{n} + \num{m})][\num{n+m}]$
\end{lem}

\begin{proof}
આ દાવો~$m$ પર આગમનથી
સાબિત કરીએ. $m = 0$ હોય, તો
$\Th{Q} \Proves
\eq[(\num{n} + \Obj 0)][\num{n}]$
બતાવવાનું છે. તે સ્વયંસિદ્ધ
વાક્ય~$!Q_4$થી મળે છે.
હવે $m$ માટે દાવો માનીને
$m+1$ માટે, એટલે કે
$\Th{Q} \Proves \eq[(\num{n} +
  \num{m+1})][\num{n+m+1}]$,
સાબિત કરીએ. ધ્યાન આપો કે
$\num{m+1}$ એ $\num{m}'$
અને $\num{n+m+1}$ એ
$\num{n+m}'$ જ છે.
સ્વયંસિદ્ધ વાક્ય $!Q_5$થી
$\Th{Q} \Proves
\eq[(\num{n} + \num{m}')][(\num{n}+\num{m})']$.
આગમન-ધારણાથી
$\Th{Q} \Proves
\eq[(\num{n} + \num{m})][\num{n+m}]$.
તેથી $\Th{Q} \Proves
\eq[(\num{n} + \num{m}')][\num{n+m}']$.
\end{proof}

\begin{proof}[\olref{prop:rep-add}ની સાબિતી]
$\Add$નું નિરૂપણ કરતું
સૂત્ર~$!A_\Add(x_0, x_1, y)$
એ $\eq[y][(x_0 + x_1)]$ છે.
પહેલાં બતાવીએ કે
$\Add(n, m) = k$ હોય, તો
$\Th{Q} \Proves
!A_\Add(\num{n}, \num{m}, \num{k})$,
એટલે કે $\Th{Q} \Proves
\eq[\num{k}][(\num{n} + \num{m})]$.
પરંતુ $k = n + m$ હોવાથી
$\num{k}$ એ $\num{n+m}$ જ છે;
અને \olref{lem:q-proves-add}માં
$\Th{Q} \Proves \eq[(\num{n} +
  \num{m})][\num{n+m}]$
પહેલેથી બતાવ્યું છે.

$\Add(n, m) = k$ હોય, તો આ પણ
બતાવવું પડે:
\[
\Th{Q} \Proves \lforall[y][(!A_\Add(\num{n}, \num{m}, y) \lif
  \eq[y][\num{k}])].
\]
માનો કે
$\eq[(\num{n} + \num{m})][y]$ છે.
કારણ કે
\[
\Th{Q} \Proves \eq[(\num{n}+\num{m})][\num{n+m}],
\]
ડાબી બાજુને $\num{n+m}$ વડે
બદલતાં કોઈપણ~$y$ માટે
$\eq[\num{n+m}][y]$ મળે છે.
\end{proof}

\begin{prop}
\ollabel{prop:rep-mult}
ગુણાકારનું વિધેય
$\Mult(x_0, x_1) = x_0 \cdot x_1$ને
$\Th{Q}$માં આ સૂત્ર નિરૂપે છે:
\[
  !A_{\Mult}(x_0, x_1, y) \ident y = (x_0 \times x_1).
\]
\end{prop}

\begin{proof} અભ્યાસપ્રશ્ન. \end{proof}

\begin{lem}
\ollabel{lem:q-proves-mult}
$\Th{Q} \Proves \eq[(\num{n} \times \num{m})][\num{n \cdot m}]$
\end{lem}

\begin{proof} અભ્યાસપ્રશ્ન. \end{proof}

\begin{prob}
\olref[inc][req][bre]{lem:q-proves-mult}
સાબિત કરો.
\end{prob}

\begin{prob}
\olref[inc][req][bre]{lem:q-proves-mult}
વાપરીને
\olref[inc][req][bre]{prop:rep-mult}
સાબિત કરો.
\end{prob}

\begin{explain}
યાદ રાખો કે અંકગણિતની
ભાષામાં વિધેય-ચિહ્ન માટે
$\times$ અને સંખ્યાઓ પરની
સામાન્ય ગુણાકારક્રિયા માટે
$\cdot$ વાપરીએ છીએ.
તેથી સંખ્યાઓ દર્શાવતી
અભિવ્યક્તિઓની વચ્ચે
$\cdot$ આવી શકે છે (જેમ કે
$m \cdot n$માં), જ્યારે
$\times$ માત્ર અંકગણિતની
ભાષાનાં પદોની વચ્ચે આવે છે
(જેમ કે $(\num{m} \times
  \num{n})$માં). વાત વધુ
ગૂંચવાય છે, કારણ કે $+$
બંને માટે વપરાય છે:
ફલન માટે પણ અને
સરવાળાની ક્રિયા માટે પણ.
તે પદોની વચ્ચે આવે—દાખલા
તરીકે $(\num{n} + \num{m})$માં—તો
તે અંકગણિતની ભાષાનું $2$-સ્થાની
ફલન છે; અને સંખ્યાઓની
વચ્ચે આવે—દાખલા તરીકે
$n+m$માં—તો તે સરવાળાની
ક્રિયા છે. $\num{n+m}$નો
કિસ્સો પણ આમાં આવે છે:
તે સંખ્યા~$n+m$ને અનુરૂપ
પ્રમાણભૂત અંકનિરૂપણ છે.
\end{explain}
% END OLP-0294
\endgroup


\begingroup
\makeatletter\def\ol@sectioncs{section}\makeatother

% BEGIN OLP-0295
\olfileid{inc}{req}{cmp}
\olsection{સંયોજન~$\Th{Q}$માં નિરૂપ્ય છે}
\phantomsection\label{guunit:OLP-0295}


માનો કે $h$ આ પ્રમાણે
વ્યાખ્યાયિત છે:
\[
h(x_0,\dots,x_{l-1}) = f(g_0(x_0,\dots,x_{l-1}), \dots,
g_{k-1}(x_0,\dots,x_{l-1})).
\]
અહીં $f$ અને અનુક્રમે $g_0$,
\dots,~$g_{k-1}$ને નિરૂપતાં
સૂત્રો $!A_f, !A_{g_0}, \dots,
!A_{g_{k-1}}$ આપણે અગાઉથી શોધ્યાં
છે. હવે~$h$ને નિરૂપતું
સૂત્ર~$!A_h$ શોધવાનું છે.

પહેલાં બધાં વિધેયો $1$-સ્થાની
હોય એવો સરળ કિસ્સો લો, એટલે
$h(x) = f(g(x))$ વિચારો.
$!A_f(y, z)$ એ~$f$ને અને
$!A_g(x, y)$ એ~$g$ને નિરૂપે,
તો~$h$ને નિરૂપતું
સૂત્ર~$!A_h(x, z)$ જોઈએ.
ધ્યાન આપો કે $h(x) = z$
ત્યારે અને માત્ર ત્યારે જ એવું~$y$
છે કે $z = f(y)$ અને
$y = g(x)$ બંને સાચાં છે.
($h(x) = z$ હોય, તો $g(x)$
આવું~$y$ છે; અને આવું $y$
હોય, તો $y = g(x)$ તથા
$z = f(y)$ હોવાથી $z
= f(g(x))$.) તેથી
$\lexists[y][(!A_g(x, y) \land !A_f(y,
  z))]$ એ~$!A_h(x, z)$ માટે
યોગ્ય ઉમેદવાર લાગે છે.
હવે સંબંધિત સૂત્રો
$\Th{Q}$માં સાબિત થાય છે
તેની ખાતરી કરવી બાકી છે.

\begin{prop}
\ollabel{prop:rep1}
$h(n) = m$ હોય, તો
$\Th{Q} \Proves !A_h(\num{n}, \num{m})$.
\end{prop}

\begin{proof}
માનો કે $h(n) = m$, એટલે કે
$f(g(n)) = m$. $k = g(n)$ લો.
તો
\begin{align*}
  \Th{Q} & \Proves !A_g(\num{n}, \num{k})
  \intertext{કારણ કે $!A_g$ એ~$g$ને નિરૂપે છે, અને}
  \Th{Q} & \Proves !A_f(\num{k}, \num{m})
  \intertext{કારણ કે $!A_f$ એ~$f$ને નિરૂપે છે. તેથી}
  \Th{Q} & \Proves !A_g(\num{n}, \num{k}) \land !A_f(\num{k}, \num{m})
  \intertext{અને પરિણામે}
  \Th{Q} & \Proves \lexists[y][(!A_g(\num{n}, y) \land !A_f(y, \num{m}))],
\end{align*}
એટલે કે $\Th{Q} \Proves
!A_h(\num{n}, \num{m})$.
\end{proof}

\begin{prop}
\ollabel{prop:rep2}
$h(n) = m$ હોય, તો
$\Th{Q} \Proves \lforall[z][(!A_h(\num{n}, z) \lif
  z = \num{m})]$.
\end{prop}

\begin{proof}
માનો કે $h(n) = m$, એટલે કે
$f(g(n)) = m$. $k = g(n)$ લો.
તો
\begin{align*}
  \Th{Q} & \Proves \lforall[y][(!A_g(\num{n}, y) \lif \eq[y][\num{k}])]
  \intertext{કારણ કે $!A_g$ એ~$g$ને નિરૂપે છે, અને}
  \Th{Q} & \Proves \lforall[z][(!A_f(\num{k}, z) \lif \eq[z][\num{m}])]
  \intertext{કારણ કે $!A_f$ એ~$f$ને નિરૂપે છે. થોડો તર્ક વાપરીને આ પણ બતાવી શકીએ:}
  \Th{Q} & \Proves \lforall[z][(\lexists[y][(!A_g(\num{n}, y) \land
      !A_f(y, z))] \lif \eq[z][\num{m}])].
\end{align*}
એટલે કે $\Th{Q} \Proves
\lforall[y][(!A_h(\num n, y) \lif \eq[y][\num m])]$.
\end{proof}

$f$ અને~$g_i$ની સ્થાનસંખ્યા
$1$ કરતાં વધુ હોય એવા જટિલ
કિસ્સામાં પણ આ જ વિચાર કામ કરે છે.

\begin{prop}
\ollabel{prop:rep-composition}
$!A_f(y_0, \dots, y_{k-1}, z)$
એ~$f(y_0, \dots, y_{k-1})$ને
$\Th{Q}$માં નિરૂપે, અને
$!A_{g_i}(x_0, \dots, x_{l-1}, y)$
એ $g_i(x_0, \dots, x_{l-1})$ને
$\Th{Q}$માં નિરૂપે, તો આ સૂત્ર
\begin{multline*}
  \lexists[y_0]\dots\lexists[y_{k-1}]
  (!A_{g_0}(x_0,\dots,x_{l-1},y_0) \land
      \dots \land {}\\
  !A_{g_{k-1}}(x_0,\dots,x_{l-1},y_{k-1}) \land !A_f(y_0,\dots,y_{k-1},z))
\end{multline*}
\sourcecorrection{OLINC-027}{સ્થિર અંગ્રેજી
મૂળના આ પ્રદર્શિત સૂત્રમાં
પહેલો પરિમાણક $y_0\dots$ને
એક જ ચલ તરીકે લે છે અને
છેલ્લો સંયોજક અંદરના પરિમાણકની
હદ બહાર રહી જાય છે. અહીં
$y_0$થી~$y_{k-1}$ સુધીના
અસ્તિત્વ-પરિમાણકો તથા આખા
સંયોજનની હદ સ્પષ્ટ કરી છે.}
આ વિધેયને નિરૂપે છે:
\[
h(x_0, \dots, x_{l-1}) = f(g_0(x_0, \dots, x_{l-1}), \dots, g_{k-1}(x_0,
\dots, x_{l-1})).
\]
\end{prop}

\begin{proof}
અભ્યાસપ્રશ્ન.
\end{proof}

\begin{prob}
\olref[inc][req][cmp]{prop:rep1}
અને \olref[inc][req][cmp]{prop:rep2}ની
સાબિતીઓને માર્ગદર્શક બનાવી
\olref[inc][req][cmp]{prop:rep-composition}ની
સાબિતી વિગતે પૂર્ણ કરો.
\sourcecorrection{OLINC-028}{સ્થિર અંગ્રેજી
મૂળના આ અભ્યાસપ્રશ્નમાં
એક જ વિધાનનો સંદર્ભ બે વાર છે;
અગાઉની બે પૂરક સાબિતીઓમાંથી
પહેલી અને બીજી બંનેનો સંદર્ભ
અહીં આપ્યો છે.}
\end{prob}
% END OLP-0295
\endgroup


\begingroup
\makeatletter\def\ol@sectioncs{section}\makeatother

% BEGIN OLP-0296
\olfileid{inc}{req}{min}
\olsection{નિયમિત લઘુત્તમીકરણ~$\Th{Q}$માં નિરૂપ્ય છે}
\phantomsection\label{guunit:OLP-0296}


હવે અમર્યાદિત શોધનો વિચાર કરીએ.
માનો કે $g(x, z)$ નિયમિત છે અને
સૂત્ર~$!A_g(x, z, y)$ વડે
$\Th{Q}$માં નિરૂપ્ય છે.
માનો કે $f$ને
$f(z) = \umin{x}{[g(x, z) = 0]}$
દ્વારા વ્યાખ્યાયિત કરીએ.
$f$ને નિરૂપતું
સૂત્ર~$!A_f(z, y)$ શોધવું છે.
$f(z)$નું મૂલ્ય એવી સંખ્યા~$x$
છે જે (a) $g(x, z) = 0$ સંતોષે
છે અને (b) એવી સંખ્યાઓમાં સૌથી
નાની છે, એટલે દરેક $w < x$
માટે $g(w, z) \neq 0$ છે.
તેથી આ સ્વાભાવિક પસંદગી છે:
\[
!A_f(z,y) \ident !A_g(y, z, \Obj 0) \land \lforall[w][(w < y \lif \lnot
  !A_g(w, z, \Obj 0))].
\]
સામાન્ય કિસ્સામાં, અલબત્ત,
$z$ને $z_0$, \dots, $z_k$થી
બદલવું પડશે.

સાબિતીમાં ફરીથી એ બાબતનાં
કેટલાંક લેમ્મા આવશે કે
$\Th{Q}$ શું સાબિત કરવા જેટલું
સમર્થ છે.

\begin{lem}
\ollabel{lem:succ} દરેક
અચળ~$a$ અને દરેક
પ્રાકૃતિક સંખ્યા~$n$ માટે,
\[
\Th{Q} \Proves \eq[(a' + \num n)][(a + \num n)'].
\]
\end{lem}

\begin{proof}
હંમેશની જેમ, $n$ પર અનુમાન
વડે સાબિતી કરીએ. આરંભિક કિસ્સામાં
$n = 0$ છે; એટલે $\Th{Q}$માં
$\eq[(a' + \Obj 0)][(a + \Obj
0)']$ સાબિત કરવાનું છે.
આપણી પાસે આ છે:
\begin{align}
  \Th{Q} & \Proves \eq[(a' + \Obj 0)][a'] \quad \text{સ્વયંસિદ્ધિ $Q_4$ વડે}
  \ollabel{step1}\\
  \Th{Q} & \Proves  \eq[(a + \Obj 0)][a] \quad \text{સ્વયંસિદ્ધિ $Q_4$ વડે}
  \ollabel{step2} \\
  \Th{Q} & \Proves \eq[(a + \Obj 0)'][a'] \quad \text{\olref{step2} વડે}
  \ollabel{step3} \\
  \Th{Q} & \Proves \eq[(a' + \Obj 0)][(a + \Obj 0)'] \quad
  \text{\olref{step1} અને \olref{step3} વડે}\notag
\end{align}
અનુમાનના પગલામાં ધારો કે
$\Th{Q} \Proves \eq[(a' + \num n)][(a + \num n)']$
બતાવી ચૂક્યા છીએ. કારણ કે
$\num{n+1}$ એટલે $\num{n}'$,
હવે $\Th{Q}$માં
$\eq[(a' + \num{n}')][(a + \num{n}')']$
સાબિત કરવાનું છે. આપણે મેળવીએ:
\begin{align}
  \Th{Q} & \Proves \eq[(a' + \num n')][(a' + \num n)'] \quad
  \text{સ્વયંસિદ્ધિ $!Q_5$ વડે} \ollabel{step5}\\
  \Th{Q} & \Proves \eq[(a' + \num n)'][(a + \num n')'] \quad
  \text{અનુમાનની ધારણા અને સ્વયંસિદ્ધિ $!Q_5$ વડે} \ollabel{step6}\\
  \Th{Q} & \Proves \eq[(a' + \num n')][(a + \num n')'] \quad
  \text{\olref{step5} અને \olref{step6} વડે.} \notag
\end{align}
\sourcecorrection{OLINC-029}{સ્થિર અંગ્રેજી
મૂળના અનુમાન-પગલાની બીજી પંક્તિમાં
સાબિત કરવાનું અંતિમ સમીકરણ જ
અનુમાનની ધારણા તરીકે લખાયું છે;
છેલ્લી પંક્તિમાં મધ્યવર્તી સમીકરણ
આવી જાય છે. અહીં અનુમાનની
ધારણા તથા $Q_5$ વડે મળતું
મધ્યવર્તી સમીકરણ બીજી પંક્તિમાં
અને માંગેલું તારણ છેલ્લી પંક્તિમાં
આપ્યું છે.}
\end{proof}

ફરી નોંધવું જોઈએ કે આ દાવો
$\Th{Q}$ આ
$\lforall[x][\lforall[y][(x' + y) = (x + y)']]$
સાબિત કરે છે એવા દાવા કરતાં
નબળો છે. આ વાક્ય
$\Struct{N}$માં સાચું હોવા છતાં
$\Th{Q}$ તેને સાબિત કરતું નથી.

\begin{lem}
\ollabel{lem:less-zero}
$\Th{Q} \Proves \lforall[x][\lnot x < \Obj 0]$.
\end{lem}

\begin{proof}
સાબિતી અનૌપચારિક રીતે આપીએ
(એટલે ઔપચારિક નિષ્પત્તિ
કેવી રીતે રચવું તેના માત્ર સંકેતો
આપીએ).

કોઈપણ~$a$ માટે $\lnot a < \Obj 0$
સાબિત કરવાનું છે. $<$ની
વ્યાખ્યા પ્રમાણે $\Th{Q}$માં
$\lnot \lexists[y][\eq[(y'+a)][\Obj 0]]$
સાબિત કરવાનું છે. આપણે
$\lexists[y][\eq[(y'+a)][\Obj 0]]$
ધારીને વિરોધાભાસ મેળવીશું.
માનો કે $\eq[(b'+a)][\Obj 0]$.
$!Q_3$ વડે
$\eq[a][\Obj 0] \lor \lexists[y][\eq[a][y']]$
મળે છે. હવે બે કિસ્સા જોઈએ.

કિસ્સો 1: $\eq[a][\Obj 0]$.
$\eq[(b'+a)][\Obj 0]$માંથી
$\eq[(b' + \Obj 0)][\Obj 0]$
મળે છે. $\Th{Q}$ની સ્વયંસિદ્ધિ
$!Q_4$ વડે
$\eq[(b' + \Obj 0)][b']$,
એટલે $\eq[b'][\Obj 0]$ મળે છે.
પણ સ્વયંસિદ્ધિ $!Q_2$ વડે
$\eq/[b'][\Obj 0]$ પણ મળે છે;
આ વિરોધાભાસ છે.

કિસ્સો 2: કોઈ~$c$ માટે
$\eq[a][c']$. તો
$\eq[(b' + c')][\Obj 0]$ મળે છે.
સ્વયંસિદ્ધિ $!Q_5$ વડે
$\eq[(b' + c)'][\Obj 0]$ મળે,
જે ફરી સ્વયંસિદ્ધિ $Q_2$ની
વિરુદ્ધ છે.
\end{proof}

\begin{lem}
\ollabel{lem:less-nsucc}
દરેક પ્રાકૃતિક સંખ્યા~$n$ માટે,
  \[
  \Th{Q} \Proves
  \lforall[x][(x < \num {n+1} \lif (\eq[x][\Obj 0] \lor \dots \lor
    \eq[x][\num n]))].
  \]
\end{lem}

\begin{proof}
$n$ પર અનુમાન વાપરીએ.
પહેલાં $n = 0$નો આરંભિક
કિસ્સો જોઈએ. ત્યારે કોઈપણ~$a$
માટે $a < \num 1 \lif
\eq[a][\Obj 0]$ બતાવવું છે.
માનો કે $a < \num 1$. તો
$<$ની વ્યાખ્યા આપતી સ્વયંસિદ્ધિ
પ્રમાણે
$\lexists[y][\eq[(y'+a)][\Obj 0']]$
છે (કારણ કે $\num 1 \ident
\Obj 0'$).

માનો કે $b$ આ ગુણધર્મ
ધરાવે છે, એટલે
$\eq[(b'+a)][\Obj 0']$.
$\eq[a][\Obj 0]$ બતાવવાનું છે.
સ્વયંસિદ્ધિ $!Q_3$ પ્રમાણે
$\eq[a][\Obj 0]$ અથવા
કોઈ~$c$ માટે $\eq[a][c']$.
પહેલા કિસ્સામાં બતાવવાનું કંઈ
બાકી નથી. તેથી
$\eq[a][c']$ ધારો. તો
$\eq[(b' + c')][\Obj 0']$.
$\Th{Q}$ની સ્વયંસિદ્ધિ $!Q_5$
વડે $\eq[(b'+c)'][\Obj 0']$.
સ્વયંસિદ્ધિ $!Q_1$ વડે
$\eq[(b' + c)][\Obj 0]$.
પણ સ્વયંસિદ્ધિ $!Q_8$ પ્રમાણે
આનો અર્થ $c < \Obj 0$ છે,
જે \olref{lem:less-zero}થી
મળેલા પરિણામની વિરુદ્ધ છે.

હવે અનુમાનનું પગલું જોઈએ.
$n$ માટેનો કિસ્સો ધારીને
$n+1$ માટે સાબિત કરીએ.
માનો કે $a < \num {n+2}$.
ફરી $!Q_3$ વડે બે કિસ્સા:
$\eq[a][\Obj 0]$ અથવા
કોઈ~$b$ માટે $\eq[a][b']$.
પહેલા કિસ્સામાં
$\eq[a][\Obj 0] \lor \dots \lor
\eq[a][\num{n+1}]$ તરત મળે છે.
બીજા કિસ્સામાં
$b' < \num {n+2}$, એટલે
$b' < \num{n+1}'$.
સ્વયંસિદ્ધિ $!Q_8$ વડે કોઈ~$c$
માટે
$\eq[(c'+b')][\num{n+1}']$.
સ્વયંસિદ્ધિ $!Q_5$ વડે
$\eq[(c'+b)'][\num{n+1}']$.
સ્વયંસિદ્ધિ $!Q_1$ વડે
$\eq[(c'+b)][\num{n+1}]$,
અને તેથી સ્વયંસિદ્ધિ $!Q_8$
વડે $b < \num{n+1}$.
અનુમાનની ધારણા પ્રમાણે
$\eq[b][\Obj 0] \lor \dots \lor
\eq[b][\num{n}]$. તર્કથી
$\eq[b'][\Obj 0'] \lor \dots \lor
\eq[b'][\num{n}']$ મળે છે;
અને $\eq[a][b']$ હોવાથી
$\eq[a][\num{1}] \lor \dots \lor
\eq[a][\num{n+1}]$ મળે છે.
\end{proof}

\begin{lem}
  \ollabel{lem:trichotomy} દરેક પ્રાકૃતિક સંખ્યા~$m$ માટે,
  \[
  \Th{Q} \Proves
  \lforall[y][((y < \num{m} \lor \num{m} < y) \lor \eq[y][\num{m}])].
  \]
\end{lem}

\begin{proof}
$m$ પર અનુમાન કરીએ.
પહેલાં $m=0$નો કિસ્સો લો.
$!Q_3$ વડે
$\Th{Q} \Proves
\lforall[y][(\eq[y][\Obj 0] \lor \lexists[z][\eq[y][z']])]$.
કોઈપણ~$a$ લો. તો
$\eq[a][\Obj 0]$ અથવા
કોઈ~$b$ માટે $\eq[a][b']$.
પહેલા કિસ્સામાં
$(a < \Obj 0 \lor \Obj
0 < a) \lor \eq[a][\Obj 0]$
પણ મળે છે. જો $\eq[a][b']$,
તો~$\eq$ માટેના તર્ક વડે
$\eq[(b' + \Obj 0)][(a + \Obj 0)]$.
$!Q_4$ વડે
$\eq[(a + \Obj 0)][a]$;
તેથી $\eq[(b' + \Obj 0)][a]$,
અને પરિણામે
$\lexists[z][\eq[(z' + \Obj 0)][a]]$.
$!Q_8$ની $<$ની વ્યાખ્યા
પ્રમાણે $\Obj 0 < a$.
જો $\Obj 0 < a$, તો
$(\Obj 0 < a \lor a
< \Obj 0) \lor \eq[a][\Obj 0]$
પણ મળે છે.

હવે ધારો કે આ સાબિત થયું છે:
\begin{align*}
  \Th{Q} & \Proves \lforall[y][((y < \num{m} \lor \num{m} < y) \lor
    \eq[y][\num{m}])]
  \intertext{અને આપણે આ બતાવવું છે:}
  \Th{Q} & \Proves \lforall[y][((y < \num{m+1} \lor \num{m+1} < y) \lor
    \eq[y][\num{m+1}])]
\end{align*}
કોઈપણ~$a$ લો. $!Q_3$ પ્રમાણે
$\eq[a][\Obj 0]$ અથવા
કોઈ~$b$ માટે $\eq[a][b']$.
પહેલા કિસ્સામાં $!Q_4$ વડે
$\eq[\num{m}' +
a][\num{m+1}]$; એટલે $!Q_8$
વડે $a < \num{m+1}$.

હવે બીજો કિસ્સો લો:
$\eq[a][b']$. અનુમાનની
ધારણા પ્રમાણે
$(b < \num{m} \lor \num{m} < b) \lor
    \eq[b][\num{m}]$.

પહેલો વિકલ્પ $b < \num{m}$
$!Q_8$ પ્રમાણે
$\lexists[z][\eq[(z' + b)][\num{m}]]$
ની સમકક્ષ છે. માનો કે~$c$ આ
ગુણધર્મ ધરાવે છે. જો
$\eq[(c' + b)][\num{m}]$,
તો $\eq[(c' + b)'][\num{m}']$
પણ મળે છે. $!Q_5$ વડે
$\eq[(c' + b')][(c' + b)']$.
એટલે $\eq[(c' +
b')][\num{m}']$. હવે~$c$
પર અસ્તિત્વ-પરિમાણક લાવીને
અને $\num{m}' \ident
\num{m+1}$ ધ્યાનમાં રાખીને
$\lexists[u][\eq[(u' + b')][\num{m+1}]]$
મળે છે. તેથી $b < \num{m}$
હોય, તો $b' < \num{m+1}$;
અને $a < \num{m+1}$.
\sourcecorrection{OLINC-030}{સ્થિર અંગ્રેજી
મૂળ અહીં $Q_5$ વડે નિરર્થક
સ્વસમાનતા લખે છે અને પછી
$c'$ પર અસ્તિત્વ-પરિમાણક
લાવવાનું કહે છે. જરૂરી સમીકરણ
$(c'+b')=(c'+b)'$ છે;
તેમાં સાક્ષી~$c$ જ છે, જેને
પરિમાણિત કરતાં $b'<\num{m+1}$
મળે છે.}

હવે $\num{m} < b$ માનો,
એટલે
$\lexists[z][\eq[(z' +
\num{m})][b]]$. માનો કે~$c$
આવું~$z$ છે, એટલે
$\eq[(c' + \num{m})][b]$.
તર્કથી
$\eq[(c' + \num{m})'][b']$.
$!Q_5$ વડે
$\eq[(c' + \num{m}')][b']$.
$\eq[a][b']$ અને
$\num{m}' \ident \num{m+1}$
હોવાથી
$\eq[(c' + \num{m+1})][a]$.
$!Q_8$ વડે $\num{m+1} < a$.

છેવટે $\eq[b][\num{m}]$
ધારો. તો તર્કથી
$\eq[b'][\num{m}']$;
એટલે $\eq[a][\num{m+1}]$.

આથી $m$ અને~$b$ માટેના
દરેક વિકલ્પમાંથી $m+1$
અને~$a$ માટેનો અનુરૂપ
વિકલ્પ મળે છે.
\end{proof}

\begin{prop}
\ollabel{prop:rep-minimization}
$!A_g(x, z, y)$ એ~$g(x, z)$ને
$\Th{Q}$માં નિરૂપે, તો
\[
!A_f(z,y) \ident !A_g(y, z, \Obj 0) \land \lforall[w][(w < y \lif \lnot
  !A_g(w, z, \Obj 0))]
\]
એ $f(z) = \umin{x}{[g(x, z) = 0]}$ને
નિરૂપે છે.
\end{prop}

\begin{proof}
પહેલાં બતાવીએ કે $f(n) = m$
હોય, તો
$\Th{Q} \Proves !A_f(\num n, \num m)$,
એટલે કે
\begin{align}
  \Th{Q} & \Proves !A_g(\num{m}, \num{n}, \Obj 0) \land \lforall[w][(w <
    \num{m} \lif \lnot !A_g(w, \num{n}, \Obj 0))]. \notag
  \intertext{કારણ કે $!A_g(x, z, y)$ એ $g(x, z)$ને નિરૂપે છે અને $f(n) = m$ હોય તો $g(m, n) = 0$, તેથી}
\Th{Q} & \Proves !A_g(\num{m}, \num{n}, \Obj 0). \notag
\intertext{$f(n) = m$ હોય, તો દરેક $k < m$ માટે $g(k, n) \neq 0$. તેથી}
\Th{Q} & \Proves \lnot !A_g(\num{k}, \num{n}, \Obj 0). \notag
\intertext{આ પરથી મળે છે:}
\Th{Q} & \Proves \lforall[w][(w < \num{m} \lif \lnot
  !A_g(w, \num{n}, \Obj 0))]. \ollabel{rep-less}
\end{align}
$m = 0$ હોય તો
\olref{lem:less-zero} વડે અને
અન્યથા \olref{lem:less-nsucc}
વડે છેલ્લું તારણ મળે છે.

હવે $f(n) = m$ હોય, તો
$\Th{Q} \Proves
\lforall[y][(!A_f(\num{n}, y) \lif
\eq[y][\num{m}])]$
તે બતાવીએ. દલીલ ફરી
અનૌપચારિક રીતે દર્શાવીએ અને
ઔપચારિક રૂપાંતર વાચક પર છોડીએ.

માનો કે $!A_f(\num{n}, b)$.
ત્યારે (a) $!A_g(b, \num{n},
\Obj 0)$ અને (b)
$\lforall[w][(w < b \lif \lnot !A_g(w, \num{n}, \Obj
  0))]$ મળે છે.
\olref{lem:trichotomy} પ્રમાણે
$(b < \num{m} \lor \num{m} < b)
\lor \eq[b][\num{m}]$.
$b < \num{m}$ અને $\num{m}
< b$ બંને વિરોધાભાસ આપે છે
તે બતાવીએ.

જો $\num{m} < b$, તો~(b)
પરથી $\lnot !A_g(\num{m},
\num{n}, \Obj 0)$ મળે છે.
પરંતુ $m = f(n)$ હોવાથી
$g(m, n) = 0$; અને~$!A_g$
એ~$g$ને નિરૂપે છે, એટલે
$\Th{Q} \Proves !A_g(\num{m},
\num{n}, \Obj 0)$.
આ વિરોધાભાસ છે.

હવે $b < \num{m}$ માનો.
\olref{rep-less} પ્રમાણે
$\Th{Q} \Proves \lforall[w][(w <
  \num{m} \lif \lnot !A_g(w, \num{n}, \Obj 0))]$;
તેથી $\lnot !A_g(b, \num{n},
\Obj 0)$ મળે છે. આ ફરી~(a)
ની વિરુદ્ધ છે.
\end{proof}
% END OLP-0296
\endgroup


\begingroup
\makeatletter\def\ol@sectioncs{section}\makeatother

% BEGIN OLP-0297
\olfileid{inc}{req}{crq}
\olsection{સંગણનીય વિધેયો~$\Th{Q}$માં નિરૂપ્ય છે}
\phantomsection\label{guunit:OLP-0297}


\begin{thm}
દરેક સંગણનીય વિધેય
$\Th{Q}$માં નિરૂપ્ય છે.
\end{thm}

\begin{proof}
ચોક્કસતા માટે, ચર્ચ--ટ્યુરિંગ
માન્યતાનો ઉપયોગ કરીને એવું
કહીએ કે વિધેય સંગણનીય છે
ત્યારે અને માત્ર ત્યારે જ
તે સામાન્ય પુનરાવર્તી છે.
સામાન્ય પુનરાવર્તી વિધેયો
એવાં છે જે શૂન્ય વિધેય~$\Zero$,
ઉત્તરગામી વિધેય~$\Succ$ અને
પ્રક્ષેપણ વિધેય~$\Proj{n}{i}$માંથી
સંયોજન, આદિમ પુનરાવર્તન અને
નિયમિત લઘુત્તમીકરણ વડે
વ્યાખ્યાયિત થઈ શકે છે.
\olref[pri]{lem:prim-rec} પ્રમાણે,
$f$ અને~$g$માંથી વ્યાખ્યાયિત
કરી શકાય એવું કોઈપણ વિધેય~$h$,
$f$, $g$ અને $\Zero$,
$\Succ$, $\Proj{n}{i}$,
$\Add$, $\Mult$, $\Char{=}$માંથી
સંયોજન અને નિયમિત લઘુત્તમીકરણ
વડે પણ વ્યાખ્યાયિત થઈ શકે છે.
પરિણામે, વિધેય સામાન્ય પુનરાવર્તી
હોય ત્યારે અને માત્ર ત્યારે જ
તે $\Zero$, $\Succ$,
$\Proj{n}{i}$, $\Add$,
$\Mult$, $\Char{=}$માંથી
સંયોજન અને નિયમિત લઘુત્તમીકરણ
વડે વ્યાખ્યાયિત થઈ શકે છે.

વધુમાં, સંબંધિત મૂળભૂત વિધેયો
$\Th{Q}$માં નિરૂપ્ય છે તે આપણે
બતાવ્યું છે
(\cref{inc:req:bre:prop:rep-zero,inc:req:bre:prop:rep-succ,inc:req:bre:prop:rep-proj,inc:req:bre:prop:rep-id,inc:req:bre:prop:rep-add,inc:req:bre:prop:rep-mult}).
નિરૂપ્ય વિધેયોમાંથી સંયોજન
અથવા નિયમિત લઘુત્તમીકરણ
વડે વ્યાખ્યાયિત કોઈપણ વિધેય
પણ નિરૂપ્ય છે
(\olref[cmp]{prop:rep-composition},
\olref[min]{prop:rep-minimization}).
તેથી દરેક સામાન્ય પુનરાવર્તી
વિધેય~$\Th{Q}$માં નિરૂપ્ય છે.
\end{proof}

\begin{explain}
આપણે બતાવ્યું છે કે સંગણનીય
વિધેયોનો ગણ એટલે $\Th{Q}$માં
નિરૂપ્ય વિધેયોનો ગણ.
હકીકતમાં આ સાબિતી વધુ સામાન્ય
છે. નિરૂપ્યતાની વ્યાખ્યામાંથી
જોવું મુશ્કેલ નથી કે
$\Th{Q}$નું વિસ્તરણ કરતો
કોઈપણ સિદ્ધાંત (અથવા જેમાં
$\Th{Q}$નું અર્થઘટન કરી શકાય)
સંગણનીય વિધેયોને નિરૂપી શકે છે.
પણ બીજી દિશામાં, જે
નિષ્પત્તિ પદ્ધતિમાં
નિષ્પત્તિની વિભાવના સંગણનીય
છે, તેમાં દરેક નિરૂપ્ય વિધેય
સંગણનીય છે. આથી, દાખલા તરીકે,
સંગણનીય વિધેયોનો ગણ પીઆનો
અંકગણિતમાં, અથવા ઝર્મેલો--ફ્રેન્કેલ
સમૂહસિદ્ધાંતમાં પણ નિરૂપ્ય
વિધેયોના ગણ તરીકે વર્ણવી શકાય
છે. ગડેલે નોંધ્યું તેમ, આ કંઈક
આશ્ચર્યજનક છે. આગળ જોઈશું
કે સાબિતિક્ષમતાના પ્રશ્નો માટે
કયો સિદ્ધાંત લેવાય છે તે ખૂબ
મહત્ત્વનું છે: સ્થૂળ રીતે કહીએ
તો સ્વયંસિદ્ધિઓ જેટલી વધુ
શક્તિશાળી, તેટલું વધુ સાબિત
કરી શકાય. પરંતુ અનેક પ્રકારના
સ્વયંસિદ્ધિમૂલક સિદ્ધાંતોમાં
નિરૂપ્ય વિધેયો બરાબર સંગણનીય
વિધેયો જ છે; સિદ્ધાંતોને
અસરકારક રીતે સ્વયંસિદ્ધિઓથી વર્ણવી શકાય
ત્યાં સુધી વધુ શક્તિશાળી
સિદ્ધાંત વધુ વિધેયોને નિરૂપતો
નથી.
\sourcecorrection{OLINC-031}{સ્થિર અંગ્રેજી
મૂળના અંતિમ વાક્યમાં માત્ર
``સ્વયંસિદ્ધીકરણ શક્ય'' હોવાની
શરત લખાઈ છે. અગાઉની દલીલ
સાબિતીઓ તપાસી શકાય એવી
નિગમન-પદ્ધતિ પર આધાર રાખે છે;
તેથી અહીં અસરકારક સ્વયંસિદ્ધીકરણની
શરત સ્પષ્ટ કરી છે.}
\end{explain}
% END OLP-0297
\endgroup


\begingroup
\makeatletter\def\ol@sectioncs{section}\makeatother

% BEGIN OLP-0298
\olfileid{inc}{req}{rel}

\olsection{સંબંધોનું નિરૂપણ}
\phantomsection\label{guunit:OLP-0298}


\emph{સંબંધ} નિરૂપ્ય હોય તેનો
અર્થ શું છે તે કહીએ.

\begin{defn}
\ollabel{defn:representing-relations}
પ્રાકૃતિક સંખ્યાઓ પરનો સંબંધ
$R(x_0,\dots,x_k)$ {\em
$\Th{Q}$માં નિરૂપ્ય} કહેવાય,
જો એવું સૂત્ર $!A_R(x_0,\dots,x_k)$
હોય કે $R(n_0,\dots,n_k)$
સાચું હોય ત્યારે
$\Th{Q}$માં
$!A_R(\num{n_0},\dots,\num{n_k})$
સાબિત થાય, અને
$R(n_0,\dots,n_k)$ ખોટું હોય
ત્યારે $\Th{Q}$માં
$\lnot !A_R(\num{n_0},
\dots, \num{n_k})$ સાબિત થાય.
\end{defn}

\begin{thm}
\ollabel{thm:representing-rels}
સંબંધ $\Th{Q}$માં નિરૂપ્ય
હોય ત્યારે અને માત્ર ત્યારે જ
તે સંગણનીય હોય છે.
\end{thm}

\begin{proof}
પહેલી દિશા માટે માનો કે
$R(x_0,\dots,x_k)$નું નિરૂપણ
$!A_R(x_0,\dots,x_k)$ સૂત્ર
કરે છે. $R$ ગણવા માટેની
રીત આ છે: ઇનપુટ
$n_0$, \dots,~$n_k$ મળે ત્યારે
$!A_R(\num{n_0}, \dots, \num{n_k})$
ની સાબિતી અને
$\lnot !A_R(\num{n_0}, \dots, \num{n_k})$
ની સાબિતી એકસાથે શોધો.
આપણી ધારણા પ્રમાણે એ બેમાંથી
એક સાબિતી જરૂર મળશે; પહેલી
મળે તો ``હા'' કહો, અન્યથા
``ના'' કહો.

બીજી દિશામાં, માનો કે
$R(x_0, \dots, x_k)$ સંગણનીય
છે. વ્યાખ્યા પ્રમાણે આનો અર્થ
તેનું લાક્ષણિક વિધેય
$\Char{R}(x_0, \dots, x_k)$
સંગણનીય છે. \olref[int]{thm:representable-iff-comp}
પ્રમાણે, $\Char{R}$ને કોઈ
સૂત્ર, કહો કે
$!A_{\Char{R}}(x_0, \dots, x_k,
y)$, નિરૂપે છે.
$!A_R(x_0, \dots, x_k)$ તરીકે
$!A_{\Char{R}}(x_0,
\dots, x_k, \num{1})$ સૂત્ર
લો. હવે કોઈપણ $n_0$,
\dots,~$n_k$ માટે,
$R(n_0, \dots, n_k)$ સાચું
હોય તો $\Char{R}(n_0, \dots,
n_k) = 1$; આ કિસ્સામાં
$\Th{Q}$માં
$!A_{\Char{R}}(\num{n_0}, \dots, \num{n_k},
\num{1})$ સાબિત થાય છે,
એટલે $\Th{Q}$માં
$!A_R(\num{n_0}, \dots,
\num{n_k})$ સાબિત થાય છે.
બીજી તરફ, $R(n_0, \dots,
n_k)$ ખોટું હોય તો
$\Char{R}(n_0, \dots, n_k) = 0$.
આથી $\Th{Q}$માં આ સાબિત થાય છે:
\[
\lforall[y][(!A_{\Char{R}}(\num{n_0}, \dots, \num{n_k}, y) \lif y =
  \num{0})].
\]
$\Th{Q}$માં
$\eq/[\num{0}][\num{1}]$
સાબિત થાય છે, તેથી
$\Th{Q}$માં
$\lnot !A_{\Char{R}}(\num{n_0}, \dots, \num{n_k}, \num{1})$
સાબિત થાય છે; પરિણામે
$\lnot !A_R(\num{n_0}, \dots, \num{n_k})$
પણ સાબિત થાય છે.
\end{proof}

\begin{prob}
$R$ એ~$\Th{Q}$માં નિરૂપ્ય
હોય, તો~$\Char{R}$ પણ તેમાં
નિરૂપ્ય છે તે બતાવો.
\end{prob}
% END OLP-0298
\endgroup


\begingroup
\makeatletter\def\ol@sectioncs{section}\makeatother

% BEGIN OLP-0299
\olfileid{inc}{req}{und}
\olsection{અનિર્ણેયતા}
\phantomsection\label{guunit:OLP-0299}


સિદ્ધાંત~$\Th{T}$ને આપણે
\emph{અનિર્ણેય} કહીએ, જો એવી
કોઈ સંગણકીય પ્રક્રિયા ન હોય
જે $\Th{T}$ની ભાષાના કોઈપણ
વાક્ય~$!A$ માટે ``શું
$\Th{T}$માં~$!A$ સાબિત થાય
છે?'' તેનો સાચો જવાબ
અચૂક રીતે અને સાન્ત પગલાંમાં
આપે. તેથી $\Th{Q}$ નિર્ણેય
હોય ત્યારે અને માત્ર ત્યારે જ
અંકગણિતની ભાષાનું વાક્ય~$!A$
આપતાં, $\Th{Q} \Proves !A$
છે કે નહીં તે નક્કી કરતી
સંગણકીય પ્રક્રિયા હોય.
આને વધુ ચોક્કસ રીતે પૂછી
શકીએ: સંબંધ~$\Prov[\Th{Q}](y)$
$y$ માટે ત્યારે અને માત્ર ત્યારે જ
સાચો હોય કે $y$
$\Th{Q}$માં સાબિત થઈ શકે
એવા વાક્યની ગડેલ સંખ્યા
હોય, તો શું આ સંબંધ પુનરાવર્તી
છે? જવાબ છે: ના.

\begin{thm}
$\Th{Q}$ અનિર્ણેય છે, એટલે
આ સંબંધ
\[
\Prov[\Th{Q}](y) \defiff \fn{Sent}(y) \land
\lexists[x][\Prf[\Th{Q}](x, y)]
\]
પુનરાવર્તી નથી.
\end{thm}

\begin{proof}
ધારો કે તે પુનરાવર્તી હોય.
તો અટકવાની સમસ્યા આ રીતે
ઉકેલી શકીએ: $e$ અને $n$
આપેલાં હોય, તો આપણે જાણીએ
છીએ કે $\cfind{e}(n) \fdefined$
ત્યારે અને માત્ર ત્યારે જ
કોઈ~$s$ માટે $T(e, n, s)$ છે;
અહીં $T$ એ
\olref[cmp][rec][nft]{thm:kleene-nf}માંથી
ક્લીનીનું વિધેયધર્મ છે.
$T$ આદિમ પુનરાવર્તી હોવાથી
તેને સૂત્ર~$!B_T$ વડે
$\Th{Q}$માં નિરૂપી શકાય છે,
એટલે $T(e, n, s)$ હોય
ત્યારે અને માત્ર ત્યારે જ
$\Th{Q} \Proves
!B_T(\num{e}, \num{n}, \num{s})$.
જો $\Th{Q} \Proves
!B_T(\num{e}, \num{n}, \num{s})$,
તો $ \Th{Q} \Proves
\lexists[y][!B_T(\num{e}, \num{n}, y)]$
પણ છે. જો આવું કોઈ~$s$
ન હોય, તો દરેક~$s$ માટે
$\Th{Q} \Proves \lnot
!B_T(\num{e}, \num{n}, \num{s})$.
પરંતુ $\Th{Q}$
$\omega$-સુસંગત છે, એટલે
દરેક~$n \in \Nat$ માટે
$\Th{Q} \Proves \lnot
!A(\num{n})$ હોય, તો
$\Th{Q} \Proves/
\lexists[y][!A(y)]$.
આપણે આ જાણીએ છીએ, કારણ કે
$\Th{Q}$ની સ્વયંસિદ્ધિઓ
પ્રમાણભૂત નમૂના~$\Struct{N}$માં
સાચી છે. તેથી
$\Th{Q} \Proves/
\lexists[y][!B_T(\num{e}, \num{n}, y)]$.
બીજા શબ્દોમાં,
$\Th{Q} \Proves
\lexists[y][!B_T(\num{e}, \num{n}, y)]$
ત્યારે અને માત્ર ત્યારે જ
કોઈ~$s$ માટે $T(e, n, s)$,
એટલે ત્યારે અને માત્ર ત્યારે જ
$\cfind{e}(n) \fdefined$.
$e$ અને~$n$ પરથી આપણે
$\Gn{\lexists[y][!B_T(\num{e},
    \num{n}, y)]}$ ગણી શકીએ;
આ કામ કરતું આદિમ પુનરાવર્તી
વિધેય $g(e, n)$ લો. તો
\[
h(e, n) =
\begin{cases}
1 & \text{જો $\Prov[\Th{Q}](g(e, n))$}\\
0 & \text{અન્યથા}.
\end{cases}
\]
આથી $\Prov[\Th{Q}]$
પુનરાવર્તી હોય તો~$h$ પણ
પુનરાવર્તી હોવાનું નીકળે.
પરંતુ
\olref[cmp][rec][hlt]{thm:halting-problem}
પ્રમાણે~$h$ પુનરાવર્તી નથી;
તેથી $\Prov[\Th{Q}]$ પણ
પુનરાવર્તી ન હોઈ શકે.
\end{proof}

\begin{cor}
પ્રથમ-ક્રમનો તર્ક અનિર્ણેય છે.
\end{cor}

\begin{proof}
પ્રથમ-ક્રમનો તર્ક નિર્ણેય
હોય, તો $\Th{Q}$માં
સાબિતિક્ષમતા પણ નિર્ણેય
હોત, કારણ કે
$\Th{Q} \Proves !A$ ત્યારે અને
માત્ર ત્યારે જ
$\Proves !T \lif !A$;
અહીં~$!T$ એ $\Th{Q}$ની
સ્વયંસિદ્ધિઓનું સંયોજન છે.
\end{proof}
% END OLP-0299
\endgroup


\begingroup
\makeatletter\def\ol@sectioncs{section}\makeatother

% BEGIN OLP-0300
\olfileid{inc}{inp}{s1c}
\olsection{\texorpdfstring{$\Sigma_1$}{Sigma-1} પૂર્ણતા}
\phantomsection\label{guunit:OLP-0300}


$\Th{Q}$ અને તેના સુસંગત,
સ્વયંસિદ્ધિગત વિસ્તરણો અપૂર્ણ
હોવા છતાં, આપણે જોયું છે કે
$\Th{Q}$ અંકો વિશેનાં ઘણાં
મૂળભૂત તથ્યો સાબિત કરે છે.
હકીકતમાં આ વાતને ખાસ્સી
આગળ લઈ જઈ શકાય છે.
$\Th{Q}$માં શું સાબિત કરી
શકાય છે તેનો વ્યાપ સમજવા
$\Delta_0$, $\Sigma_1$ અને
$\Pi_1$ સૂત્રોની
વિભાવનાઓ રજૂ કરીએ.
સ્થૂળ રીતે કહીએ તો
$\Sigma_1$ સૂત્ર એ
$\lexists[x][!B(x)]$ આકારનું
છે, જ્યાં~$!B$ માત્ર વિધાનાત્મક
સંયોજકો અને મર્યાદિત પરિમાણકો
વડે રચાયેલું છે. આપણે બતાવીશું
કે $!A$ એ $\Struct{N}$માં
સાચું $\Sigma_1$ વાક્ય
હોય, તો $\Th{Q} \Proves !A$
(\olref{thm:sigma1-completeness}).

\begin{defn}
\ollabel{defn:bd-quant}
\emph{મર્યાદિત અસ્તિત્વવાચક
સૂત્ર} એ
$\lexists[x][(x < t \land !A(x))]$
આકારનું છે, જ્યાં~$t$ કોઈપણ
પદ છે; પરંપરાગત રીતે તેને
$\bexists{x < t}{!A(x)}$
લખીએ છીએ.
\emph{મર્યાદિત સર્વવાચક
સૂત્ર} એ
$\lforall[x][(x < t \lif !A(x))]$
આકારનું છે, જ્યાં~$t$ કોઈપણ
પદ છે; પરંપરાગત રીતે તેને
$\bforall{x < t}{!A(x)}$
લખીએ છીએ.
\end{defn}

\begin{defn}
\ollabel{defn:delta0-sigma1-pi1-frm}
સૂત્ર~$!B$ એ $\Delta_0$
ત્યારે કહેવાય જ્યારે તે
પરમાણુ સૂત્રોમાંથી
માત્ર વિધાનાત્મક સંયોજકો
અને મર્યાદિત પરિમાણીકરણ
વડે રચાયેલું હોય.
સૂત્ર~$!A$ એ $\Sigma_1$
ત્યારે કહેવાય જ્યારે
$!A \ident \lexists[x][!B(x)]$
અને~$!B$ એ $\Delta_0$ હોય.
સૂત્ર~$!A$ એ $\Pi_1$
ત્યારે કહેવાય જ્યારે
$!A \ident \lforall[x][!B(x)]$
અને~$!B$ એ $\Delta_0$ હોય.
\end{defn}

\begin{lem}
\ollabel{lem:q-proves-clterm-id}
માનો કે~$t$ બંધ પદ છે અને
$\Value{t}{N} = n$. તો
$\Th{Q} \Proves \eq[t][\num n]$.
\end{lem}

\begin{proof}
$t$ની જટિલતા પર અનુમાન
વડે આ સાબિત કરીએ.
આરંભિક કિસ્સામાં
$\Value{\Obj 0}{N} = 0$,
અને $\num 0 \ident \Obj 0$
હોવાથી
$\Th{Q} \Proves \eq[\Obj 0][\num 0]$.
અનુમાનના કિસ્સામાં માનો કે
$t_1$ અને~$t_2$ એવાં પદો છે કે
$\Value{t_1}{N} = n_1$,
$\Value{t_2}{N} = n_2$,
$\Th{Q} \Proves \eq[t_1][\num n_1]$,
અને $\Th{Q} \Proves
\eq[t_2][\num n_2]$.

તો $\Value{(t_1')}{N} = n_1 + 1$.
અનુમાનની ધારણા અને
સૂત્ર~$\eq[\num{n_1}'][\num{n_1}']$
પર સમતા માટેના પ્રથમ-ક્રમના
નિયમો લગાડવાથી
$\Th{Q} \Proves
\eq[t_1'][{\num n_1}']$.
તેથી અંકોની વ્યાખ્યા પ્રમાણે
$\Th{Q} \Proves
\eq[t_1'][\num{n_1 + 1}]$.

સરવાળા માટે
\[
      \Value{(t_1 + t_2)}{N}
    = \Value{t_1}{N} + \Value{t_2}{N}
    = n_1 + n_2.
\]
અનુમાનની ધારણા તથા
સમતાના નિયમો વડે પહેલાં
$\Th{Q} \Proves
\eq[t_1 + t_2][\num{n_1} + t_2]$
અને તે નિયમો બીજી વાર
લગાડતાં
$\Th{Q} \Proves
\eq[t_1 + t_2][\num{n_1} + \num{n_2}]$
મળે છે.
\olref[inc][req][bre]{lem:q-proves-add}
પ્રમાણે
$\Th{Q} \Proves
\eq[\num{n_1} + \num{n_2}][\num{n_1 + n_2}]$;
તેથી $\Th{Q} \Proves
\eq[t_1 + t_2][\num{n_1 + n_2}]$.

$\times$ માટે પણ
\olref[inc][req][bre]{lem:q-proves-mult}
વાપરીને આવી જ દલીલ થાય છે.
અંકગણિતનાં બધાં બંધ પદો
આ કિસ્સાઓમાં આવી જાય છે;
તેથી $\Value{t}{N} = n$
હોય એવા દરેક બંધ પદ~$t$
માટે $\Th{Q} \Proves
\eq[t][\num n]$.
\end{proof}
 
\begin{prob}
$\times$ સંબંધિત
\olref{lem:q-proves-clterm-id}ની
સાબિતી વિગતે આપો.
\end{prob}

\begin{lem}
\ollabel{lem:atomic-completeness}
માનો કે $t_1$ અને~$t_2$
બંધ પદો છે. તો
\begin{enumerate}
\item જો $\Value{t_1}{N} = \Value{t_2}{N}$,
    તો $\Th{Q} \Proves \eq[t_1][t_2]$.
\item જો $\Value{t_1}{N} \neq \Value{t_2}{N}$,
    તો $\Th{Q} \Proves \eq/[t_1][t_2]$.
\item જો $\Value{t_1}{N} < \Value{t_2}{N}$,
    તો $\Th{Q} \Proves t_1 < t_2$.
\item જો $\Value{t_2}{N} \leq \Value{t_1}{N}$,
    તો $\Th{Q} \Proves \lnot(t_1 < t_2)$.
\end{enumerate}
\end{lem}

\begin{proof}
$t_1$ અને~$t_2$ આપેલાં હોય,
તો $n = \Value{t_1}{N}$ અને
$m = \Value{t_2}{N}$ નક્કી કરીએ.

માનો કે $!A \ident t_1 = t_2$.
\olref{lem:q-proves-clterm-id}
પ્રમાણે $\Th{Q} \Proves
\eq[t_1][\num n]$ અને
$\Th{Q} \Proves \eq[t_2][\num m]$.
જો $n = m$, તો
$\Th{Q} \Proves \eq[\num n][\num m]$;
અને સમતાની સંક્રામકતા વડે
$\Th{Q} \Proves \eq[t_1][t_2]$.
જો $n \neq m$, તો
$\Th{Q} \Proves \eq/[\num n][\num m]$;
સમતાની સંક્રામકતા ફરીથી
લગાડતાં $\Th{Q} \Proves
\eq/[t_1][t_2]$.
\sourcecorrection{OLINC-032}{સ્થિર અંગ્રેજી
મૂળ અહીં $t_2$નું મૂલ્ય~$m$
વ્યાખ્યાયિત કરે છે છતાં
$t_2=\num n$ લખે છે.
આ સમીકરણ $t_2=\num m$
કર્યું છે, જેથી $n=m$ અને
$n\neq m$ બંને કિસ્સા
યોગ્ય રીતે ચાલે.}

હવે $!A \ident t_1 < t_2$
લો. બંને કિસ્સામાં સ્વયંસિદ્ધિ
$!Q_8$ વાપરીશું; તે દરેક
$x,y$ માટે
$x < y \liff
\lexists[z][\eq[z' + x][y]]$
કહે છે.

માનો કે $\Sat{N}{t_1 < t_2}$.
તો કોઈ~$k \in \Nat$ માટે
$n + k + 1 = m$.
\olref{lem:q-proves-clterm-id}
પ્રમાણે $\Th{Q} \Proves
\eq[t_1][\num n]$ અને
$\Th{Q} \Proves \eq[t_2][\num m]$.
આ લેમ્માના પહેલા ભાગથી
$\Th{Q} \Proves
\eq[{\num k}' + \num n][\num m]$.
સમતાની સંક્રામકતા વડે
$\Th{Q} \Proves
\eq[{\num k}' + t_1][t_2]$;
એટલે $\Th{Q} \Proves
\lexists[z][\eq[z' + t_1][t_2]]$.
$!Q_8$ની જમણેથી ડાબી
દિશા પ્રમાણે $\Th{Q} \Proves
t_1 < t_2$.
\sourcecorrection{OLINC-033}{સ્થિર અંગ્રેજી
મૂળ પહેલા $\num n+\num k'$
નું સમીકરણ આપે છે, પરંતુ
પછીના અસ્તિત્વવાચક સાક્ષી માટે
$\num k'+\num n$ જરૂરી છે.
$\Th{Q}$ સામાન્ય સરવાળાની
અદલાબદલી સાબિત કરતું નથી;
અહીં નિશ્ચિત અંકોનું જરૂરી
સમીકરણ લેમ્માના પહેલા
ભાગથી સીધું લીધું છે.}

તેના બદલે $\Sat/{N}{t_1 < t_2}$,
એટલે $m \leq n$, માનો.
$\Th{Q}$માં રહીને
$t_1 < t_2$ ધારો. $!Q_8$ની
ડાબેથી જમણી દિશા પ્રમાણે
કોઈ~$z$ માટે
$\eq[z' + t_1][t_2]$.
$\Th{Q} \Proves
\eq[t_1][\num n]$ અને
$\Th{Q} \Proves \eq[t_2][\num m]$
હોવાથી
$\eq[z' + \num n][\num m]$.
હવે $!Q_5$ વાપરીને~$m$
પર બાહ્ય અનુમાન કરીએ, તો
$\eq[z' + \num{n - m}][\Obj 0]$.
જો $m = n$, તો $!Q_4$
વડે $\eq[z'][\Obj 0]$ મળે,
જે $!Q_2$ને વિરુદ્ધ છે.
જો $m < n$, તો ફરી~$!Q_5$
વડે
$\eq[(z' + \num{n - m - 1})'][\Obj 0]$,
જે $!Q_2$ને વિરુદ્ધ છે.
આથી $\Th{Q} \Proves
\lnot(t_1 < t_2)$.
\sourcecorrection{OLINC-034}{સ્થિર અંગ્રેજી
મૂળ આ બંને વિરોધાભાસને
$Q_3$ સાથે જોડે છે; ઉત્તરગામી
શૂન્ય ન હોય તે સ્વયંસિદ્ધિ
$Q_2$ છે. $m=n$ના
કિસ્સામાં સરવાળામાંથી શૂન્ય
દૂર કરવા $Q_4$ પણ જરૂરી છે.}
\end{proof}

\begin{lem}
\ollabel{lem:bounded-quant-equiv}
માનો કે $!A$ એ સૂત્ર,
$t$ એ બંધ પદ અને
$k=\Value{t}{N}$. તો
\begin{enumerate}
\item $\Th{Q} \Proves \bforall{x<t}{!A(x)}$
ત્યારે અને માત્ર ત્યારે જ
$\Th{Q} \Proves
    !A(\num 0) \land \dots \land !A(\num{k-1})$.
\item $\Th{Q} \Proves \bexists{x<t}{!A(x)}$
ત્યારે અને માત્ર ત્યારે જ
$\Th{Q} \Proves
    !A(\num 0) \lor \dots \lor !A(\num{k-1})$.
\end{enumerate}
\end{lem}

\begin{proof}
    મર્યાદિત સર્વવાચક પરિમાણકનો
    કિસ્સો સાબિત કરીએ.
    $\Value{t}{N} = 0$ હોય,
    તો સમકક્ષતાની ડાબી બાજુ
    $\Th{Q}$માં સાબિત થાય છે,
    કારણ કે \olref[inc][req][min]{lem:less-zero}
    પ્રમાણે $x<\num 0$
    એવું શક્ય નથી.
    તેવી જ રીતે ખાલી સંયોજનને
    માત્ર $\ltrue$ લઈ શકીએ,
    જે પણ $\Th{Q}$માં
    સાબિત થાય છે.
    \sourcecorrection{OLINC-035}{સ્થિર અંગ્રેજી
    મૂળ આ સર્વવાચક કિસ્સામાં
    ખાલી વિયોજનને $\ltrue$
    કહે છે. અહીં સાન્ત
    સર્વવાચક દાવાને અનુરૂપ
    ખાલી સંયોજન જ $\ltrue$
    છે; ખાલી વિયોજનનો અર્થ
    જુદો થાય છે.}
    તેથી કોઈ પ્રાકૃતિક સંખ્યા~$k$
    માટે $\Value{t}{N} = k+1$
    ધારો. \olref{lem:q-proves-clterm-id}
    વડે આપણે
    $\bforall{x<\num{k+1}}{!A(x)}$
    આકારના સૂત્ર સાથે
    કામ કરીએ તેમ માની શકીએ.
    
    માનો કે $\Th{Q} \Proves
    \bforall{x<\num{k+1}}{!A(x)}$,
    અને $n \leq k$ લો.
    \olref{lem:atomic-completeness}
    પ્રમાણે $\Th{Q} \Proves
    \num n < \num{k+1}$;
    તેથી તર્કથી $\Th{Q} \Proves
    !A(\num n)$. દરેક
    $n \leq k$ માટે આ વાત
    $k+1$ વાર લગાડવાથી
    માંગ્યા પ્રમાણે
    $\Th{Q} \Proves !A(\num 0)
    \land \dots \land !A(\num k)$.
    
    બીજી દિશા માટે
    $\Th{Q} \Proves
    !A(\num 0) \land \dots \land !A(\num k)$
    ધારો. $\Th{Q}$માં
    $x < \num{k+1}$ માનો.
    \olref[inc][req][min]{lem:less-nsucc}
    વડે
    $x = \num 0 \lor \dots \lor x = \num k$;
    તેથી તર્કથી~$!A(x)$
    મળે છે અને પરિણામે
    સર્વવાચક દાવો
    $\bforall{x<\num{k+1}}{!A(x)}$
    મળે છે.
    
    મર્યાદિત અસ્તિત્વવાચક
    પરિમાણિત સૂત્રોની
    સમકક્ષતાની સાબિતી પણ આવી જ છે.
\end{proof}

\begin{prob}
\olref{lem:bounded-quant-equiv}માં
અસ્તિત્વવાચક કિસ્સાની વિગતવાર
સાબિતી આપો.
\end{prob}

\begin{lem}
\ollabel{lem:delta0-completeness}
$!A$ એ $\Struct{N}$માં સાચું
$\Delta_0$ વાક્ય હોય,
તો $\Th{Q} \Proves !A$.
\end{lem}

\begin{proof}
સૂત્રની જટિલતા પર
અનુમાન વડે સાબિત કરીએ.
આરંભિક કિસ્સો
\olref{lem:atomic-completeness}
આપે છે; તેથી અનુમાનના
પગલા તરફ જઈએ. સરળતા માટે,
નિષેધ લાગુ પડે તે
સૂત્રની રચના મુજબ
નિષેધના કિસ્સાને પેટાકિસ્સામાં
વહેંચીએ.

\begin{enumerate}
\item માનો કે $(!A \land !B)$
$\Struct{N}$માં સાચું છે;
ત્યારે $!A$ અને~$!B$ બંને
$\Struct{N}$માં સાચાં છે.
અનુમાનની ધારણા પ્રમાણે
$\Th{Q} \Proves !A$ અને
$\Th{Q} \Proves !B$;
તેથી તર્કથી
$\Th{Q} \Proves (!A \land !B)$.
\item માનો કે
$\lnot (!A \land !B)$ એ
$\Struct{N}$માં સાચું છે;
ત્યારે $\lnot !A$ અથવા
$\lnot !B$ એ $\Struct{N}$માં
સાચું છે. વ્યાપકતા ગુમાવ્યા
વિના પહેલું સાચું છે એમ
ધારો. અનુમાનની ધારણા પ્રમાણે
$\Th{Q} \Proves \lnot !A$;
તેથી તર્કથી $\Th{Q} \Proves
\lnot (!A \land !B)$.
\item માનો કે $(!A \lor !B)$
$\Struct{N}$માં સાચું છે;
ત્યારે $!A$ એ $\Struct{N}$માં
સાચું છે અથવા~$!B$ એ
$\Struct{N}$માં સાચું છે.
વ્યાપકતા ગુમાવ્યા વિના પહેલું
સાચું છે એમ ધારો.
અનુમાનની ધારણા પ્રમાણે
$\Th{Q} \Proves !A$;
તેથી તર્કથી $\Th{Q} \Proves
(!A \lor !B)$.
\item માનો કે
$\lnot(!A \lor !B)$ એ
$\Struct{N}$માં સાચું છે;
ત્યારે $\lnot !A$ અને
$\lnot !B$ બંને $\Struct{N}$માં
સાચાં છે. અનુમાનની ધારણા
પ્રમાણે $\Th{Q} \Proves
\lnot !A$ અને
$\Th{Q} \Proves \lnot !B$.
પરિણામે તર્કથી
$\Th{Q} \Proves
\lnot(!A \lor !B)$.
\item માનો કે
$\bforall{x<t}{!A(x)}$
$\Struct{N}$માં સાચું છે,
જ્યાં~$t$ બંધ પદ અને
$k=\Value{t}{N}$. અનુમાનની
ધારણા તથા તર્ક પ્રમાણે
દરેક $n < \Value{t}{N}$
માટે $!A(\num n)$ એ
$\Struct{N}$માં સાચું હોય,
તો $\Th{Q} \Proves
!A(\num 0) \land \dots \land !A(\num{k-1})$.
\olref{lem:bounded-quant-equiv}
પ્રમાણે $\Th{Q} \Proves
\bforall{x<t}{!A(x)}$.
\item મર્યાદિત અસ્તિત્વવાચક
પરિમાણકનો કિસ્સો, જેમાં
$\bexists{x < t}{!A(x)}$
આકારનું વાક્ય હોય,
તે મર્યાદિત સર્વવાચક
પરિમાણકના કિસ્સા જેવો જ છે.
\item માનો કે
$\lnot \bforall{x<t}{!A(x)}$
$\Struct{N}$માં સાચું છે,
જ્યાં~$t$ બંધ પદ છે.
આ વાક્ય એ
વાક્ય
$\bexists{x<t}{\lnot !A(x)}$
ની સમકક્ષ છે, અને સમકક્ષતા
$\Th{Q}$માં નિગમ્ય છે;
તેથી મર્યાદિત અસ્તિત્વવાચક
પરિમાણકની દલીલ વાપરી શકીએ.
\item તેવી જ રીતે માનો કે
$\lnot \bexists{x<t}{!A(x)}$
$\Struct{N}$માં સાચું છે,
જ્યાં~$t$ બંધ પદ છે.
આ વાક્ય એ
$\Th{Q}$માં
$\bforall{x<t}{\lnot!A(x)}$
ની સમકક્ષ છે; તેથી
મર્યાદિત સર્વવાચક પરિમાણકની
દલીલ વાપરી શકીએ.
\sourcecorrection{OLINC-036}{સ્થિર અંગ્રેજી
મૂળના આ કિસ્સામાં
$\bexists$ માટે બીજો
કૌંસબંધ દલીલ-ઘટક છૂટી
ગયો છે. મેક્રોની વ્યાખ્યા
પ્રમાણે અહીં $!A(x)$ને
બીજા ઘટકમાં મૂક્યું છે.}
\item છેલ્લે માનો કે
$\lnot !A$ એ $\Struct{N}$માં
સાચું છે. હવે માત્ર બે
કિસ્સા બાકી છે: $!A$
પરમાણુ છે, અથવા કોઈ
$\Delta_0$ વાક્ય~$!B$
માટે
$\lnot !A \ident \lnot\lnot !B$.
જો $!A$ પરમાણુ હોય, તો
\olref{lem:atomic-completeness}
પ્રમાણે $\Th{Q} \Proves
\lnot !A$. જો
$\lnot !A \ident \lnot\lnot !B$,
તો તર્કથી તે~$!B$ની
$\Th{Q}$માં સાબિત કરી શકાય
એવી સમકક્ષતા ધરાવે છે.
તે $\Struct{N}$માં સાચું છે,
કારણ કે $\lnot !A$ પણ
$\Struct{N}$માં સાચું છે.
તેથી અનુમાનની ધારણા પ્રમાણે
$\Th{Q} \Proves \lnot !A$.
\end{enumerate}
\end{proof}

\begin{prob}
\olref{lem:delta0-completeness}માં
અસ્તિત્વવાચક કિસ્સાની વિગતવાર
સાબિતી આપો.
\end{prob}

\begin{thm}
\ollabel{thm:sigma1-completeness}
$!A$ એ $\Struct{N}$માં સાચું
$\Sigma_1$ વાક્ય હોય,
તો $\Th{Q} \Proves !A$.
\end{thm}

\begin{proof}
$\lexists[x][!A(x)]$ એ
$\Struct{N}$માં સાચું
$\Sigma_1$ વાક્ય
હોય, તો કોઈ પ્રાકૃતિક
સંખ્યા~$n$ અને ચલ-નિયુક્તિ~$s$
એવાં છે કે $s(x) = n$ અને
$\Sat{N}{!A(x)}[s]$.
સંતોષસંબંધનાં પ્રમાણભૂત
તથ્યો વડે $\Sat{N}{!A(\num n)}$
મળે છે. પરંતુ $!A(\num n)$
$\Delta_0$ સૂત્ર છે;
તેથી \olref{lem:delta0-completeness}
પ્રમાણે $\Th{Q} \Proves
!A(\num n)$ અને પરિણામે
તર્કથી $\Th{Q} \Proves
\lexists[x][!A(x)]$ પણ છે.
\sourcecorrection{OLINC-037}{સ્થિર અંગ્રેજી
મૂળ અહીં $\lexists{x}!A(x)$
લખે છે; આ મેક્રો ચલ અને
વ્યાપ માટે ચોરસ કૌંસવાળા
વૈકલ્પિક ઘટકો લે છે.
વ્યાખ્યા અને સાબિતીના
છેલ્લા સૂત્ર પ્રમાણે
$\lexists[x][!A(x)]$
આકાર સ્પષ્ટ કર્યો છે.}
\end{proof}
% END OLP-0300
\endgroup


\OLEndChapterHook


\begin{editorial}
આગળનો સંબંધિત અંશ મૂળના સક્રિય આયાતક્રમમાં નથી. તેને અહીં સ્પષ્ટ વૈકલ્પિક સંદર્ભમાં જાળવ્યો છે.
\end{editorial}
\begingroup
\makeatletter\def\ol@sectioncs{section}\makeatother

% BEGIN OLP-0647
\olfileid{inc}{req}{cre}
\olsection{$C$નાં વિધેયોનું $\Th{Q}$માં પ્રતિનિધિત્વ કરી શકાય છે}
\phantomsection\label{guunit:OLP-0647}


$C$માંના દરેક વિધેયનું
$\Th{Q}$માં પ્રતિનિધિત્વ
કરી શકાય છે તે
બતાવવાનું છે. અંતે,
$C$માંના દરેક $k$-સ્થાની
વિધેય $f(x_0,\dots,x_{k-1})$
માટે તેનું પ્રતિનિધિત્વ
કરતું સૂત્ર
$!A_f(x_0,\dots,x_{k-1},y)$
કેવી રીતે આપવું
તે બતાવવું પડશે.

\begin{lem}
પ્રાકૃતિક સંખ્યાઓ $n$
અને $m$ માટે, જો
$n \neq m$, તો
$Q \Proves \num
n \neq \num m$.
\end{lem}

\begin{proof}
$n$ પરના આગમનથી
બતાવો કે દરેક $m$
માટે, જો $n \neq m$,
તો $Q \Proves \eq/[\num n][\num m]$.

પાયાના કિસ્સામાં
$n = 0$. જો $m$
એ $0$ નથી, તો
કોઈ પ્રાકૃતિક
સંખ્યા $k$ માટે
$m = k + 1$.
આપણી પાસે
$\lforall[x][\eq/[0][x']]$
એવું સ્વયંસિદ્ધિ છે.
પરિમાણક-સ્વયંસિદ્ધિથી
$x$ની જગ્યાએ
$\num k$ મૂકતાં
$\eq/[0][\num k']$
મળે છે. પરંતુ
$\num k'$ એ
$\num m$ જ છે.

આગમન-પગલામાં
દાવો $n$ માટે સાચો
છે એમ માની $n+1$
વિચારો. $m$ કોઈ
પણ પ્રાકૃતિક
સંખ્યા હોય. બે
શક્યતાઓ છે:
કાં તો $m = 0$,
અથવા કોઈ $k$
માટે $m = k+1$.
પહેલો કિસ્સો
ઉપર મુજબ છે.
બીજા કિસ્સામાં
$n+1 \neq k+1$
માનો; ત્યારે
$n \neq k$.
$n$ માટેની
આગમન-પરિકલ્પનાથી
$\Th{Q} \Proves \eq/[\num n][\num k]$.
આપણી પાસે
$\lforall[x][\lforall[y][\eq[x'][y'] \lif \eq[x][y]]]$
સ્વયંસિદ્ધિ છે.
પરિમાણક-સ્વયંસિદ્ધિ
વાપરી
$\eq[\num n'][\num k'] \lif \eq[\num n][\num
  k]$ મળે. વિધાનાત્મક
તર્કથી $\Th{Q}$માં
$\eq/[\num n][\num k] \lif \eq/[\num n'][\num k']$
મળે છે. મોડસ
પોનેન્સથી ઇચ્છિત
$\eq/[\num n'][\num k']$
મળે છે, કારણ કે
$\num k'$ એ
$\num m$ છે.
\end{proof}

આ લેમ્મા બહુ મોટો
દાવો કરતું નથી:
મૂળભૂત રીતે, જુદા
સંખ્યાચિહ્નો જુદી
વસ્તુઓ નિર્દેશે છે
એ $\Th{Q}$ સાબિત
કરી શકે છે. દાખલા
તરીકે, $\Th{Q}$
$0'' \neq 0'''$
સાબિત કરે છે.
પરંતુ સામાન્ય કિસ્સો
બતાવવા થોડી
કાળજી જોઈએ.
આગમન વાપરીએ
છીએ, પણ તે
$\Th{Q}$ની
\emph{બહારનું}
આગમન છે.

શૂન્ય, ઉત્તરગામી વિધેય, સરવાળો, ગુણાકાર, સમાનતા માટેનું
લાક્ષણિક વિધેય અને પ્રક્ષેપોનું પ્રતિનિધિત્વ કરી શકાશે.
દરેક કિસ્સામાં પ્રતિનિધિત્વ કરતું યોગ્ય સૂત્ર સીધું છે.
દાખલા તરીકે, શૂન્યનું પ્રતિનિધિત્વ આ સૂત્ર કરે છે:
\[
y = 0,
\]
ઉત્તરગામી વિધેયનું આ સૂત્ર કરે છે:
\[
x_0' = y,
\]
અને સરવાળાનું આ સૂત્ર કરે છે:
\[
x_0 + x_1 = y.
\]
અહીં કામ એ સાબિત કરવાનું છે કે $\Th{Q}$ સંબંધિત વિધાનો
સાબિત કરી શકે છે. દાખલા તરીકે, ઉપરનું સૂત્ર સરવાળાનું
પ્રતિનિધિત્વ કરે છે એમ કહેવા માટે દરેક પ્રાકૃતિક સંખ્યા
$m$ અને $n$ માટે $\Th{Q}$ નીચેનાં બંને સૂત્રો સાબિત કરે છે
એ બતાવવું પડે:
\[
\eq[\num n + \num m][\num {n+m}]
\]
અને
\[
\lforall[y][(\eq[(\num n + \num m)][y] \lif \eq[y][\num {n+m}])].
\]

હવે સંયોજનનું શું? ધારો કે $h$ આમ વ્યાખ્યાયિત છે:
\[
h(x_0,\dots,x_{l-1}) = f(g_0(x_0,\dots,x_{l-1}), \dots,
g_{k-1}(x_0,\dots,x_{l-1})).
\]
અહીં અનુક્રમે $f,g_0,\dots,g_{k-1}$નું પ્રતિનિધિત્વ
કરતાં સૂત્રો $!A_f, !A_{g_0}, \dots,
!A_{g_{k-1}}$ આપણે પહેલેથી મેળવ્યાં છે. ત્યારે
$h$નું પ્રતિનિધિત્વ કરતું સૂત્ર $!A_h$ મેળવવા માટે
$!A_h(x_0,\dots,x_{l-1},y)$ને આમ વ્યાખ્યાયિત કરી શકીએ:
\begin{multline*}
  \lexists[z_0][\dots \lexists[z_{k-1}][(!A_{g_0}(x_0,\dots,x_{l-1},z_0) \land
  \dots \land !A_{g_{k-1}}(x_0,\dots,x_{l-1},z_{k-1}) \land\\
  !A_f(z_0,\dots,z_{k-1},y))]].
\end{multline*}
\sourcecorrection{OLMD-181}{મૂળ સૂત્રમાં અંતિમ $!A_f$ સંયોજક
અસ્તિત્વવાચક પરિમાણકોના વ્યાપની બહાર પડી જાય છે; અહીં બધા
મધ્યવર્તી $z_i$ અને અંતિમ સંયોજકને એ જ વ્યાપમાં રાખ્યાં છે.}

છેવટે અનિયત મર્યાદા સુધીની શોધ વિચારીએ. ધારો કે
$g(x,\vec z)$ નિયમિત છે અને $\Th{Q}$માં, માનો કે
$!A_g(x,\vec z,y)$ સૂત્રથી, તેનું પ્રતિનિધિત્વ કરી
શકાય છે. $f$ને $f(\vec z) = \mu x \; g(x,\vec z)$
વડે વ્યાખ્યાયિત કરો. $f$નું પ્રતિનિધિત્વ કરતું
$!A_f(\vec z,y)$ સૂત્ર શોધવું છે. સ્વાભાવિક પસંદગી આ છે:
\[
!A_f(\vec z,y) \ident !A_g(y,\vec z,0) \land \lforall[w][(w < y
\lif \lnot !A_g(w,\vec z,0))].
\]

નીચેના જેવાં વધારાનાં લેમ્મા વાપરી આ પસંદગી કામ કરે છે
એ બતાવી શકાય:

\begin{lem}
  દરેક ચલ $x$ અને દરેક પ્રાકૃતિક સંખ્યા $n$ માટે
  $\Th{Q}$ એ $\eq[(x'
  + \num n)][(x + \num n)']$ સાબિત કરે છે.
\end{lem}

ફરી નોંધવું જોઈએ કે આ દાવો $\Th{Q}$ એ
$\lforall[x][\lforall[y][(x' + y) = (x +y)']]$
સાબિત કરે છે એવા દાવા કરતાં નબળો છે; પછીનો દાવો ખોટો છે.

\begin{proof}
રાબેતા મુજબ સાબિતી $n$ પરના આગમનથી મળે છે.
પાયાના કિસ્સામાં $n = 0$ છે અને $\Th{Q}$ એ
$\eq[(x'+0)][(x+0)']$ સાબિત કરે છે એમ બતાવવાનું છે.
પરંતુ આપણી પાસે આ છે:
\begin{eqnarray*}
& & \eq[(x' + 0)][x'] \quad \text{સ્વયંસિદ્ધિ 4થી} \\
& & \eq[(x + 0)][x] \quad \text{સ્વયંસિદ્ધિ 4થી} \\
& & \eq[(x+0)'][x'] \quad \text{બીજી પંક્તિથી} \\
& & \eq[(x' + 0)][(x + 0)'] \quad \text{પહેલી અને ત્રીજી પંક્તિથી}
\end{eqnarray*}
આગમન-પગલામાં $\Th{Q}$માં
$\eq[(x' + \num n)][(x + \num n)']$ નિષ્પન્ન થયું
છે એમ ધારીએ. $\num{n+1}$ એ $\num n'$ હોવાથી,
$\Th{Q}$ એ $\eq[(x' + \num n')][(x + \num n')']$
સાબિત કરે છે એમ બતાવવાનું છે. $\Th{Q}$માં
સમાનતાઓની નીચેની સાંકળ નિષ્પન્ન કરી શકાય છે:
\begin{eqnarray*}
(x' + \num n') & \eq & (x' + \num n)' \quad \text{સ્વયંસિદ્ધિ 5થી} \\
& \eq & (x + \num n')' \quad \text{આગમન-પરિકલ્પનાથી}
\end{eqnarray*}
\sourcecorrection{OLMD-182}{બીજી સમાનતા માટે આગમન-પરિકલ્પના
ઉપરાંત $(x+\num n')=(x+\num n)'$ આપતું સ્વયંસિદ્ધિ 5
અને સમાનતાની અવેજી પણ જોઈએ; મૂળમાં આ મધ્યવર્તી
પગલાં સ્પષ્ટ લખ્યાં નથી.}
\end{proof}

\begin{lem}
\begin{enumerate}
\item $\Th{Q}$ એ $\lnot (x < \num 0)$ સાબિત કરે છે.
\item દરેક પ્રાકૃતિક સંખ્યા $n$ માટે $\Th{Q}$ એ
નીચેનું સાબિત કરે છે:
\[
x < \num {n+1} \lif (\eq[x][\num 0] \lor \dots \lor \eq[x][\num n]).
\]
\end{enumerate}
\end{lem}

\begin{proof}
પહેલો દાવો અને બીજા દાવાનો અમુક ભાગ અનૌપચારિક રીતે
સાબિત કરીએ; એટલે ઔપચારિક નિષ્પત્તિ કેવી રીતે
બનાવવી તેના માત્ર સંકેતો આપીશું.

પહેલા દાવા માટે $<$ની વ્યાખ્યા મુજબ $\Th{Q}$માં
$\lnot \lexists[y][\eq[(y'+x)][0]]$ સાબિત કરવાનું છે.
પ્રથમ-ક્રમ તર્કના સ્વયંસિદ્ધિઓ અને નિયમો વડે આ
$\lforall[y][\eq/[(y' + x)][0]]$ની સમકક્ષ છે.
વિચાર આ છે: $\eq[(y'+x)][0]$ ધારો. જો $x$ એ
શૂન્ય હોય, તો $\eq[(y' + 0)][0]$ મળે. પરંતુ
$\Th{Q}$ની સ્વયંસિદ્ધિ 4થી $\eq[(y' + 0)][y']$
મળે છે અને સ્વયંસિદ્ધિ 2થી $\eq/[y'][0]$
મળે છે; આ વિરોધાભાસ છે. જો $x$ એ $0$
ન હોય, તો સ્વયંસિદ્ધિ 3થી કોઈ $z$ માટે
$\eq[x][z']$ મળે છે. ત્યારે
$\eq[(y' + z')][0]$ મળે. સ્વયંસિદ્ધિ~5થી
$\eq[(y' + z)'][0]$ મળે, જે ફરી
સ્વયંસિદ્ધિ~2નો વિરોધ કરે છે. બંને કિસ્સાઓ
પૂરા થતાં $\lforall[y][\eq/[(y' +x)][0]]$ મળે છે.
\sourcecorrection{OLMD-183}{મૂળ ગદ્યમાં બંને
કિસ્સાઓ પૂરા થાય તે પહેલાં જ સર્વવાચી નિષ્કર્ષ
લખાયો છે; અહીં તેને બીજા કિસ્સા પછી મૂક્યો છે.}

બીજા દાવા માટે $n$ પર આગમન વાપરો. પાયાનો
કિસ્સો $n =0$ છે. ત્યારે
$x < \num 1 \lif \eq[x][\num 0]$ બતાવવાનું છે.
$x < \num 1$ ધારો. $<$ માટેની વ્યાખ્યાત્મક
સ્વયંસિદ્ધિથી $\lexists[y][\eq[(y'+x)][0']]$
મળે છે. એવો ગુણધર્મ ધરાવતું $y$ લો;
તેથી $\eq[y'+x][0']$ મળે છે.

$\eq[x][0]$ બતાવવાનું છે. સ્વયંસિદ્ધિ 3થી,
જો $x$ એ શૂન્ય ન હોય, તો કોઈ $z$ માટે તે
$z'$ને સમાન છે. તેથી $\eq[(y' + z')][0']$
મળે છે. $\Th{Q}$ની સ્વયંસિદ્ધિ~5થી
$\eq[(y'+z)'][0']$ મળે. સ્વયંસિદ્ધિ 1થી
$\eq[(y' + z)][0]$ મળે. પરંતુ વ્યાખ્યા મુજબ
તેનો અર્થ $z < 0$ થાય છે, જે પહેલા દાવાનો
વિરોધ કરે છે.
\sourcecorrection{OLMD-184}{મૂળ પાયાના કિસ્સામાં
$x<\num 1$થી $x=0$ દિશા જ સમજાવી છે. વિપરીત
દિશા $x=0$ હોય ત્યારે $<$ની વ્યાખ્યામાં
$y=0$ સાક્ષી લઈને મળે છે; આગળનું આગમન-પગલું
પણ આ જૂના સંક્ષિપ્ત વિભાગમાં લખાયું નથી.}
\end{proof}

\begin{explain}
આપણે બતાવ્યું છે કે સંગણનીય વિધેયોનો ગણ
$\Th{Q}$માં પ્રતિનિધિત્વ કરી શકાય એવાં
વિધેયોના ગણ તરીકે ઓળખી શકાય છે. હકીકતમાં
સાબિતી વધુ સામાન્ય છે. પ્રતિનિધિત્વની
વ્યાખ્યાથી સમજાય છે કે $\Th{Q}$નું વિસ્તરણ
હોય એવી કોઈ પણ થિયરી, અથવા જેમાં $\Th{Q}$નું
અર્થઘટન કરી શકાય એવી થિયરી, સંગણનીય
વિધેયોનું પ્રતિનિધિત્વ કરી શકે છે. બીજી
બાજુ, સાબિતીની સંકલ્પના સંગણનીય હોય એવી
કોઈ પણ નિષ્પત્તિ પ્રણાલીમાં દરેક
પ્રતિનિધિત્વ કરી શકાય એવું વિધેય
સંગણનીય છે. તેથી સંગણનીય વિધેયોના
ગણને પેયાનો અંકગણિતમાં, અથવા ઝેર્મેલો--ફ્રેન્કેલ
ગણસિદ્ધાંતમાં પણ, પ્રતિનિધિત્વ થતાં
વિધેયોના ગણ તરીકે ઓળખી શકાય છે.
ગેડેલે નોંધ્યું તેમ, આ કંઈક આશ્ચર્યજનક છે.
સાબિત કરી શકાય તેવા વિધાનો વિશેના પ્રશ્નો
પસંદ કરેલી થિયરી પર ખૂબ આધાર રાખે છે
એ આપણે જોઈશું: સામાન્ય રીતે સ્વયંસિદ્ધિઓ
વધુ પ્રબળ હોય તેમ વધુ સાબિત કરી શકાય.
પરંતુ સ્વયંસિદ્ધિ-આધારિત થિયરીઓના
વિશાળ વર્ગમાં પ્રતિનિધિત્વ કરી શકાય
એવાં વિધેયો ચોક્કસ સંગણનીય વિધેયો જ છે.
\end{explain}
% END OLP-0647
\endgroup

\begin{editorial}
આગળનો સંબંધિત અંશ મૂળના સક્રિય આયાતક્રમમાં નથી. તેને અહીં સ્પષ્ટ વૈકલ્પિક સંદર્ભમાં જાળવ્યો છે.
\end{editorial}
\begingroup
\makeatletter\def\ol@sectioncs{section}\makeatother

% BEGIN OLP-0648
\olfileid{inc}{req}{cee}
\olsection{$C$નાં વિધેયો}
\phantomsection\label{guunit:OLP-0648}


$C$ એ નીચેના વિધેયો ધરાવતો સૌથી નાનો વિધેય-ગણ હોય:
\begin{enumerate}
\item $0$,
\item ઉત્તરગામી વિધેય,
\item સરવાળો,
\item ગુણાકાર,
\item પ્રક્ષેપો, અને
\item સમાનતા માટેનું લાક્ષણિક વિધેય, $\Char{=}$;
\end{enumerate}
અને જે નીચેની ક્રિયાઓ હેઠળ બંધ હોય:
\begin{enumerate}
\item સંયોજન, અને
\item નિયમિત વિધેયો પર લાગુ પડતી અનિયત મર્યાદા સુધીની શોધ.
\end{enumerate}
યાદ રાખો કે છેલ્લી મર્યાદાનો અર્થ એટલો જ છે કે
પરિણામી વિધેય સંપૂર્ણ હોય ત્યારે જ $\mu$ ક્રિયા
વાપરી શકાય. તેની સરખામણી \emph{સામાન્ય પુનરાવર્તી}
વિધેયોની વ્યાખ્યા સાથે કરો: અહીં આપણે સરવાળો, ગુણાકાર
અને $\Char{=}$ ઉમેર્યાં છે, પરંતુ આદિમ પુનરાવર્તન
દૂર કર્યું છે.

$C$માંનું દરેક વિધેય પુનરાવર્તી છે એ સ્પષ્ટ છે,
કારણ કે સરવાળો, ગુણાકાર અને $\Char{=}$ પણ
પુનરાવર્તી છે. વિપરીત દિશા પણ સાચી છે એ આપણે
બતાવીશું; એટલે $C$ની અન્ય સામગ્રી વડે આદિમ
પુનરાવર્તન કરી શકાય છે.
% END OLP-0648
\endgroup
% END OLP-0289
\endgroup


\begingroup
\makeatletter\def\ol@sectioncs{section}\makeatother

% BEGIN OLP-0301
\begin{editorial}
  આ પ્રકરણ સંગણનક્ષમતા-સિદ્ધાંતના
  પ્રકરણની સામગ્રી પર આધાર રાખે છે;
  તે સામગ્રી ભણાવવામાં ન આવી
  હોય તો આ પ્રકરણ છોડી શકાય.
  હાલમાં તે જેરેમી અવિગાડની
  નોંધોનું મૂળભૂત રૂપાંતર છે;
  તેનું પુનરીક્ષણ થયું નથી અને
  અભ્યાસપ્રશ્નો પણ ખૂટે છે.
\end{editorial}

\olchapter{inc}{tcp}{સિદ્ધાંતો અને સંગણનક્ષમતા}
\olchapterid{inc}{tcp}
\phantomsection\label{guunit:OLP-0301}


\begingroup
\makeatletter\def\ol@sectioncs{section}\makeatother

% BEGIN OLP-0302
\olfileid{inc}{tcp}{int}

\olsection{પરિચય}
\phantomsection\label{guunit:OLP-0302}


\begin{editorial}
આ વિભાગ ફરીથી લખવો જોઈએ.
\end{editorial}

આપણી પાસે નીચેની બાબતો છે:
\begin{enumerate}
\item $\Th{Q}$માં વિધેય નિરૂપ્ય
  હોવાનો અર્થ આપતી વ્યાખ્યા
  (\olref[req][int]{defn:representable-fn})
\item $\Th{Q}$માં સંબંધ નિરૂપ્ય
  હોવાનો અર્થ આપતી વ્યાખ્યા
  (\olref[req][rel]{defn:representing-relations})
\item $\Th{Q}$માં નિરૂપ્ય વિધેયો
  બરાબર સંગણનીય વિધેયો જ છે
  એવું કહેતો પ્રમેય
  (\olref[req][int]{thm:representable-iff-comp})
\item $\Th{Q}$માં નિરૂપ્ય સંબંધો
  બરાબર સંગણનીય સંબંધો જ છે
  એવું કહેતો પ્રમેય
  (\olref[req][rel]{thm:representing-rels})
  \sourcecorrection{OLINC-038}{સ્થિર અંગ્રેજી
  મૂળની આ છેલ્લી યાદી-બાબતમાં
  સંદર્ભની શરૂઆતનો ગોળ કૌંસ
  ગાયબ છે; અગાઉની ત્રણ
  બાબતોની જેમ અહીં બંને
  કૌંસ આપ્યા છે.}
\end{enumerate}

\emph{સિદ્ધાંત} એ નિગમન હેઠળ
બંધ વાક્યોનો ગણ છે; એટલે
$T$માંથી $!A$ સાબિત થાય
ત્યારે $!A$ એ~$T$માં જ
હોય છે. સિદ્ધાંતને વાક્યોનો
એવો સંગ્રહ સમજવો સારું છે,
જેમાં તે વાક્યોમાંથી નીકળતાં
બધાં તારણો પણ સમાવિષ્ટ છે.
હવેથી $\Th{Q}$ વડે
\olref[req][int]{sec}ની
આઠ સ્વયંસિદ્ધિઓમાંથી નિગમ્ય
વાક્યોના ગણરૂપ~\emph{સિદ્ધાંત}
દર્શાવીશું. યાદ રાખો કે
$\Th{Q}$ની ભાષાનાં સૂત્રોને
સંખ્યાઓ વડે સંકેતિત કરી
શકાય છે; $!A$ એવું સૂત્ર
હોય, તો $\Gn{!A}$ એ~$!A$ને
સંકેતિત કરતી સંખ્યા છે.
આ સંકેતીકરણને ધ્યાનમાં રાખીને
હવે વિવિધ સૂત્રગણો સંગણનીય
છે કે નહીં તે પૂછી શકીએ.
% END OLP-0302
\endgroup


\begingroup
\makeatletter\def\ol@sectioncs{section}\makeatother

% BEGIN OLP-0303
\olfileid{inc}{tcp}{qce}

\olsection{$\Th{Q}$ એ સંગણકીય રીતે પરિગણનીય-પૂર્ણ છે}
\phantomsection\label{guunit:OLP-0303}



\begin{thm}
$\Th{Q}$ સંગણકીય રીતે પરિગણનીય છે,
પણ નિર્ણેય નથી. હકીકતમાં તે
પૂર્ણ સંગણકીય રીતે પરિગણનીય ગણ છે.
\end{thm}

\begin{proof}
$\Th{Q}$ સંગણકીય રીતે પરિગણનીય છે
તે જોવું મુશ્કેલ નથી, કારણ કે
તે એવા વાક્યોના (સંકેતોના)
ગણ~$y$ છે કે $\Th{Q}$માં
$y$ની કોઈ સાબિતી~$x$ છે:
\[
Q = \Setabs{y}{\lexists[x][\Prf[\Th{Q}](x,y)]}.
\]
પરંતુ આપણે જાણીએ છીએ કે
$\Prf[\Th{Q}](x,y)$ સંગણનીય
છે (હકીકતમાં આદિમ પુનરાવર્તી),
અને ઉપરના આકારમાં લખી શકાય
એવો દરેક ગણ સંગણકીય રીતે
પરિગણનીય છે.

તે પૂર્ણ સંગણકીય રીતે
પરિગણનીય ગણ છે એમ કહેવું
$K \leq_m Q$ કહેવાની સમકક્ષ
છે, જ્યાં $K = \Setabs{x}{!A_x(x) \downarrow}$.
તેથી $K$નું $\Th{Q}$માં
ન્યૂનીકરણ થાય છે તે બતાવીએ.
ક્લીનીનું વિધેયધર્મ
$T(e,x,s)$ આદિમ પુનરાવર્તી
હોવાથી તે $\Th{Q}$માં,
કહો કે~$!A_T$ વડે, નિરૂપ્ય
છે. તેથી દરેક~$x$ માટે
\begin{align*}
x \in K & \lif \lexists[s][T(x,x,s)] \\
& \lif \lexists[s][(\Th{Q} \Proves !A_T(\num x,\num x, \num s))]
  \\
& \lif \Th{Q} \Proves \lexists[s][!A_T(\num x, \num x, s)].
\end{align*}
વિપરીત દિશામાં, જો
$\Th{Q} \Proves \lexists[s][!A_T(\num x, \num x, s)]$,
તો હકીકતમાં કોઈ પ્રાકૃતિક
સંખ્યા~$n$ માટે સૂત્ર
$!A_T(\num x, \num x, \num n)$
સાચું હોવું જ પડે. હવે
$T(x,x,n)$ ખોટું હોય, તો
$!A_T$ એ~$T$ને નિરૂપે છે
એટલે $\Th{Q}$માં
$\lnot !A_T(\num x, \num x, \num n)$
સાબિત થાય. પણ ત્યારે
$\Th{Q}$માં ખોટું સૂત્ર
સાબિત થતું હોત; આ
વિરોધાભાસ છે. તેથી
$T(x,x,n)$ સાચું જ હોવું
જોઈએ, અને પરિણામે
$!A_x(x) \downarrow$.

ટૂંકમાં, દરેક~$x$ માટે
$x$ એ~$K$માં છે ત્યારે અને
માત્ર ત્યારે જ $\Th{Q}$માં
$\lexists[s][!A_T(\num x,\num x,s)]$
સાબિત થાય છે. તેથી~$x$ને
વાક્ય $\lexists[s][!A_T(\num x,
  \num x, s)]$નો સંકેત આપતું
વિધેય~$f$ એ~$K$નું
$\Th{Q}$માં ન્યૂનીકરણ છે.
\sourcecorrection{OLINC-039}{સ્થિર અંગ્રેજી
મૂળ આ અંતિમ બે વાક્યોમાં
અંકગણિતની ભાષામાં વપરાતું
નિરૂપક સૂત્ર~$!A_T$ બદલે
બાહ્ય ક્લીની વિધેયધર્મ~$T$
સૂત્રની અંદર લખે છે.
ન્યૂનીકરણનું વાક્ય અગાઉની
સાબિતી પ્રમાણે $!A_T$
વડે જ રચ્યું છે.}
\end{proof}
% END OLP-0303
\endgroup


\begingroup
\makeatletter\def\ol@sectioncs{section}\makeatother

% BEGIN OLP-0304
\olfileid{inc}{tcp}{oqn}

\olsection{$\Th{Q}$નાં $\omega$-સુસંગત વિસ્તરણો અનિર્ણેય છે}
\phantomsection\label{guunit:OLP-0304}


\begin{explain}
$\Th{Q}$ સંગણકીય રીતે
પરિગણનીય ગણોમાં પૂર્ણ છે
તેની સાબિતી આ તથ્ય પર
આધાર રાખતી હતી કે
$\Th{Q}$માં સાબિત થતું
દરેક વાક્ય પ્રાકૃતિક
સંખ્યાઓ વિશે ``સાચું'' છે.
આગળની વ્યાખ્યા અને પ્રમેય
ઉપરની સાબિતી માટે જરૂરી
``સત્ય''ના ચોક્કસ પાસાઓ
ઓળખીને તે પરિણામને વધુ
મજબૂત બનાવે છે. જોકે આ
પ્રમેય પર વધારે સમય રોકાવાની
જરૂર નથી, કારણ કે ટૂંકમાં
તેને હજુ વધુ મજબૂત કરીશું.
અહીં તેને મુખ્યત્વે ઐતિહાસિક
હેતુથી સામેલ કર્યો છે:
ગડેલના મૂળ લેખમાં
$\omega$-સુસંગતતાની
વિભાવના વપરાઈ હતી, પરંતુ
થોડા સમય પછી સામાન્ય
સુસંગતતા વડે
$\omega$-સુસંગતતાને
બદલીને તેનું પરિણામ વધુ
મજબૂત બનાવાયું.
\end{explain}

\begin{defn}
\ollabel{thm:oconsis-q}
સિદ્ધાંત~$\Th{T}$ને
$\omega$-સુસંગત ત્યારે
કહીએ જ્યારે નીચેની શરત હોય:
$\lexists[x][!A(x)]$ કોઈપણ
વાક્ય હોય અને $\Th{T}$માં
$\lnot !A(\num 0)$,
$\lnot !A(\num 1)$,
$\lnot !A(\num 2)$,
\dots સાબિત થાય, તો
$\Th{T}$માં
$\lexists[x][!A(x)]$
સાબિત ન થાય.
\end{defn}

\begin{thm}
  $\Th{T}$ એ $\Th{Q}$ને
  સમાવતા કોઈપણ
  $\omega$-સુસંગત સિદ્ધાંત
  હોય. તો $\Th{T}$
  નિર્ણેય નથી.
\end{thm}

\begin{proof}
$\Th{T}$ એ $\Th{Q}$ને
સમાવતું હોય, તો
$\Th{T}$ સંગણનીય વિધેયો
અને સંબંધોને નિરૂપે છે.
અગાઉની સાબિતીમાં માત્ર
થોડો ફેરફાર કરવો પડે.
ઉપરની જેમ $x \in K$ હોય,
તો $\Th{T}$માં
$\lexists[s][!A_T(\num
  x,\num x, s)]$ સાબિત થાય.
વિપરીત રીતે માનો કે
$\Th{T}$માં
$\lexists[s][!A_T(\num x, \num x, s)]$
સાબિત થાય છે. તો $x$ એ
$K$માં હોવું જ પડે: અન્યથા
$x$ મશીનની $x$ ઇનપુટ
પર અટકતી સંગણના જ ન હોય.
$!A_T$ ક્લીનીના~$T$
સંબંધને નિરૂપે છે, એટલે
$\Th{T}$માં
$\lnot !A_T(\num x, \num x, \num 0)$,
$\lnot !A_T(\num x, \num
x, \num 1)$, \dots સાબિત
થાય. આથી $\Th{T}$
$\omega$-અસુસંગત થઈ જાય,
જે ધારણાની વિરુદ્ધ છે.
\end{proof}
% END OLP-0304
\endgroup


\begingroup
\makeatletter\def\ol@sectioncs{section}\makeatother

% BEGIN OLP-0305
\olfileid{inc}{tcp}{cqn}

\olsection{$\Th{Q}$નાં સુસંગત વિસ્તરણો અનિર્ણેય છે}
\phantomsection\label{guunit:OLP-0305}


\begin{explain}
યાદ રાખો કે સિદ્ધાંત
\emph{સુસંગત} ત્યારે કહેવાય
જ્યારે તે કોઈપણ
સૂત્ર~$!A$ માટે
$!A$ અને $\lnot !A$
બંને સાબિત ન કરે.
વિરોધાભાસમાંથી કંઈપણ તારવી
શકાય છે, એટલે અસુસંગત
સિદ્ધાંત તુચ્છ હોય છે:
દરેક વાક્ય સાબિત થાય છે.
સ્પષ્ટ છે કે સિદ્ધાંત
$\omega$-સુસંગત હોય તો
સુસંગત પણ હોય. પરંતુ
સુસંગતતા નબળી શરત છે
(એટલે સુસંગત છતાં
$\omega$-અસુસંગત સિદ્ધાંતો
હોય છે). વધુ મજબૂત પ્રમેય
મેળવવા
\olref[oqn]{thm:oconsis-q}માં
લેવાયેલી ધારણાને સામાન્ય
સુસંગતતા સુધી નબળી કરી
શકીએ.
\end{explain}

\begin{lem}
``સાર્વત્રિક સંગણનીય સંબંધ''
હોતો નથી. એટલે એવો કોઈ
દ્વિચલ સંગણનીય સંબંધ
$R(x,y)$ નથી કે જે નીચેનો
ગુણધર્મ ધરાવે: $S(y)$
કોઈપણ એકચલ સંગણનીય
સંબંધ હોય, તો એવું~$k$
હોય કે દરેક~$y$ માટે
$S(y)$ સાચું ત્યારે અને
માત્ર ત્યારે જ $R(k,y)$
સાચું હોય.
\end{lem}

\begin{proof}
માનો કે $R(x,y)$ સાર્વત્રિક
સંગણનીય સંબંધ છે.
$S(y)$ તરીકે $\lnot R(y,y)$
સંબંધ લો. $S(y)$ સંગણનીય
હોવાથી કોઈ~$k$ માટે
$S(y)$ એ $R(k,y)$ની
સમકક્ષ છે. પરંતુ ત્યારે
$S(k)$ એ $R(k,k)$ અને
$\lnot R(k,k)$ બંનેની
સમકક્ષ થાય છે; આ
વિરોધાભાસ છે.
\end{proof}


\begin{thm}
$\Th{T}$ એ $\Th{Q}$ને
સમાવતો કોઈપણ સુસંગત
સિદ્ધાંત હોય. તો $\Th{T}$
નિર્ણેય નથી.
\end{thm}


\begin{proof}
માનો કે $\Th{T}$ એ
$\Th{Q}$નું સુસંગત,
નિર્ણેય વિસ્તરણ છે.
$\Th{T}$ વડે સાર્વત્રિક
સંગણનીય સંબંધ વ્યાખ્યાયિત
કરીને વિરોધાભાસ મેળવીશું.

$R(x,y)$ નીચેની શરત હોય
ત્યારે અને માત્ર ત્યારે જ
સાચો લો:
\begin{quote}
$x$ એ સૂત્ર $!D(u)$નો
સંકેત છે અને $\Th{T}$માં
$!D(\num y)$ સાબિત થાય છે.
\end{quote}
$\Th{T}$ નિર્ણેય છે એવી
ધારણા હોવાથી $R$ સંગણનીય
છે. હવે $R$ સાર્વત્રિક છે
તે બતાવીએ. $S(y)$ કોઈપણ
સંગણનીય સંબંધ હોય, તો
તેને $\Th{Q}$માં, એટલે
$\Th{T}$માં પણ,
$!D_S(u)$ સૂત્ર નિરૂપે છે.
તેથી દરેક~$n$ માટે
\begin{eqnarray*}
S(\num n) & \lif & T \vdash !D_S(\num n) \\
& \lif & R(\Gn{!D_S(u)}, n)
\end{eqnarray*}
અને
\begin{eqnarray*}
\lnot S(\num n) & \lif & T \vdash \lnot !D_S(\num n) \\
& \lif & T \not\vdash !D_S(\num n) \quad \text{(કારણ કે $\Th{T}$ સુસંગત છે)} \\
& \lif & \lnot R(\Gn{!D_S(u)}, n).
\end{eqnarray*}
એટલે દરેક~$y$ માટે
$S(y)$ સાચું ત્યારે અને
માત્ર ત્યારે જ
$R(\Gn{!D_S(u)}, y)$
સાચું છે. તેથી $R$
સાર્વત્રિક છે, અને માંગેલો
વિરોધાભાસ મળી ગયો.
\end{proof}

``સાચું અંકગણિત'' એટલે
સિદ્ધાંત~$\Setabs{!A}{ \Sat{\Nat}{!A}}$,
અર્થાત્ પ્રમાણભૂત અર્થઘટનમાં
સાચાં અંકગણિતીય વાક્યોનો ગણ.

\begin{cor}
સાચું અંકગણિત નિર્ણેય નથી.
\end{cor}
% END OLP-0305
\endgroup


\begingroup
\makeatletter\def\ol@sectioncs{section}\makeatother

% BEGIN OLP-0306
\olfileid{inc}{tcp}{cax}

\olsection{સ્વયંસિદ્ધીકરણીય સિદ્ધાંતો}
\phantomsection\label{guunit:OLP-0306}


સિદ્ધાંત~$\Th{T}$ પાસે
સ્વયંસિદ્ધિઓનો સંગણનીય ગણ~$A$
હોય તો તેને
\emph{સ્વયંસિદ્ધીકરણીય} કહીએ.
($A$ એ~$\Th{T}$નો
સ્વયંસિદ્ધિઓનો ગણ છે એમ
કહેવાનો અર્થ
$T = \Setabs{!A}{A \Proves !A}$.)
પ્રાકૃતિક સંખ્યાઓના કોઈપણ
``વાજબી'' સ્વયંસિદ્ધીકરણમાં
આ ગુણધર્મ હશે. ખાસ કરીને
સાન્ત સંખ્યામાં સ્વયંસિદ્ધિઓ
ધરાવતો દરેક સિદ્ધાંત
સ્વયંસિદ્ધીકરણીય છે.

\begin{lem}
માનો કે $\Th{T}$
સ્વયંસિદ્ધીકરણીય છે.
તો $\Th{T}$
સંગણકીય રીતે પરિગણનીય છે.
\end{lem}

\begin{proof}
માનો કે~$A$ એ~$\Th{T}$
માટે સ્વયંસિદ્ધિઓનો
સંગણનીય ગણ છે.
$!A \in T$ હોય તેની
પુષ્ટિ માટે સ્વયંસિદ્ધિઓમાંથી
$!A$નું નિષ્પત્તિ
શોધતા રહો.
\sourcecorrection{OLINC-040}{સ્થિર અંગ્રેજી
મૂળ અહીં શોધથી $!A\in T$
છે કે નહીં ``નક્કી''
થઈ જાય એવું કહે છે.
આ શોધ સભ્યતા હોય ત્યારે
સાબિતી શોધી આપે છે, પરંતુ
સભ્યતા ન હોય ત્યારે અટકવી
જરૂરી નથી; લેમ્માનો દાવો
પરિગણનીયતાનો છે, નિર્ણેયતાનો નહીં.}

થોડું જુદી રીતે કહીએ તો
$!A$ એ $\Th{T}$માં છે
ત્યારે અને માત્ર ત્યારે જ
$A$માં $!B_1$, \dots,
$!B_k$ સ્વયંસિદ્ધિઓની
સાન્ત યાદી અને પ્રથમ-ક્રમના
તર્કમાં
$(!B_1 \land \dots \land !B_k) \lif !A$
નું નિષ્પત્તિ છે.
પરંતુ આપણે અગાઉથી જાણીએ
છીએ કે ``એવા \dots છે કે
\dots'' આકારની વ્યાખ્યા
ધરાવતો દરેક ગણ સંગણકીય રીતે પરિગણનીય
છે, જો બીજું ``\dots''
સંગણનીય હોય.
\end{proof}
% END OLP-0306
\endgroup


\begingroup
\makeatletter\def\ol@sectioncs{section}\makeatother

% BEGIN OLP-0307
\olfileid{inc}{tcp}{cdc}

\olsection{સ્વયંસિદ્ધીકરણીય પૂર્ણ સિદ્ધાંતો નિર્ણેય છે}
\phantomsection\label{guunit:OLP-0307}


સિદ્ધાંતને {\em પૂર્ણ} ત્યારે કહીએ
જ્યારે દરેક વાક્ય~$!A$ માટે
$!A$ અથવા $\lnot !A$
બેમાંથી એક સાબિત થાય.

\begin{lem}
માનો કે સિદ્ધાંત~$\Th{T}$
પૂર્ણ અને સ્વયંસિદ્ધીકરણીય છે.
તો $\Th{T}$ નિર્ણેય છે.
\end{lem}

\begin{proof}
માનો કે $\Th{T}$ પૂર્ણ છે
અને $A$ તેના સ્વયંસિદ્ધિઓનો
સંગણનીય ગણ છે.
જો $\Th{T}$ અસુસંગત હોય,
તો તે સ્પષ્ટપણે સંગણનીય છે.
(કલનવિધિ: ``હા'' કહો.)
આથી માની શકીએ કે
$\Th{T}$ સુસંગત પણ છે.

વાક્ય~$!A$ એ~$\Th{T}$માં
છે કે નહીં તે નક્કી કરવા,
$\Th{T}$માંથી~$!A$નું
નિષ્પત્તિ અને
$\lnot !A$નું નિષ્પત્તિ
બંને એકસાથે શોધો.
$\Th{T}$ પૂર્ણ હોવાથી
બેમાંથી એક મળી જ જશે;
અને $\Th{T}$ સુસંગત હોવાથી,
જો $\lnot !A$નું
નિષ્પત્તિ મળે તો
$!A$નું નિષ્પત્તિ
હોઈ શકે નહીં.

જુદી રીતે કહીએ તો
$\Th{T}$ સંગણકીય રીતે પરિગણનીય છે
તે આપણે જાણીએ છીએ;
અગાઉ સાબિત કરેલા પ્રમેય
મુજબ $\Th{T}$નો પૂરક પણ
સંગણકીય રીતે પરિગણનીય છે એટલું
બતાવવું પૂરતું છે.
પરંતુ સૂત્ર~$!A$ એ~$\Th{\bar
T}$માં ત્યારે અને માત્ર
ત્યારે જ હોય જ્યારે
$\lnot !A$ એ~$\Th{T}$માં
હોય; તેથી
$\Th{\bar T} \leq_m
\Th{T}$.
\end{proof}
% END OLP-0307
\endgroup


\begingroup
\makeatletter\def\ol@sectioncs{section}\makeatother

% BEGIN OLP-0308
\olfileid{inc}{tcp}{inc}

\olsection{$\Th{Q}$ને પૂર્ણ, સુસંગત અને સ્વયંસિદ્ધીકરણીય
  વિસ્તરણો નથી}
\phantomsection\label{guunit:OLP-0308}


\begin{thm}
\ollabel{thm:first-incompleteness}
$\Th{Q}$નું પૂર્ણ, સુસંગત અને
સ્વયંસિદ્ધીકરણીય વિસ્તરણ
હોતું નથી.
\end{thm}

\begin{proof}
$\Th{Q}$નું કોઈ સુસંગત,
નિર્ણેય વિસ્તરણ નથી
તે આપણે પહેલેથી જાણીએ છીએ.
પરંતુ $\Th{T}$ પૂર્ણ અને
સ્વયંસિદ્ધીકૃત હોય,
તો તે નિર્ણેય છે.
\end{proof}

\begin{explain}
આ પ્રમેય પ્રથમ અપૂર્ણતા
પ્રમેયની ગોડેલે ૧૯૩૧માં
આપેલી મૂળ રજૂઆતથી
ખૂબ દૂર નથી.
વધુ આધુનિક પરિભાષા સિવાય,
મુખ્ય તફાવત આ છે:
ગોડેલે ``સુસંગત''ને બદલે
``$\omega$-સુસંગત''
વાપર્યું હતું; અને સંગણનીયતાનો
ઔપચારિક ખ્યાલ ત્યારે હજી
ઉપલબ્ધ ન હોવાથી તેઓ
સંપૂર્ણ સામાન્ય અર્થમાં
``સ્વયંસિદ્ધીકરણીય'' કહી
શકતા નહોતા. (પછીના
દાયકામાં, ગોડેલ સહિતના
સંશોધકોએ સંગણનીયતાના
ઔપચારિક નમૂનાઓ વિકસાવ્યા.
તેનું મોટું કારણ ગોડેલની
અપૂર્ણતાને આધીન થતા
સિદ્ધાંતોના પ્રકારો ઓળખવાનું
પણ હતું.)

પ્રમેય કહે છે કે પૂર્ણતા,
સુસંગતતા અને
સ્વયંસિદ્ધીકરણીયતા
ત્રણેય એકસાથે મળી શકતાં
નથી. છતાં આમાંથી કોઈ
એક છોડીએ તો બાકીનાં બે
મળી શકે છે: $\Th{Q}$
સુસંગત અને સંગણનીય રીતે
સ્વયંસિદ્ધિત છે, પરંતુ
પૂર્ણ નથી; અસુસંગત
સિદ્ધાંત પૂર્ણ છે અને
સંગણનીય રીતે સ્વયંસિદ્ધિત
છે (દાખલા તરીકે
$\{ \eq/[0][0] \}$ વડે),
પરંતુ સુસંગત નથી; અને
અંકગણિતનાં સાચાં વાક્યોનો
ગણ પૂર્ણ અને સુસંગત છે,
પરંતુ સંગણનીય રીતે
સ્વયંસિદ્ધિત નથી.
\end{explain}
% END OLP-0308
\endgroup


\begingroup
\makeatletter\def\ol@sectioncs{section}\makeatother

% BEGIN OLP-0309
\olfileid{inc}{tcp}{ins}

\olsection{$\Th{Q}$માં સાબિત અને નકારી શકાય એવાં વાક્યો
  સંગણનીય રીતે અવિભાજ્ય છે}
\phantomsection\label{guunit:OLP-0309}



$\Th{\bar Q}$ એ એવાં
વાક્યોનો ગણ માનો કે જેના
{\em નિષેધો} $\Th{Q}$માં
સાબિત થાય છે; એટલે
$\Th{\bar Q} = \Setabs{!A}{\Th{Q} \Proves
  \lnot !A}$.
યાદ રાખો કે વિચ્છિન્ન
ગણો $A$ અને $B$ને
સંગણનીય રીતે અવિભાજ્ય
ત્યારે કહીએ જ્યારે એવો
સંગણનીય ગણ~$C$ ન હોય
કે $A \subseteq C$ અને
$B \subseteq \Complement{C}$.

\begin{lem}
$\Th{Q}$ અને
$\Th{\bar Q}$ સંગણનીય
રીતે અવિભાજ્ય છે.
\end{lem}

\begin{proof}
માનો કે $C$ એવો સંગણનીય
ગણ છે કે $\Th{Q} \subseteq C$
અને $\Th{\bar Q} \subseteq \Complement{C}$.
$R(x,y)$ સંબંધ નીચે મુજબ લો:
\begin{quote}
$x$ એ સૂત્ર $!D(u)$નો
સંકેત છે અને $!D(\num y)$
એ $C$માં છે.
\end{quote}
$R(x,y)$ સાર્વત્રિક
સંગણનીય સંબંધ છે એમ
બતાવીને વિરોધાભાસ
મેળવીશું.

માનો કે $S(y)$ સંગણનીય
છે અને તેને $\Th{Q}$માં
$!D_S(u)$ નિરૂપે છે.
ત્યારે
\begin{eqnarray*}
S(\num n) & \lif & \Th{Q} \Proves !D_S(\num n) \\
& \lif & !D_S(\num n) \in C
\end{eqnarray*}
અને
\begin{eqnarray*}
\lnot S(\num n) & \lif & \Th{Q} \Proves \lnot !D_S(\num n) \\
& \lif & !D_S(\num n) \in \Th{\bar Q} \\
& \lif & !D_S(\num n) \not\in C
\end{eqnarray*}
આથી $S(y)$ એ
$R(\#(!D_S(u)),y)$ની
સમકક્ષ છે.
\sourcecorrection{OLINC-041}{સ્થિર અંગ્રેજી
મૂળના અંતિમ સંકેતમાં
$!D_S(\num u)$ લખાયું છે.
અહીં $R$નો પહેલો ઘટક
મુક્ત ચલવાળા સૂત્ર
$!D_S(u)$નો સંકેત છે;
ચલને અંકરૂપ આપતું
ચિહ્ન તેથી દૂર કર્યું છે.}
\end{proof}
% END OLP-0309
\endgroup


\begingroup
\makeatletter\def\ol@sectioncs{section}\makeatother

% BEGIN OLP-0310
\olfileid{inc}{tcp}{con}

\olsection{$\Th{Q}$ સાથે સુસંગત સિદ્ધાંતો અનિર્ણેય છે}
\phantomsection\label{guunit:OLP-0310}


નીચેનું પ્રમેય કહે છે કે
$\Th{Q}$ પોતે તો અનિર્ણેય
છે જ, પરંતુ $\Th{Q}$ સાથે
વિરોધ ન ધરાવતો કોઈપણ
સિદ્ધાંત અનિર્ણેય છે.

\begin{thm}
$\Th{T}$ એ અંકગણિતની
ભાષામાં એવો કોઈપણ સિદ્ધાંત
માનો કે જે $\Th{Q}$ સાથે
સુસંગત છે (એટલે
$\Th{T} \cup \Th{Q}$
સુસંગત છે). તો
$\Th{T}$ અનિર્ણેય છે.
\end{thm}

\begin{proof}
યાદ રાખો કે $\Th{Q}$ને
સ્વયંસિદ્ધિઓનો સાન્ત ગણ
$!Q_1$, \dots,
$!Q_8$ છે. આપણે તેમને
એક જ સ્વયંસિદ્ધિ
$!E = !Q_1 \land
\dots \land !Q_8$ વડે
બદલી પણ શકીએ.

માનો કે $\Th{T}$ એ
$\Th{Q}$ સાથે સુસંગત
નિર્ણેય સિદ્ધાંત છે.
આવો ગણ લો:
\[
C = \Setabs{!A}{\Th{T} \Proves !E \lif !A}.
\]
$C$ એ $\Th{Q}$ અને
$\Th{\bar Q}$ને અલગ પાડતો
સંગણનીય ગણ છે એવું બતાવીશું,
જે વિરોધાભાસ છે. પ્રથમ,
જો $!A$ એ $\Th{Q}$માં
હોય, તો $\Th{Q}$ની
સ્વયંસિદ્ધિઓમાંથી $!A$
સાબિત થાય છે; નિગમન
પ્રમેયથી પ્રથમ-ક્રમના
તર્કમાં $!E \lif !A$નું
નિષ્પત્તિ મળે છે.
આથી $!A$ એ~$C$માં છે.

બીજી બાજુ, જો $!A$ એ
$\Th{\bar Q}$માં હોય,
તો પ્રથમ-ક્રમના તર્કમાં
$!E \lif \lnot !A$ની
સાબિતી છે. જો $\Th{T}$
$!E \lif !A$ પણ સાબિત
કરે, તો $\Th{T}$ એ
$\lnot !E$ સાબિત કરે,
અને તેથી $\Th{T}
\cup \Th{Q}$ અસુસંગત બને.
પરંતુ $\Th{T} \cup \Th{Q}$
સુસંગત છે એવી આપણી
ધારણા છે. તેથી $\Th{T}$
$!E \lif !A$ સાબિત કરતું
નથી અને $!A$ એ~$C$માં
નથી.

અમે બતાવ્યું કે $!A$ જો
$\Th{Q}$માં હોય તો $C$માં
છે, અને જો $!A$ એ
$\Th{\bar Q}$માં હોય તો
$\Complement{C}$માં છે.
આથી $C$ સંગણનીય વિભાજક
છે, અને આપણને જોઈતો
વિરોધાભાસ મળે છે.
\end{proof}

આ પ્રમેય ઘણું શક્તિશાળી છે.
દાખલા તરીકે તેનાથી આ
અનુવિધાન મળે છે:

\begin{cor}
  અંકગણિતની ભાષામાં
  પ્રથમ-ક્રમનો તર્ક (એટલે
  $\Setabs{!A}{\text{પ્રથમ-ક્રમના તર્કમાં $!A$ સાબિત થાય છે}}$)
  અનિર્ણેય છે.
\end{cor}

\begin{proof}
પ્રથમ-ક્રમનો તર્ક એ
$\emptyset$નાં નિષ્કર્ષોનો
ગણ છે, અને તે
$\Th{Q}$ સાથે સુસંગત છે.
\end{proof}
% END OLP-0310
\endgroup


\begingroup
\makeatletter\def\ol@sectioncs{section}\makeatother

% BEGIN OLP-0311
\olfileid{inc}{tcp}{itp}

\olsection{$\Th{Q}$નું અર્થઘટન થઈ શકે એવા સિદ્ધાંતો અનિર્ણેય છે}
\phantomsection\label{guunit:OLP-0311}


આ પરિણામોને વધુ મજબૂત
બનાવી શકીએ. અનૌપચારિક રીતે,
એક ભાષા~$\Lang{L_1}$નું
બીજી ભાષા~$\Lang{L_2}$માં
અર્થઘટન એટલે
$\Lang{L_1}$ના વિશ્વ,
સંબંધ-ચિહ્નો અને
વિધેય-ચિહ્નોને
$\Lang{L_2}$નાં
સૂત્ર વડે વ્યાખ્યાયિત
કરવાં. અહીં એટલો સમય
લેતા નથી, છતાં આ
વ્યાખ્યાને ચોક્કસ બનાવી
શકાય છે.

\begin{thm}
  માનો કે $\Th{T}$ એવી
  ભાષાનો સિદ્ધાંત છે જેમાં
  અંકગણિતની ભાષાનું એવું
  અર્થઘટન શક્ય છે કે
  $\Th{T}$ એ $\Th{Q}$ના
  અર્થઘટન સાથે સુસંગત છે.
  તો $\Th{T}$ અનિર્ણેય છે.
  જો $\Th{T}$ એ
  $\Th{Q}$ની સ્વયંસિદ્ધિઓના
  અર્થઘટનને સાબિત કરે,
  તો $\Th{T}$નું કોઈ
  સુસંગત વિસ્તરણ નિર્ણેય
  નથી.
\end{thm}

સાબિતી અગાઉના પ્રમેયની
સાબિતીમાં નાનો ફેરફાર
કરીને મળે છે: વિરોધી
ઉદાહરણ હોય તો તેનાથી
$\Th{Q}$ અને
$\Th{\bar Q}$ને અલગ
પાડી શકીએ. પસંદગીના
સ્વયંસિદ્ધિ સહિતનો
ઝર્મેલો--ફ્રેન્કલ ગણસિદ્ધાંત
$\Th{ZFC}$ સામાન્ય ગણિતનો
મોટો ભાગ વિકસાવી શકાય
એવો શક્તિશાળી
સ્વયંસિદ્ધીય આધાર છે.
$\Th{ZFC}$માં પ્રાકૃતિક
સંખ્યાઓ વ્યાખ્યાયિત કરી
શકાય છે, અને આ અર્થઘટન
હેઠળ $\Th{Q}$ની
સ્વયંસિદ્ધિઓ સાચી છે.
આથી આપણને મળે છે:

\begin{cor}
$\Th{ZFC}$નું કોઈ
સુસંગત વિસ્તરણ નિર્ણેય નથી.
\sourcecorrection{OLINC-042}{સ્થિર અંગ્રેજી
મૂળના અનુવિધાનમાં
``સુસંગત'' શરત છૂટી છે.
અસુસંગત વિસ્તરણ દરેક
વાક્ય સાબિત કરે છે અને
તેથી નિર્ણેય હોય છે.
ઉપરના પ્રમેય પરથી મળતા
અનુવિધાનમાં સુસંગતતા
જરૂરી છે.}
\end{cor}

\begin{cor}
$\Th{ZFC}$નું પૂર્ણ,
સુસંગત અને સંગણનીય
રીતે સ્વયંસિદ્ધીકરણીય
વિસ્તરણ હોતું નથી.
\end{cor}

$\Th{ZFC}$ની ભાષામાં
માત્ર એક દ્વિચલ સંબંધ
$\in$ છે. (હકીકતમાં
સમતા પણ જરૂરી નથી.)
આથી મળે છે:

\begin{cor}
દ્વિચલ સંબંધ-ચિહ્ન
ધરાવતી કોઈપણ ભાષાનો
પ્રથમ-ક્રમનો તર્ક
અનિર્ણેય છે.
\end{cor}

આ પરિણામ બે એકચલ
વિધેય-ચિહ્ન ધરાવતી કોઈપણ
ભાષા માટે પણ લાગુ પડે છે,
કારણ કે તેમની મદદથી
દ્વિચલ સંબંધ-ચિહ્નનું
અનુકરણ કરી શકાય છે.
હમણાં આપેલાં પરિણામો
ચોક્કસ સીમાએ છે: માત્ર
\emph{એકચલ} સંબંધ-ચિહ્નો
અને વધુમાં વધુ એક
\emph{એકચલ} વિધેય-ચિહ્ન
ધરાવતી ભાષાનો પ્રથમ-ક્રમનો
તર્ક નિર્ણેય છે.

એક વધારાની માહિતી.
પ્રમાણભૂત નમૂનામાં સાચાં
હોય એવાં, ભાષા
$\Obj 0$, $'$, $+$,
$\times$, $<$નાં વાક્યોનો
ગણ અનિર્ણેય છે તે આપણે
જાણીએ છીએ. હકીકતમાં
$<$ને બાકીનાં ચિહ્નો
વડે વ્યાખ્યાયિત કરી
શકાય છે, અને પછી
$+$ને $\times$ અને $'$
વડે વ્યાખ્યાયિત કરી
શકાય છે. તેથી ભાષા
$\Obj 0$, $'$, $\times$નાં
સાચાં વાક્યોનો ગણ
અનિર્ણેય છે. બીજી બાજુ,
પ્રેસબર્ગરે બતાવ્યું કે
અંકગણિતની ભાષામાં
સાચાં હોય એવાં, ભાષા
$\Obj 0$, $'$, $+$નાં
વાક્યોનો ગણ નિર્ણેય છે.
જોકે આ નિર્ણય-પ્રક્રિયા
ગણતરી માટે વ્યવહારુ નથી.
% END OLP-0311
\endgroup


\OLEndChapterHook
% END OLP-0301
\endgroup


\begingroup
\makeatletter\def\ol@sectioncs{section}\makeatother

% BEGIN OLP-0312
\olchapter{inc}{inp}{અપૂર્ણતા અને સાબિતિક્ષમતા}
\olchapterid{inc}{inp}
\phantomsection\label{guunit:OLP-0312}


\begingroup
\makeatletter\def\ol@sectioncs{section}\makeatother

% BEGIN OLP-0313
\olfileid{inc}{inp}{int}

\olsection{પરિચય}
\phantomsection\label{guunit:OLP-0313}


હિલ્બર્ટ માનતા હતા કે
પ્રાકૃતિક સંખ્યાઓ જેવી
ગણિતીય રચના માટેની
સ્વયંસિદ્ધિ-પદ્ધતિ તે
રચના વિશેનાં બધાં સાચાં
વિધાનો તારવી ન શકે તો
અપૂરતી છે. પછી તેમને
ઔપચારિક નિગમન-પદ્ધતિઓમાં
રસ પડ્યો, તેની સાથે આ
વિચાર જોડીએ તો એમ લાગે
છે કે પ્રાકૃતિક સંખ્યાઓ
વિશે તર્ક કરવા વપરાતી
ઔપચારિક પદ્ધતિ માત્ર
સુસંગત જ નહીં, પણ
\emph{પૂર્ણ} પણ હોવી જોઈએ
એવું તેઓ માનતા હતા:
એટલે તેની ભાષામાંનું દરેક
વિધાન અથવા તેનો નિષેધ
નિષ્પન્ન કરી શકાય એવું હોય.
ગોડેલનું પ્રથમ અપૂર્ણતા
પ્રમેય બતાવે છે કે આવી
સ્વયંસિદ્ધિ-પદ્ધતિ નથી:
અંકગણિત માટે પૂર્ણ,
સુસંગત અને
સ્વયંસિદ્ધીકરણીય
ઔપચારિક પદ્ધતિ નથી.
હકીકતમાં ``પૂરતા મજબૂત'',
સુસંગત અને
સ્વયંસિદ્ધીકરણીય
ગણિતીય સિદ્ધાંતોમાંથી
કોઈપણ પૂર્ણ નથી.

હિલ્બર્ટનું એક વધુ મહત્વનું
ધ્યેય, આધુનિક
(``શાસ્ત્રીય'') ગણિતને
વાજબી ઠેરવવાના તેમના
કાર્યક્રમનું કેન્દ્ર, શાસ્ત્રીય
તર્કને રજૂ કરતી ઔપચારિક
પદ્ધતિઓ માટે સાન્તવાદી
સુસંગતતા-સાબિતીઓ શોધવાનું
હતું. તેથી હિલ્બર્ટના
કાર્યક્રમ માટે ગોડેલનું
બીજું અપૂર્ણતા પ્રમેય
ઘણો મોટો આઘાત હતો.
બીજા અપૂર્ણતા પ્રમેયને પણ
પહેલાની જેમ અચોક્કસ
શબ્દોમાં કહી શકાય.
આશરે તે કહે છે કે
અંકગણિતનો કોઈ પૂરતો
મજબૂત સિદ્ધાંત પોતાની
સુસંગતતા સાબિત કરી
શકતો નથી. ``પૂરતો
મજબૂત''માં $\Th{Q}$થી
થોડું વધુ આવરી લેવું
પડશે.

અપૂર્ણતા પ્રમેયની ગોડેલની
મૂળ સાબિતીનો વિચાર
એપિમેનિડીસના વિરોધાભાસમાં
દેખાય છે. ક્રિટનો વતની
એપિમેનિડીસે કહ્યું કે
ક્રિટના બધા લોકો જૂઠા છે;
વિરોધાભાસનું વધુ સીધું
રૂપ ``આ વાક્ય અસત્ય છે''
એવું વિધાન છે. સત્યને
નિષ્પન્નક્ષમતાથી બદલીને
ગોડેલ એવા વાક્યને
ઔપચારિક બનાવી શક્યા જે
આડકતરી રીતે પોતે
નિષ્પન્ન કરી શકાય એવું નથી એમ
કહે છે. જો તે વાક્ય
નિષ્પન્ન કરી શકાય એવું હોય, તો
સિદ્ધાંત અસુસંગત બને.
ગોડેલે બતાવ્યું કે તેઓ
વિચારતા હતા એવી
સ્વયંસિદ્ધિ-પદ્ધતિમાંથી
તે વાક્યનો નિષેધ
પણ નિષ્પન્ન કરી શકાય એવું નથી.
(આ બીજા ભાગ માટે ગોડેલે
સિદ્ધાંત~$\Th{T}$ને
``$\omega$-સુસંગત''
હોવાની ધારણા કરવી પડી.
$\omega$-સુસંગતતા
સુસંગતતા સાથે સંબંધિત છે,
પરંતુ વધુ મજબૂત ગુણધર્મ
છે.\footnote{એટલે દરેક
$\omega$-સુસંગત સિદ્ધાંત
સુસંગત હોય છે, પરંતુ
દરેક સુસંગત સિદ્ધાંત
એવો હોય જ એવું
નથી.} ગોડેલના થોડાં વર્ષો
પછી રોસરે બતાવ્યું કે
માત્ર~$\Th{T}$ની
સુસંગતતા ધારવી પૂરતી છે.)

પહેલો પડકાર એવો
વાક્ય કેવી રીતે
રચવું તે સમજવાનો છે,
જે પોતાનો જ ઉલ્લેખ કરે.
$\Th{Q}$ની ભાષાના દરેક
સૂત્ર~$!A$ માટે
~$\gn{!A}$ એ
~$\Gn{!A}$ને અનુરૂપ
અંકપદ દર્શાવે એમ માનો.
આનો અર્થ વિચારો:
$!A$ એ~$\Th{Q}$ની
ભાષાનું સૂત્ર છે,
$\Gn{!A}$ એ પ્રાકૃતિક
સંખ્યા છે, અને
$\gn{!A}$ એ~$\Th{Q}$ની
ભાષાનું \emph{પદ} છે.
આથી $\Th{Q}$ની ભાષાના
દરેક સૂત્ર~$!A$ને
એક \emph{નામ},
~$\gn{!A}$, મળે છે,
જે $\Th{Q}$ની ભાષાનું
પદ છે. આમ એક વૈચારિક
માળખું મળે છે જેમાં
$\Th{Q}$ની ભાષાનાં
સૂત્ર બીજા
સૂત્ર વિશે કંઈક
``કહી'' શકે છે.
નીચેનું લેમ્મા નિશ્ચિત-બિંદુ
લેમ્મા કહેવાય છે.

\begin{lem}
$\Th{T}$ એ~$\Th{Q}$નું
કોઈપણ વિસ્તરણ હોય અને
$!B(x)$ એવું કોઈપણ
સૂત્ર હોય જેમાં
માત્ર ચલ~$x$ મુક્ત હોય.
તો એવું વાક્ય~$!A$
છે કે
$\Th{T} \Proves!A \liff !B(\gn{!A})$.
\end{lem}

લેમ્મા કહે છે કે કોઈપણ
ગુણધર્મ~$!B(x)$ આપેલો
હોય, તો એવું
વાક્ય~$!A$ છે
જે ``$!B(x)$ મારા માટે
સાચું છે'' એમ કહે છે,
અને $\Th{T}$ પણ આ
``જાણે'' છે.

આવું વાક્ય કેવી
રીતે રચી શકીએ?
ક્વાઇનનું નીચેનું
એપિમેનિડીસ-વિરોધાભાસનું
રૂપ વિચારો:
\begin{quote}
``પોતાનું અવતરણ પહેલાં
મૂકવાથી અસત્ય મળે છે''
પોતાનું અવતરણ પહેલાં
મૂકવાથી અસત્ય મળે છે.
\end{quote}
આ વાક્ય સીધો પોતાનો
ઉલ્લેખ કરતું નથી.
તે અવતરણચિહ્નો વચ્ચેના
વાક્યખંડ નામના વાક્યરચનાત્મક
પદાર્થ વિશે માત્ર વિધાન
કરે છે; આમ તે નીચેનાં
વાક્યોની હરોળમાં આવે છે:
\begin{enumerate}
\item ``રોબર્ટ'' સુંદર નામ છે.
\item ``હું દોડ્યો.'' ટૂંકું વાક્ય છે.
\item ``ત્રણ શબ્દો છે''માં ત્રણ શબ્દો છે.
\end{enumerate}
પરંતુ ``પોતાનું અવતરણ
પહેલાં મૂકવાથી અસત્ય મળે
છે'' વાક્યખંડની પહેલાં
એ જ વાક્યખંડને
અવતરણચિહ્નોમાં મૂકીશું
તો શું થશે? આપણને
મૂળ વાક્ય જ મળશે!{}
ટૂંકમાં, આ વાક્ય પોતે
અસત્ય છે એમ કહે છે.
% END OLP-0313
\endgroup


\begingroup
\makeatletter\def\ol@sectioncs{section}\makeatother

% BEGIN OLP-0314
\olfileid{inc}{inp}{fix}

\olsection{નિશ્ચિત-બિંદુ લેમ્મા}
\phantomsection\label{guunit:OLP-0314}


\begin{explain}
નિશ્ચિત-બિંદુ લેમ્મા કહે છે કે
કોઈપણ સૂત્ર~$!B(x)$ માટે
એવું વાક્ય~$!A$ છે કે
$\Th{T} \Proves !A \liff !B(\gn{!A})$,
જો $\Th{T}$ એ~$\Th{Q}$નું
વિસ્તરણ હોય. જૂઠા વાક્યના
કિસ્સામાં, આપણે $!A$ એ
~``$\gn{!A}$ અસત્ય છે''
એની સાથે (~$\Th{T}$માં
સાબિતિક્ષમ રીતે) સમકક્ષ
હોય એવું ઇચ્છીએ; એટલે
$\Gn{!A}$ એ અસત્ય
વાક્યની ગોડેલ સંખ્યા
છે એવું વિધાન.
સાબિતીનો વિચાર સમજવા,
$!A$ માટે ક્વાઇનનું
અનૌપચારિક રૂપ સરખાવવું
ઉપયોગી છે: ``{}`પોતાનું
અવતરણ પહેલાં મૂકવાથી
અસત્ય મળે છે' પોતાનું
અવતરણ પહેલાં મૂકવાથી
અસત્ય મળે છે.'' કોઈ
અભિવ્યક્તિની પહેલાં તેનું
પોતાનું અવતરણ મૂકીને
વાક્ય રચવાની ક્રિયાને તે
અભિવ્યક્તિનું
\emph{વિકર્ણીકરણ} કહી શકાય;
મળેલું વાક્ય તેનું
વિકર્ણિત રૂપ છે.
એટલે `પોતાનું અવતરણ
પહેલાં મૂકવાથી અસત્ય મળે
છે'નું વિકર્ણિત રૂપ
``{}`પોતાનું અવતરણ પહેલાં
મૂકવાથી અસત્ય મળે છે'
પોતાનું અવતરણ પહેલાં
મૂકવાથી અસત્ય મળે છે'' છે.
હવે નોંધો કે ક્વાઇનનું
જૂઠું વાક્ય `અસત્ય મળે છે'નું
વિકર્ણિત રૂપ નથી, પણ
`પોતાનું અવતરણ પહેલાં
મૂકવાથી અસત્ય મળે છે'નું
છે. તેથી જૂઠું વાક્ય
મેળવવા જે ગુણધર્મનું
વિકર્ણીકરણ થાય છે તેમાં
જ વિકર્ણીકરણ સમાયેલું છે!{}

અંકગણિતની ભાષામાં એક
મુક્ત ચલવાળા સૂત્રની
ગોડેલ સંખ્યા ગણીએ અને
પછી તે ચલને એ ગોડેલ
સંખ્યાના પ્રમાણભૂત અંકપદથી
બદલીએ; આમ તેના અવતરણ
રચાય છે.
~$!E(x)$નું વિકર્ણિત
રૂપ $!E(\num{n})$ છે,
જ્યાં $n = \Gn{!E(x)}$.
(હવેથી $\num{\Gn{!E(x)}}$ને
$\gn{!E(x)}$ લખીશું.)
તેથી $!B(x)$ જો
``અસત્ય છે'' હોય,
તો ``પોતાનું અવતરણ પહેલાં
મૂકતાં અસત્ય મળે છે''
એનો અર્થ ``પોતાના
વિકર્ણિત રૂપની ગોડેલ
સંખ્યા પર લાગુ કરતાં
અસત્ય મળે છે'' થાય.
ધારો કે સૂત્રની
ગોડેલ સંખ્યા~$n$ પરથી
તેના વિકર્ણિત રૂપની
ગોડેલ સંખ્યા ગણતા વિધેય
$\fn{diag}(n)$ માટે
આપણી પાસે
~$\Obj{diag}$ ચિહ્ન હોય.
તો $!E(x)$ને
$!B(\Obj{diag}(x))$
લખી શકીએ. અને ક્વાઇનના
જૂઠા વાક્યનું રૂપ તેનું
વિકર્ણીકરણ, એટલે
$!E(\gn{!E(x)})$ અથવા
$!B(\Obj{diag}(\gn{!B(\Obj{diag}(x))}))$,
થાય. અલબત્ત, હવે
$!B(x)$ બીજો કોઈપણ
ગુણધર્મ હોઈ શકે અને એ
જ રચના કામ કરશે.
અપૂર્ણતા પ્રમેય માટે
$!B(x)$ તરીકે
``$x$ એ~$\Th{T}$માં
નિષ્પન્ન કરી શકાય એવું નથી'' લઈશું.
ત્યારે $!E(x)$નો અર્થ
``પોતાના વિકર્ણિત રૂપની
ગોડેલ સંખ્યા પર લાગુ
કરતાં $\Th{T}$માં
નિષ્પન્ન કરી શકાય એવું ન હોય એવું
વાક્ય મળે છે''
થાય.

આને~$\Th{T}$માં ઔપચારિક
બનાવવા $\fn{diag}$ને
ઔપચારિક કરવાની રીત
જોઈએ. વિધેય
$\fn{diag}(n)$ સંગણનીય,
હકીકતમાં આદિમ પુનરાવર્તી
છે: $n$ જો
સૂત્ર~$!E(x)$ની ગોડેલ
સંખ્યા હોય, તો
$\fn{diag}(n)$ એ
~$!E(\gn{!E(x)})$ની
ગોડેલ સંખ્યા આપે છે.
(યાદ રાખો,
$\gn{!E(x)}$ એ
~$!E(x)$ની ગોડેલ સંખ્યાનું
પ્રમાણભૂત અંકપદ છે,
એટલે $\num{\Gn{!E(x)}}$.)
જો $\Obj{diag}$ એ
$\fn{diag}$નું નિરૂપણ
કરતું $\Th{T}$માંનું
વિધેય-ચિહ્ન હોય,
તો $!A$ તરીકે સૂત્ર
$!B(\Obj{diag}(\gn{!B(\Obj{diag}(x))}))$
લઈ શકીએ. નોંધો કે
\begin{align*}
\fn{diag}(\Gn{!B(\Obj{diag}(x))}) & =
\Gn{!B(\Obj{diag}(\gn{!B(\Obj{diag}(x))}))} \\
& = \Gn{!A}.
\end{align*}
માનો કે $\Th{T}$ નીચેનું
નિષ્પન્ન કરી શકે:
\[
\Obj{diag}(\gn{!B(\Obj{diag}(x))}) = \gn{!A},
\]
તો તે
$!B(\Obj{diag}(\gn{!B(\Obj{diag}(x))}))
\liff !B(\gn{!A})$
પણ નિષ્પન્ન કરી શકે.
પરંતુ ડાબી બાજુ
વ્યાખ્યા મુજબ~$!A$ છે.

અલબત્ત, સામાન્ય રીતે
$\Obj{diag}$ એ
$\Th{T}$નું વિધેય-ચિહ્ન
હશે નહીં; ખાસ કરીને
~$\Th{Q}$નું તો નથી જ.
પરંતુ $\fn{diag}$ સંગણનીય
હોવાથી કોઈ સૂત્ર
$!D_{\fn{diag}}(x,y)$ તેને
~$\Th{Q}$માં
\emph{નિરૂપે} છે.
તેથી $!B(\Obj{diag}(x))$
લખવાને બદલે
$\lexists[y][(!D_{\fn{diag}}(x,y) \land !B(y))]$
લખી શકીએ. બાકી ઉપર
રેખાંકિત કરેલી સાબિતી
ચાલે છે; ખરેખર, તે
~$\Th{Q}$માં જ ચાલે છે.
\end{explain}

\begin{lem}
\ollabel{lem:fixed-point}
$!B(x)$ એ માત્ર એક
મુક્ત ચલ~$x$ ધરાવતું
કોઈપણ સૂત્ર હોય.
તો એવું વાક્ય~$!A$ છે કે
$\Th{Q} \Proves
!A \liff !B(\gn{!A})$.
\end{lem}


\begin{proof}
$!B(x)$ આપેલું હોય ત્યારે
$!E(x)$ એ સૂત્ર
$\lexists[y][(!D_{\fn{diag}}(x,y) \land !B(y))]$
લો અને $!A$ એ તેનું
વિકર્ણિત રૂપ, એટલે
સૂત્ર $!E(\gn{!E(x)})$ લો.

$!D_{\fn{diag}}$ એ
$\fn{diag}$નું નિરૂપણ કરે
છે અને
$\fn{diag}(\Gn{!E(x)}) = \Gn{!A}$
હોવાથી $\Th{Q}$ નીચેનું
નિષ્પન્ન કરી શકે છે:
\begin{align}
  & !D_{\fn{diag}}(\gn{!E(x)}, \gn{!A}) \ollabel{repdiag1} \\
  & \lforall[y][(!D_{\fn{diag}}(\gn{!E(x)},y) \lif
  \eq[y][\gn{!A}])]. \ollabel{repdiag2}
\end{align}
હવે બતાવીએ કે
$\Th{Q} \Proves !A \liff !B(\gn{!A})$.
માત્ર તર્ક અને
~$\Th{Q}$માં નિષ્પન્ન કરી શકાય એવું
હકીકતો વાપરીને
અનૌપચારિક દલીલ કરીશું.

પ્રથમ $!A$, એટલે
$!E(\gn{!E(x)})$ માનો.
$!E(x)$ની વ્યાખ્યા પરથી
$!E(\gn{!E(x)})$ એટલે
\[
\lexists[y][(!D_{\fn{diag}}(\gn{!E(x)},y) \land !B(y))].
\]
આવો~$y$ લો.
$!D_{\fn{diag}}(\gn{!E(x)},y)$
હોવાથી \olref{repdiag2}
મુજબ $y = \gn{!A}$.
તેથી $!B(y)$ પરથી
~$!B(\gn{!A})$ મળે છે.

હવે $!B(\gn{!A})$ માનો.
\olref{repdiag1} મુજબ
\begin{align*}
& !D_{\fn{diag}}(\gn{!E(x)}, \gn{!A}) \land !B(\gn{!A}).
\intertext{આથી મળે છે કે}
& \lexists[y][(!D_{\fn{diag}}(\gn{!E(x)},y) \land !B(y))].
\end{align*}
પરંતુ એ જ
$!E(\gn{!E(x)})$,
એટલે~$!A$ છે.
\end{proof}

\begin{digress}
આ સાબિતીને સંગણનીયતા
સિદ્ધાંતના નિશ્ચિત-બિંદુ
લેમ્માની સાબિતી સાથે
સરખાવો. તફાવત એ છે કે
અહીં પોતાને આધારે
\emph{વિધાન} વ્યાખ્યાયિત
કરવું છે, જ્યારે ત્યાં
પોતાને આધારે
\emph{વિધેય} વ્યાખ્યાયિત
કરવું હતું. આ ભેદ સિવાય
વિચાર ખરેખર એક જ છે.
\end{digress}

\begin{prob}
  સૂત્ર~$!A(x)$ને
  \emph{સત્યની વ્યાખ્યા}
  ત્યારે કહીએ જ્યારે બધાં
  વાક્ય~$!B$ માટે
  $\Th{Q}
  \Proves !B \liff !A(\gn{!B})$
  હોય. નિશ્ચિત-બિંદુ લેમ્મા
  વાપરીને બતાવો કે કોઈ
  સૂત્ર સત્યની
  વ્યાખ્યા નથી.
\end{prob}
% END OLP-0314
\endgroup


\begingroup
\makeatletter\def\ol@sectioncs{section}\makeatother

% BEGIN OLP-0315
\olfileid{inc}{inp}{1in}

\olsection{પ્રથમ અપૂર્ણતા પ્રમેય}
\phantomsection\label{guunit:OLP-0315}


હવે પ્રથમ અપૂર્ણતા પ્રમેયની
ગોડેલની મૂળ સાબિતી
વર્ણવી શકીએ. $\Th{T}$
એ અંકગણિતની ભાષાનું
વિસ્તરણ કરતી ભાષામાંનો
કોઈપણ સંગણનીય રીતે
સ્વયંસિદ્ધિત સિદ્ધાંત હોય
અને $\Th{T}$માં
$\Th{Q}$નાં સ્વયંસિદ્ધિઓ
સમાવેલાં હોય. ખાસ કરીને
તેનો અર્થ એ છે કે
$\Th{T}$ સંગણનીય
વિધેયો અને સંબંધોનું
નિરૂપણ કરે છે.

અગાઉ દલીલ કરી છે કે
સૂત્રો અને સાબિતીઓને
સંખ્યાઓ વડે યોગ્ય રીતે
સંકેતિત કરીએ તો સંબંધ
$\Prf[\Th{T}](x,y)$
સંગણનીય છે. અહીં
$\Prf[\Th{T}](x,y)$
ત્યારે અને માત્ર ત્યારે જ
સાચો છે જ્યારે $x$ એ
$\Th{T}$માં ગોડેલ સંખ્યા
~$y$ ધરાવતા સૂત્રના
નિષ્પત્તિની ગોડેલ
સંખ્યા હોય. હકીકતમાં
ગોડેલે ધ્યાનમાં રાખેલા
ચોક્કસ સિદ્ધાંત માટે
પોતાના લેખની ૪૫ વિધેયો
અને સંબંધોની યાદીથી
તેમણે આ સંબંધ આદિમ
પુનરાવર્તી છે એવું બતાવ્યું.
૪૫મો સંબંધ, $x B y$,
તેમની $\Th{T}$ની ખાસ
પસંદગી માટેનો
$\Prf[\Th{T}](x,y)$
જ છે. યાદ રાખો કે
ગોડેલ પોતાના લેખમાં
``પુનરાવર્તી'' કહે છે
ત્યાં આજે આપણે
``આદિમ પુનરાવર્તી''
કહીશું.

$\Prf[\Th{T}](x,y)$ સંગણનીય
હોવાથી $\Th{T}$માં તેનું
નિરૂપણ થાય છે. તેને
નિરૂપતા સૂત્ર માટે
$\OPrf[\Th{T}](x,y)$
લખીશું. $\OProv[\Th{T}](y)$
એ સૂત્ર
$\lexists[x][\OPrf[\Th{T}](x,y)]$
હોય. આ ગોડેલની યાદીનો
૪૬મો સંબંધ,
$\fn{Bew}(y)$, વર્ણવે છે.
ગોડેલ નોંધે છે તેમ,
આ એકમાત્ર એવો સંબંધ
છે જે ``પુનરાવર્તી છે''
એવું કહી શકાતું નથી.
કદાચ તેમનો અર્થ આ હતો:
વ્યાખ્યા પરથી તે સંગણનીય
છે એવું સ્પષ્ટ થતું નથી;
અને પછીનો વિકાસ બતાવે
છે કે હકીકતમાં તે
સંગણનીય નથી.

$\Th{T}$ એ~$\Th{Q}$ને
સમાવતો કોઈ
સ્વયંસિદ્ધીકરણીય સિદ્ધાંત
હોય. ત્યારે
$\Prf[\Th{T}](x, y)$
નિર્ણેય છે અને તેથી
કોઈ સૂત્ર
~$\OPrf[\Th{T}](x, y)$
વડે~$\Th{Q}$માં નિરૂપાય
છે. ઉપર વર્ણવેલું
$\OProv[\Th{T}](y)$
સૂત્ર લો. નિશ્ચિત-બિંદુ
લેમ્માથી એવું સૂત્ર
$!G_\Th{T}$ છે કે
$\Th{Q}$ (અને તેથી
$\Th{T}$ પણ) નીચેનું
નિષ્પન્ન કરે છે:
\begin{equation}
\ollabel{eqn:qpf}
!G_\Th{T} \liff \lnot \OProv[\Th{T}](\gn{!G_\Th{T}}).
\end{equation}
નોંધો કે $!G_\Th{T}$નો
મૂળ અર્થ ``$!G_\Th{T}$
એ~$\Th{T}$માં
નિષ્પન્ન કરી શકાય એવું નથી'' છે.

\begin{lem}\ollabel{lem:cons-G-unprov}
જો $\Th{T}$ એ~$\Th{Q}$નું
સુસંગત,
સ્વયંસિદ્ધીકરણીય વિસ્તરણ
હોય, તો
$\Th{T} \Proves/ !G_\Th{T}$.
\end{lem}

\begin{proof}
માનો કે $\Th{T}$ એ
$!G_\Th{T}$ને નિષ્પન્ન
કરે છે. તો ખરેખર
નિષ્પત્તિ \emph{છે};
તેથી કોઈ સંખ્યા~$m$ માટે
સંબંધ $\Prf[\Th{T}](m,
\Gn{!G_\Th{T}})$ સાચો છે.
પણ તેથી $\Th{Q}$ વાક્ય
$\OPrf[\Th{T}](\num m, \gn{!G_\Th{T}})$ને
નિષ્પન્ન કરે છે.
આથી $\Th{Q}$ એ
$\lexists[x][\OPrf[\Th{T}](x,\gn{!G_\Th{T}})]$
પણ નિષ્પન્ન કરે છે,
જે વ્યાખ્યા મુજબ
$\OProv[\Th{T}](\gn{!G_\Th{T}})$
છે. \olref{eqn:qpf}
મુજબ $\Th{Q}$ એ
$\lnot !G_\Th{T}$ને
નિષ્પન્ન કરે છે;
અને $\Th{T}$ એ
$\Th{Q}$નું વિસ્તરણ
હોવાથી $\Th{T}$ પણ
એમ કરે છે. બતાવ્યું કે
જો $\Th{T}$ એ
$!G_\Th{T}$ને નિષ્પન્ન
કરે, તો $\lnot !G_\Th{T}$ને
પણ નિષ્પન્ન કરે છે;
તેથી તે અસુસંગત બને.
\end{proof}

\begin{defn}
\ollabel{thm:oconsis-q}
સિદ્ધાંત $\Th{T}$ને
\emph{$\omega$-સુસંગત}
ત્યારે કહીએ જ્યારે આ
શરત સંતોષાય: જો
$\lexists[x][!A(x)]$
કોઈપણ વાક્ય હોય અને
$\Th{T}$ એ
$\lnot
!A(\num 0)$,
$\lnot !A(\num 1)$,
$\lnot !A(\num 2)$,
\dots બધાંને
નિષ્પન્ન કરે,
તો $\Th{T}$ એ
$\lexists[x][!A(x)]$ને
સાબિત કરતું નથી.
\end{defn}

દરેક $\omega$-સુસંગત
સિદ્ધાંત સુસંગત પણ છે.
કારણ કે જો $\Th{T}$
અસુસંગત હોય, તો દરેક
~$!A$ માટે
$\Th{T} \Proves !A$.
ખાસ કરીને, જો
$\Th{T}$ અસુસંગત હોય,
તો દરેક~$n$ માટે
$\lnot !A(\num n)$ને
નિષ્પન્ન કરે છે અને
$\lexists[x][!A(x)]$ને
પણ નિષ્પન્ન કરે છે.
તેથી $\Th{T}$ અસુસંગત
હોય તો તે
$\omega$-અસુસંગત છે.
વિપરિત અનુમાનથી,
$\Th{T}$ જો
$\omega$-સુસંગત હોય
તો તે સુસંગત હોવું જ
જોઈએ.

\begin{lem}\ollabel{lem:omega-cons-G-unref}
જો $\Th{T}$ એ~$\Th{Q}$નું
$\omega$-સુસંગત,
સ્વયંસિદ્ધીકરણીય વિસ્તરણ
હોય, તો
$\Th{T} \Proves/ \lnot !G_\Th{T}$.
\end{lem}

\begin{proof}
જો $\Th{T}$ એ
$\lnot !G_\Th{T}$ને
નિષ્પન્ન કરે તો તે
$\omega$-અસુસંગત છે
એવું બતાવીએ. માનો કે
$\Th{T}$ એ
$\lnot !G_\Th{T}$ને
નિષ્પન્ન કરે છે.
જો $\Th{T}$ અસુસંગત
હોય તો તે
$\omega$-અસુસંગત છે,
એટલે કામ પૂરું.
નહીં તો $\Th{T}$
સુસંગત છે, તેથી
\olref{lem:cons-G-unprov}
મુજબ તે $!G_\Th{T}$ને
નિષ્પન્ન કરતું નથી.
$\Th{T}$માં
$!G_\Th{T}$નું
કોઈ નિષ્પત્તિ નથી;
તેથી $\Th{Q}$ નીચેનું
નિષ્પન્ન કરે છે:
\[
\lnot \OPrf[\Th{T}](\num 0, \gn{!G_\Th{T}}), \lnot \OPrf[\Th{T}](\num 1,
\gn{!G_\Th{T}}), \lnot \OPrf[\Th{T}](\num 2, \gn{!G_\Th{T}}), \dots
\]
અને $\Th{T}$ પણ એમ
કરે છે. બીજી બાજુ,
\olref{eqn:qpf} મુજબ
$\lnot
!G_\Th{T}$ એ
$\lexists[x][\OPrf[\Th{T}](x,\gn{!G_\Th{T}})]$
ની સમકક્ષ છે. તેથી
$\Th{T}$
$\omega$-અસુસંગત છે.
\end{proof}

\begin{prob}
  દરેક $\omega$-સુસંગત
  સિદ્ધાંત સુસંગત હોય છે.
  વિપરીત વિધાન ખોટું છે,
  એટલે સુસંગત છતાં
  $\omega$-અસુસંગત
  સિદ્ધાંતો છે, તે બતાવો.
  આ માટે $\Th{Q} \cup
  \{\lnot !G_\Th{Q}\}$
  સુસંગત પરંતુ
  $\omega$-અસુસંગત છે
  તે સાબિત કરો.
\end{prob}

\begin{thm}
\ollabel{thm:first-incompleteness}
$\Th{T}$ એ~$\Th{Q}$નું
કોઈપણ $\omega$-સુસંગત,
સ્વયંસિદ્ધીકરણીય વિસ્તરણ
હોય. તો $\Th{T}$ પૂર્ણ
નથી.
\end{thm}

\begin{proof}
  જો $\Th{T}$
  $\omega$-સુસંગત છે,
  તો તે સુસંગત છે; તેથી
  \olref{lem:cons-G-unprov}
  મુજબ $\Th{T}
  \Proves/ !G_\Th{T}$.
  \olref{lem:omega-cons-G-unref}
  મુજબ
  $\Th{T} \Proves/ \lnot !G_\Th{T}$.
  આથી $\Th{T}$ અપૂર્ણ
  છે: તે $!G_\Th{T}$
  કે $\lnot !G_\Th{T}$,
  બેમાંથી એકેયને
  નિષ્પન્ન કરતું નથી.
\end{proof}
% END OLP-0315
\endgroup


\begingroup
\makeatletter\def\ol@sectioncs{section}\makeatother

% BEGIN OLP-0316
\olfileid{inc}{inp}{ros}

\olsection{રોસરનું પ્રમેય}
\phantomsection\label{guunit:OLP-0316}


``$\omega$-સુસંગત''ને
બદલે માત્ર ``સુસંગત''
રાખીને વધુ મજબૂત
પરિણામ મેળવવા ગોડેલની
સાબિતી બદલી શકીએ?
રોસરે શોધેલી યુક્તિથી
જવાબ ``હા'' છે.
રોસરની યુક્તિ
$\OProv[\Th{T}](y)$ને
બદલે ``બદલાયેલું''
નિષ્પન્નક્ષમતા વિધેયધર્મ
$\ORProv_T(y)$ વાપરે છે.

\begin{thm}
\ollabel{thm:rosser}
$\Th{T}$ એ~$\Th{Q}$નું
કોઈપણ સુસંગત,
સ્વયંસિદ્ધીકરણીય વિસ્તરણ
હોય. તો $\Th{T}$
પૂર્ણ નથી.
\end{thm}

\begin{proof}
યાદ રાખો કે
$\OProv[\Th{T}](y)$ની
વ્યાખ્યા
$\lexists[x][\OPrf[\Th{T}](x,
  y)]$ છે; અહીં
$\OPrf[\Th{T}](x, y)$
એ એવો નિર્ણેય સંબંધ
નિરૂપે છે કે $x$ એ
ગોડેલ સંખ્યા~$y$
ધરાવતા વાક્યના
નિષ્પત્તિની ગોડેલ
સંખ્યા હોય ત્યારે અને
માત્ર ત્યારે જ સંબંધ
સાચો છે. $x$ અને~$y$
વચ્ચેનો એવો સંબંધ પણ
નિર્ણેય છે જેમાં $x$ એ
ગોડેલ સંખ્યા~$y$ ધરાવતા
વાક્યના \emph{ખંડન}ની
ગોડેલ સંખ્યા હોય ત્યારે
તે સાચો થાય. $\fn{not}(x)$
નીચે મુજબનું આદિમ
પુનરાવર્તી વિધેય હોય:
$x$ જો સૂત્ર~$!A$નો
સંકેત હોય, તો
$\fn{not}(x)$ એ
$\lnot
!A$નો સંકેત છે.
ત્યારે $\Refut[\Th{T}](x, y)$
ત્યારે અને માત્ર ત્યારે જ
સાચો છે જ્યારે
$\Prf[\Th{T}](x, \fn{not}(y))$
સાચો હોય. તેને
$\ORefut[\Th{T}](x, y)$
નિરૂપે એમ માનો. પછી
જો $\Th{T} \Proves \lnot !A$
અને $\delta$ એ અનુરૂપ
નિષ્પત્તિ હોય, તો
$\Th{Q} \Proves
\ORefut[\Th{T}](\gn{\delta}, \gn{!A})$.
$\ORProv[\Th{T}](y)$ની
વ્યાખ્યા આ રાખીએ:
\[
\lexists[x][(\OPrf[\Th{T}](x,y) \land \lforall[z][(z < x \lif \lnot
  \ORefut[\Th{T}](z,y))])].
\]
આશરે, $\ORProv[\Th{T}](y)$
કહે છે કે ``$\Th{T}$માં
$y$ની સાબિતી છે, અને
$y$નું તેની સાબિતીની
સંકેતસંખ્યા કરતાં નાની
સંકેતસંખ્યા ધરાવતું
ખંડન નથી.''
\sourcecorrection{OLINC-043}{સ્થિર અંગ્રેજી
મૂળમાં ``ટૂંકું ખંડન''
કહ્યું છે, પરંતુ વ્યાખ્યાનું
ક્રમ-માપ સાબિતી અને
ખંડનની સંકેતસંખ્યાઓની
સરખામણી છે. સૂત્રમાં
લખાયેલી નાની સંકેતસંખ્યા
અનુસાર અર્થ સ્પષ્ટ કર્યો છે.}
$\Th{T}$ સુસંગત
છે એમ માનીએ તો
$\ORProv[\Th{T}](y)$ અને
$\OProv[\Th{T}](y)$
એક જ સંખ્યાઓ માટે
સાચાં છે; પરંતુ
~$\Th{T}$માં
\emph{સાબિતિક્ષમતા}ની
દૃષ્ટિએ (સત્ય અને
સાબિતિક્ષમતા વચ્ચે તફાવત
છે તે હવે જાણીએ છીએ!)
બંનેના ગુણધર્મો જુદા છે.
જો $\Th{T}$
\emph{અ}સુસંગત હોય,
તો બંને એક જ સંખ્યાઓ
માટે સાચાં \emph{નથી}!{}
($\ORProv[\Th{T}](y)$ને
ઘણી વાર ``$y$ રોસર-સાબિતિક્ષમ
છે'' એમ વાંચે છે.
હમણાં ચર્ચ્યું તેમ,
રોસર-સાબિતિક્ષમતા કોઈ
ખાસ પ્રકારની સાબિતિક્ષમતા
નથી: અસુસંગત સિદ્ધાંતોમાં
કેટલાંક વાક્ય
સાબિતિક્ષમ હોય છતાં
રોસર-સાબિતિક્ષમ નથી.
એટલે આ નામ ગૂંચવી શકે.
ગૂંચવણ ટાળવા તેને
``$y$ શ્મૂવેબલ છે'' પણ
વાંચી શકો.)

નિશ્ચિત-બિંદુ લેમ્માથી એવું
સૂત્ર $!R_\Th{T}$ છે કે
\begin{equation}
  \Th{Q} \Proves !R_\Th{T} \liff \lnot \ORProv[\Th{T}](\gn{!R_\Th{T}}).
  \ollabel{RT}
\end{equation}
\olref[1in]{thm:first-incompleteness}ની
સાબિતીથી વિપરીત, અહીં
દાવો એવો છે કે
$\Th{T}$ સુસંગત હોય તો
$\Th{T}$ એ $!R_\Th{T}$ને
નિષ્પન્ન કરતું નથી,
અને $\Th{T}$ એ
$\lnot !R_\Th{T}$ને પણ
નિષ્પન્ન કરતું નથી.
(બીજા શબ્દોમાં,
$\omega$-સુસંગતતાની
ધારણા જરૂરી નથી.)

પ્રથમ બતાવીએ કે
$\Th{T} \Proves/ !R_{\Th{T}}$.
માનો કે તે સાબિત કરે;
તો~$T$માંથી
~$!R_{\Th{T}}$નું
નિષ્પત્તિ છે;
તેની ગોડેલ સંખ્યા~$n$
માનો. તેથી
$\Th{Q} \Proves \OPrf[\Th{T}](\num{n}, \gn{!R_{\Th{T}}})$,
કારણ કે $\OPrf[\Th{T}]$
એ~$\Th{Q}$માં
$\Prf[\Th{T}]$નું નિરૂપણ
કરે છે. ઉપરાંત દરેક
$k < n$ માટે $k$ એ
$\lnot !R_{\Th{T}}$ના
નિષ્પત્તિની ગોડેલ
સંખ્યા નથી, કારણ કે
$\Th{T}$ સુસંગત છે.
તેથી દરેક $k < n$ માટે
$\Th{Q} \Proves \lnot
\ORefut[\Th{T}](\num{k}, \gn{!R_{\Th{T}}})$.
\olref[req][min]{lem:less-nsucc}
મુજબ
$\Th{Q} \Proves \lforall[z][(z < \num{n} \lif \lnot \ORefut[\Th{T}](z,
  \gn{!R_{\Th{T}}}))]$.
આથી
\[
\Th{Q} \Proves \lexists[x][(\OPrf[\Th{T}](x,\gn{!R_{\Th{T}}}) \land \lforall[z][(z
    < x \lif \lnot \ORefut[\Th{T}](z,\gn{!R_{\Th{T}}}))])],
\]
પણ આ તો
$\ORProv[\Th{T}](\gn{!R_{\Th{T}}})$
જ છે. \olref{RT} મુજબ
$\Th{Q}
\Proves \lnot !R_{\Th{T}}$.
$\Th{T}$ એ $\Th{Q}$નું
વિસ્તરણ હોવાથી
$\Th{T}
\Proves \lnot !R_{\Th{T}}$
પણ. પરંતુ
$\Th{T} \Proves !R_{\Th{T}}$
એવી ધારણા કરી હતી;
તેથી $\Th{T}$ અસુસંગત
બને, જે પ્રમેયની
ધારણાથી વિરુદ્ધ છે.

હવે બતાવીએ કે
$\Th{T} \Proves/ \lnot !R_{\Th{T}}$.
ફરી માનો કે તે
સાબિત કરે અને
~$\lnot !R_{\Th{T}}$ના
નિષ્પત્તિની ગોડેલ
સંખ્યા~$n$ હોય. ત્યારે
$\Refut[\Th{T}](n, \Gn{!R_{\Th{T}}})$
સાચો છે; અને
$\ORefut[\Th{T}]$ એ
$\Th{Q}$માં
$\Refut[\Th{T}]$નું
નિરૂપણ કરે છે, તેથી
$\Th{Q} \Proves
\ORefut[\Th{T}](\num{n}, \gn{!R_{\Th{T}}})$.
અહીં પણ $\Th{T}$
અસુસંગત બનશે એવું
બતાવીશું, કારણ કે તે
~$!R_{\Th{T}}$ને પણ
નિષ્પન્ન કરશે. કારણ કે
\begin{align*}
\Th{Q} & \Proves !R_{\Th{T}} \liff \lnot \ORProv[\Th{T}](\gn{!R_{\Th{T}}}),
\intertext{અને $\Th{T}$ એ~$\Th{Q}$નું વિસ્તરણ હોવાથી, આ બતાવવું પૂરતું છે કે}
\Th{Q} & \Proves \lnot \ORProv[\Th{T}](\gn{!R_{\Th{T}}}).
\end{align*}
વાક્ય $\lnot
\ORProv[\Th{T}](\gn{!R_{\Th{T}}})$,
એટલે
\begin{align*}
  \lnot & \lexists[x][(\OPrf[\Th{T}](x,\gn{!R_{\Th{T}}}) \land
    \lforall[z][(z < x \lif \lnot \ORefut[\Th{T}](z,\gn{!R_{\Th{T}}}))])],
  \intertext{તર્કની દૃષ્ટિએ આને સમકક્ષ છે:}
  & \lforall[x][(\OPrf[\Th{T}](x,\gn{!R_{\Th{T}}}) \lif
    \lexists[z][(z < x \land \ORefut[\Th{T}](z,\gn{!R_{\Th{T}}}))])].
\end{align*}
~$\Th{Q}$ જે હકીકતો
નિષ્પન્ન કરે છે તે
વાપરીને અનૌપચારિક
તાર્કિક દલીલ કરીએ.
$x$ ગમે તે હોય અને
$\OPrf[\Th{T}](x,
\gn{!R_{\Th{T}}})$ માનો.
આપણે પહેલેથી જાણીએ
છીએ કે
$\Th{T} \Proves/ !R_{\Th{T}}$;
તેથી દરેક $k$ માટે
$\Th{Q} \Proves \lnot \OPrf[\Th{T}](\num{k}, \gn{!R_{\Th{T}}})$.
આથી દરેક $k$ માટે
$\eq/[x][\num{k}]$ મળે
છે. ખાસ કરીને
(a) $\eq/[x][\num{n}]$.
ઉપરાંત
$\lnot(\eq[x][\num{0}]
\lor \eq[x][\num{1}] \lor \dots \lor \eq[x][\num{n-1}])$
મળે છે; તેથી
\olref[req][min]{lem:less-nsucc}
મુજબ (b)
$\lnot(x < \num{n})$.
\olref[req][min]{lem:trichotomy}
મુજબ $\num{n} < x$.
અને
$\Th{Q} \Proves
\ORefut[\Th{T}](\num{n}, \gn{!R_{\Th{T}}})$
હોવાથી
$\num{n} < x \land
\ORefut[\Th{T}](\num{n}, \gn{!R_{\Th{T}}})$,
અને તેથી
$\lexists[z][(z < x
\land \ORefut[\Th{T}](z,\gn{!R_{\Th{T}}}))]$
મળે છે. $x$ ગમે તે
હોવાથી માંગેલું વિધાન
મળે છે:
\[
\lforall[x][(\OPrf[\Th{T}](x,\gn{!R_{\Th{T}}}) \lif
  \lexists[z][(z < x \land \ORefut[\Th{T}](z,\gn{!R_{\Th{T}}}))])].
\]
\end{proof}

\begin{prob}
પ્રાકૃતિક સંખ્યાઓના બે
ગણો $A$ અને $B$ને
\emph{સંગણનીય રીતે
અવિભાજ્ય} ત્યારે કહીએ
જ્યારે એવો કોઈ
નિર્ણેય ગણ~$X$
ન હોય કે
$A \subseteq X$ અને
$B \subseteq \Complement{X}$
(અહીં $\Complement{X}$
એ~$X$નો પૂરક ગણ
$\Nat \setminus X$ છે).
$\Th{T}$ એ~$\Th{Q}$નું
સુસંગત,
સ્વયંસિદ્ધીકરણીય વિસ્તરણ
હોય. માનો કે $A$ એ
~$\Th{T}$માં સાબિત
થતાં વાક્યની
ગોડેલ સંખ્યાઓનો ગણ છે,
અને $B$ એ~$\Th{T}$માં
નકારી શકાય એવાં
વાક્યોની ગોડેલ
સંખ્યાઓનો ગણ છે.
$A$ અને~$B$ સંગણનીય
રીતે અવિભાજ્ય છે એમ
સાબિત કરો.
\end{prob}
% END OLP-0316
\endgroup


\begingroup
\makeatletter\def\ol@sectioncs{section}\makeatother

% BEGIN OLP-0317
\olfileid{inc}{inp}{gop}

\olsection{ગોડેલના મૂળ લેખ સાથે સરખામણી}
\phantomsection\label{guunit:OLP-0317}


ગોડેલના ૧૯૩૧ના લેખને
થોડો સમય આપવો ઉપયોગી છે.
તેની પ્રસ્તાવનામાં હમણાં
ચર્ચેલા વિચારોનું રૂપરેખાંકન
છે. લેખ પર ઉપરછલ્લી
નજર નાખીએ તોય દરેક
તબક્કે શું થાય છે તે
સરળતાથી દેખાય છે:
પહેલાં ગોડેલ ઔપચારિક
પદ્ધતિ~$P$ વર્ણવે છે
(વાક્યરચના, સ્વયંસિદ્ધિઓ,
સાબિતીના નિયમો); પછી
આદિમ પુનરાવર્તી વિધેયો
અને સંબંધો વ્યાખ્યાયિત
કરે છે; પછી $x B y$
આદિમ પુનરાવર્તી છે એવું
બતાવે છે અને આદિમ
પુનરાવર્તી વિધેયો અને
સંબંધોનું $\Th{P}$માં
નિરૂપણ થાય છે એવી
દલીલ કરે છે. પછી ઉપર
જણાવ્યા મુજબ અપૂર્ણતા
પ્રમેય સાબિત કરે છે.
વિભાગ ૩માં તેઓ બતાવે છે
કે અસાબિતિક્ષમ વિધાનને
અંકગણિતની ભાષાનું વાક્ય
લઈ શકાય છે. અહીંથી
$\beta$-લેમ્માનો ઉદ્ભવ
થાય છે. પુનરાવર્તી
વિધેયોનું $\Th{Q}$માં
નિરૂપણ થાય છે તે
બતાવતાં અનુક્રમોને
સંભાળવા આપણે પણ આ
લેમ્મા વાપર્યું હતું.
ગોડેલ આ માટે પૂરતાં હોય
એવાં ન્યૂનતમ સ્વયંસિદ્ધિઓ
અલગ તારવતા નથી; પરંતુ
હવે જાણીએ છીએ કે
$\Th{Q}$ પૂરતું છે.
છેલ્લે, વિભાગ ૪માં
તેઓ બીજા અપૂર્ણતા
પ્રમેયની સાબિતીની
રૂપરેખા આપે છે.
% END OLP-0317
\endgroup


\begingroup
\makeatletter\def\ol@sectioncs{section}\makeatother

% BEGIN OLP-0318
\olfileid{inc}{inp}{prc}

\olsection{$\Th{PA}$ માટેની નિષ્પન્નક્ષમતા શરતો}
\phantomsection\label{guunit:OLP-0318}


પેનો અંકગણિત, એટલે
$\Th{PA}$, એ
$\Th{Q}$ને દરેક
સૂત્ર માટેની
ગાણિતિક અનુમાનની
સ્વયંસિદ્ધિઓથી વિસ્તરાવતો
સિદ્ધાંત છે. એટલે
$\Th{Q}$માં આ આકારની
સ્વયંસિદ્ધિઓ ઉમેરીએ:
\[
(!A(0) \land \lforall[x][(!A(x) \lif !A(x'))]) \lif \lforall[x][!A(x)]
\]
દરેક સૂત્ર~$!A$
માટે. નોંધો કે આ ખરેખર
એક \emph{યોજના} છે,
એટલે અનંત સ્વયંસિદ્ધિઓ
(અને $\Th{PA}$ને સાન્ત
સ્વયંસિદ્ધિ-ગણથી
વ્યાખ્યાયિત કરી શકાતું
નથી). પરંતુ ચિહ્નોની
કોઈ શ્રેણી ગાણિતિક
અનુમાનની સ્વયંસિદ્ધિનું
ઉદાહરણ છે કે નહીં તે
અસરકારક રીતે નક્કી કરી
શકાય છે; તેથી
$\Th{PA}$નો સ્વયંસિદ્ધિ-ગણ
સંગણનીય છે. $\Th{PA}$
એ~$\Th{Q}$ કરતાં ઘણો
વધુ શક્તિશાળી સિદ્ધાંત
છે. દાખલા તરીકે, સામાન્ય
રીતે ગાણિતિક અનુમાન
વાપરીને સરવાળો અને ગુણાકાર
ક્રમવિનિમયી છે એ સરળતાથી
સાબિત થાય છે. હકીકતમાં
મોટા ભાગની સાન્તવાદી
સંખ્યાસિદ્ધાંત અને
સંયોજનશાસ્ત્રની દલીલો
~$\Th{PA}$માં કરી શકાય છે.

$\Th{PA}$ સંગણનીય રીતે
સ્વયંસિદ્ધિત હોવાથી
નિષ્પન્નક્ષમતા વિધેયધર્મ
$\Prf[\Th{PA}](x,y)$
સંગણનીય છે; તેથી તેનું
~$\Th{Q}$માં (અને તેથી
~$\Th{PA}$માં પણ)
નિરૂપણ થાય છે. પહેલાંની
જેમ સંબંધને નિરૂપતા
સૂત્ર માટે
$\OPrf[\Th{PA}](x,y)$
લખીશું. $\OProv[\Th{PA}](y)$
એ સૂત્ર
$\lexists[x][\OPrf[\Th{PA}](x,y)]$
હોય. તેનો સાહજિક અર્થ
``$y$ એ $\Th{PA}$ની
સ્વયંસિદ્ધિઓમાંથી
નિષ્પન્ન કરી શકાય એવું છે'' થાય છે.
\sourcecorrection{OLINC-044}{સ્થિર અંગ્રેજી
મૂળમાં આ અસ્તિત્વવાચક
સૂત્રની અંદર બહારનો
સંબંધ મૂકાયો છે. અહીં
અગાઉ વ્યાખ્યાયિત
અંકગણિતીય નિરૂપક સૂત્ર
મૂક્યું છે જેથી આખી
અભિવ્યક્તિ સિદ્ધાંતની
ભાષાનું સૂત્ર રહે.}
$\Th{Q}$ની સ્વયંસિદ્ધિઓ
કરતાં થોડું વધારે
શા માટે જોઈએ?
કારણ કે વપરાતો સિદ્ધાંત
આ નિષ્પન્નક્ષમતા
વિધેયધર્મ વિશેની કેટલીક
મૂળભૂત હકીકતો
નિષ્પન્ન કરી શકે એટલો
મજબૂત હોવો જોઈએ.
જરૂરી હકીકતો આ છે:
\begin{enumerate}
\item[P1.] જો
  $\Th{PA} \Proves !A$,
  તો $\Th{PA} \Proves
  \OProv[\Th{PA}](\gn{!A})$.
\item[P2.] બધાં
  સૂત્ર $!A$ અને
  $!B$ માટે,
  \[
  \Th{PA} \Proves \OProv[\Th{PA}](\gn{!A \lif !B}) \lif
  (\OProv[\Th{PA}](\gn{!A}) \lif \OProv[\Th{PA}](\gn{!B})).
  \]
\item[P3.] દરેક
  સૂત્ર~$!A$ માટે,
  \[
  \Th{PA} \Proves \OProv[\Th{PA}](\gn{!A})
  \lif \OProv[\Th{PA}](\gn{\OProv[\Th{PA}](\gn{!A})}).
  \]
\end{enumerate}
આ ત્રણ ગુણધર્મો સાચા
છે તે ચકાસવા માટે
સૂત્ર
$\OProv[\Th{PA}](y)$નું
કાળજીથી વર્ણન કરવું અને
સંબંધિત ઔપચારિક
નિષ્પત્તિ વર્ણવવા
$\Th{PA}$ની
સ્વયંસિદ્ધિઓ વાપરવી
પડે. શરતો (૧) અને
~(૨) સરળ છે; ખરેખર
કામ શરત~(૩)માં છે.
(વિચારો કે તેમાં કેવું
કામ કરવું પડે \dots)
વિગતો આપવી કંટાળાજનક
અને અહીં રસપ્રદ નથી;
તેથી $\Th{PA}$ ઉપરના
ત્રણ ગુણધર્મો ધરાવે છે
એ વાત અહીં માન્ય રાખવા
કહીએ છીએ. $\OProv[\Th{PA}](y)$ની
યોગ્ય પસંદગી આ પણ
સંતોષે છે:
\begin{enumerate}
\item[P4.] જો
  $\Th{PA} \Proves \OProv[\Th{PA}](\gn{!A})$,
  તો $\Th{PA} \Proves !A$.
\end{enumerate}
પરંતુ આ હકીકતની
આપણને જરૂર પડશે નહીં.

\begin{digress}
એક વાત વધુ: ગોડેલ પણ
અત્યારે આપણે છીએ એટલા
જ વિગતોમાં આળસુ હતા.
૧૯૩૧ના લેખના અંતે તેમણે
બીજા અપૂર્ણતા પ્રમેયની
સાબિતીની રૂપરેખા આપી
અને પછીના લેખમાં વિગતો
આપવાનું વચન આપ્યું.
પરંતુ તેમણે એ વિગતો
ક્યારેય આપી નહીં;
દલીલ સમજતા સૌ માનતા
હતા કે તેને પૂરું કરી
શકાય છે (એટલે તેમને
વિગતો ભરવાની જરૂર
પડી નહીં).
\end{digress}
% END OLP-0318
\endgroup


\begingroup
\makeatletter\def\ol@sectioncs{section}\makeatother

% BEGIN OLP-0319
\olfileid{inc}{inp}{2in}

\olsection{બીજું અપૂર્ણતા પ્રમેય}
\phantomsection\label{guunit:OLP-0319}


$\Th{PA}$ પોતાની સુસંગતતા સાબિત કરતું નથી એવો દાવો કેવી રીતે વ્યક્ત કરવો? $\Th{PA}$ અસુસંગત છે એમ કહેવું એટલે $\Th{PA} \Proves \eq[0][1]$ એમ કહેવું. તેથી તેની સુસંગતતાનું વિધાન $\OCon[\Th{PA}]$ આપણે વાક્ય $\lnot\OProv[\Th{PA}](\gn{\lfalse})$ લઈ શકીએ: અહીં બંને વિરોધાભાસ સમકક્ષ છે. પછીનું પ્રમેય આપણને જોઈતું પરિણામ આપે છે.
\sourcecorrection{OLINC-045}{સ્થિર અંગ્રેજી મૂળ શરૂઆતમાં સુસંગતતાનું વિધાન $\lnot\OProv[\Th{PA}](\gn{\eq[0][1]})$ લે છે, પણ નીચેની ઔપચારિક દલીલમાં $\OCon \ident \lnot\OProv(\gn{\lfalse})$ વાપરે છે. અહીં દલીલ સાથે સુસંગત રહેવા બીજું નિરૂપણ પસંદ કર્યું છે; $\Th{PA}$માં બંને વિરોધાભાસોની સાબિતિક્ષમતા સમકક્ષ છે.}

\begin{thm}
\ollabel{thm:second-incompleteness}
જો $\Th{PA}$ સુસંગત હોય, તો $\Th{PA}$ એ $\OCon[\Th{PA}]$ને નિષ્પન્ન કરતું નથી.
\end{thm}

આ પ્રમેય $\OCon[\Th{PA}]$ના ચોક્કસ નિરૂપણ પર, એટલે $\OProv[\Th{PA}](y)$ના ચોક્કસ નિરૂપણ પર, આધાર રાખે છે એ નોંધવું જરૂરી છે. આપણે માત્ર એટલું જ વાપરીશું કે $\OProv[\Th{PA}](y)$નું નિરૂપણ ત્રણ નિષ્પન્નક્ષમતા શરતો સંતોષે છે. તેથી આ ગુણધર્મો ધરાવતા નિષ્પન્નક્ષમતા વિધેયધર્મવાળા કોઈપણ સિદ્ધાંત માટે પ્રમેય વ્યાપક થાય છે.

ગોડેલે આપેલી દલીલની રૂપરેખા વાંચવા જેવી છે: પ્રમેય દલીલના અસરકારક અંતિમ પરિણામની જેમ નીકળે છે. દલીલ આ છે. \olref[1in]{thm:first-incompleteness}ની સાબિતીમાં રચેલું ગોડેલનું વાક્ય $!G_\Th{PA}$ લો. આપણે બતાવ્યું છે: ``જો $\Th{PA}$ સુસંગત હોય, તો $\Th{PA}$ એ $!G_\Th{PA}$ને નિષ્પન્ન કરતું નથી.'' આ વાતને $\Th{PA}$ની \emph{અંદર} ઔપચારિક કરીએ તો નીચેની સાબિતી મળે છે:
\[
\OCon[\Th{PA}] \lif \lnot \OProv[\Th{PA}](\gn{!G_\Th{PA}}).
\]
હવે માનો કે $\Th{PA}$ એ $\OCon[\Th{PA}]$ને નિષ્પન્ન કરે છે. ત્યારે તે $\lnot\OProv[\Th{PA}](\gn{!G_\Th{PA}})$ને પણ નિષ્પન્ન કરે છે. પરંતુ $!G_\Th{PA}$ ગોડેલનું વાક્ય હોવાથી આ $!G_\Th{PA}$ની સમકક્ષ છે. તેથી $\Th{PA}$ એ $!G_\Th{PA}$ને નિષ્પન્ન કરે છે.
\sourcecorrection{OLINC-046}{સ્થિર અંગ્રેજી મૂળમાં અહીં બહારના $\Prov[\Th{PA}]$નો ઉપયોગ છે. અગાઉ વ્યાખ્યાયિત અને દલીલમાં જરૂરી અંકગણિતીય સૂત્ર $\OProv[\Th{PA}]$ મૂક્યું છે.}

પરંતુ આપણે જાણીએ છીએ કે $\Th{PA}$ સુસંગત હોય તો તે $!G_\Th{PA}$ને નિષ્પન્ન કરતું નથી!{} તેથી $\Th{PA}$ સુસંગત હોય તો તે $\OCon[\Th{PA}]$ને નિષ્પન્ન કરી શકતું નથી.

દલીલને વધુ ચોક્કસ બનાવવા $!G_\Th{PA}$ને $\Th{PA}$ માટેનું ગોડેલ વાક્ય લઈએ અને નિષ્પન્નક્ષમતા શરતો (P1)--(P3) વડે બતાવીએ કે $\Th{PA}$ એ $\OCon[\Th{PA}] \lif !G_\Th{PA}$ને નિષ્પન્ન કરે છે. તેથી $\Th{PA}$ એ $\OCon[\Th{PA}]$ને નિષ્પન્ન કરતું નથી એવું નીકળશે. હવે $\Th{PA}$ની \emph{અંદર} સાબિતીની રૂપરેખા જોઈએ. સરળતા માટે $\Th{PA}$ના નીચલા સૂચક છોડીએ છીએ.
\begin{align}
& !G \liff \lnot \OProv(\gn{!G}) \ollabel{G2-1}\\
& \qquad\text{$!G$ ગોડેલનું વાક્ય છે}\notag \\
& !G \lif \lnot \OProv(\gn{!G}) \ollabel{G2-2}\\
  & \qquad\text{\olref{G2-1} પરથી} \notag\\
& !G \lif
  (\OProv(\gn{!G}) \lif \lfalse) \ollabel{G2-3}\\
  & \qquad\text{\olref{G2-2} અને તર્ક પરથી}\notag\\
& \OProv(\gn{
    !G \lif
    (\OProv(\gn{!G}) \lif \lfalse)
  }) \ollabel{G2-4}\\
  & \qquad\text{\olref{G2-3} અને શરત P1 પરથી} \notag\\
& \OProv(\gn{!G}) \lif
  \OProv(\gn{
    (\OProv(\gn{!G}) \lif \lfalse)
    }) \ollabel{G2-5}\\
  & \qquad\text{\olref{G2-4} અને શરત P2 પરથી} \notag\\
& \OProv(\gn{!G}) \lif (\OProv(\gn{\OProv(\gn{!G})}) \lif \OProv(\gn{\lfalse})) \ollabel{G2-6}\\
  & \qquad\text{\olref{G2-5}, શરત P2 અને તર્ક પરથી} \notag\\
& \OProv(\gn{!G}) \lif
  \OProv(\gn{\OProv(\gn{!G})}) \ollabel{G2-7}\\
   & \qquad\text{શરત P3 પરથી} \notag\\
& \OProv(\gn{!G}) \lif \OProv(\gn{\lfalse}) \ollabel{G2-8}\\
  & \qquad \text{\olref{G2-6}, \olref{G2-7} અને તર્ક પરથી}\notag\\
& \OCon \lif \lnot \OProv(\gn{!G}) \ollabel{G2-9}\\
  & \qquad\text{\olref{G2-8}નું પ્રતિધનાત્મક રૂપ અને $\OCon \ident \lnot \OProv(\gn{\lfalse})$ પરથી}\notag \\
& \OCon \lif !G \notag\\
  & \qquad\text{\olref{G2-1}, \olref{G2-9} અને તર્ક પરથી}\notag
\end{align}
ઉપર તર્કનો ઉપયોગ માત્ર વિધાનતર્કની પ્રાથમિક હકીકતો માટે છે. દાખલા તરીકે, \olref{G2-3}માં $\Proves \lnot!A \liff (!A\lif\lfalse)$ વપરાય છે; અને \olref{G2-8}માં $!A \lif (!B \lif !C), !A \lif !B \Proves !A \lif !C$ વપરાય છે. \olref{G2-5} અને \olref{G2-6}માં શરત~P2નો ઉપયોગ તેના ઉદાહરણો $\OProv(\gn{!A \lif !B}) \lif (\OProv(\gn{!A}) \lif \OProv(\gn{!B}))$ પર આધાર રાખે છે. પહેલા ઉદાહરણમાં $!A \ident !G$ અને $!B \ident \OProv(\gn{!G}) \lif \lfalse$ છે; બીજામાં $!A \ident \OProv(\gn{!G})$ અને $!B \ident \lfalse$ છે.
\sourcecorrection{OLINC-047}{સ્થિર અંગ્રેજી મૂળની છેલ્લી સમજણમાં બીજા ઉદાહરણ માટે $\OProv(\gn{G})$ છે, જ્યારે ઉપરની દલીલ તથા પહેલા ઉદાહરણમાં વાક્ય $!G$ છે. અહીં $\OProv(\gn{!G})$ રાખ્યું છે.}

બીજા અપૂર્ણતા પ્રમેયનું વધુ અમૂર્ત સ્વરૂપ આ છે:

\begin{thm}
\ollabel{thm:second-incompleteness-gen}
$\Th{T}$ એ $\Th{Q}$નું કોઈ સુસંગત, સ્વયંસિદ્ધીકૃત વિસ્તરણ હોય અને $\OProv[\Th{T}](y)$ એ $\Th{T}$ માટે નિષ્પન્નક્ષમતા શરતો P1--P3 સંતોષતું કોઈ સૂત્ર હોય. તો $\Th{T}$ એ $\OCon[\Th{T}]$ને નિષ્પન્ન કરતું નથી.
\end{thm}
\sourcecorrection{OLINC-048}{સ્થિર અંગ્રેજી મૂળમાં છેલ્લે $\OCon[T]$ છે. આ પ્રમેય અને અગાઉની વ્યાખ્યા સાથે સુસંગત સિદ્ધાંત-ચિહ્ન $\OCon[\Th{T}]$ રાખ્યું છે.}

\begin{prob}
બતાવો કે $\Th{PA}$ એ $!G_{\Th{PA}} \lif \OCon[\Th{PA}]$ને નિષ્પન્ન કરે છે.
\end{prob}

\begin{digress}
આમાંથી બોધ એ મળે છે કે ગણિતનો કોઈપણ ``વાજબી'' સુસંગત સિદ્ધાંત પોતાની સુસંગતતાનું વિધાન નિષ્પન્ન કરી શકતો નથી. માનો કે $\Th{T}$ એ ગણિતનો એવો સિદ્ધાંત છે જેમાં $\Th{Q}$ તથા હિલ્બર્ટની ``સાન્તવાદી'' દલીલો (જેનો અર્થ જે કંઈ હોય તે) સામેલ છે. તો આખો $\Th{T}$ એ $\Th{T}$નું સુસંગતતાનું વિધાન નિષ્પન્ન કરી શકતો નથી; તેથી તેની સાન્તવાદી ભાગ-રચના પણ $\Th{T}$નું એ વિધાન નિષ્પન્ન કરી શકતી નથી. આ અર્થમાં ``આખા ગણિત'' માટે સાન્તવાદી સુસંગતતા-સાબિતી હોઈ શકતી નથી.

``સાન્તવાદી'' શબ્દનો અર્થ સમજવામાં થોડો અવકાશ છે. ગોડેલે ૧૯૩૧ના લેખમાં એવી શક્યતા સ્વીકારી હતી કે જેને આપણે ``સાન્તવાદી'' ગણીએ તે હિલ્બર્ટ ઔપચારિક કરવા માગતા હતા તે પ્રકારના ગણિતની બહાર હોઈ શકે. પરંતુ ગોડેલ અહીં ઉદાર હતા: આજે સાન્તવાદી કહેવા યોગ્ય, છતાં દાખલા તરીકે પસંદગીના સ્વયંસિદ્ધિ સાથેના ઝર્મેલો--ફ્રેન્કેલ ગણસિદ્ધાંત $\Th{ZFC}$માં ઔપચારિક ન કરી શકાય એવું કંઈ શોધવું મુશ્કેલ છે.
\end{digress}
% END OLP-0319
\endgroup


\begingroup
\makeatletter\def\ol@sectioncs{section}\makeatother

% BEGIN OLP-0320
\olfileid{inc}{inp}{lob}

\olsection{લોબનું પ્રમેય}
\phantomsection\label{guunit:OLP-0320}


સિદ્ધાંત~$\Th{T}$ માટેનું ગોડેલ વાક્ય $\lnot\OProv[\Th{T}](y)$નું નિશ્ચિત-બિંદુ છે; એટલે એવું વાક્ય~$!G$ કે
\[
\Th{T} \Proves \lnot \OProv[\Th{T}](\gn{!G}) \liff !G.
\]
સુસંગત સિદ્ધાંતમાં તે નિષ્પન્ન કરી શકાય એવું નથી. કારણ કે જો $\Th{T} \Proves !G$, તો (a) નિષ્પન્નક્ષમતા શરત~(1)થી $\Th{T} \Proves \OProv[\Th{T}](\gn{!G})$; અને (b) $\Th{T} \Proves !G$ તથા $\Th{T} \Proves \lnot \OProv[\Th{T}](\gn{!G}) \liff !G$ મળીને $\Th{T} \Proves \lnot \OProv[\Th{T}](\gn{!G})$ આપે. તેથી $\Th{T}$ અસુસંગત બને. હવે $\OProv[\Th{T}](y)$ના નિશ્ચિત-બિંદુની સ્થિતિ શું છે એ પૂછવું સ્વાભાવિક છે. એટલે એવું વાક્ય~$!H$ કે
\[
\Th{T} \Proves \OProv[\Th{T}](\gn{!H}) \liff !H.
\]
જો તે નિષ્પન્ન કરી શકાય એવું હોય, તો શરત~(1)થી $\Th{T} \Proves \OProv[\Th{T}](\gn{!H})$; પણ ઉપરની સમકક્ષતા પર મોડસ પોનેન્સ લાગુ કરતાં પણ એ જ તારણ મળે. તેથી, ઓછામાં ઓછું ગોડેલ વાક્યના કિસ્સામાં વપરાયેલી દલીલથી, $\Th{T}$ અસુસંગત છે એવું નીકળતું નથી. અલબત્ત, આથી $\Th{T}$ ખરેખર $!H$ને નિષ્પન્ન કરે છે એવું પણ સાબિત થતું નથી.
\sourcecorrection{OLINC-049}{સ્થિર અંગ્રેજી મૂળ ગોડેલ વાક્ય અસાબિત્ય છે એવું નિઃશરત કહે છે, પરંતુ પછીની દલીલ માત્ર તેની સાબિતીથી સિદ્ધાંત અસુસંગત બનશે એવું બતાવે છે. તેથી અહીં સુસંગતતાની શરત સ્પષ્ટ કરી છે; અસુસંગત સિદ્ધાંતમાં વાક્ય સાબિત થઈ શકે.}

આ પ્રશ્નને થોડો વ્યાપક કરીએ તો આગળ વધી શકીએ. નિશ્ચિત-બિંદુની સમકક્ષતામાં ડાબેથી જમણી દિશા, $\OProv[\Th{T}](\gn{!H}) \lif !H$, \emph{પ્રતિબિંબન સિદ્ધાંત} કહેવાતી સામાન્ય યોજનાનું ઉદાહરણ છે: $\OProv[\Th{T}](\gn{!A}) \lif !A$. તેને એવું નામ એટલા માટે છે કે, એક અર્થમાં, તે $\Th{T}$ને પોતે શું નિષ્પન્ન કરી શકે છે તે અંગે ``પ્રતિબિંબિત'' કરવા દે છે. મૂળભૂત રીતે તે દરેક~$!A$ માટે કહે છે: ``જો $\Th{T}$ એ $!A$ને નિષ્પન્ન કરી શકે, તો $!A$ સાચું છે.'' અલબત્ત, આ ફક્ત યથાર્થ સિદ્ધાંતો માટે સાચું છે; તેથી સામાન્ય રીતે સિદ્ધાંતો તેના દરેક ઉદાહરણને નિષ્પન્ન કરતા નથી એવું સૂચાય છે. તો પૂરતી શક્તિ અને નિષ્પન્નક્ષમતા શરતો ધરાવતો સિદ્ધાંત તેના કયા ઉદાહરણો નિષ્પન્ન કરી શકે? નિશ્ચિતપણે એવા બધાં, જ્યાં $!A$ પોતે નિષ્પન્ન કરી શકાય એવું છે. પછીના પરિણામ મુજબ, બસ એટલાં જ.

\begin{thm}\ollabel{thm:lob}
$\Th{T}$ એ $\Th{Q}$નું સ્વયંસિદ્ધીકરણીય વિસ્તરણ હોય અને $\OProv[\Th{T}](y)$ એ \olref[2in]{sec}ની શરતો P1--P3 સંતોષતું સૂત્ર હોય. જો $\Th{T}$ એ $\OProv[\Th{T}](\gn{!A}) \lif !A$ને નિષ્પન્ન કરે, તો હકીકતમાં $\Th{T}$ એ $!A$ને નિષ્પન્ન કરે છે.
\end{thm}

બીજી રીતે કહીએ તો, જો $\Th{T} \Proves/ !A$, તો $\Th{T} \Proves/ \OProv[\Th{T}](\gn{!A}) \lif !A$. આ પરિણામ લોબના પ્રમેય તરીકે ઓળખાય છે.

\begin{explain}
લોબના પ્રમેયની સાબિતી પાછળની યુક્તિ સાન્તા ક્લોઝ અસ્તિત્વમાં છે એવી ચતુર સાબિતી છે. (તમને આ તારણ ન ગમતું હોય તો ગમતું કોઈ બીજું તારણ તેની જગ્યાએ મૂકી શકો.) દલીલ આ છે:
\begin{enumerate}
\item $X$ એ આ વાક્ય હોય: ``જો $X$ સાચું હોય, તો સાન્તા ક્લોઝ અસ્તિત્વમાં છે.''
\item માનો કે $X$ સાચું છે.
\item તો તે જે કહે છે તે સાચું છે; એટલે જો $X$ સાચું હોય, તો સાન્તા ક્લોઝ અસ્તિત્વમાં છે.
\item $X$ સાચું છે એમ માન્યું હોવાથી (2) અને~(3) પરથી મોડસ પોનેન્સ વડે સાન્તા ક્લોઝ અસ્તિત્વમાં છે એવું તારવી શકીએ.
\item ધારણા~(2), ``$X$ સાચું છે,'' પરથી આપણે (4), ``સાન્તા ક્લોઝ અસ્તિત્વમાં છે,'' તારવ્યું. શરતી સાબિતી વડે બતાવ્યું: ``જો $X$ સાચું હોય, તો સાન્તા ક્લોઝ અસ્તિત્વમાં છે.''
\item પરંતુ આ તો વાક્ય~$X$ જ છે. એટલે $X$ સાચું છે એવું બતાવ્યું.
\item પછી ઉપરની (2)--(4) દલીલથી સાન્તા ક્લોઝ અસ્તિત્વમાં છે.
\end{enumerate}
આ વિચારને ઔપચારિક કરતાં ``સાચું છે''ની જગ્યાએ ``નિષ્પન્ન કરી શકાય એવું છે'' અને ``સાન્તા ક્લોઝ અસ્તિત્વમાં છે''ની જગ્યાએ~$!A$ મૂકીએ તો લોબના પ્રમેયની સાબિતી મળે છે. યુક્તિ સૂત્ર~$\OProv[\Th{T}](y) \lif !A$ પર નિશ્ચિત-બિંદુ લેમ્મા લાગુ કરવાની છે. તેનું નિશ્ચિત-બિંદુ ઉપરની રૂપરેખાના વાક્ય~$X$ને અનુરૂપ છે.
\end{explain}

\begin{proof}[\olref{thm:lob}ની સાબિતી]
માનો કે $!A$ એવું વાક્ય છે કે $\Th{T}$ એ $\OProv[\Th{T}](\gn{!A}) \lif !A$ને નિષ્પન્ન કરે છે. $!B(y)$ એ સૂત્ર~$\OProv[\Th{T}](y) \lif !A$ હોય. નિશ્ચિત-બિંદુ લેમ્માથી એવું વાક્ય~$!D$ શોધો કે $\Th{T}$ એ $!D \liff !B(\gn{!D})$ને નિષ્પન્ન કરે. હવે નીચેની દરેક અભિવ્યક્તિ $\Th{T}$માં નિષ્પન્ન કરી શકાય એવું છે:
\begin{align}
  & !D \liff (\OProv[\Th{T}](\gn{!D}) \lif !A) \ollabel{L-1}\\
  & \qquad \text{$!D$ એ~$!B(y)$નું નિશ્ચિત-બિંદુ છે}\notag \\
  & !D \lif (\OProv[\Th{T}](\gn{!D}) \lif !A) \ollabel{L-2}\\
  & \qquad\text{\olref{L-1} પરથી}\notag\\
  & \OProv[\Th{T}](\gn{!D \lif (\OProv[\Th{T}](\gn{!D}) \lif !A)}) \ollabel{L-3}\\
  & \qquad \text{\olref{L-2} અને શરત P1 પરથી}\notag \\
  & \OProv[\Th{T}](\gn{!D}) \lif \OProv[\Th{T}](\gn{\OProv[\Th{T}](\gn{!D}) \lif !A})
  \ollabel{L-4}\\
  &\qquad \text{\olref{L-3} અને શરત P2 પરથી}\notag \\
  & \OProv[\Th{T}](\gn{!D}) \lif (\OProv[\Th{T}](\gn{\OProv[\Th{T}](\gn{!D})}) \lif \OProv[\Th{T}](\gn{!A})) \ollabel{L-5}\\
  &\qquad \text{\olref{L-4} અને ફરી P2 પરથી} \notag\\
& \OProv[\Th{T}](\gn{!D}) \lif \OProv[\Th{T}](\gn{\OProv[\Th{T}](\gn{!D})}) \ollabel{L-6}\\
  & \qquad\text{નિષ્પન્નક્ષમતા શરત P3 પરથી} \notag\\
  & \OProv[\Th{T}](\gn{!D}) \lif \OProv[\Th{T}](\gn{!A}) \ollabel{L-7} \\
  &\qquad\text{\olref{L-5} અને \olref{L-6} પરથી}\notag\\
  & \OProv[\Th{T}](\gn{!A}) \lif !A \ollabel{L-8}\\
  &\qquad\text{પ્રમેયની ધારણા પરથી} \notag\\
  & \OProv[\Th{T}](\gn{!D}) \lif !A \ollabel{L-9}\\
  &\qquad\text{\olref{L-7} અને \olref{L-8} પરથી}\notag\\
  & (\OProv[\Th{T}](\gn{!D}) \lif !A) \lif !D \ollabel{L-10}\\
  & \qquad \text{\olref{L-1} પરથી}\notag \\
  & !D \ollabel{L-11}\\
  & \qquad\text{\olref{L-9} અને \olref{L-10} પરથી}\notag \\
  & \OProv[\Th{T}](\gn{!D}) \ollabel{L-12}\\
  & \qquad\text{\olref{L-11} અને શરત~P1 પરથી}\notag \\
  & !A \qquad\qquad\text{\olref{L-8} અને \olref{L-12} પરથી}\notag
\end{align}
\end{proof}

લોબનું પ્રમેય મળ્યા પછી બીજું અપૂર્ણતા પ્રમેય ટૂંકમાં સાબિત કરી શકાય છે (જ્યાં નિષ્પન્નક્ષમતા વિધેયધર્મ P1--P3 શરતો સંતોષે છે): જો $\Th{T} \Proves \OProv[\Th{T}](\gn{\lfalse}) \lif \lfalse$, તો $\Th{T} \Proves \lfalse$. જો $\Th{T}$ સુસંગત હોય, તો $\Th{T} \Proves/ \lfalse$. તેથી $\Th{T} \Proves/ \OProv[\Th{T}](\gn{\lfalse}) \lif \lfalse$, એટલે $\Th{T} \Proves/ \OCon[\Th{T}]$. આ પ્રમેયથી $\OProv[\Th{T}](x)$ના નિશ્ચિત-બિંદુ~$!H$ને પણ નિષ્પન્ન કરી શકાય એવું બતાવી શકાય છે. કારણ કે
\begin{align*}
  \Th{T} & \Proves \OProv[\Th{T}](\gn{!H}) \liff !H\\
  \intertext{ખાસ કરીને}
    \Th{T} & \Proves \OProv[\Th{T}](\gn{!H}) \lif !H
\end{align*}
અને તેથી લોબના પ્રમેયથી $\Th{T} \Proves !H$.


\begin{prob}
$\Th{T}$ એ સંગણનીય રીતે સ્વયંસિદ્ધિત સિદ્ધાંત હોય અને $\OProv[\Th{T}]$ એ $\Th{T}$ માટેનો નિષ્પન્નક્ષમતા વિધેયધર્મ હોય. નીચેના ચાર દાવા વિચારો:
\begin{enumerate}
\item જો $T \Proves !A$, તો $T \Proves \OProv[\Th{T}](\gn{!A})$.
\item $T \Proves !A \lif \OProv[\Th{T}](\gn{!A})$.
\item જો $T \Proves \OProv[\Th{T}](\gn{!A})$, તો $T \Proves !A$.
\item $T \Proves \OProv[\Th{T}](\gn{!A}) \lif !A$.
\end{enumerate}
કઈ શરતો હેઠળ આમાંથી દરેક દાવો સાચો છે?
\end{prob}
% END OLP-0320
\endgroup


\begingroup
\makeatletter\def\ol@sectioncs{section}\makeatother

% BEGIN OLP-0321
\olfileid{inc}{inp}{tar}

\olsection{સત્યની અવ્યાખ્યેયતા}
\phantomsection\label{guunit:OLP-0321}


\emph{વ્યાખ્યેયતા}નો ખ્યાલ અંકગણિતની ભાષા માટે ઔપચારિક અર્થવિજ્ઞાન હોવા પર આધાર રાખે છે. આપણે અંકગણિતની ભાષામાં સૂત્રો અને વાક્યોનો ગણ વર્ણવ્યો છે. ``અભિપ્રેત અર્થઘટન'' મુજબ એ વાક્યો કુદરતી સંખ્યાઓ વિશે દાવા કરે છે; આવા દાવા સાચા કે ખોટા હોઈ શકે. $\Struct{N}$ એ એવો સંરચના હોય જેનો પ્રદેશ $\Nat$ છે અને જેમાં અંકગણિતની ભાષાના ચિહ્નોનું પ્રમાણભૂત અર્થઘટન છે. ત્યારે $\Sat{N}{!A}$નો અર્થ ``પ્રમાણભૂત અર્થઘટનમાં $!A$ સાચું છે'' થાય છે.

\begin{defn}
કુદરતી સંખ્યાઓનો સંબંધ $R(x_1,\dots,x_k)$, $\Struct{N}$માં \emph{વ્યાખ્યેય} ત્યારે અને માત્ર ત્યારે જ છે જ્યારે અંકગણિતની ભાષામાં એવું સૂત્ર $!A(x_1,\dots,x_k)$ હોય કે બધાં $n_1,\dots,n_k$ માટે $R(n_1,\dots,n_k)$ ત્યારે અને માત્ર ત્યારે જ હોય જ્યારે $\Sat{N}{!A(\num n_1,\dots,\num n_k)}$ હોય.
\end{defn}

બીજી રીતે કહીએ તો, સંબંધ $\Struct{N}$માં વ્યાખ્યેય ત્યારે અને માત્ર ત્યારે જ છે જ્યારે તે સિદ્ધાંત $\Th{TA}$માં નિરૂપણક્ષમ હોય. અહીં $\Th{TA} = \Setabs{!A}{\Sat{N}{!A}}$ અંકગણિતનાં સાચાં વાક્યોનો ગણ છે. (આ વાત તરત સ્પષ્ટ ન લાગે તો વ્યાખ્યાઓ ફરી જોઈને પોતાને ખાતરી કરાવો.)

\begin{lem}
દરેક સંગણનીય સંબંધ $\Struct{N}$માં વ્યાખ્યેય છે.
\end{lem}

\begin{proof}
$\Th{Q}$માં કોઈ સંબંધને નિરૂપતું સૂત્ર $\Struct{N}$માં એ જ સંબંધ વ્યાખ્યાયિત કરે છે એ ચકાસવું સહેલું છે.
\end{proof}

હવે વિપરીત દિશા પણ સાચી છે કે નહીં એવું પૂછીએ: શું $\Struct{N}$માં વ્યાખ્યેય દરેક સંબંધ સંગણનીય છે? જવાબ ના છે. દાખલા તરીકે:

\begin{lem}
અટકવાનો સંબંધ $\Struct{N}$માં વ્યાખ્યેય છે.
\end{lem}

\begin{proof}
યાદ કરો કે ક્લીનીનું સામાન્ય સ્વરૂપ પ્રમેય કહે છે: દરેક આંશિક સંગણનીય વિધેય~$f$ માટે એવો સૂચક~$e$ છે કે બધાં $x \in \Nat$ માટે $f(x) = U(\umin{s}{T(e,x,s)})$; અહીં $U$ અને $T$ આદિમ પુનરાવર્તી અને તેથી સર્વત્ર વ્યાખ્યાયિત છે. આથી $f(x)$ વ્યાખ્યાયિત હોય (એટલે સંગણના અટકે) ત્યારે અને માત્ર ત્યારે જ કોઈ~$s$ માટે $T(e,x,s)$ સાચું હોય.

હવે $H$ અટકવાનો સંબંધ હોય, એટલે
\[
H = \Setabs{\tuple{e,x}}{\lexists[s][T(e, x, s)]}.
\]
$!D_T$ એ $\Struct{N}$માં $T$ને વ્યાખ્યાયિત કરતું સૂત્ર હોય. તો
\[
H = \Setabs{\tuple{e,x}}{\Sat{N}{\lexists[s][!D_T(\num e, \num x, s)]}},
\]
અને તેથી $\lexists[s][!D_T(z, x, s)]$ એ $\Struct{N}$માં $H$ને વ્યાખ્યાયિત કરે છે.
\end{proof}

\begin{prob}
બતાવો કે $Q(n) \defiff n \in \Setabs{\Gn{!A}}{\Th{Q} \Proves !A}$ અંકગણિતમાં વ્યાખ્યેય છે.
\end{prob}

પોતે $\Th{TA}$ વિશે શું? શું તે અંકગણિતમાં વ્યાખ્યેય છે? એટલે, શું ગણ $\Setabs{\Gn{!A}}{\Sat{N}{!A}}$ અંકગણિતમાં વ્યાખ્યેય છે? ટાર્સ્કીનું પ્રમેય તેનો નકારાત્મક જવાબ આપે છે.

\begin{thm}
\ollabel{thm:tarski}
અંકગણિતનાં સાચાં વાક્યનો ગણ અંકગણિતમાં વ્યાખ્યેય નથી.
\end{thm}

\begin{proof}
માનો કે $!D(x)$ તેને વ્યાખ્યાયિત કરે છે, એટલે $\Sat{N}{!A}$ ત્યારે અને માત્ર ત્યારે જ જ્યારે $\Sat{N}{!D(\gn{!A})}$. નિશ્ચિત-બિંદુ લેમ્માથી એવું સૂત્ર $!A$ છે કે $\Th{Q} \Proves !A \liff \lnot !D(\gn{!A})$ અને તેથી $\Sat{N}{!A \liff \lnot !D(\gn{!A})}$. પરંતુ ત્યાર પછી $\Sat{N}{!A}$ ત્યારે અને માત્ર ત્યારે જ જ્યારે $\Sat{N}{\lnot !D(\gn{!A})}$; આ વાત $!D(y)$ અંકગણિતનાં સાચાં વિધાનોનો ગણ વ્યાખ્યાયિત કરે છે એવી ધારણાને વિરોધે છે.
\end{proof}

ટાર્સ્કીએ આ વિશ્લેષણને સત્યના વધુ વ્યાપક દાર્શનિક ખ્યાલ પર લાગુ કર્યું. કોઈપણ ભાષા $L$ માટે તેમણે દલીલ કરી કે $L$ માટે સત્યનો યોગ્ય ખ્યાલ દરેક વાક્ય $X$ માટે આ શરત સંતોષે:
\begin{quote}
`$X$' ત્યારે અને માત્ર ત્યારે જ સાચું છે જ્યારે $X$.
\end{quote}
અંગ્રેજી માટેનું ટાર્સ્કીનું વારંવાર ઉલ્લેખાતું ઉદાહરણ ગુજરાતી અર્થમાં આ છે:
\begin{quote}
`બરફ સફેદ છે' ત્યારે અને માત્ર ત્યારે જ સાચું છે જ્યારે બરફ સફેદ છે.
\end{quote}
પરંતુ કર્ણીય વિધેયનું નિરૂપણ કરી શકે એટલી શક્તિશાળી કોઈપણ ભાષા અને કોઈપણ ભાષાકીય વિધેયધર્મ $T(x)$ માટે આપણે એવું વાક્ય $X$ રચી શકીએ જે ``$X$ ત્યારે અને માત્ર ત્યારે જ જ્યારે $T(\text{`$X$'})$ નથી'' સંતોષે. સત્યનો વિધેયધર્મ કોઈ વાક્યને એકસાથે સાચું અને ખોટું ન ઠરાવે એવી આપણી માગણી હોવાથી ટાર્સ્કીએ તારવ્યું કે ભાષાની મર્યાદાની બહાર કોઈ રીતે ગયા વિના તે ભાષાનાં બધાં વાક્યો માટે સત્યનો વિધેયધર્મ નિર્દેશી શકાતો નથી. બીજા શબ્દોમાં, ભાષા માટેનો સત્યનો વિધેયધર્મ એ ભાષામાં જ વ્યાખ્યાયિત થઈ શકતો નથી.
% END OLP-0321
\endgroup


\OLEndChapterHook
% END OLP-0312
\endgroup


\OLEndPartHook
% END OLP-0274
